\documentclass[11pt]{article}
\usepackage{amsmath,amssymb,amsthm,mathtools}
\usepackage{enumitem}
\usepackage{booktabs,array,longtable}
\usepackage[margin=2.4cm]{geometry}
\usepackage[hidelinks]{hyperref}
\setlist{itemsep=1pt,topsep=3pt}
\allowdisplaybreaks
\emergencystretch=2em

% The handbook's macros, so the family's notation is available.
\newcommand{\cE}{\mathcal E}
\newcommand{\cB}{\mathcal B}
\newcommand{\cH}{\mathcal H}
\newcommand{\one}{\mathbf 1}
\newcommand{\mvec}{\mathbf m}
\newcommand{\bmu}{\boldsymbol\mu}
\newcommand{\ip}[2]{\langle #1,#2\rangle}
\newcommand{\oE}{\mathbb E}
\newcommand{\sG}{\mathsf G}
\newcommand{\sZ}{\mathsf Z}
\newcommand{\sB}{\mathsf B}
\newcommand{\sT}{\mathsf T}
\newcommand{\Gst}[4]{\Gamma_{\substack{#1\\#2}}\!\begin{bmatrix}#3\\#4\end{bmatrix}}
\newcommand{\IZ}{\mathbb I_{\mathbb Z}}
\newcommand{\Ilat}[1]{\mathbb I_{#1}}
\newcommand{\Tx}[3]{(\sT_{#1;(1)}[x]\,#2)((#3,\infty))}
\DeclareMathOperator{\Reg}{Reg}
\DeclareMathOperator{\re}{Re}
\DeclareMathOperator{\im}{Im}
\DeclareMathOperator*{\CT}{CT}
% Typing aids only; each prints the full expression it stands for.
\newcommand{\Pc}[1]{P_{\cB}(#1\,\one;#1,#1)}
\newcommand{\Clo}[1]{[\alpha^0]\big(K(\alpha)\Psi(\alpha)/\alpha^{\cE}\big)^{#1}}
\newcommand{\Cor}{\sum_\pi\prod_{k=2}^nc(U_k)^{-1}}
\newcommand{\Pth}{P_{\binom{[3]}1}(\one;1,1)}
\newcommand{\Pfour}{P_{\binom{[4]}1}(\one;1,1)}
\newcommand{\Lch}{\mathbb I_{\mathbb R^{\cE-e_0}}}
\newcommand{\Zt}{\Tx3{\IZ}0}

\theoremstyle{plain}
\newtheorem{proposition}{Proposition}[section]
\newtheorem{theorem}[proposition]{Theorem}
\newtheorem{corollary}[proposition]{Corollary}
\theoremstyle{definition}
\newtheorem{definition}[proposition]{Definition}
\newtheorem{procedure}[proposition]{Procedure}
\newtheorem{question}[proposition]{Question}
\newtheorem{hypothesis}[proposition]{Hypothesis}
\newtheorem{example}[proposition]{Example}
\theoremstyle{remark}
\newtheorem{remark}[proposition]{Remark}
\renewcommand{\thesection}{15}
\renewcommand{\theproposition}{15.\arabic{proposition}}
\newcommand{\derived}{\textup{[derived]}}
\newcommand{\computed}{\textup{[computed]}}
\newcommand{\recalled}{\textup{[recalled]}}
\newcommand{\notyet}{\textup{[not yet computed]}}
\setcounter{tocdepth}{1}

\title{Foreman notebook, \S\S15--22\\[4pt]\large The centre ray read off the list; the five-mark lists along their hinge trees; the masters of $\binom{[5]}2$ and two marks doubled, closed; the reading of a ray through the origin; the three-mark centre, set as an examination; the law of the stride; the gauge end; the torus. With a ledger of what is sorted, open and corrected}
\date{22--24 October 2026}
\author{}

\begin{document}
\maketitle

\noindent\emph{This file brings together the entries of 22--24 October 2026 and replaces them: \S\S15--19 with the questions of \S20 (22--23 October), the answers to those questions (23 October), \S21 ``Which negatives are absolute'' and \S22 ``The negatives hashed out'' (24 October). \S15 follows \S14.17 of the running notebook. Where two entries proved or computed the same thing, one statement stands: \S20 prints each question at the head of its answer, with the later readings of the five-mark frame, the transform, the six-mark list and the merit placed where they answer it; \S21 is the gauge end, its first form, its form for every number of marks and the readings of both entries merged; \S22 is the torus. Corrections made by later entries are applied in place. The ongoing things, the sortings of negatives and hypotheses, the open lists, the corrections and the receipts, are merged in the Ledger at the end, which also gives the old numbers against the new; numbers $15.x$--$19.x$ are unchanged. Handbook citations are by printed number and page of the handbook PDF of 20 October 2026; ``the note'' is the tree-coaction note of 22 October 2026. Every object is a member of the serif family or a construction of the handbook, written with its slots; classical names appear only as aliases, in the last section.}

{\small\tableofcontents}

\section{The centre ray read off the list}

The neighbours of the five-mark centre are the values of the centre ray of the list (Def.~35.9, p.~309), the period $P_{\cB}(\sigma\one;\sigma,\sigma)$. The interest of the list is one Gramma (Thm.~8.3, p.~54), the period is its residue on the balance wall, and the residue is the Gramma of the letter columns (Thm.~2.9, p.~15). Its share side (Prop.~4.2, p.~25) is the working form of this entry: a fibre integral of Gammas, one per basis for each letter, pinned by equal incidence. The pin is what the list contributes. Closing the fibre at the negative strengths (Thm.~4.5, p.~26) lands on the pairs of $k$-tuples of bases with equal incidence; the residue at zero shares is the corner of Thm.~38.5(i); the native shares (Thm.~4.3, p.~25) are uniform, the share side of the table's gauge; and the banana is the list whose pin identifies the two letters mark by mark, which turns its share side into the Mellin read of a power of one Gramma.

\subsection*{15.1\quad The centre period on its share side}

\begin{proposition}[The share side of the centre]\label{a:p:share}\derived; receipts \computed\
Let $(\cE,\cB)$ be a list, $n=|\cE|\ge2$, $e_0\in\cE$, and $\re\sigma>0$. Then
\[
\Gamma(\sigma)^2\,\Pc\sigma=\int_{\substack{|\mathbf a|=\sigma,\ |\mathbf b|=\sigma\\ \sum_Ta_T\one_T=\sum_Tb_T\one_T}}\ \prod_{T\in\cB}\Gamma(a_T)\,\Gamma(b_T),
\]
the fibre integral of Def.~4.1, p.~25, with one share $a_T$ for the letter $K$ and one share $b_T$ for the letter $\Psi$ at every basis, pinned by the two letter totals and by the incidence at the marks $e\neq e_0$; it has dimension $2|\cB|-n-1$, and the uniform shares $a_T=b_T=\sigma/|\cB|$ are an admissible real point.
\end{proposition}

\begin{proof}
By Thm.~8.3, p.~54, $\Gamma(s)\Gamma(t)\,J_{\cB}(\mvec;s,t)=\sG_{(\bmu;s;t);t_{\cB}}(\one)$ with the unit columns $(\mathbf e_e;0;0)$ and the letter columns $(\chi_T/2;1;0)$, $(-\chi_T/2;0;1)$. The linear form $\ell(x;v;w)=\sum_ex_e-\frac{2r-n}2(v-w)$ is $1$ on the unit columns and $0$ on the letter columns, and at the index it is $\kappa_{\cE}$ (Thm.~8.2(iii), p.~54): the balance wall is the first wall of this Gramma along $\ell$, with the letter columns as facet terms. By Thm.~8.10, p.~56, the period is the residue of $J_{\cB}$ there, and by Thm.~2.9, p.~15, at $p=0$ that residue is the Gramma of the facet terms. Its share side (Prop.~4.2) has one share per letter column; the pins of the two letter scales are $|\mathbf a|=s$, $|\mathbf b|=t$, and those of the marks at $\bmu=0$ are $\sum_T(a_T-b_T)\chi_T/2\in\mathbb C\one$, which with $|\mathbf a|=|\mathbf b|$ is $\sum_T(a_T-b_T)\one_T=0$. At the centre $s=t=\sigma$. The uniform shares satisfy the pins and have positive real parts.
\end{proof}

\noindent Receipts. For $\binom{[2]}1$ the fibre has dimension one and the integral is four Gammas on a line (Cor.~4.7, p.~26),
\[
\frac1{2\pi i}\int\Gamma(a)^2\,\Gamma(\sigma-a)^2\,da=\Gst41{\sigma,\ \sigma,\ \sigma,\ \sigma}{2\sigma},\qquad\text{so}\qquad \Pc\sigma=\Gst21{\sigma,\ \sigma}{2\sigma}\quad\derived.
\] For $\binom{[3]}1$ at $\sigma=1$ the two-dimensional fibre integral is $2.343907238689463\ldots$, the value of the receipts table below, to fourteen digits \computed.

\begin{proposition}[The banana's pin]\label{a:p:bananapin}\derived; receipt \computed\
For $\binom{\cE}1$ the pin is $a_e=b_e$ at every mark, so each mark carries $\Gamma(a_e)^2$, which is the share kernel of the Gramma $\sG_{0;(1,-1)}(1,y)=\int_0^\infty x^{-1}e^{-x-y/x}dx$ (Prop.~4.2), and
\[
\Gamma(\sigma)^2\,\Pc\sigma=\int_{|\mathbf a|=\sigma}\prod_{e\in\cE}\Gamma(a_e)^2=\int_0^\infty y^{\sigma-1}\,\sG_{0;(1,-1)}(1,y)^n\,dy .
\]
At $n=5$, $\sigma=1$ the right side is $39.945588011996245699\ldots$ \computed.
\end{proposition}

\begin{proof}
With $T=\{e\}$ the pins are $a_e-b_e=0$ for $e\neq e_0$ and the two totals, so $a=b$ and the fibre is the share integral $|\mathbf a|=\sigma$ (Def.~4.1). The share side of $\sG_{0;(1,-1)}(1,y)$ is $\frac1{2\pi i}\int\Gamma(c)^2y^{-c}dc$, and the read $\int_0^\infty y^{\sigma-1}(\cdot)\,dy$ of a product of $n$ such fibre integrals pins the sum of their shares at $\sigma$.
\end{proof}

So the ``Bessel moments'' are the banana's share side: the pin of a single mark turns the two letters into one Gramma per mark (alias $\sG_{0;(1,-1)}(1,y)=2K_0(2\sqrt y)$). A general list keeps a share per basis and per letter and couples them only through the incidence; that coupling is the broader arena.

\begin{proposition}[The corner is the residue at zero shares]\label{a:p:corner}\derived\
With $U_k=\{\pi(k),\dots,\pi(n)\}$ the tails of an order $\pi$ and $c$ the width (Def.~7.10, p.~49),
\[
\lim_{\sigma\downarrow0}\sigma^{n-1}\Pc\sigma=\Cor ,
\]
and $\sigma^{n-1}\Pc\sigma$ has no term of order $\sigma$. This is Thm.~38.5(i), p.~330, at the centre question, whose walls are $\sigma c(S)$ (Prop.~8.21(ii), p.~59); the term of order $\sigma$ there is a sum of $\log\theta$ at greedy bases, zero for a list. On the share side, $\sigma^{n+1}\Gamma(\sigma)^2\Pc\sigma$ tends to the same number: it is the residue of the fibre integral of Proposition~\ref{a:p:share} at the origin, where all $2|\cB|$ Gammas have their first pole. For $\binom{\cE}1$ the corner is $n!$, the $5!$ of \S14.17; for $\binom{[5]}2$ the tails of every order have widths $1,2,2,1$ and it is $30$.
\end{proposition}

\begin{proposition}[Closing the centre at the negative strengths]\label{a:p:negcorner}\derived\
For $k\in\mathbb N$, at $\sigma=-k$ the share side closes (Thm.~4.5, p.~26) onto the lattice of pairs $(\boldsymbol\alpha,\boldsymbol\beta)\in\mathbb N^{\cB}\times\mathbb N^{\cB}$ with $|\boldsymbol\alpha|=|\boldsymbol\beta|=k$ and $\sum_T\alpha_T\one_T=\sum_T\beta_T\one_T$, where every Gamma is at a pole, $a_T=-\alpha_T$, $b_T=-\beta_T$, with the residues $\prod_T(-1)^{\alpha_T+\beta_T}/(\alpha_T!\,\beta_T!)$ (Prop.~1.3(iv), p.~11). The continued period has a pole of order at most $n-1$ at $\sigma=-k$, and
\begin{align*}
\lim_{\sigma\to-k}(\sigma+k)^{n-1}\,\Pc\sigma&=\Cor\cdot\sum_{(\boldsymbol\alpha,\boldsymbol\beta)}\frac{(k!)^2}{\prod_T\alpha_T!\,\beta_T!}\\
&=\Cor\cdot\Clo k ,
\end{align*}
the sum over that lattice: it counts the pairs of $k$-tuples of bases with equal incidence.
\end{proposition}

\begin{proof}
On the balance wall the period is $\sum_\pi\int_{[0,1]^{n-1}}\prod_{k'=2}^nq_{k'}^{\kappa_{U_{k'}}-1}h_\pi(q)\,dq$ in the sector coordinates $\alpha_{\pi(k')}=\rho\,q_2\cdots q_{k'}$ (Thm.~8.7(ii), p.~55, and Thm.~8.10). At the centre $\kappa_{U_{k'}}=\sigma c(U_{k'})$ and $h_\pi=(k_\pi\psi_\pi)^{-\sigma}$, and $k_\pi\psi_\pi=\prod_{k'}q_{k'}^{c(U_{k'})}K(\alpha)\Psi(\alpha)/\alpha^{\cE}$ is a polynomial in $q$ with nonnegative coefficients and constant term one (Lemma~8.5, p.~55, weights one). By Lemma~8.6, p.~55, each cube integral continues with at most a simple pole in each $\kappa_{U_{k'}}$ at $-\mathbb N$, and at $\sigma=-k$ the $n-1$ poles meet only through the term $q^{kc}$, $c=(c(U_{k'}))_{k'}$, of $(k_\pi\psi_\pi)^k(k_\pi\psi_\pi)^{-(\sigma+k)}$, with coefficient $\prod_{k'}c(U_{k'})^{-1}[q^{kc}](k_\pi\psi_\pi)^k$. The map $\alpha^v\mapsto q^{(v(U_{k'}))_{k'}}$ is injective on exponents of degree zero, so $[q^{kc}](k_\pi\psi_\pi)^k=\Clo k$. Expanding each letter's $k$-th power (the share kernel of Cor.~4.4, p.~26, closed at the exponent $-k$) gives the lattice sum. On the share side the factor $(k!)^{-2}$ is the residue of $\Gamma(\sigma)^2$ at $-k$, and the common factor of the lattice points is the residue at zero shares, the corner (Proposition~\ref{a:p:corner}).
\end{proof}

\begin{corollary}[One law for the period and its closing]\label{a:c:law}\derived\
If polynomials $c_0,\dots,c_d$ satisfy
\[
\sum_{i=0}^dc_i(\sigma)\,\Pc{(\sigma+i)}=0
\]
identically, then for every $k\in\mathbb N$
\[
\sum_{i=0}^{\min(d,k)}c_i(-k)\,\Clo{k-i}=0 .
\]
So the least law of the period in its strength is at least as long as the least recurrence of its closings, and the closings obey the reflected law from $k=0$, with nothing below.
\end{corollary}

\begin{proof}
Multiply by $(\sigma+k)^{n-1}$ and let $\sigma\to-k$: the terms with $i\le k$ tend to $c_i(-k)$ times the corner times the closing at $k-i$ by Proposition~\ref{a:p:negcorner}, and the others are regular there.
\end{proof}

The laws come from the family's own relations. By the letter presentation (Prop.~3.6, p.~20) the centre period is a Breta with one bunch for the shape and one singleton bunch per letter; its letter splits and scale exchanges (Thm.~3.5, p.~19) are the letter and mark relations of the interest (Thm.~8.3(i)--(iii)), and along the centre ray they are the relations of Thm.~35.3, p.~308, whose eliminations are the laws of Cor.~35.6. The converse of Corollary~\ref{a:c:law}, that the least recurrence of the closings, read at $k=-\sigma$, is a law of the period, is not derived here; for the banana through seven marks it is not needed, the laws being derived by adding marks (Theorem~\ref{a:t:cert}). For every list below the minimal solution it gives agrees with an independent quadrature to fourteen digits, and the closings obey it from $k=0$ with nothing below \computed.

\begin{procedure}[The centre ray from the list]\label{a:p:proc}
(i) Close the share side at $\sigma=-k$: $x_k=\Clo k$, $k\le k_{\max}$, the letter $x_k$ abbreviating them in this procedure. (ii) Find their least recurrence $\sum_{j=0}^dp_j(k)x_{k+j}=0$ by exact linear algebra and check it on the surplus equations. (iii) Read it at $k=-\sigma$: $c_i(\sigma)=p_{d-i}(-\sigma-d)$, $\sum_ic_i(\sigma)\Pc{(\sigma+i)}=0$. (iv) The period is its minimal solution: $\Pc{(\sigma+1)}/\Pc\sigma\to K(\one)^{-2}$ (Prop.~35.7, p.~309), the native value of least modulus. (v) Scale it by the residue at zero shares: $c_0$ has a zero of order $n-1$ at $0$, and $[\sigma^{n-1}]c_0\cdot\Cor+\sum_{i\ge1}c_i(0)\Pc i=0$.
\end{procedure}

\noindent Receipts \computed, no quadrature in the procedure:
\begin{center}
\begin{tabular}{lll}
\toprule
list & $P_{\cB}(\one;1,1)$ & reading\\ \midrule
$\binom{[3]}1$ & $2.343907238689458890601562\ldots$ & $3L(\chi_{-3},2)$ (alias)\\
$\binom{[4]}1$ & $8.414398322117159997798167\ldots$ & $\Tx{3}{\Ilat{\frac12+\mathbb Z}}{0}=7\,\Tx{3}{\IZ}{0}$\\
$\binom{[4]}2$ & $2.343907238689458890601562\ldots$ & as $\binom{[3]}1$: Cor.~19.17 (p.~143) and duality\\
$\binom{[5]}1$ & $39.94558801199624569894402\ldots$ & $M_5(1)$ of \S14.17\\
\bottomrule
\end{tabular}
\end{center}
The four-mark values are those of the note (Ex.~6.3 and 6.2, p.~39). There they are closings of the share side at positive strength: the Gammas of a basis meet in double poles, whose Laurent data are Trailmoments of the counting input and of lattice inputs (Prop.~5.14, p.~33; Lemma~35.20, p.~314).

\subsection*{15.2\quad The masters are the signs}

A native gauge of the centre ray is a point where its direction $(\one;1,1)$ is native (Def.~35.4, p.~308): the defect $1-R_e(\alpha)-R^\ast_e(\alpha)$ vanishes at every mark, and since $R^\ast_e(\alpha)=1-R_e(\alpha^{-1})$ (Thm.~7.14, p.~50) this reads
\[
R_e(\alpha)=R_e(\alpha^{-1})\qquad(e\in\cE):
\]
inverting the gauge leaves the occupation unchanged. Its native value is $\big(K(\alpha)K(\alpha^{-1})\big)^{-1}$.

\begin{proposition}[The banana has only signs]\label{a:p:onlysigns}\derived\
The native gauges of the centre ray of $\binom{\cE}1$ are exactly the signs $\eta\in\{\pm1\}^{\cE}$, $\eta_{e_0}=1$, with $K(\eta)=\sum_e\eta_e\neq0$; their native values are $K(\eta)^{-2}$ (Thm.~35.12, p.~310).
\end{proposition}

\begin{proof}
For $K=\sum_e\alpha_e$ the occupation is $R_e(\alpha)=\alpha_e/K(\alpha)$, so $R_e(\alpha)=R_e(\alpha^{-1})$ for every $e$ says $\alpha_e^2=K(\alpha)/K(\alpha^{-1})$ for every $e$.
\end{proof}

So along the banana sequence the native values are $(n-2j)^{-2}$, $j$ the number of minus signs, $\lceil n/2\rceil$ of them, all simple. At $n=5$ they are $\frac1{25},\frac19,1$, the symbol of the ray law vanishes there (Cor.~35.6, p.~308), the order is three, and the corner leaves two masters. Every list has the signs with $K(\eta)\neq0$ among the native gauges of its centre ray; what is special about $\binom{\cE}1$ is that it has nothing else.

\begin{remark}[The lowered settings of \S14.17]\label{a:r:lowered}
With $g=\sG_{0;(1,-1)}(1,\cdot)$ and $h(y)=\sum_my^m/\Gamma(1+m)^2$, the settings $x^{(j)}(s)=\int_0^\infty y^{s-1}g^{5-j}h^j\,dy$, $j=0,1,2$, satisfy the ray law in the normalization $M(s)=\Gamma(s)^2\Pc s$, and $x^{(0)}=M$ is Proposition~\ref{a:p:bananapin}. Their corners are $120,0,0$ \computed, and $x^{(j)}(s+1)/\big(s^2x^{(j)}(s)\big)\to(5-2j)^{-2}$: they are the sign classes with $j$ minus signs. ``The $g$'s outnumber the $h$'s'' is $K(\eta)\neq0$; the tie at four marks is $K(\eta)=0$, which is not a native gauge.
\end{remark}

\begin{proposition}[The lowered-setting block is the frame]\label{a:p:frame}\derived; receipt \computed\
Let $\Phi(s)=\det\big[x^{(i)}(s+l)\big]_{i,l=0}^2$. By Prop.~35.8, p.~309, applied to the ray law of $M$ (end coefficients $s^6$ and $-225$), $\Phi(s+1)=\frac{s^6}{225}\Phi(s)$ and $\Phi(s)=\Phi(1)\,\Gamma(s)^6\,225^{1-s}$. The amplitudes at the three sign classes, $x^{(j)}(s)\sim2\pi^{(5-2j)/2}2^{-j}\rho_j^{5/2}\,\Gamma(2s-\frac52)\,\rho_j^{-2s}$ with $\rho_j=2(5-2j)$, give
\[
\Phi(1)=\frac{128\sqrt{15}}{16875}\,\pi^3=0.910880564798203032890397187975\ldots,
\]
which the settings computed directly (quadrature, as a check) reproduce to thirty digits. The corners reduce the third column, so
\[
\Phi(1)=\frac{120}{225}\Big(x^{(1)}(1)\,x^{(2)}(2)-x^{(1)}(2)\,x^{(2)}(1)\Big),\qquad x^{(1)}(1)x^{(2)}(2)-x^{(1)}(2)x^{(2)}(1)=\frac{16\sqrt{15}}{1125}\,\Gamma(\tfrac12)^6 :
\]
the block of \S14.17 is the frame of the ray over its corner. The masters do not enter $\Phi$; their coefficients are minors that the zero corners of the lowered settings kill. So the frame does not shorten the last minor of \S14.17.
\end{proposition}

\subsection*{15.3\quad A parallel class closes through the banana}

\begin{proposition}\label{a:p:parallel}\derived\
Let $\cB$ be obtained from a list $\cH$ on the marks $1,\dots,m$ by replacing mark $i$ with a parallel class $P_i$ of $s_i$ marks, so that $K_{\cB}(\alpha)=K_{\cH}\big((\alpha(P_i))_i\big)$ with $\alpha(P_i)=\sum_{e\in P_i}\alpha_e$. Then $[\alpha^v]K_{\cB}^k=\big([A^w]K_{\cH}^k\big)\prod_i\big(w_i!/\prod_{e\in P_i}v_e!\big)$ with $w_i=v(P_i)$, and the closings of $\cB$ are
\[
\sum_w\big([A^w]K_{\cH}(A)^k\big)^2\prod_i[\alpha^0]\Big(K_{P_i}(\alpha)\Psi_{P_i}(\alpha)/\alpha^{P_i}\Big)^{w_i},
\]
with $K_{P_i}=\sum_{e\in P_i}\alpha_e$ and $\Psi_{P_i}$ the letters of the banana $\binom{P_i}1$: its closings at $-w_i$.
\end{proposition}

The banana enters the closing only through its own closings, as the parallel node enters the network through its full leans (Cor.~25.9, p.~209).

\subsection*{15.4\quad The five-mark centres}

All connected lists on five marks, up to duality, which leaves the centre unchanged (Prop.~8.14, p.~57); ``doubled'' and ``tripled'' are parallel classes. \computed
\begin{center}
\small
\begin{tabular}{@{}l c l l c c c@{}}
\toprule
list & $|\cB|$ & $K(\eta)^2$ at the signs (mult.) & roots of the symbol & order & corner & masters\\ \midrule
$\binom{[5]}1$ & 5 & $25\,(1),\,9\,(5),\,1\,(10)$ & $25,9,1$ & 3 & 120 & 2\\
$\binom{[4]}2$, a mark doubled & 9 & $81\,(1),\,9\,(6),\,1\,(9)$ & $81,9,1$ & 3 & 36 & 2\\
$\binom{[5]}2$ & 10 & $100\,(1),\,4\,(15)$ & $100,4,1,-8$ & 4 & 30 & 3\\
$\binom{[3]}2$, two marks doubled & 8 & $64\,(1),\,16\,(2),\,4\,(8),\,0\,(5)$ & $64,16,4,4,1$ & 5 & 44 & $3^\dagger$\\
$\binom{[3]}2$, a mark tripled & 7 & $49\,(1),\,25\,(1),\,9\,(3),\,1\,(11)$ & $49,25,9,1,1,1$ & 6 & 60 & $4^\dagger$\\
\bottomrule
\end{tabular}

\medskip
\begin{tabular}{@{}l l@{}}
\toprule
list & $P_{\cB}(\one;1,1)$\\ \midrule
$\binom{[5]}1$ & $39.94558801199624569894402849842399589184978292477\ldots$\\
$\binom{[4]}2$, a mark doubled & $5.4939011435808440124442576078206911958536386406801\ldots$\\
$\binom{[5]}2$ & $4.0297141707110496073206957344284069331398675636396\ldots$\\
$\binom{[3]}2$, two marks doubled & $7.8658373683640215924219746974461682109823828094581\ldots$\\
$\binom{[3]}2$, a mark tripled & $12.848722141915555957312339931253534636619123934709\ldots$\\
\bottomrule
\end{tabular}
\end{center}
The roots of the symbol are those of the recurrence of the closings, the values $K(\alpha)K(\alpha^{-1})$ at the native gauges; the ray law of the period has their reciprocals. ``Masters'' is $m-1$ for the least $m$ with a rational relation among $1,P_{\cB}(\one;1,1),\dots,\Pc m$ (PSLQ at $120$ digits, coefficients below $10^{15}$, values good to $150$ digits): over $\mathbb Q$, for the symmetric insertion, on the centre ray. In the first three rows the one relation is the corner. Over $\mathbb Q+\mathbb Q\,\Pth$ the list $\binom{[5]}2$ has two masters, and over $\mathbb Q+\mathbb Q\,\Pth+\mathbb Q\,\Gamma(\frac12)^4$ the triangle with two marks doubled has one (Theorems~\ref{a:t:u25masters} and~\ref{a:t:r4masters}).

$^\dagger$\,In the last two rows $c_0$ and $c_1$ both carry $(\sigma+1)^2$, so the ray law at $\sigma=-1$ reaches the period at neither $-1$ nor $0$, and it is a second relation, a \emph{second corner}: the boundary term, the constant term at $\sigma=-1$ (Thm.~1.6, p.~12),
\[
\CT_{\sigma=-1}\Big(c_0(\sigma)\Pc\sigma+c_1(\sigma)\Pc{(\sigma+1)}\Big)=-\sum_{i\ge2}c_i(-1)\,\Pc{(i-1)},
\]
is rational \computed. In lowest terms the two relations read
\begin{align*}
&2503-566\,P_{\cB}(\one;1,1)+88248\,\Pc2-908872\,\Pc3+1425600\,\Pc4=0,\\
&13800002-1848045\,P_{\cB}(\one;1,1)+227950815\,\Pc2-3601110627\,\Pc3\\
&\qquad\qquad+11177654265\,\Pc4-8565102000\,\Pc5=0,
\end{align*}
so the masters are $5-2=3$ and $6-2=4$. Write $c_0=(\sigma+1)^2A$ and $c_1=(\sigma+1)^2B$. Proposition~\ref{a:p:negcorner} at $\sigma=-1$ requires $A(-1)\,|\cB|+B(-1)=0$, the closing at $1$ being $|\cB|$; it holds, $B(-1)/A(-1)=-8$ and $-7$ \computed. There the recursion of Corollary~\ref{a:c:law} leaves the corner at $-1$ free, so this is a check of Proposition~\ref{a:p:negcorner} that the recursion does not supply.

$\binom{[5]}2$ has native gauges off the signs: on the pattern $\alpha\propto(a,a,b,b,1)$ the native equations have solutions with $K(\alpha)K(\alpha^{-1})=1$ and $-8$. There is no rational relation among $1$, the two masters of $\binom{[5]}1$ and the three of $\binom{[5]}2$ (coefficients below $10^8$, 150 digits), nor with $\pi^2$ adjoined.

\subsection*{15.5\quad The native shares, the roots, and an exact pairing}

Every root of the symbol of the closings' recurrence is the reciprocal native value $K(\alpha_c)K(\alpha_c^{-1})$ of a native gauge $\alpha_c$ of the centre ray (Cor.~35.6, p.~308).

\begin{proposition}[Both reads sit at the native shares]\label{a:p:edge}\derived; receipts \computed\
Let the list be connected, and let $H=2\big(\kappa_{\one}(X_e,X_{e'})\big)_{e,e'\neq e_0}$, twice the covariances of the indicators under the weight law (Def.~9.2, p.~61), the Hessian of Thm.~35.11, p.~310.
\begin{enumerate}[label=(\roman*)]
\item On the real fibre of Proposition~\ref{a:p:share} the native shares of Thm.~4.3(ii), p.~25, are the uniform shares $a_T=b_T=\sigma/|\cB|$: every basis takes the same share from each letter. This is the table's gauge $\one$ read on the share side.
\item The period reads the fibre along its imaginary directions. By Thm.~4.3(iii), which at the period is Thm.~38.5(ii), p.~330,
\[
\Pc k\sim K(\one)^{-2k}\Big(\frac{2\pi}k\Big)^{\frac{n-1}2}(\det H)^{-\frac12}.
\]
\item The closing reads the same fibre on its integer points. Dropping the pin, the multinomial theorem gives $\sum_{|\boldsymbol\alpha|=|\boldsymbol\beta|=k}(k!)^2/\prod_T\alpha_T!\beta_T!=K(\one)^k\Psi(\one)^k=|\cB|^{2k}$; the pin keeps the Gaussian of the fibre at the native shares scaled by $k$, and
\[
\Clo k\sim K(\one)^{2k}(2\pi k)^{-\frac{n-1}2}(\det H)^{-\frac12}.
\]
\item So $k^{n-1}\,\Pc k\,\Clo k\to(\det H)^{-1}$: one Gaussian, read once along a contour and once on a lattice, so that the powers of $2\pi$ cancel and the determinant appears twice.
\end{enumerate}
\end{proposition}

\begin{proof}
(i) The uniform point satisfies the pins, and $\nabla S=-\log\mathbf w$ is constant there, a combination of the two rows of the letter totals; $S$ is strictly concave on the real fibre (Thm.~4.3(ii)). (ii) Thm.~38.5(ii) at the centre: $X(0)=-2\log K(\one)$ and $-\nabla_z^2X(0)=8\kappa_{\one}=4H$ on the chart, $y=2z$; Thm.~35.11 is the same edge read on the generating function of the ray. (iii) The terms $(k!)^2/\prod\alpha_T!\beta_T!$ on the pinned lattice peak at the uniform shares scaled by $k$; their Gaussian there is the quadratic form of $\sum\log\Gamma$ at the native shares restricted to the fibre, the same form as in (ii), now summed over the integer points of the fibre; $\nabla^2_zE(0)=8\kappa_{\one}$ (Thm.~7.14, p.~50) identifies it with $H$.
\end{proof}

\noindent Receipts \computed: $(\det H)^{-1}=\frac{3125}{16},\frac{729}{16},\frac{3125}{81},\frac{512}9,\frac{16807}{192}$ in the order of the table of \S15.4. At $k=70$ the product agrees with these to about $10^{-4}$, the error falling like $k^{-2}$. For $\binom{[2]}1$ the product is exact at every $k$: $k\,\Gst21{k,\ k}{2k}\binom{2k}k=2=(\det H)^{-1}$ \derived.

\begin{remark}[The signs are the pin read modulo two]\label{a:auto:15.12}
Averaging over the signs detects the incidence modulo two: for every $k$,
\[
2^{-n}\sum_{\eta\in\{\pm1\}^{\cE}}K(\eta)^{2k}=\sum_{\substack{|\boldsymbol\alpha|=|\boldsymbol\beta|=k\\ \sum_T\alpha_T\one_T\equiv\sum_T\beta_T\one_T\ (2)}}\frac{(k!)^2}{\prod_T\alpha_T!\,\beta_T!}\ \ \ge\ \Clo k\quad\derived,
\]
since $K(\eta)\Psi(\eta)/\eta^{\cE}=K(\eta)^2$. So the sign values $K(\eta)^2$ are the growth rates of the pin's reduction modulo two, with $K(\one)^2$ the largest. The full pin keeps them as roots (Thm.~35.12) and adds the native gauges off the signs, below.
\end{remark}

\begin{proposition}[An exact pairing of the period with its closing for the five-mark banana]\label{a:p:pairing}\derived\ from the ray law; receipt \computed.\ For $\binom{[5]}1$ the ray law $\sum_{i=0}^3c_i(\sigma)\Pc{(\sigma+i)}=0$ has
\begin{gather*}
c_0=-\sigma^4,\qquad c_1=35\sigma^4+70\sigma^3+63\sigma^2+28\sigma+5,\\
c_2=-(\sigma+1)^2(259\sigma^2+518\sigma+285),\qquad c_3=225(\sigma+1)^2(\sigma+2)^2,
\end{gather*}
and it is self-adjoint under the reflection of the strength: $c_i(-\sigma-i)=c_i(\sigma)$ for every $i$. So the closings, read as a function of $\sigma$, solve the adjoint law, and for every $\sigma\in\{1,2,3,\dots\}$
\[
\sum_{i=1}^3\sum_{j=1}^ic_i(\sigma-j)\ \Clo{\sigma-j}\ \Pc{(\sigma-j+i)}=120,
\]
the residue at zero shares, with the closings at negative indices read as $0$.
\end{proposition}

\begin{proof}
The left side, $B(\sigma)$ in this proof, is the bilinear concomitant of the law and its adjoint: $B(\sigma+1)-B(\sigma)$ is the law applied to the period times the closing at $\sigma$, less the period at $\sigma$ times the adjoint law applied to the closings. The first vanishes by the law. The adjoint law at $\sigma$ is $\sum_ic_i(\sigma-i)x_{\sigma-i}$, which by $c_i(\sigma-i)=c_i(-\sigma)$ is the recurrence of Corollary~\ref{a:c:law} at $k=\sigma$, zero from $k=0$ on. At $\sigma=1$ only $j=1$ survives, and $B(1)=\sum_ic_i(0)\Pc i=120$ is the corner relation of Procedure~\ref{a:p:proc}(v). The self-adjointness is an identity of the four polynomials; $B(\sigma)=120$ holds to seventy digits for $\sigma\le55$ \computed.
\end{proof}

At large $\sigma$ the closings and the period move by the factors $25$ and $\frac1{25}$ per step, so $B(\sigma)\to\lambda\chi'(\lambda)\cdot\lim_kk^4\,\Pc k\,\Clo k$, with $\chi(\lambda)=\sum_i[\sigma^4]c_i\,\lambda^i=-1+35\lambda-259\lambda^2+225\lambda^3$ the symbol and $\lambda=\frac1{25}$. With Proposition~\ref{a:p:edge}:
\[
120\cdot\det H=\lambda\chi'(\lambda)\big|_{\lambda=1/25}=\frac{384}{625},\qquad \det H=\frac{16}{3125}\quad\derived.
\]
The residue at zero shares (the low end of strength, Thm.~38.5(i)) and the Gaussian at the native shares (the high end, Thm.~38.5(ii)) are the two readings of one exact identity between the fibre and its lattice. For the other four five-mark lists the closings solve no adjoint of the ray law in the forms tried: no combination $\sum_{l<d}r_l(\sigma)\Clo{\sigma+l}$ with polynomial $r_l$ of degree at most $20$, nor with the apparent factors of $c_0$ as denominators up to degree $14$, solves it \computed. In that framework the pairing of the period with its closing is not a single product there; it must carry other solution classes.

\begin{proposition}[Signs through parallel classes]\label{a:p:signsparallel}\derived\
In the setting of Proposition~\ref{a:p:parallel}, $K_{\cB}(\eta)=K_{\cH}\big(s_1(\eta),\dots,s_m(\eta)\big)$ with $s_i(\eta)=\sum_{e\in P_i}\eta_e\in\{s_i,s_i-2,\dots,-s_i\}$: the sign values of the bananas $\binom{P_i}1$ feed the letter of $\cH$.
\end{proposition}

For the triangle $\cH=\binom{[3]}2$, $K_{\cH}=e_2$, and with one mark tripled, $s=s_1(\eta)\in\{3,1,-1,-3\}$:
\[
e_2(s,1,1)=2s+1\in\{7,3,-1,-5\},\qquad e_2(s,1,-1)=-1\quad\text{for every }s,
\]
which is the census $49\,(1),\,25\,(1),\,9\,(3),\,1\,(11)$ of \S15.4. The second form does not see $s$.

\begin{proposition}[Letters that forget a class]\label{a:p:families}\derived\
Write $K_{\cB}=A_i\,K_{\cH/i}(A)+K_{\cH\setminus i}(A)$ with $A_i=\alpha(P_i)$ (Def.~7.3, p.~48), and $A_i'=\sum_{e\in P_i}\alpha_e^{-1}$.
\begin{enumerate}[label=(\roman*)]
\item Tripled triangle. On $\alpha_3=-\alpha_2$ the contraction letter $K_{\cH/1}=A_2+A_3$ vanishes, so every copy has occupation zero at $\alpha$ and at $\alpha^{-1}$, and $K(\alpha)=-\alpha_2^2$, $K(\alpha^{-1})=-\alpha_2^{-2}$ whatever the copies are. The native gauges there are
\[
\big\{\alpha_3=-\alpha_2,\ A_1=\alpha_2^2A_1'\big\},
\]
a family of dimension two modulo scale, all of native value $1$, containing the eight sign classes with $\eta_3=-\eta_2$.
\item Two marks doubled. On $P_2=(t,-t)$ the class letter vanishes, $A_2=A_2'=0$, and $K=A_1\alpha_3$ is the letter of $\cH\setminus2$, in which both remaining classes are coloops. The native gauges there are
\[
\Big\{P_1=(c,c),\ P_2=(t,-t),\ t^2=\frac{c\,\alpha_3\,(2\alpha_3+c)}{2c+\alpha_3}\Big\},
\]
a curve of native value $\frac14$ which on the torus $|\alpha_e|=1$ is a circle through the signs with $(s_1,s_2)=(\pm2,0)$.
\end{enumerate}
\end{proposition}

\begin{proof}
Nativity is $R_e(\alpha)=R_e(\alpha^{-1})$ at every mark. (i) At a copy $x\in P_1$, $R_x(\alpha)=\alpha_x(A_2+A_3)/K(\alpha)$ and $R_x(\alpha^{-1})=\alpha_x^{-1}(A_2'+A_3')/K(\alpha^{-1})$, both zero on $\alpha_3=-\alpha_2$; at $\alpha_2$ and $\alpha_3$ the two sides differ by $\mp(A_1/\alpha_2-\alpha_2A_1')$. (ii) At $\alpha_3$ both occupations are $1$ on $A_2=A_2'=0$; at $P_1$ the equation holds when $\alpha_{1a}^2=\alpha_{1b}^2=A_1/A_1'$, and $\alpha_{1b}=-\alpha_{1a}$ gives $K=0$; at $P_2$ it holds exactly at the stated $t^2$. When $|c|=|\alpha_3|=1$, $|2\alpha_3+c|=|2c+\alpha_3|$, so $|t|=1$.
\end{proof}

\begin{remark}[Local exponents]\label{a:auto:15.16}
Put $x_k\approx\lambda^kk^\beta$ into the recurrence of the closings \computed. At every simple root of every row, $\beta=-2=-\frac{n-1}2$, the exponent of a nondegenerate native gauge (Thm.~35.5(i), p.~308). At the repeated roots the indicial polynomials are
\[
(\beta+1)(\beta+2)(\beta+3)\quad(\text{tripled triangle},\ \lambda=1),\qquad(2\beta+3)(2\beta+5)\quad(\text{two marks doubled},\ \lambda=4).
\]
A family of real dimension $d$ on the torus $|\alpha_e|=1$, transversally nondegenerate, gives $\beta=-\frac{n-1-d}2$. The Hessian of $\log\big(K(\alpha)K(\alpha^{-1})\big)$ there has rank two along the family of (i) and rank three along the circle of (ii) \computed, so these are $-1$ and $-\frac32$. The isolated sign orbit at $\lambda=1$ (one copy plus, two minus, $\eta_2=\eta_3$) has full rank and gives $-2$. The exponents $-3$ and $-\frac52$: \notyet. These families are integral directions in the exceptional set of Thm.~35.19(i), p.~313: the order of the ray law counts native values with these multiplicities, and the families raise the order, and with it the masters, past the number of native values.
\end{remark}

\begin{remark}[Native gauges at complex points off the signs \computed]\label{a:auto:15.17}
$\binom{[5]}2$: $K(\alpha)K(\alpha^{-1})=1$ at $(\omega^2,\omega^2,\omega,\omega,1)$, $\omega=e^{2\pi\mathrm i/3}$, where $K=K(\alpha^{-1})=1$; and $-8$ at $(1,1,b,b,1)$, $b=-(2+\sqrt3)$, and at its inverse, where $K=-2-2\sqrt3$ and $K(\alpha^{-1})=-2+2\sqrt3$. Two marks doubled: $1$ at $(\omega,\omega^2,\omega,\omega^2,1)$. The signs are the points of order two of the torus; points of order three are the next layer that occurs here.
\end{remark}

\subsection*{15.6\quad Closing the banana at strength one}

\begin{proposition}[Four marks close on the half-lattice]\label{a:p:fourclose}\derived\
For $\binom{[4]}1$, merging the marks in pairs by four Gammas on a line (Cor.~4.7, p.~26), $\int_{a_1+a_2=c}\Gamma(a_1)^2\Gamma(a_2)^2=\Gst41{c,\ c,\ c,\ c}{2c}$, leaves a fibre of dimension one, and
\begin{align*}
P_{\cB}(\one;1,1)&=\frac1{2\pi i}\int_{(1/2)}\Gst41{c,\ c,\ c,\ c}{2c}\Gst41{1-c,\ 1-c,\ 1-c,\ 1-c}{2-2c}dc\\
&=\Tx{3}{\Ilat{\frac12+\mathbb Z}}{0}=7\,\Tx{3}{\IZ}{0}.
\end{align*}
\end{proposition}

\begin{proof}
By Lemma~5.16, p.~34, and the shift and reflection of the stacked Gamma (Prop.~1.3(iii),(v), p.~11) the integrand is
\[
\Gst40{c,\ c,\ 1-c,\ 1-c}{}^2\,\frac{\sin2\pi c}{\pi(1-2c)}=\frac{2\pi^3\cos\pi c}{(1-2c)\sin^3\pi c}.
\]
On the lines $\re c=N+\frac12$ it is $O\big(N^{-1}e^{-2\pi|\im c|}\big)$, so it closes to the right (Thm.~4.5, p.~26). At $c=m\ge1$, $2\pi^3\cos\pi c/\sin^3\pi c=2(c-m)^{-3}+O(c-m)$, with no parts of orders one and two, as in Lemma~5.16; the half-integers are zeros. So the residue is $2\,[\varepsilon^2](1-2m-2\varepsilon)^{-1}=-8/(2m-1)^3$, and the closing is $\sum_{m\ge1}8(2m-1)^{-3}=\sum_{x\in\frac12+\mathbb Z,\,x>0}x^{-3}$. The last equality is the coset relation at the step $2$ (Thm.~5.12(i), p.~32). The half-lattice is the doubled share $2c$ of a merged pair, the lower entry of its stacked Gamma.
\end{proof}

\noindent This is the note's value (Ex.~6.2, p.~39), now read as a closing of the share side.

\begin{proposition}[A mark added is a critical-line kernel]\label{a:p:addmark}\derived; receipts \computed\
For every $n\ge1$,
\[
P_{\binom{[n+1]}1}(\one;1,1)=-\sum_{m\ge1}\partial_\sigma P_{\binom{[n]}1}(\sigma\one;\sigma,\sigma)\big|_{\sigma=m}.
\]
\end{proposition}

\begin{proof}
In the share side of $\binom{[n+1]}1$ (Proposition~\ref{a:p:bananapin}) merge the first $n$ marks: their fibre at the total $d$ is $\Gamma(d)^2P_{\binom{[n]}1}(d\one;d,d)$, by the same proposition at $n$ marks. The last mark carries $\Gamma(1-d)^2$, and $\Gamma(d)^2\Gamma(1-d)^2=\Gst40{d,\ d,\ 1-d,\ 1-d}{}$, whose right closing reads the derivative of the rest at the positive integers (Lemma~5.16(i)). The rest is holomorphic on $\re d>0$, bounded on each vertical line by its value at the real part, and decays as $\re d\to\infty$, so the closing applies (Thm.~4.5).
\end{proof}

\noindent Receipts \computed. At $n=2$, with $\partial_\sigma\Gst21{\sigma,\ \sigma}{2\sigma}=2\Gst21{\sigma,\ \sigma}{2\sigma}\big((\log\Gamma)'(\sigma)-(\log\Gamma)'(2\sigma)\big)$ and $(\log\Gamma)'(j+1)=-\Tx{1}{\IZ}{j}$ (Prop.~5.14, p.~33),
\[
P_{\binom{[3]}1}(\one;1,1)=2\sum_{m\ge1}\Gst21{m,\ m}{2m}\,(\sT_{1;(1)}[x]\,\IZ)((m-1,2m-1])=2.343907238689458890601562\ldots,
\]
the period of two marks times a head of the counting input, to $25$ digits. At $n=3$ the sum is $7\Tx{3}{\IZ}{0}$ and at $n=4$ it is $39.945588011996245699\ldots$, both to $20$ digits.

\begin{remark}[The five-mark centre is a Trailmoment of a live input]\label{a:r:live}
By Proposition~\ref{a:p:addmark},
\[
P_{\binom{[5]}1}(\one;1,1)=\big(\sT_{0;(1)}[x]\,f\,\IZ\big)((0,\infty)),\qquad f(m)=-\partial_\sigma P_{\binom{[4]}1}(\sigma\one;\sigma,\sigma)\big|_{\sigma=m},
\]
a live input (Chapter~6): its density is read from the four-mark centre period, and the four-mark ray law differentiated in $\sigma$ gives it a recurrence of order two whose inhomogeneity is the values. At $\sigma=0$ the differentiated law is a relation among the first two slopes, the corner and the values, the term of order $\sigma$ of Proposition~\ref{a:p:corner} being zero. So every $f(m)$ lies in $\mathbb Q+\mathbb Q\Tx{3}{\IZ}{0}+\mathbb Q f(1)$ \derived. For example $P_{\binom{[4]}1}(2\,\one;2,2)=(7\zeta(3)-6)/16$ and $f(2)=-\big(36+21\zeta(3)-2f(1)\big)/32$ (alias $\zeta(3)=\Tx{3}{\IZ}{0}$) \computed. The first slope is
\[
-f(1)=48\,\mathrm{Li}_4(\tfrac12)+2\log^42-2\pi^2\log^22-\tfrac{17}{30}\pi^4=-39.38158212312216562493980934508\ldots
\]
by PSLQ at $60$ digits \computed. In the family these aliases are Trailmoments of index at most four of the counting input at the frequencies $0$ and $\frac12$ and of the lattice input of density $2^{-x}$ (Def.~5.2, p.~28). So the four-mark slope sits on the half-lattice of Proposition~\ref{a:p:fourclose}, at index four. The masters of the five-mark centre are values of this live Trailmoment, and this is where the closing stops landing on plain lattice inputs.
\end{remark}

\subsection*{15.7\quad Closing the live input of five marks}

Throughout this subsection $p(\sigma)=P_{\binom{[4]}1}(\sigma\one;\sigma,\sigma)$ for real $\sigma>0$ and $f=-p'$; $p\,\IZ$ and $f\,\IZ$ are the \emph{values input} and the \emph{slope input} (Def.~5.4, p.~29). The four-mark ray law is
\[
\sum_{i=0}^2c_i(\sigma)\,p(\sigma+i)=0,\qquad c_0=\sigma^3,\quad c_1=-2(2\sigma+1)(5\sigma^2+5\sigma+2),\quad c_2=64(\sigma+1)^3,
\]
read off the closings by Procedure~\ref{a:p:proc} and derived in Theorem~\ref{a:t:cert} below. For a function $R$ on a coset of $\mathbb Z$ put $(L^\ast R)(x)=\sum_ic_i(x-i)\,R(x-i)$, and let $L'^\ast$ be the same with the derivatives $c_i'$ in place of the $c_i$.

\begin{proposition}[The live recurrence]\label{a:p:liverec}\derived\
The slope input has a recurrence of order two whose inhomogeneity is the values input:
\[
\sum_{i=0}^2c_i(\sigma)\,f(\sigma+i)=\sum_{i=0}^2c_i'(\sigma)\,p(\sigma+i),\qquad c_0'=3\sigma^2,\quad c_1'=-6(10\sigma^2+10\sigma+3),\quad c_2'=192(\sigma+1)^2 .
\]
\end{proposition}

\begin{proof}
The law holds identically in $\sigma$ (Theorem~\ref{a:t:cert}); differentiate it.
\end{proof}

\begin{proposition}[A mark added, at every strength]\label{a:p:cosetmark}\derived; receipts \computed\
For real $s>0$ let $W_s(d)=\Gst22{d,\ d}{1+d-s,\ 1+d-s}$. Then
\begin{align*}
\Gamma(s)^2\,P_{\binom{[5]}1}(s\one;s,s)&=-\big(\sT_{0;(1)}[x]\,(W_sp)'\,\Ilat{s+\mathbb Z}\big)((s-1,\infty)),\\
x^{(1)}(s)&=\phantom{-}\big(\sT_{0;(1)}[x]\,W_sp\,\Ilat{s+\mathbb Z}\big)((s-1,\infty)),
\end{align*}
with $x^{(1)}$ the lowered setting with one $h$ of Remark~\ref{a:r:lowered}. The period reads the slopes of the word $W_sp$ on the coset of its strength, and the lowered setting reads its values. At $s\in\{1,2,\dots\}$ the word $W_s(x)=\prod_{j=1}^{s-1}(x-j)^2$ is a polynomial with double zeros at $1,\dots,s-1$, the cut moves to $0$, and
\[
\Gamma(s)^2P_{\binom{[5]}1}(s\one;s,s)=\Tx{0}{(W_sf-W_s'p)\,\IZ}{0},\qquad x^{(1)}(s)=\Tx{0}{W_sp\,\IZ}{0}.
\]
\end{proposition}

\begin{proof}
Merge the first four marks in the share side of five (Proposition~\ref{a:p:bananapin}): $\Gamma(s)^2P_{\binom{[5]}1}(s\one;s,s)=\frac1{2\pi i}\int_{(c)}\Gamma(d)^2p(d)\,\Gamma(s-d)^2\,dd$ with $0<c<s$, as in the proof of Proposition~\ref{a:p:addmark}. By the reflection (Prop.~1.3(v), p.~11), $\Gamma(s-d)\Gamma(1-s+d)=\pi/\sin\pi(s-d)$, so
\[
\Gamma(d)^2\,\Gamma(s-d)^2=W_s(d)\,\Gst40{s-d,\ s-d,\ 1-s+d,\ 1-s+d}{},
\]
the stacked Gamma of Lemma~5.16, p.~34, at $a=s-d$: double poles at $d\in s+\mathbb Z$, with no simple parts. The double zeros of $W_s$ at $d=s-1,s-2,\dots$ remove the poles between the contour and $s$. The rest is as in Proposition~\ref{a:p:addmark}: $W_sp$ is holomorphic on $\re d>0$, of polynomial growth along vertical lines, where the kernel decays like $e^{-2\pi|\im d|}$, and it decays like $16^{-\re d}$ to the right; closing to the right (Thm.~4.5, p.~26) reads $-(W_sp)'$ at $d=s,s+1,\dots$. For $x^{(1)}$, expand $h=\sum_ky^k/\Gamma(1+k)^2$ against $g^4$ (Proposition~\ref{a:p:bananapin} at four marks): $x^{(1)}(s)=\sum_{k\ge0}\Gamma(s+k)^2p(s+k)/\Gamma(1+k)^2$, and $\Gamma(s+k)^2/\Gamma(1+k)^2=W_s(s+k)$.
\end{proof}

\noindent Receipts \computed: at $s=\frac32$ both sides of the first identity are $2.27912905076377$, and at $s=\frac94$ both are $0.129204079141924$ (the left side by one-dimensional quadrature of the share side; fifteen digits).

\begin{proposition}[Summation by parts at the cut: what lands on lattice inputs]\label{a:p:sbp}\derived\
For every polynomial $R$,
\begin{gather*}
\Tx{0}{(L^\ast R)\,p\,\IZ}{0}=-24\,R(0),\\
\Tx{0}{(L^\ast R)\,f\,\IZ}{0}=\Tx{0}{(L'^\ast R)\,p\,\IZ}{0}.
\end{gather*}
\end{proposition}

\begin{proof}
By the shifted inputs (Thm.~6.2, p.~42) a shift of the input is a shift of the word, and the cut relation (Prop.~5.3, p.~29) returns the points that cross the cut: for a density $\varphi$,
\[
\sum_{x\ge1}R(x)\sum_ic_i(x)\,\varphi(x+i)=\Tx{0}{(L^\ast R)\,\varphi\,\IZ}{0}-R(0)\sum_{i\ge1}c_i(0)\,\varphi(i)-R(-1)\,c_2(-1)\,\varphi(1),
\]
and $c_2(-1)=c_2'(-1)=0$. For $\varphi=p$ the left side vanishes by the law. For $\varphi=f$ it is $\sum_xR(x)\sum_ic_i'(x)\,p(x+i)$ by Proposition~\ref{a:p:liverec}, to which the same move applies. So the boundaries are
\[
\sum_{i\ge1}c_i(0)\,p(i)=-4p(1)+64p(2),\qquad \sum_{i\ge1}\big(c_i(0)f(i)-c_i'(0)p(i)\big)=-4f(1)+64f(2)+18p(1)-192p(2).
\]
The first is $-[\sigma^3]c_0$ times the corner $24$ (Procedure~\ref{a:p:proc}(v)); the second is the term of order $\sigma$ of $\sigma^3p(\sigma)$, which vanishes (Proposition~\ref{a:p:corner}). With the values of Remark~\ref{a:r:live} they are $-24$ and $0$: the lattice Trailmoments in them, $p(1)=7\Tx{3}{\IZ}{0}$, $f(1)$ and $f(2)$, cancel identically.
\end{proof}

So what the closing of the live Trailmoment puts on lattice inputs is the residue at zero shares, and nothing else: neither $\Tx{3}{\IZ}{0}$ nor the four-mark slope $f(1)$ survives.

\begin{proposition}[The fixed block]\label{a:p:block}\derived; receipts \computed\
$L^\ast$ raises the degree of a polynomial by three, with leading factor $1-20+64=45$, the symbol of the law at $\lambda=1$, which is not a native value of four marks; so the polynomials are $L^\ast(\mathbb Q[x])\oplus\mathrm{span}\{1,\,x-1,\,(x-1)^2\}$. By Proposition~\ref{a:p:sbp}, every Trailmoment at index zero of a polynomial word against the values input or the slope input, cut at $0$, is a rational number plus a rational combination of the six
\[
\Tx{0}{(x-1)^j\,p\,\IZ}{0},\qquad \Tx{0}{(x-1)^j\,f\,\IZ}{0}\qquad(j=0,1,2).
\]
The share side adds one relation:
\[
5\,\Tx{0}{p\,\IZ}{0}-5\,\Tx{0}{(x-1)\,f\,\IZ}{0}=\Tx{0}{f\,\IZ}{0}.
\]
\end{proposition}

\begin{proof}
The splitting: $\deg L^\ast R=\deg R+3$, so $L^\ast$ is injective on polynomials and its image misses exactly three dimensions in each degree. The relation says that each mark takes on average a fifth of the strength. On the fibre of Proposition~\ref{a:p:bananapin} the shares sum to $\sigma$, and the integrand and the pins are symmetric in the five marks; so the fibre integral with one share $a_e$ inserted is the same for every $e$, and the five sum to $\sigma\,\Gamma(\sigma)^2P_{\binom{[5]}1}(\sigma\one;\sigma,\sigma)$. Merging the first four marks, the fifth share is $\sigma-d$, and at $\sigma=1$ the closing of Proposition~\ref{a:p:cosetmark} reads its insertion as $-\sum_{x\ge1}\partial_d\big((1-d)p(d)\big)\big|_{d=x}=\Tx{0}{p\,\IZ}{0}-\Tx{0}{(x-1)f\,\IZ}{0}$; the right side is $\frac15P_{\binom{[5]}1}(\one;1,1)=\frac15\Tx{0}{f\,\IZ}{0}$.
\end{proof}

\noindent Receipts \computed. At the strengths $2,\dots,5$ the same insertion reduces to rational multiples of this relation. There is no rational relation among $1$, $\Tx{3}{\IZ}{0}$, $f(1)$ and the five Trailmoments other than $\Tx{0}{(x-1)f\,\IZ}{0}$ with coefficients below $10^{30}$ (PSLQ, $360$ digits): in that framework, over $\mathbb Q$, the block is five-dimensional. Its six values:
\begin{center}
\small
\begin{tabular}{@{}lll@{}}
\toprule
$j$ & $\Tx{0}{(x-1)^j\,p\,\IZ}{0}$ & $\Tx{0}{(x-1)^j\,f\,\IZ}{0}$\\ \midrule
$0$ & $8.570280443374461268149695824246537286\ldots$ & $39.94558801199624569894402849842399589\ldots$\\
$1$ & $0.161069702276674858851310666248514198\ldots$ & $0.581162840975212128360890124561738107\ldots$\\
$2$ & $0.171874581383536904847683491575293916\ldots$ & $0.616835511971624849092924541153028126\ldots$\\
\bottomrule
\end{tabular}
\end{center}

\begin{theorem}[A mark added carries the ray law]\label{a:t:cert}\derived; receipts \computed\
Normalize the ray law of $\binom{[n]}1$, $\sum_{i=0}^{d_n}c^{(n)}_i(\sigma)\,P_{\binom{[n]}1}((\sigma+i)\one;\sigma+i,\sigma+i)=0$, so that $c_0^{(n)}=(-1)^n\sigma^{n-1}$, as for four marks above and for five marks in Proposition~\ref{a:p:pairing}, and let $L_n^\ast$ be its adjoint. For $n=1,\dots,6$,
\[
\sum_{i=0}^{d_{n+1}}c^{(n+1)}_i(s)\binom{x-1}{s-1+i}^{\!2}=\big(L_n^\ast A_s\big)(x),\qquad A_s(x)=\big((n+1)x-(n+2)s\big)\binom xs^{\!2},
\]
an identity of rational functions of $x$ and $s$ once both sides are divided by $\binom xs^2$. Read one way it produces the law: the words $\binom{x-1}{s-1+i}^2$ have distinct degrees, so the coefficients of $L_n^\ast A_s$ in them are the $c^{(n+1)}_i(s)$. Read the other way it proves it: for $2\le n\le6$, if the ray law of $\binom{[n]}1$ holds identically in $\sigma$, so does that of $\binom{[n+1]}1$, and the lowered setting $x^{(1)}$ of $n+1$ marks, divided by $\Gamma(s)^2$, obeys it too. The two-mark period $P_{\binom{[2]}1}(\sigma\one;\sigma,\sigma)=\Gst21{\sigma,\ \sigma}{2\sigma}$ obeys $\sigma P(\sigma)-2(2\sigma+1)P(\sigma+1)=0$ by the shift of the stacked Gamma (Prop.~1.3(iii), p.~11). So the ray laws of $\binom{[n]}1$, $n\le7$, are derived from the stacked Gamma, one mark at a time.
\end{theorem}

\begin{proof}
Write $p_n$ for the centre period of $\binom{[n]}1$, and let $s>0$ be non-integral. Proposition~\ref{a:p:cosetmark}, with $n$ marks merged in place of four, gives $P_{\binom{[n+1]}1}((s+i)\one;s+i,s+i)=-\sum_{d\in s+\mathbb N}(w_{s,i}\,p_n)'(d)$ with $w_{s,i}=W_{s+i}/\Gamma(s+i)^2$, the terms $d=s,\dots,s+i-1$ being zero by the double zeros of $W_{s+i}$. On the coset $w_{s,i}(d)$ is $\binom{d-1}{s-1+i}^2$ written with Gammas, so the displayed identity holds there. On $s+\mathbb N$ the move of Proposition~\ref{a:p:sbp} has no boundary for the words $A_s$ and $A_s'$, which vanish at $d=s-1,s-2,\dots$: $\sum_d(L_n^\ast A)\,\varphi=\sum_dA\,(L_n\varphi)$, and likewise with $L_n'$. With $(L_n^\ast A)'=L_n'^\ast A+L_n^\ast(A')$ and $L_np_n'=-L_n'p_n$ (the law differentiated),
\begin{multline*}
\sum_ic^{(n+1)}_i(s)\,P_{\binom{[n+1]}1}((s+i)\one;s+i,s+i)=-\sum_{d\in s+\mathbb N}\Big((L_n'^\ast A_s)\,p_n+A_s'\,L_np_n+A_s\,L_np_n'\Big)(d)\\
=-\sum_{d}(L_n'^\ast A_s)\,p_n+\sum_dA_s\,L_n'p_n=0,
\end{multline*}
and $\sum_ic^{(n+1)}_i(s)\,\Gamma(s+i)^{-2}x^{(1)}(s+i)=\sum_d(L_n^\ast A_s)\,p_n=\sum_dA_s\,L_np_n=0$. Both sides are holomorphic in $s$, so the laws hold identically. The displayed identity is a finite computation with rational functions of $(x,s)$.
\end{proof}

\noindent Receipts \computed. The certificate $A_s$ was found by exact linear algebra at each $n\le6$ and $s\le5$, the same at every $n$ up to the factor $(n+1)x-(n+2)s$, and the identity was then checked with $s$ indeterminate. At $n=1$ ($p_1=1$, $c^{(1)}=(-1,1)$) it holds as an identity, but the sums of the proof diverge. The five-mark law at $s=\frac12$ and $s=\frac3{10}$, from the share side by quadrature: residuals $4.6\cdot10^{-14}$ and $2.3\cdot10^{-14}$ against terms of size $10^2$. At $s=0$, with the cut at $0$, the same reduction returns the corner relation of $n+1$ marks as $A_0'(0)=n+1$ times the boundary of Proposition~\ref{a:p:sbp}: the corners $(n+1)!=(n+1)\cdot n!$ of Proposition~\ref{a:p:corner}, here $120=5\cdot24$.

\begin{corollary}[Two masters over the corner, two lowered numbers]\label{a:c:masters}\derived\
For every $s\in\{1,2,\dots\}$ there are rationals $a_s,b_s,q_s$ with
\begin{align*}
\Gamma(s)^2P_{\binom{[5]}1}(s\one;s,s)&=q_s+a_s\,P_{\binom{[5]}1}(\one;1,1)+b_s\,P_{\binom{[5]}1}(2\one;2,2),\\
x^{(1)}(s)&=a_s\,x^{(1)}(1)+b_s\,x^{(1)}(2),
\end{align*}
and in the block
\begin{gather*}
P_{\binom{[5]}1}(\one;1,1)=\Tx{0}{f\,\IZ}{0},\\
P_{\binom{[5]}1}(2\one;2,2)=\Tx{0}{(x-1)^2f\,\IZ}{0}-2\,\Tx{0}{(x-1)\,p\,\IZ}{0},\\
x^{(1)}(1)=\Tx{0}{p\,\IZ}{0},\qquad x^{(1)}(2)=\Tx{0}{(x-1)^2p\,\IZ}{0}.
\end{gather*}
Modulo $L^\ast(x\,\mathbb Q[x])$ every word $W_s$ lies in $\mathrm{span}\{1,(x-1)^2\}$: the ray reaches the direction $x-1$ of the block only through the value term $-2\Tx{0}{(x-1)p\,\IZ}{0}$ of its second master. For example $(a_3,b_3,q_3)=\big(-\frac1{45},\frac{19}{15},\frac8{15}\big)$ and $(a_4,b_4,q_4)=\big(-\frac{113}{1125},\frac{1907}{375},\frac{944}{375}\big)$.
\end{corollary}

\begin{proof}
The four block identities are Proposition~\ref{a:p:cosetmark} at $s=1,2$, with $W_1=1$, $W_2=(x-1)^2$, $W_2'=2(x-1)$. Write $W_s/\Gamma(s)^2=L^\ast B_s+r_s$ with $\deg r_s\le2$. Since the display of Theorem~\ref{a:t:cert} is an identity of polynomials and the splitting of Proposition~\ref{a:p:block} is direct, the $r_s$ obey the five-mark recurrence for $s\ge1$, and at $s=0$, where $A_0=5x$, they obey $\sum_{i\ge1}c^{(5)}_i(0)\,r_i=0$. As $r_1=1$ and $r_2=(x-1)^2$, every $r_s$ lies in their span, and the same recursion, from $B_1=B_2=0$, gives $B_s(0)=0$. Then $x^{(1)}(s)=\Gamma(s)^2\Tx{0}{r_s\,p\,\IZ}{0}$, and by Proposition~\ref{a:p:sbp} $P_{\binom{[5]}1}(s\one;s,s)=\Tx{0}{(r_sf-r_s'p)\,\IZ}{0}+24B_s'(0)$, where $\sum_{i\ge1}c^{(5)}_i(0)B_i'(0)=A_0'(0)=5$ carries the corner.
\end{proof}

\begin{proposition}[The lowered half is a stacked Gamma]\label{a:p:stackedhalf}\computed\
\[
x^{(1)}(1)=\Tx{0}{p\,\IZ}{0}=\frac1{120}\Gst84{\tfrac1{15},\ \tfrac2{15},\ \tfrac4{15},\ \tfrac8{15},\ \tfrac15,\ \tfrac25,\ \tfrac35,\ \tfrac45}{\tfrac12,\ \tfrac12,\ \tfrac12,\ \tfrac12}
\]
to $364$ digits, and
\begin{gather*}
x^{(1)}(2)=\frac{65\,x^{(1)}(1)-32\,\Gamma(\frac12)^8/x^{(1)}(1)}{1125}\quad(360\text{ digits}),\\
x^{(2)}(1)=\frac{\sqrt{15}}{4\,\Gamma(\frac12)^2}\,x^{(1)}(1),\qquad x^{(2)}(2)=\frac{\sqrt{15}}{4\,\Gamma(\frac12)^2}\cdot\frac{65\,x^{(1)}(1)+32\,\Gamma(\frac12)^8/x^{(1)}(1)}{1125}\quad(100\text{ digits}).
\end{gather*}
So the block of Proposition~\ref{a:p:frame} is $x^{(1)}(1)x^{(2)}(2)-x^{(1)}(2)x^{(2)}(1)=\frac{\sqrt{15}}{4\Gamma(1/2)^2}\cdot\frac{64\,\Gamma(1/2)^8}{1125}=\frac{16\sqrt{15}}{1125}\Gamma(\frac12)^6$: the stacked Gamma cancels from the frame, which is why the frame saw only $\Gamma(\frac12)$.
\end{proposition}

The values half that the ray reaches is two stacked-Gamma numbers, $x^{(1)}(1)$ and $\Gamma(\frac12)^8/x^{(1)}(1)$. The direction $x-1$ is not among them: $\Tx{0}{(x-1)\,p\,\IZ}{0}$ has no rational relation with $1$, $\Tx{3}{\IZ}{0}$, $\Gamma(\frac12)^4$, $x^{(1)}(1)$ and $\Gamma(\frac12)^8/x^{(1)}(1)$, nor with ten products of $x^{(1)}(1)^{\pm1}$, powers of $\Gamma(\frac12)^2$ and $\sqrt{15}$, with coefficients below $10^{12}$ \computed.

\begin{proposition}[The block on the four-mark chart and its pencil]\label{a:p:chart}\derived; receipts \computed\
On the chart $y_{e_0}=0$ of four marks let $X=K(\alpha)K(\alpha^{-1})$, the reciprocal of the stride of the centre ray (Def.~35.9, p.~309), so that $p(\sigma)=\int X^{-\sigma}dy$, the value along the ray from the real origin $(\sigma\one;\sigma,\sigma)$ (Def.~35.1, p.~307), with $X\ge16$ on the chart (Thm.~35.10(i), p.~309). For a polynomial $R$,
\begin{gather*}
\Tx{0}{R\,p\,\IZ}{0}=\int\rho_R(X)\,dy,\qquad\Tx{0}{R\,f\,\IZ}{0}=\int\log X\,\rho_R(X)\,dy,\\
\rho_R(X)=\sum_{m\ge1}R(m)X^{-m},
\end{gather*}
a rational function of $X$ whose only pole is $X=1$. In particular
\begin{gather*}
P_{\binom{[5]}1}(\one;1,1)=\int\frac{\log X}{X-1}\,dy,\qquad P_{\binom{[5]}1}(2\one;2,2)=\int\frac{(X+1)\log X-2(X-1)}{(X-1)^3}\,dy,\\
x^{(1)}(1)=\int\frac{dy}{X-1},\qquad x^{(1)}(2)=\int\frac{X+1}{(X-1)^3}\,dy,\qquad \Tx{0}{(x-1)\,p\,\IZ}{0}=\int\frac{dy}{(X-1)^2}.
\end{gather*}
The first is the left closing of Lemma~5.16(ii), p.~34, read on the chart. With the generating function of the four-mark ray from the origin $(\varepsilon\one;\varepsilon,\varepsilon)$, $G_\varepsilon(q)=\sum_{k\ge1}p(k+\varepsilon)\,q^{k-1}=\int X^{-\varepsilon}\,dy/(X-q)$ for $|q|<16$, carried by the pencil $K\Psi-q\alpha^{\one}$ (Thm.~35.10(ii)),
\begin{gather*}
\Tx{0}{(x-1)^j\,p\,\IZ}{0}=(q\partial_q)^jG_0(q)\big|_{q=1},\\
\Tx{0}{(x-1)^j\,f\,\IZ}{0}=-\partial_\varepsilon(q\partial_q)^jG_\varepsilon(q)\big|_{q=1,\ \varepsilon=0}:
\end{gather*}
the block is the jet at $q=1$ of the four-mark generating function and of its slope in the origin.
\end{proposition}

\begin{proof}
Def.~35.1 along the ray from $(\sigma\one;\sigma,\sigma)$, the positivity of every term, and $\sum_{m\ge1}R(m)X^{-m}$ summed for $X>1$: $\rho_1=1/(X-1)$, $\rho_{x-1}=1/(X-1)^2$, $\rho_{(x-1)^2}=(X+1)/(X-1)^3$. For the masters use Corollary~\ref{a:c:masters}; for the pencil, Thm.~35.10(ii) with $k_0=1$ and the origin form $X^{-\varepsilon}dy$.
\end{proof}

\noindent Receipts \computed, by three-dimensional quadrature on the chart, to twelve digits:
\[
\int\frac{dy}X=8.414398322115,\qquad\int\frac{\log X}{X-1}\,dy=39.945588011991,\qquad\int\frac{dy}{X-1}=8.570280443372,
\]
against $7\Tx{3}{\IZ}{0}$, the five-mark centre at strength one, and $x^{(1)}(1)$.

\begin{remark}[What the closing found]\label{a:r:found}
(i) The live Trailmoment lands on lattice inputs only through the residue at zero shares (Proposition~\ref{a:p:sbp}).

(ii) The words of the five-mark centre ray reach a plane of the five-dimensional block on each side: on the values side the lowered setting, two stacked-Gamma numbers (Proposition~\ref{a:p:stackedhalf}); on the slopes side the period, whose two masters are
\begin{gather*}
P_{\binom{[5]}1}(\one;1,1)=\Tx{0}{f\,\IZ}{0},\\
P_{\binom{[5]}1}(2\one;2,2)=\Tx{0}{(x-1)^2f\,\IZ}{0}-2\,\Tx{0}{(x-1)\,p\,\IZ}{0}:
\end{gather*}
the slope in the origin of the four-mark generating function at $q=1$, and its second jet less twice the first jet of the generating function itself. They have no rational relation with $1$, $\Tx{3}{\IZ}{0}$, $f(1)$, $\Gamma(\frac12)^4$, $\Gamma(\frac12)^8$, $\Tx{5}{\IZ}{0}$, $\Gamma(\frac12)^4\Tx{3}{\IZ}{0}$, the two stacked-Gamma numbers and the direction $x-1$ of the block, with coefficients below $10^{20}$ (PSLQ, $360$ digits) \computed: in that framework, over $\mathbb Q$, they are new numbers.

(iii) The last minor of \S14.17 is now written in the masters over the stacked Gamma (Proposition~\ref{a:p:stackedhalf}):
\begin{multline*}
P_{\binom{[5]}1}(\one;1,1)\,x^{(1)}(2)-P_{\binom{[5]}1}(2\one;2,2)\,x^{(1)}(1)\\
=\frac{65\,x^{(1)}(1)-32\,\Gamma(\frac12)^8/x^{(1)}(1)}{1125}\,P_{\binom{[5]}1}(\one;1,1)-x^{(1)}(1)\,P_{\binom{[5]}1}(2\one;2,2)\\
=4.34000293153566203407\ldots,
\end{multline*}
and it is not a rational combination of $1$, $\Gamma(\frac12)^4$, $\Gamma(\frac12)^8$, $\Tx{3}{\IZ}{0}$ and seven products of $x^{(1)}(1)^{\pm1}$ or $x^{(1)}(1)^{\pm2}$ with powers of $\Gamma(\frac12)^4$, with coefficients below $10^{15}$ \computed. Its scale is the masters' scale; nothing smaller is missing.

(iv) The rational words and the pole columns of Chapter~6 did not enter: every word of the centre ray is a polynomial, read at index zero. What closed the live Trailmoment is the shifted inputs (Thm.~6.2), the cut relation (Prop.~5.3), summation by parts in its adjoint form (Thm.~6.5, p.~43), and a fixed block at nonpositive index, the counterpart of Thm.~6.9, p.~45.
\end{remark}

\renewcommand{\thesection}{16}
\setcounter{proposition}{0}
\renewcommand{\theproposition}{16.\arabic{proposition}}

\section{The five-mark lists along their hinge trees}

A connected list is glued from prime tables and series and parallel nodes at its bunches of width one (Thm.~20.9, p.~166), and its canonical hinge tree (Def.~22.14, p.~189) is read off the Skeleton: the width $c(S)=r(S)+r(\cE\setminus S)-r(\cE)$ of Def.~7.10, p.~49, is a count of the circuit and cocircuit layers, and a bunch of width one is where the list composes itself out of two pieces. The period is read from the leaves up (Thm.~22.15, p.~190). A cell enters only through its hinge wall and its law (Thm.~22.7, p.~185), a series node multiplies the transforms of its cells and a parallel node their dual transforms (Thm.~22.13, p.~188), and a mark left in place is a one-mark cell whose two transforms at the centre are the Gramma $\sG_{0;(1,-1)}(1,y)$ (Def.~22.17(1), p.~191). This entry reads the five lists of \S15.4 that way. Three are series-parallel, one carries a prime node that acts as a mark, and one is prime.

\begin{center}
\small
\begin{tabular}{@{}l l@{}}
\toprule
list & canonical hinge tree \derived\\ \midrule
$\binom{[5]}1$ & one parallel node\\
$\binom{[4]}2$, a mark doubled & the prime node $\binom{[4]}2$, a parallel node of three marks at one marker\\
$\binom{[5]}2$ & one prime node\\
$\binom{[3]}2$, two marks doubled & a series node of three marks, parallel nodes of three marks at two markers\\
$\binom{[3]}2$, a mark tripled & a series node of three marks, a parallel node of four marks at one marker\\
\bottomrule
\end{tabular}
\end{center}
The bunches of width one with two marks on each side are the parallel classes, of width $1+r(\cE\setminus P)-r(\cE)=1$; in $\binom{[4]}2$ and $\binom{[5]}2$ every set of two marks has width $2$, so both are prime.

\subsection*{16.1\quad A prime node that is a mark}

\begin{proposition}[The doubled mark of $\binom{[4]}2$]\label{a:p:u24}\derived; receipts \computed\
Replace a mark $i$ of $\binom{[4]}2$ by a parallel class of $k\ge1$ marks. Then
\[
P_{\binom{[4]}2[i\to k]}(\sigma\one;\sigma,\sigma)=P_{\binom{[3]}1}(\sigma\one;\sigma,\sigma)\;P_{\binom{[k+1]}1}(\sigma\one;\sigma,\sigma).
\]
At $k=1$ this is $P_{\binom{[4]}2}=P_{\binom{[3]}1}\,\Gst21{\sigma,\ \sigma}{2\sigma}$, the agreement of the table of \S15.1 at every strength; at $k=2$ it is the second row of \S15.4,
\[
P_{\binom{[4]}2,\ \text{a mark doubled}}(\sigma\one;\sigma,\sigma)=P_{\binom{[3]}1}(\sigma\one;\sigma,\sigma)^2 .
\]
\end{proposition}

\begin{proof}
$\binom{[4]}2$ is identically self-dual. The cell of $\binom{[4]}2$ hinged at $i$ has the hinge minors $\binom{[4]}2\setminus i=\binom{[3]}2$ and $\binom{[4]}2/i=\binom{[3]}1$, dual with the identity and the same weights, so its ratio is identically one, and in every gluing it acts as $\phi(\Delta)$ times one mark, with $\phi(\Delta)=P_{\binom{[3]}2}(\sigma\one;\kappa,2\sigma-\kappa)$ (Cor.~22.8(i),(ii), p.~186). Its hinge wall is $\kappa=\sigma$, balanced against the parallel cell of the class, whose hinge wall is $k\sigma-(k-1)\sigma=\sigma$ (Thm.~22.13(ii)). By Prop.~8.14, p.~57, $P_{\binom{[3]}2}(\sigma\one;\sigma,\sigma)=P_{\binom{[3]}1}(\sigma\one;\sigma,\sigma)$. What is left is the class glued to one mark, the banana $\binom{[k+1]}1$.
\end{proof}

\noindent Receipts \computed: Procedure~\ref{a:p:proc} for $\binom{[4]}2$ with a mark doubled agrees with $P_{\binom{[3]}1}(\sigma\one;\sigma,\sigma)^2$ to $60$ digits at $\sigma=1,\dots,4$, and for $\binom{[4]}2$ with $P_{\binom{[3]}1}\Gst21{\sigma,\ \sigma}{2\sigma}$ at $\sigma=1,2,3$.

\begin{corollary}\label{a:c:u24}\derived\
The ray law of $\binom{[4]}2$ with a mark doubled is the symmetric square of the three-mark law: its native values are the products of $\frac19$ and $1$ in pairs, the roots $81,9,1$ of \S15.4, and its corner is $6^2=36$. Every value lies in $\mathbb Q+\mathbb Q\,P_{\binom{[3]}1}(\one;1,1)+\mathbb Q\,P_{\binom{[3]}1}(\one;1,1)^2$: its two masters are the three-mark centre and its square.
\end{corollary}

The same argument at the centre of any identically self-dual list splits off a mark (Cor.~22.8(iii) and Prop.~8.14) \derived:
\[
P_{\binom{[2r]}r}(\sigma\one;\sigma,\sigma)=\Gst21{\sigma,\ \sigma}{2\sigma}\,P_{\binom{[2r-1]}{r-1}}(\sigma\one;\sigma,\sigma).
\]
So $\binom{[6]}3$ carries exactly the masters of $\binom{[5]}2$.

\subsection*{16.2\quad A series root over banana leaves}

\begin{proposition}[The triangle rows]\label{a:p:triangles}\derived; receipts \computed\
Let $f_{P_k}$ be the transform at the centre of the parallel cell of a class of $k$ marks; $\int_0^\infty y^{w-1}f_{P_k}(y)\,dy=\Gamma(w)^2P_{\binom{[k]}1}((\sigma-w)\one;\sigma-w,\sigma-w)$ for $0<\re w<\sigma$. Then
\begin{gather*}
\Gamma(\sigma)^2P_{\binom{[3]}2,\ \text{a mark tripled}}(\sigma\one;\sigma,\sigma)=\int_0^\infty y^{\sigma-1}\,\sG_{0;(1,-1)}(1,y)^2\,f_{P_3}(y)\,dy\\
\qquad=\frac1{2\pi i}\int_{(c)}\Gamma(d)^2\Gamma(\sigma-d)^2\,P_{\binom{[4]}2}(d\one;d,d)\,dd,\\
\Gamma(\sigma)^2P_{\binom{[3]}2,\ \text{two marks doubled}}(\sigma\one;\sigma,\sigma)=\int_0^\infty y^{\sigma-1}\,\sG_{0;(1,-1)}(1,y)\,f_{P_2}(y)^2\,dy .
\end{gather*}
So the tripled triangle is the closing of Proposition~\ref{a:p:cosetmark} with the ray of the four-mark banana replaced by the ray of $\binom{[4]}2$.
\end{proposition}

\begin{proof}
The triangle is a series root (Def.~22.17(5), p.~191): its cells are the parallel cells of the classes and the marks left in place, the latter with the transform $\sG_{0;(1,-1)}(1,y)$. The dual transform of a parallel cell is the product of the $k$ Grammas of its marks (Thm.~22.13(ii)), whose Mellin read is $\Gamma(z)^2P_{\binom{[k]}1}(z\one;z,z)$ (Proposition~\ref{a:p:bananapin}); its hinge wall is $\sigma$, and at $s=t=\kappa=\sigma$ the change of transform reflects $z\mapsto\sigma-z$ (Prop.~22.5(iii),(v), p.~184). The Mellin read of a product pins the sum of the shares, as in the proof of Proposition~\ref{a:p:bananapin}, and $P_{\binom{[2]}1}P_{\binom{[3]}1}=P_{\binom{[4]}2}$ is Proposition~\ref{a:p:u24} at $k=1$.
\end{proof}

\noindent Receipts \computed: the closing $-\sum_{d\ge1}(W_s\,P_{\binom{[4]}2})'(d)/\Gamma(s)^2$ of Proposition~\ref{a:p:cosetmark} agrees with Procedure~\ref{a:p:proc} for the tripled triangle to $60$ digits at $s=1,2,3$; for two marks doubled, opening the list at its unmarked mark (Prop.~19.19, read as in Proposition~\ref{a:p:u25split} below: the contraction is the four-mark banana, the deletion two parallel pairs) gives $7.8658373683630$ by three-dimensional quadrature, against $7.86583736836402159\ldots$.

\begin{proposition}[Closing the tripled triangle over its input]\label{a:p:tripled}\derived; receipts \computed\
Let $q(\sigma)=P_{\binom{[4]}2}(\sigma\one;\sigma,\sigma)$ and $f_q=-q'$.
\begin{enumerate}[label=(\roman*)]
\item $q$ obeys the three-mark law times the shift of the stacked Gamma,
\[
\sigma^3q(\sigma)-2(2\sigma+1)(10\sigma^2+10\sigma+3)\,q(\sigma+1)+36(\sigma+1)(2\sigma+1)(2\sigma+3)\,q(\sigma+2)=0,
\]
with corner $12=2\cdot6$ and no term of order $\sigma$, and its symbol at $\lambda=1$ is $105\neq0$.
\item Summation by parts at the cut $0$, as in Proposition~\ref{a:p:sbp}: for every polynomial $R$,
\begin{gather*}
\Tx{0}{(L_q^\ast R)\,q\,\IZ}{0}=-12\,R(0),\\
\Tx{0}{(L_q^\ast R)\,f_q\,\IZ}{0}=\Tx{0}{(L_q'^\ast R)\,q\,\IZ}{0}+36\,R(-1)\,P_{\binom{[3]}1}(\one;1,1).
\end{gather*}
The last coefficient of the law has only a simple zero at $\sigma=-1$: the shift of the stacked Gamma, $P_{\binom{[2]}1}(\sigma+2)/P_{\binom{[2]}1}(\sigma+1)=(\sigma+1)/(2(2\sigma+3))$, cancels one of the two factors $\sigma+1$ of the three-mark coefficient $9(\sigma+1)^2$. So the slope boundary there reads $q(1)=P_{\binom{[3]}1}(\one;1,1)$.
\item Every value of the tripled triangle is a rational number, plus a rational multiple of $P_{\binom{[3]}1}(\one;1,1)$, plus a rational combination of the three slope Trailmoments $-\Tx{0}{(Rq)'\,\IZ}{0}$, $R\in\{1,\,x,\,x^2/2\}$. The words $W_s$, $s\le8$, reach all three; the five-mark banana reached two (Corollary~\ref{a:c:masters}).
\item In these coordinates the rational relation of \S15.4 for this row, the second corner, is the relation among $1,P(1),\dots,P(5)$, coordinate by coordinate.
\end{enumerate}
\end{proposition}

\begin{proof}
(i) $\sigma\,P_{\binom{[2]}1}(\sigma)=2(2\sigma+1)\,P_{\binom{[2]}1}(\sigma+1)$ (Prop.~1.3(iii)) multiplies the three-mark law coefficient by coefficient into $\sigma$ times the display. The corner is the product of the corners $\lim\sigma P_{\binom{[2]}1}=2$ and $6$; $\sigma\Gst21{\sigma,\ \sigma}{2\sigma}=2\Gamma(1+\sigma)^2/\Gamma(1+2\sigma)$ and $\sigma^2P_{\binom{[3]}1}$ have no terms of order $\sigma$ (Proposition~\ref{a:p:corner}). (ii) The move of Proposition~\ref{a:p:sbp}: the value boundary is $-6q(1)+108q(2)=-12$; for the slopes the boundary at the cut $0$ is the term of order $\sigma$, zero, and the one at $-1$ is $-R(-1)\,c_2'(-1)\,q(1)$ with $c_2'(-1)=-36$. (iii) As in Proposition~\ref{a:p:block}, with $\mathrm{span}\{1,x,x^2\}$ as the complement of $L_q^\ast(\mathbb Q[x])$, and Proposition~\ref{a:p:cosetmark} for the values. (iv) Exact arithmetic on the eight coordinates.
\end{proof}

\noindent Receipts \computed. The coordinates reproduce Procedure~\ref{a:p:proc} to $60$ digits at $s\le5$. There is no rational relation among $1$, the three-mark centre $P_{\binom{[3]}1}(\one;1,1)$, its slope $\partial_\sigma P_{\binom{[3]}1}(\sigma\one;\sigma,\sigma)|_{\sigma=1}$ and the six Trailmoments
\[
\Tx{0}{x^jq\,\IZ}{0},\qquad \Tx{0}{x^jf_q\,\IZ}{0}\qquad(j\le2)
\]
with coefficients below $10^{25}$, and none among $1$, the three-mark centre and the three slope Trailmoments of (iii) below $10^{40}$ (PSLQ, $300$ digits). So the four masters of the tripled triangle are the three-mark centre, entering at the boundary $\sigma=-1$, and the three slope Trailmoments.

\subsection*{16.3\quad The prime node at a mark}

\begin{proposition}[$\binom{[5]}2$ opened at a mark]\label{a:p:u25split}\derived; receipt \computed\
For a mark $e$ of $\binom{[5]}2$ the deletion is $\binom{[4]}2$ and the contraction $\binom{[4]}1$. Write $X_{\setminus}=X_{\binom{[4]}2}$ and $X_{/}=X_{\binom{[4]}1}$ for their strides on the chart of the other four marks. Then
\[
P_{\binom{[5]}2}(\sigma\one;\sigma,\sigma)=\int dy\int_0^1w^{\sigma-1}(1-w)^{\sigma-1}\big(w\,X_{\setminus}+(1-w)\,X_{/}\big)^{-\sigma}dw,
\]
and on the share side of the shape,
\[
\Gamma(\sigma)^2P_{\binom{[5]}2}(\sigma\one;\sigma,\sigma)=\frac1{2\pi i}\int_{(c)}\Gamma(d)^2\Gamma(\sigma-d)^2\,\Phi_\sigma(d)\,dd,\qquad \Phi_\sigma(d)=\int X_{/}^{-d}X_{\setminus}^{-(\sigma-d)}\,dy,
\]
$0<c<\sigma$, with $\Phi_\sigma(d)=P_F\big(\sigma\one;(d,\sigma-d),(d,\sigma-d)\big)$ the centre of the family $F$ of the two children (Def.~8.1, p.~53). At $\sigma=1$,
\[
P_{\binom{[5]}2}(\one;1,1)=\int\frac{\log X_{\setminus}-\log X_{/}}{X_{\setminus}-X_{/}}\,dy .
\]
\end{proposition}

\begin{proof}
The mark is the Breta of its children's even letters (Prop.~19.19, p.~144): at the centre the prefactor and the stacked Gamma are one, and $E_{\setminus e}=\alpha^{\cE\setminus e}X_{\setminus}$, $E_{/e}=\alpha^{\cE\setminus e}X_{/}$. Splitting the coefficient $wX_{\setminus}+(1-w)X_{/}$ (Cor.~4.4, p.~26) and integrating the shape gives the Gammas, one pair for each child.
\end{proof}

\noindent Receipt \computed: the last integral is $4.0297141707108$ by three-dimensional quadrature, against $4.02971417071104960732\ldots$ of \S15.4.

\begin{proposition}[The contraction under the deletion]\label{a:p:under}\derived\
On the chart $X_{\setminus}-X_{/}=(e_2^2-e_1e_3)/e_4$ with
\[
e_2^2-e_1e_3=\sum_{i<j}\alpha_i^2\alpha_j^2+\sum_i\sum_{\{j,k\}\not\ni i}\alpha_i^2\alpha_j\alpha_k+2\,\alpha_1\alpha_2\alpha_3\alpha_4,
\]
a polynomial with positive coefficients, so $X_{/}<X_{\setminus}$ everywhere. The walls of $F$ at the strengths $(d,1-d)$ are $1$ at $|S|$ odd and $2-d$ at $|S|=2$, so $\Phi=\Phi_1$ is holomorphic on $\re d<2$, and closing to the left (Thm.~4.5, Lemma~5.16),
\[
P_{\binom{[5]}2}(\one;1,1)=\sum_{m\ge0}\Phi'(-m),\qquad \Phi(-m)=\int X_{/}^{\,m}\,X_{\setminus}^{-1-m}\,dy :
\]
the slopes of the deletion's centre ray with the contraction's even letter inserted $m$ times.
\end{proposition}

\begin{proof}
$X_{/}=e_1e_3/e_4$, $X_{\setminus}=e_2^2/e_4$, and the polynomial is read off the monomials of $e_2^2$ and $e_1e_3$. The widths are $1$ for every $S$ in $\binom{[4]}1$ and $1,2,1$ at $|S|=1,2,3$ in $\binom{[4]}2$. On the line $\re d=-N-\frac12$,
\[
|\Phi(d)|\le\int\big(X_{/}/X_{\setminus}\big)^{N+1/2}X_{\setminus}^{-1}\,dy\to0,
\]
so the left closing applies; at $d=-m$ the kernel $\Gamma(d)^2\Gamma(1-d)^2$ of Lemma~5.16 has the principal part $(d+m)^{-2}$ and no simple part. Equivalently, expand $1/(X_{\setminus}-X_{/})=\sum_mX_{/}^mX_{\setminus}^{-m-1}$ in the last display of Proposition~\ref{a:p:u25split}.
\end{proof}

The ratio $X_{/}/X_{\setminus}$ is largest at the table's gauge, where it is $\frac{16}{36}=\frac49$, so the closing converges like $(\frac49)^m$ (alias: Newton's inequality $e_1e_3\le\frac49e_2^2$) \recalled.

\begin{proposition}[The values at the lattice points are leans of the three-mark banana]\label{a:p:leans}\derived; receipts \computed\
For $m\in\mathbb N$, with $(e_1e_3)^m=\sum_vc_v\alpha^v$,
\[
\Phi(-m)=\sum_vc_v\,P_{\binom{[4]}2}\big(\one+v;1+m,1+m\big),
\]
and for masses $\bmu$ on the balance wall and every mark $g$,
\[
P_{\binom{[4]}2}(\bmu;s,s)=\frac1{\Gamma(2s)}\int_0^\infty y^{\mu_g-1}\prod_{e\neq g}\sG_{2s-\mu_g-\mu_e;(1,-1)}(1,y)\,dy :
\]
$\Phi(-m)$ is a finite sum of leans of the three-mark banana, Grammas of integer index read three at a time. The first values are
\begin{gather*}
\Phi(0)=P_{\binom{[3]}1}(\one;1,1),\qquad \Phi(-1)=\frac{2+2P_{\binom{[3]}1}(\one;1,1)}9,\\
\Phi(-2)=\frac{71+56P_{\binom{[3]}1}(\one;1,1)}{810}\quad\computed .
\end{gather*}
\end{proposition}

\begin{proof}
On the chart $X_{/}^mX_{\setminus}^{-1-m}=\alpha^{\one}(e_1e_3)^m\,e_2^{-2(1+m)}$, and $e_2$ is both letters of $\binom{[4]}2$; each monomial $\alpha^v$ is an insertion moving the masses (Lemma~35.2(iii), p.~308). Then Cor.~22.8(iii) at $g$, $P_{\binom{[4]}2}(\bmu;s,s)=\Gst21{\mu_g,\ 2s-\mu_g}{2s}P_{\binom{[3]}2}(\bmu|;2s-\mu_g,\mu_g)$; Prop.~8.14 turns $\binom{[3]}2$ into $\binom{[3]}1$ at the exponents $(\mu_g,2s-\mu_g)$; and the share side of the three-mark banana at a lean, $\Gamma(a)\Gamma(b)P_{\binom{[3]}1}(\bmu';a,b)=\int_0^\infty y^{a-1}\prod_e\sG_{b-\mu'_e;(1,-1)}(1,y)\,dy$ (Prop.~4.2 with the pins $a_e-b_e=\mu'_e-b$, as in Proposition~\ref{a:p:bananapin}), whose Gammas cancel the Beta corner.
\end{proof}

\noindent Receipts \computed: $\Phi(-1)$ by three-dimensional quadrature is $0.7430904974869$, against $\frac29(1+P_{\binom{[3]}1}(\one;1,1))=0.74309049748654642013\ldots$; the leans themselves by one-dimensional quadrature to $40$ digits, with PSLQ for the rational coefficients; on the way, $\int_0^\infty\sG_{0;(1,-1)}(1,y)\,\sG_{1;(1,-1)}(1,y)^2\,dy=\frac13$.

\begin{remark}[What the prime node reads]\label{a:auto:16.8}
The prime node reads the three-mark banana at leans. Its values at the lattice points land on $\mathbb Q+\mathbb Q\,P_{\binom{[3]}1}(\one;1,1)$, and its period at strength one is the sum of the slopes there: the division of Propositions~\ref{a:p:sbp} and~\ref{a:p:block}, values on lattice inputs and slopes carrying the masters, now along the deletion's ray with the contraction's even letter as its stride, of native value $\frac49$ at the table's gauge. Two of the three masters of $\binom{[5]}2$ are slope Trailmoments of that ray, as the two of $\binom{[5]}1$ were of the four-mark ray, and the third is $\Pth$ (Theorem~\ref{a:t:u25masters}); the law in $m$ and its closing are Propositions~\ref{a:p:u25law} and~\ref{a:p:u25clos}.
\end{remark}

\begin{remark}[What the tree decided]\label{a:auto:16.9}
A parallel node multiplies Grammas: the banana. An identically self-dual prime node is a mark of ratio one, and the period factors: the doubled mark of $\binom{[4]}2$ is the square of the three-mark ray, with no new number. A series root over parallel cells multiplies transforms; its closing has the banana's shape over a product input, and the stacked Gamma in that input's law opens the boundary at $\sigma=-1$, which is where the second corner of the tripled triangle comes from (for two marks doubled it is the boundary at zero shares, Remark~\ref{a:r:closingfound}(ii)). A prime node that is not self-dual is opened at a mark into its children, weighed by the Breta of their even letters, and reads leans. The three-mark centre enters every row of \S15.4 but the banana, row four included, through the contraction of $H$ at its single mark. The row with two marks doubled has only pairs and one mark at its leaves; its masters are $\Pth$, $\Gamma(\frac12)^4$ and $P_R(1)$ (Theorem~\ref{a:t:r4masters}).
\end{remark}

\renewcommand{\thesection}{17}
\setcounter{proposition}{0}
\renewcommand{\theproposition}{17.\arabic{proposition}}

\section{The masters of $\binom{[5]}2$; two marks doubled, closed}

Both lists are read at a mark (Prop.~19.19, p.~144): the period is the Breta of the even letters of the two children, its share side is a Mellin--Barnes integral over the split $d$ with the kernel of Lemma~5.16, p.~34, and the closing (Thm.~4.5, p.~26) reads a values input and a slope input on the lattice of the split. For $\binom{[5]}2$ the children are $\binom{[4]}2$ and $\binom{[4]}1$ (\S16.3). For the triangle with two marks doubled the mark is taken in a doubled class; the children are then those of the triangle $H$ with one class doubled, opened at its single mark: the three-mark banana $H/1$ and a pair with a coloop $H\setminus1$, and the class enters through its Mellin read. In both lists the lattice values of the input are leans of the three-mark banana. This entry finds the law of each input along its lines, the boundaries of its slope input, and the masters. Throughout, $\Pth$ is the three-mark centre and $\Pfour$ the four-mark centre.

\subsection*{17.1\quad The input of $\binom{[5]}2$ along its lines}

For $s\in\{1,2,\dots\}$ put $\Phi_s(d)=\Phi(d,s-d)=\int X_{/}^{-d}X_{\setminus}^{-(s-d)}\,dy$, the family of Proposition~\ref{a:p:u25split} on the line $s$; $\Phi_1$ is the $\Phi$ of Proposition~\ref{a:p:under}. By Proposition~\ref{a:p:leans} every lattice value is a finite sum of leans of the three-mark banana,
\[
\Phi_s(-m)=r_s(m)+c_s(m)\,\Pth,\qquad r_s(m),\,c_s(m)\in\mathbb Q\quad(m\in\mathbb N),
\]
and $c_s$ is the \emph{three-mark coordinate} of the values input.

\begin{proposition}[The closing at every integer strength]\label{a:p:u25clos}\derived\
For $s\in\{1,2,\dots\}$ let $V_s(m)=\big((m+1)(m+2)\cdots(m+s-1)\big)^2$. Then
\[
\Gamma(s)^2P_{\binom{[5]}2}(s\one;s,s)=\sum_{m\ge0}\Big(V_s(m)\,\Phi_s'(-m)-V_s'(m)\,\Phi_s(-m)\Big),
\]
and $x_s=\Gamma(s)^{-2}\sum_{m\ge0}V_s(m)\,\Phi_s(-m)$ is the lowered setting, the values half of the closing.
\end{proposition}

\begin{proof}
Proposition~\ref{a:p:u25split} and the left closing of Proposition~\ref{a:p:under} at the strength $s$. By the reflection (Prop.~1.3(v), p.~11), $\Gamma(d)^2\Gamma(s-d)^2=V_s(-d)\,\Gamma(d)^2\Gamma(1-d)^2$, with double poles at $d=-m$ and no simple parts; the residue at $d=-m$ is $\big(V_s(-d)\Phi_s(d)\big)'$ there.
\end{proof}

At $s=1$ this is Proposition~\ref{a:p:under}.

\begin{proposition}[The law on a line]\label{a:p:u25law}\derived\ in Proposition~\ref{a:p:certs}\
On the line $s=1$,
\[
\sum_{i=0}^4c_i(m)\,\Phi_1(-m-i)=0,
\]
\begin{align*}
c_0&=-8(m+1)^3(m+2)(m+3)(7m+16),\\
c_1&=(m+2)(m+3)(630m^4+4583m^3+12351m^2+14584m+6336),\\
c_2&=-(m+3)(2m+5)(1295m^4+11241m^3+35869m^2+49386m+24408),\\
c_3&=3(2m+5)(2m+7)(385m^4+3855m^3+13913m^2+21153m+11286),\\
c_4&=-54(m+2)(m+4)(2m+5)(2m+7)(2m+9)(7m+9),
\end{align*}
and $r_1$, $c_1$ obey it separately. Its symbol is
\[
216\rho^4-330\rho^3+185\rho^2-45\rho+4=(2\rho-1)(9\rho-4)(3\rho-1)(4\rho-1).
\]
The lines $s=2,3$ have laws of order four and degree eight with the same symbol.
\end{proposition}

\noindent Receipts \computed. Found by exact linear algebra on the values $\Phi_s(-m)$, $m\le80$, of Proposition~\ref{a:p:leans} (at $s=1$: $138$ equations, $35$ unknowns, one solution). At $m=-1$, outside the data, $c_0$ has a triple zero and the law holds exactly, $\sum_{i\ge1}c_i(-1)\Phi_1(1-i)=0$ with $(c_1,\dots,c_4)(-1)=(300,-5670,25920,-34020)$. It was read off first; Proposition~\ref{a:p:certs} derives it from the relations along a ray.

\begin{proposition}[The roots are the native values of the ratio of the strides]\label{a:p:u25roots}\derived\
On the chart $X_{/}/X_{\setminus}=e_1e_3/e_2^2$. At the table's gauge its value is $\frac49$, the maximum. On the pattern $(1,1,x,x)$ its native equation is $(1-x)(1+x^2)=0$, and at $(1,1,i,i)$ the value is $\frac12$. Where one mark goes to zero it becomes $e_1e_3/e_2^2$ of three marks, of maximum $\frac13$; where one mark goes to zero and another to infinity it becomes $\alpha_1\alpha_2/(\alpha_1+\alpha_2)^2$, the reciprocal stride of the two-mark banana, of maximum $\frac14$. These four values are the roots of the symbol of Proposition~\ref{a:p:u25law}.
\end{proposition}

The values input decays like $(\frac49)^m$; its coordinates $r_1$ and $c_1$ decay like $(\frac12)^m$, and
\[
\Pth=-\lim_{m\to\infty}\frac{r_1(m)}{c_1(m)}\quad\computed:
\]
the three-mark centre is the multiplier that removes the component along the native value $\frac12$ of the complex gauge.

\begin{proposition}[Lines tied by a first-order relation]\label{a:p:u25mix}\derived\ for $s\le4$ in Proposition~\ref{a:p:certs}\
With $F(m,n)=\Phi(-m,n)$,
\begin{multline*}
4(m+1)(3m+2n+4)\,F(m,n+1)-12(n+1)(2n+3)\,F(m+1,n+2)-(m+1)(2n-m)\,F(m,n)\\
+(6n^2+4mn-3m^2+16n-4m+5)\,F(m+1,n+1)-6(n+1)(2n+3)\,F(m+2,n+2)=0.
\end{multline*}
\end{proposition}

On the lines $n=m+s$ it ties the line $s+1$ to the line $s$, so with the law every line is read off the line $1$. Receipts \computed: found at $s\le4$, the coefficients polynomial in $s$; the lines $2$ to $5$ read off the line $1$ agree with Proposition~\ref{a:p:leans} at every $m\le80$.

\subsection*{17.2\quad The slope input and its boundaries}

The law holds as a function of the split (Proposition~\ref{a:p:certs}), so its derivative is the recurrence of the slope input,
\[
\sum_{i=0}^4c_i(m)\,\Phi_1'(-m-i)=\sum_{i=0}^4c_i'(m)\,\Phi_1(-m-i),
\]
as in Proposition~\ref{a:p:liverec}. Its four initial slopes $\Phi_1'(0),\dots,\Phi_1'(-3)$ are tied by the boundaries of the law.

\begin{proposition}[The cut]\label{a:p:u25cut}\derived\ from Proposition~\ref{a:p:u25law}\
At $m=-1$, where the four-mark centre $\Phi_1(1)=\Pfour$ sits, $c_0$ has a triple zero; neither $\Phi_1(1)$ nor its slope enters, and
\[
300\,\Phi_1'(0)-5670\,\Phi_1'(-1)+25920\,\Phi_1'(-2)-34020\,\Phi_1'(-3)=-\frac{21380}{63}+\frac{8866}{63}\,\Pth .
\]
The lines $2$ and $3$ give $\frac12$ and $-\frac59$ times the same relation.
\end{proposition}

\begin{proposition}[The wall at $d=2$]\label{a:p:u25wall}residue \derived; relation \computed\
$\Phi_1$ has a simple pole at $d=2$, the wall $2-d$ of the two-element $S$ (Proposition~\ref{a:p:under}), with residue $-6$: each of the six two-element $S$ collapses (Cor.~8.8, p.~56) to a product of two two-mark centres, $P_{\binom{[2]}1}(\one;1,1)^2=1$; the law at $m=-2$ returns the same $-6$. There the law meets the slopes only through
\[
c_2(-2)\,\Phi_1'(0)+c_3(-2)\,\Phi_1'(-1)=96\,\Phi_1'(0)-432\,\Phi_1'(-1),
\]
and
\[
6\,\Phi_1'(0)-27\,\Phi_1'(-1)=7\,\Pth-22
\]
(PSLQ, $140$ digits), which is the same as the finite part
\[
\Phi_1(d)=-\frac{6}{d-2}+\frac{21}2+\frac34\,\Pfour+O(d-2).
\]
\end{proposition}

\begin{proposition}[One free slope, along the three-mark coordinate]\label{a:p:u25free}\computed\
A third relation, $1680\,\Phi_1'(0)-24300\,\Phi_1'(-2)=2308\,\Pth-7207$ (PSLQ), with Propositions~\ref{a:p:u25cut} and~\ref{a:p:u25wall}, says that $\Phi_1'(-m)-\Phi_1'(0)\,c_1(m)\in\mathbb Q+\mathbb Q\,\Pth$ for $m\le3$. By the slope recurrence and Proposition~\ref{a:p:u25mix} it holds for every $m$ and every line:
\[
\Phi_s'(-m)=\rho_s(m)+\lambda_s(m)\,\Pth+c_s(m)\,\Phi_1'(0),\qquad \rho_s(m),\,\lambda_s(m)\in\mathbb Q .
\]
The one free direction of the slope input is the three-mark coordinate of the values input. Minimality fixes it: the slope input has no component along the native value $\frac12$.
\end{proposition}

\noindent Receipts \computed. With $\Phi_1'(0)$ from the masters below, $\Phi_1'(-m)/\Phi_1(-m)=0.81342$ at $m=600$, tending to $\log\frac94=0.81093$, where a component along $\frac12$ would grow like $(\frac98)^m$. The free slope is
\begin{gather*}
\Phi_1'(0)=\int\log\frac{X_{\setminus}}{X_{/}}\,\frac{dy}{X_{\setminus}}=2.794313984758270846154613553026697141639\ldots,\\
\Phi_1'(0)=2\,\Pth-\partial_\sigma P_{\binom{[3]}1}(\sigma\one;\sigma,\sigma)\big|_{\sigma=1}-\int\log X_{/}\,\frac{dy}{X_{\setminus}},\\
\int\log X_{/}\,\frac{dy}{X_{\setminus}}=9.9686527561522909453181622662\ldots,
\end{gather*}
by three-dimensional quadrature $2.7943139847583$; and by the same quadrature
\[
\Phi_1'(-1)=0.8280938236176\quad\text{against}\quad0.82809382361938566084\ldots .
\]
There is no rational relation among $1$, $\Pth$, its slope, $\Pfour$, the four-mark slope and $\Phi_1'(0)$ with coefficients below $10^{15}$.

\subsection*{17.3\quad The masters of $\binom{[5]}2$}

\begin{theorem}[The masters of $\binom{[5]}2$]\label{a:t:u25masters}(i) \derived\ from Propositions~\ref{a:p:u25clos} and~\ref{a:p:u25free}; (ii)--(iv) \computed\
\begin{enumerate}[label=(\roman*)]
\item Put
\begin{gather*}
A_s=\frac{\sum_m\big(V_s\rho_s-V_s'r_s\big)(m)}{\Gamma(s)^2},\quad B_s=\frac{\sum_m\big(V_s\lambda_s-V_s'c_s\big)(m)}{\Gamma(s)^2},\\
C_s=\frac{\sum_mV_s(m)c_s(m)}{\Gamma(s)^2},\quad R_s=\frac{\sum_mV_s(m)r_s(m)}{\Gamma(s)^2},
\end{gather*}
each convergent, the coordinates decaying like $(\frac12)^m$. Then
\[
P_{\binom{[5]}2}(s\one;s,s)=A_s+B_s\,\Pth+C_s\,\Phi_1'(0),\qquad x_s=R_s+C_s\,\Pth :
\]
the same $C_s$ weighs the three-mark centre in the lowered setting and the free slope in the period.
\item $C_1=\dfrac{\sqrt3}{80}\Gst44{\tfrac1{20},\ \tfrac3{20},\ \tfrac7{20},\ \tfrac9{20}}{\tfrac12,\ \tfrac12,\ \tfrac12,\ \tfrac12}$ \ ($200$ digits).
\item The ray law of \S15.4 and its corner are the boundary at the cut. As functions of the four initial slopes,
\begin{gather*}
\sum_{i=1}^4c^{\binom{[5]}2}_i(0)\,P(i)-1800=-\tfrac{55}{27}\,\beta,\qquad\sum_{i=0}^4c^{\binom{[5]}2}_i(1)\,P(1+i)=-\tfrac{898}{189}\,\beta,
\end{gather*}
with $\beta$ the left side minus the right side of Proposition~\ref{a:p:u25cut}, $P(i)=P_{\binom{[5]}2}(i\one;i,i)$, and $1800=60\cdot30$ the corner times $-[\sigma^4]c_0$. The lowered setting obeys the ray law with no corner.
\item A second corner lands on the three-mark centre:
\[
20\,P(1)-1672\,P(2)+5120\,P(3)=384-135\,\Pth .
\]
It is the ray law at $\sigma=-1$: $(c_2,c_3,c_4)(-1)=40\cdot(20,-1672,5120)$ and $\CT_{\sigma=-1}\big(c_0(\sigma)P(\sigma)+c_1(\sigma)P(\sigma+1)\big)=5400\,\Pth-15360$. In the coordinates of (i) it reads $5C_1-418C_2+1280C_3=0$, so the free slope drops out.
\end{enumerate}
Over $\mathbb Q+\mathbb Q\,\Pth$ the masters of $\binom{[5]}2$ are $P(1)$ and $P(2)$; the third master of \S15.4 is the three-mark centre.
\end{theorem}

\begin{proof}
(i) Insert Proposition~\ref{a:p:u25free} in Proposition~\ref{a:p:u25clos}; every coordinate is a combination of solutions of the law with the native values $\frac12,\frac49,\frac13,\frac14$, so each sum converges. (ii)--(iv) are identities among the sums of (i), checked to $140$ to $200$ digits.
\end{proof}

\begin{center}
\small
\begin{tabular}{@{}l l l@{}}
\toprule
& $s=1$ & $s=2$\\ \midrule
$A_s$ & $1.2617550081051847724190292271463942\ldots$ & $-0.3171370910356445997740190032938055\ldots$\\
$B_s$ & $-0.4062466940152853736690239421082401\ldots$ & $-0.0530255881622819812253694257861745\ldots$\\
$C_s$ & $1.3313334684984500811824592146018574\ldots$ & $0.1607947380347384161148704355054792\ldots$\\
$R_s$ & $0.3527480810189285064734282363418756\ldots$ & $-0.2668464215995376310247031065607639\ldots$\\
$P(s)$ & $4.0297141707110496073206957344284069\ldots$ & $0.0078868342010287786455991517685324\ldots$\\
$x_s$ & $3.4732702349419903397503533563416278\ldots$ & $0.1100415288232609996435971271184530\ldots$\\
\bottomrule
\end{tabular}
\end{center}

\noindent Receipts \computed. The four initial slopes solved from $P(1),P(2),P(3)$ and Proposition~\ref{a:p:u25cut} give $P(4)$ and $P(5)$ to $154$ digits against Procedure~\ref{a:p:proc}, and the relations of Propositions~\ref{a:p:u25wall} and~\ref{a:p:u25free} then come out of PSLQ with coefficients below $10^5$. Apart from (iv) there is no rational relation among $1$, $\Pth$, $\Phi_1'(0)$, $P(1)$, $P(2)$, $x_1$, $x_2$ with coefficients below $10^8$; $R_1$, $C_2$, $R_2$, $C_3$ are not in $\mathbb Q\,aC_1+\mathbb Q\,b\,\Gamma(\frac12)^2+\mathbb Q\,c\,\Gamma(\frac12)^4/C_1+\mathbb Q+\mathbb Q\,\Pth$ for any $a,b,c\in\{1,\sqrt3,\sqrt5,\sqrt{15}\}$, coefficients below $10^{12}$.

\subsection*{17.4\quad Two marks doubled, closed}

Let $H$ be the triangle with the class $2$ doubled. Opened at its mark $1$ (Prop.~19.19), its children on the chart of $\{2a,2b,3\}$ are the three-mark banana $H/1=\binom{[3]}1$, of stride $X_{H/1}=(B+\alpha_3)(B'+\alpha_3^{-1})$, and $H\setminus1$, the pair $\{2a,2b\}$ with $3$ a coloop, of stride $X_{H\setminus1}=BB'$, where $B=\alpha_{2a}+\alpha_{2b}$ and $B'=\alpha_{2a}^{-1}+\alpha_{2b}^{-1}$. As $X_{H/1}=BB'+B\alpha_3^{-1}+B'\alpha_3+1$, the deletion lies under the contraction. Put $\Psi(a,b)=\int X_{H/1}^{-a}X_{H\setminus1}^{-b}\,dy$ and $\Psi_s(d)=\Psi(d,s-d)$.

\begin{proposition}[Opened at a mark of a doubled class]\label{a:p:r4clos}\derived\
For the triangle $R$ with two marks doubled, $0<c<\sigma$,
\[
\Gamma(\sigma)^2P_R(\sigma\one;\sigma,\sigma)=\frac1{2\pi i}\int_{(c)}\Gamma(d)^2\Gamma(\sigma-d)^2\,P_{\binom{[2]}1}(d\one;d,d)\,\Psi_\sigma(d)\,dd,
\]
and for $s\in\{1,2,\dots\}$, with $W_s(x)=\prod_{j=1}^{s-1}(x-j)^2$ as in Proposition~\ref{a:p:cosetmark},
\[
\Gamma(s)^2P_R(s\one;s,s)=\sum_{x\ge s}\big(W_s\,f_s-W_s'\,q_s\big)(x),\qquad q_s(x)=P_{\binom{[2]}1}(x\one;x,x)\,\Psi_s(x),\quad f_s=-q_s' .
\]
At the lattice points, with $k=x-s$,
\[
q_s(x)=\frac1{\Gamma(2x)}\sum_{u=-k}^{k}\binom{2k}{k+u}\int_0^\infty y^{x-1-|u|}\,\sG_{|u|;(1,-1)}(1,y)^2\,\sG_{0;(1,-1)}(1,y)\,dy\ \in\ \mathbb Q+\mathbb Q\,\Pth .
\]
\end{proposition}

\begin{proof}
$R$ is $H$ with its mark $1$ replaced by a class of two. In the opening of $H$ at $1$ the mark contributes $\Gamma(d)^2$; the class contributes the Mellin read of the product of its two Grammas, $\Gamma(d)^2P_{\binom{[2]}1}(d\one;d,d)$ (Thm.~22.13(ii), p.~188, as in the proof of Proposition~\ref{a:p:triangles}). $\Psi_\sigma$ is holomorphic on $\re d>0$, its one wall being the coloop, and $X_{H\setminus1}/X_{H/1}<1$ while $P_{\binom{[2]}1}$ decays like $4^{-d}$: close to the right as in Proposition~\ref{a:p:cosetmark}. At $d=x$ the insertion $(BB')^k=\sum_u\binom{2k}{k+u}(\alpha_{2a}/\alpha_{2b})^u$ moves the masses (Lemma~35.2(iii), p.~308), and the share side of the three-mark banana at a lean (proof of Proposition~\ref{a:p:leans}) has the indices $(-u,u,0)$: $\sG_{-u}\sG_u\sG_0=y^{-|u|}\sG_{|u|}^2\sG_0$.
\end{proof}

\begin{proposition}[The input's law]\label{a:p:r4law}\derived\ in Proposition~\ref{a:p:certs}\
On the line $1$,
\begin{multline*}
x^2(2x-1)^2(60x^2+137x+65)\,q_1(x)-(2x+1)(1680x^5+4316x^4+2782x^3+231x^2-223x-56)\,q_1(x+1)\\
+(2x+1)(2x+3)(2940x^4+9113x^3+7350x^2-89x-1080)\,q_1(x+2)\\
-18(x+2)(2x+1)(2x+3)(2x+5)(60x^2+17x-12)\,q_1(x+3)=0,
\end{multline*}
with the symbol $36\rho^3-49\rho^2+14\rho-1=(\rho-1)(4\rho-1)(9\rho-1)$, and the lines are tied by
\begin{multline*}
-4(2x-2s-1)^2\,q_{s+1}(x)+2(2x+1)(16x-1)\,q_{s+1}(x+1)\\
+\big(4x^2-(8s+1)x+2s(2s+1)\big)\,q_s(x)-2(2x+1)(4x-1)\,q_s(x+1)=0 .
\end{multline*}
The roots are the native value $\frac14$ of the pair's period times the native values $4,1,\frac49$ of $X_{H\setminus1}/X_{H/1}$: $4$ at the sign $(1,1,-1)$ of the three-mark banana, where $BB'=4$ and $X_{H/1}=1$; $1$ where $\alpha_3\to0$ and the pair splits; $\frac49$ at the table's gauge \derived.
\end{proposition}

\noindent Receipts \computed: found on $x\le91$ ($176$ equations, $28$ unknowns) and at $s\le4$; the lines $2$ to $5$ read off the line $1$ agree with Proposition~\ref{a:p:r4clos}. The native value $1$ sits on the symbol, as it does for the five-mark law (Ledger, \ref{o:certs}): the coordinates of $q_1$ decay only like powers, $\Pth=-\lim r/c$ again, and the closing cannot be summed coordinate by coordinate. It is summed whole, with the minimality below.

\begin{proposition}[Zero shares]\label{a:p:r4zero}Laurent data and boundary \derived; second relation \computed\
At $d=0$,
\[
P_{\binom{[2]}1}(d\one;d,d)=\frac2d-\frac{\Gamma(\frac12)^4}{3}\,d+O(d^3),\qquad \Psi_1(d)=\frac2d+2+O(d) .
\]
The pole of $\Psi_1$ is the coloop's wall: each end of $\alpha_3$ gives $\int X_{H\setminus1}^{-1}dy=P_{\binom{[2]}1}(\one;1,1)=1$, and the finite part is $\int\log X_{H\setminus1}\,X_{H\setminus1}^{-1}dy=2$ on the chart of the pair. So $q_1(d)=4/d^2+4/d+O(1)$ where $c_0$ has a double zero: the law at $x=0$ returns the $4$, and its derivative is the boundary
\[
56\,f_1(1)-3240\,f_1(2)+6480\,f_1(3)=-\frac{1726}3-\frac{16}{15}\,\Pth .
\]
The next relation carries $\Gamma(\frac12)^4$, the class's term of order $d$:
\[
3\,f_1(1)-27\,f_1(2)=11+9\,\Pth-2\,\Gamma(\tfrac12)^4\qquad(\text{PSLQ},\ 130\text{ digits}).
\]
Hence $f_1(x)-f_1(1)\,c_1(x)\in\mathbb Q+\mathbb Q\,\Pth+\mathbb Q\,\Gamma(\frac12)^4$ for every $x$ and, by the lines, on every line: as for $\binom{[5]}2$ the one free slope lies along the three-mark coordinate. Minimality fixes it: no component along the native value $1$.
\end{proposition}

\noindent The free slope is
\begin{gather*}
f_1(1)=2\,\Pth-\partial_\sigma P_{\binom{[3]}1}(\sigma\one;\sigma,\sigma)\big|_{\sigma=1}-\int\log X_{H\setminus1}\,\frac{dy}{X_{H/1}},\\
f_1(1)=7.45882387041008265101442831798585979305\ldots,\\
\int\log X_{H\setminus1}\,\frac{dy}{X_{H/1}}=5.3041428705004791404583475012260708\ldots,
\end{gather*}
by two-dimensional quadrature $5.3041428705004785$ \computed.

\begin{theorem}[The masters of two marks doubled]\label{a:t:r4masters}\computed\
\begin{enumerate}[label=(\roman*)]
\item $P_R(\one;1,1)-8\,P_R(2\one;2,2)=6\,\Pth-2-\frac49\,\Gamma(\frac12)^4$.
\item $521\,P_R(\one;1,1)-12288\,P_R(2\one;2,2)+28672\,P_R(3\one;3,3)=648+1350\,\Pth$.
\item Through the closing, the corner relation of \S15.4 is the minimality alone, with the value $4432269600=-[\sigma^4]c_0\cdot44$; the second corner is $-\frac{1741824}{13}$ times the boundary of Proposition~\ref{a:p:r4zero} plus the minimality, with the value $-80640\cdot2503$. The second corner is the boundary of the input's law at zero shares, where the class and the coloop put the double pole $4/d^2+4/d$.
\item The lowered settings $x_s=\Gamma(s)^{-2}\sum_xW_s(x)q_s(x)$ obey the ray law with no corner; at $\sigma=-1$ they have the rational second corner $80640\cdot432$, and $239-148\,x_1+5244\,x_2-14336\,x_3=0$.
\end{enumerate}
With the corner and the second corner of \S15.4, every $P_R(s\one;s,s)$ lies in $\mathbb Q+\mathbb Q\,\Pth+\mathbb Q\,\Gamma(\frac12)^4+\mathbb Q\,P_R(\one;1,1)$: one master over $\mathbb Q+\mathbb Q\,\Pth+\mathbb Q\,\Gamma(\frac12)^4$, and the $3^\dagger$ of \S15.4 are $\Pth$, $\Gamma(\frac12)^4$ and $P_R(\one;1,1)$.
\end{theorem}

\noindent Receipts \computed. The three initial slopes solved from the minimality at $x=640$ and $P_R(1),P_R(2)$ give $P_R(3),P_R(4),P_R(5)$ to $155$ digits against Procedure~\ref{a:p:proc}; the slope at $x=548$ is then $1.7\cdot10^{-336}$, of the size $4^{-x}$. The values:
\begin{gather*}
P_R(\one;1,1)=7.8658373683640215924219746974461682\ldots,\\
x_1=2.5094223834929320566277011203297088\ldots .
\end{gather*}
There is no rational relation among $1$, $\Pth$, $\Gamma(\frac12)^4$, $f_1(1)$ and $P_R(\one;1,1)$, nor with $\partial_\sigma P_{\binom{[3]}1}$ and $\Pfour$ adjoined, below $10^{10}$.

\begin{remark}[What the closing found]\label{a:r:closingfound}
(i) In both lists the lattice values are three-mark leans, and the slope input has one free direction, the three-mark coordinate of the values; the other initial slopes are fixed by the boundaries of the law, at the cut and the wall for $\binom{[5]}2$, at zero shares for the triangle, and the free one by minimality along the leading native value: $\frac12$ at the complex gauge $(1,1,i,i)$ for $\binom{[5]}2$, $1$ at the sign of the three-mark banana for the triangle. The words of the masters, $V_s$ with $s\le3$ and $W_s$ with $s\le2$, have degree below what the adjoint of the law adds (six for $\binom{[5]}2$; five for the triangle, whose symbol vanishes at $\rho=1$), so summation by parts in the sense of Proposition~\ref{a:p:sbp} has nothing to reduce; the content is in the boundaries.

(ii) The corners of the ray laws are those boundaries. For $\binom{[5]}2$ the corner and the law itself are the boundary at the cut; for the triangle the corner is the minimality and the second corner is the boundary at zero shares.

(iii) Over $\mathbb Q$ the masters are $\{\Pth,P(1),P(2)\}$ for $\binom{[5]}2$ and $\{\Pth,\Gamma(\frac12)^4,P_R(1)\}$ for the triangle with two marks doubled. The three-mark centre is a master of both; in the triangle it enters through the contraction $H/1$.

(iv) The weight $C_1$ of the three-mark centre in the lowered setting of $\binom{[5]}2$ is a stacked Gamma at twentieths, and the same number weighs the free slope in the period.
\end{remark}

\begin{remark}[The leaf and the closing]\label{a:r:leaf}
In each of the four rays met so far, the three-mark ray, the four-mark ray, the input of $\binom{[5]}2$ along its lines and the input of the triangle with two marks doubled, the lattice values are $r+c\,\ell$ and the lattice slopes are $\rho+\lambda\,\ell+c\,\ell'$: $r,c,\rho,\lambda$ rational, the same $c$ in both, and $\ell$ the centre of the smallest piece, $\ell=\Pth$ for the three-mark ray and the two inputs and $\ell=\Pfour$ for the four-mark ray. For the two banana rays this is \derived: the cut relation and its derivative (the corner and its vanishing term of order $\sigma$, Proposition~\ref{a:p:corner}) have the same homogeneous part, so the one new slope $\ell'=\partial_\sigma P(\sigma\one;\sigma,\sigma)|_{\sigma=1}$ follows $c$ through the recursion (exactly, for every lattice point checked). For the two inputs it is Propositions~\ref{a:p:u25free} and~\ref{a:p:r4zero}.

In every case the coordinates $r$ and $c$ carry the leading native value of the law, and that native value sits at a gauge off the positive chart: the sign $(1,1,-1)$ of three marks (value $1$), the sign $(1,1,1,-1)$ of four marks (value $\frac14$), the complex gauge $(1,1,i,i)$ of $X_{/}/X_{\setminus}$ (value $\frac12$), the sign of the three-mark banana inside the triangle (value $1$). The leaf's centre is the multiplier that removes it,
\[
\ell=-\lim\frac rc\qquad\computed,
\]
and the true values decay like the extreme the positive chart reaches: $\frac19$ and $\frac1{16}$ at the gauges of the bananas, $\frac49$ at the table's gauge for $\binom{[5]}2$, and $\frac14$ for the triangle, from the stratum where $\alpha_3\to0$ and the pair splits.

So the closing pairs two kinds of number: Trailmoments of rational coordinate sequences against the closing's words, and the leaf's centre with its one slope. For $\binom{[5]}2$ this is Theorem~\ref{a:t:u25masters}(i). The stacked Gammas sit differently in the two prime lists. For the banana the whole lowered setting $x^{(1)}(1)=R+C\,\Pfour$ is the stacked Gamma at fifteenths (Proposition~\ref{a:p:stackedhalf}), while the weight $C=1.0705107798578621429652109212947808\ldots$ of the four-mark centre is not: there is no relation with that stacked Gamma, with $\Gamma(\frac12)^8$ over it, with $1$ and $\Pfour$, times $1,\sqrt3,\sqrt5,\sqrt{15}$, below $10^{12}$. For $\binom{[5]}2$ the weight $C_1$ of the three-mark centre is the stacked Gamma at twentieths and $x_1$ is not \computed.
\end{remark}

\subsection*{17.5\quad The laws from the relations along a ray}

The laws of \S\S17.1--17.4 were read off exact lattice values. They follow from the two relations along a ray (Thm.~35.3, p.~308), integration by parts in every mark and the shift by the stride. The relation space has finite codimension (Thm.~35.19, p.~313), so every insertion has a law that follows from it (Cor.~35.6, p.~308). Finding it is linear algebra, and what the linear algebra returns is an insertion that certifies the law.

\begin{proposition}[Certificates]\label{a:p:certs}\derived; the certificates \computed\
On the line $s$ of $\binom{[5]}2$ put $\langle g\rangle_m=\int g\,t^m\,X_{\setminus}^{-s}\,dy$ with the stride $t=X_{/}/X_{\setminus}=e_1e_3/e_2^2$, so that $\langle t^i\rangle_m=\Phi_s(-m-i)$ and $\langle t^iX_{\setminus}^{-1}\rangle_m=\Phi_{s+1}(-m-i)$. For each law and each line relation $\sum_jp_j(m)\langle g_j\rangle_m=0$ used in \S\S17.1--17.3 (the laws at $s\le3$; the first-order relations and the explicit relations of the lines at $s\le4$) there is an insertion
\[
H=\frac{P}{e_2^{\,d}},
\]
$P$ symmetric in three marks, of degree $2d$, with every exponent at most $d$ and coefficients in $\mathbb Q[m]$,
with $H_e$ its image under the exchange of the first mark with $e$, such that
\[
\sum_jp_j(m)\,g_j=\sum_{e\in\cE}\Big(\alpha_e\partial_{\alpha_e}H_e+H_e\big(m\,\alpha_e\partial_{\alpha_e}\log t-s\,\alpha_e\partial_{\alpha_e}\log X_{\setminus}\big)\Big).
\]
For the triangle with two marks doubled the same holds on the chart $\alpha_{2a}=1$ of $\{2a,2b,3\}$ for the laws of $\Psi_s$ (the law of Proposition~\ref{a:p:r4law} divided by the shift of $P_{\binom{[2]}1}$), with the stride $X_{H\setminus1}/X_{H/1}$, the origin form $X_{H\setminus1}^{-s}dy$, and $H_e=P_e/(L_1^{d_1}L_2^{d_2}L_4^{d_4})$, where $L_1,L_2$ are the letters of the three-mark banana, $L_4=\alpha_{2a}+\alpha_{2b}$, and $P_e$ lies inside the Newton polygon of the denominator.
\begin{center}
\small
\begin{tabular}{@{}l l l l@{}}
\toprule
& law & first-order relation & explicit relation\\ \midrule
$\binom{[5]}2$ & $d=7$, degree $5$ ($s=1$); $d=9$, degree $7$ ($s=2,3$) & $d=3$, degree $5$ & $d=5$, degree $5$\\
triangle & $(d_1,d_2,d_4)=(3,3,2)$, degree $6$ & $(2,2,1)$, degree $5$ & $(2,2,1)$, degree $5$\\
\bottomrule
\end{tabular}
\end{center}
The degrees are in the line coordinate.
\end{proposition}

\begin{proof}
The identity integrates to the law by Thm.~35.3, once $H_e\,t^m X_{\setminus}^{-s}$ tends to zero at the boundary of the chart. The hypothesis of Thm.~35.3 fails here. Along the line the contraction's exponent goes down as the deletion's goes up, and the direction has $\kappa_{1,S}=(d+d^\ast)_{\binom{[4]}1}(S)-(d+d^\ast)_{\binom{[4]}2}(S)$, which is $0$ for $|S|$ odd and $1$ for $|S|=2$; on the odd walls the stride does not decay but tends to the ratio of the face. The origin form carries the decay instead. Along a direction $w$ with $w_1\ge w_2\ge w_3\ge w_4$,
\begin{gather*}
\operatorname{ord}_wt=w_3-w_2\le0,\qquad \operatorname{ord}_wX_{\setminus}^{-s}=s(w_3+w_4-w_1-w_2)<0,\\
\operatorname{ord}_wH_e\le d(w_1+w_2)-d(w_1+w_2)=0,
\end{gather*}
the last because no exponent of $P$ exceeds $d$. So the boundary terms vanish for $m\ge0$. Both sides of the law are holomorphic in $m$ where the values converge ($\re m>-2$ on the line $1$), so the law holds identically in the split. For the triangle, $\operatorname{ord}_wH_e\le0$ by the Newton polygon, and $(X_{H\setminus1}/X_{H/1})^xX_{H\setminus1}^{-s}$ decays in every direction for $x>0$.
\end{proof}

\noindent Receipts \computed. Each certificate comes from exact linear algebra in under a second; the symmetric ansatz brings the law of the line $1$ of $\binom{[5]}2$ down to $444$ unknowns, and every equation is checked exactly. The certificates of the two line-$1$ laws were checked again by direct evaluation at random rational points, with the line coordinate negative and non-integral. The two banana rays, whose strides decay in every direction, have certificates of the same kind (Thm.~35.3 as stated gives the relations for large $k$; down to $k\to0^+$ it needs bounds on the pieces, which \S19, Step 5, supplies for three marks), with the denominators $e_1^2e_2$ (three marks) and $e_1^2e_3^2$ (four marks).

\begin{corollary}[What is now derived]\label{a:c:certs}\derived\
The laws of Propositions~\ref{a:p:u25law} and~\ref{a:p:r4law}, the line relations of Propositions~\ref{a:p:u25mix} and~\ref{a:p:r4law} for $s\le4$, the slope recurrences (their derivatives in the split), and the boundaries at the cut (Proposition~\ref{a:p:u25cut}) and at zero shares (Proposition~\ref{a:p:r4zero}, with its Laurent data) are derived. Still computed: the relations among the initial slopes found by PSLQ (Propositions~\ref{a:p:u25wall} and~\ref{a:p:u25free}, and the relation carrying $\Gamma(\frac12)^4$ in Proposition~\ref{a:p:r4zero}), the identities of the closings in Theorems~\ref{a:t:u25masters}(iii) and~\ref{a:t:r4masters}(iii), and the relations among the masters.
\end{corollary}

\begin{remark}[The centre read off the law]\label{a:r:paper}
With the laws derived, the exchange rate of Remark~\ref{a:r:leaf} is an exact identity at every lattice point, not only in the limit. On the three-mark centre ray $p(\sigma)=P_{\binom{[3]}1}(\sigma\one;\sigma,\sigma)$ the three-mark law (its certificate is among the receipts of Proposition~\ref{a:p:certs}) and the first two values,
\begin{gather*}
\sigma^2p(\sigma)-(10\sigma^2+10\sigma+3)\,p(\sigma+1)+9(\sigma+1)^2\,p(\sigma+2)=0,\\
p(1)=\Pth,\qquad p(2)=-\tfrac23+\tfrac13\,\Pth,
\end{gather*}
give $p(q)=r(q)+c(q)\,\Pth$ with $r,c$ rational, and $\Pth=-r(q)/c(q)+p(q)/c(q)$ for every $q$. The ratios $-r/c$ are
\[
0,\ \ 2,\ \ \tfrac{23}{10},\ \ \tfrac{145}{62},\ \ \tfrac{1331}{568},\ \ \tfrac{242353}{103400},\ \ \tfrac{5495507}{2344600},\ \ \tfrac{46422941}{19805800},\ \ \dots,
\]
all below $\Pth$ at every $q$, since $c(q)>0$ at every $q$ (\S19, Step 8), the error shrinking by $0.1112$ per step at $q=30$, towards $\frac19$, the native value of the gauge over that of the sign $(1,1,-1)$ \computed: about a digit of the three-mark centre per step of exact arithmetic. On the line $1$ of $\binom{[5]}2$, $\Phi_1(-m)=\int t^m X_{\setminus}^{-1}\,dy>0$ and $\Phi_1(0)=\Pth$, so $\Pth=-r_1(m)/c_1(m)+\Phi_1(-m)/c_1(m)$ \derived; the exact $c_1(m)$ is positive for $m\le80$, so these ratios are lower bounds as well, with error $6.7\cdot10^{-4}$ at $m=80$, shrinking per step towards $\frac89=\frac49\big/\frac12$ \computed. With $V_1=1$, Theorem~\ref{a:t:u25masters}(ii) is a series of the rationals that the law of Proposition~\ref{a:p:u25law} produces:
\[
\sum_{m\ge0}c_1(m)=1+\frac29+\frac{28}{405}+\frac{208}{8505}+\frac{304}{32805}+\frac{16672}{4546773}+\cdots=\frac{\sqrt3}{80}\Gst44{\tfrac1{20},\ \tfrac3{20},\ \tfrac7{20},\ \tfrac9{20}}{\tfrac12,\ \tfrac12,\ \tfrac12,\ \tfrac12}.
\]
\end{remark}

\renewcommand{\thesection}{18}
\setcounter{proposition}{0}
\renewcommand{\theproposition}{18.\arabic{proposition}}

\section{The reading of a ray through the origin}

Remark~\ref{a:r:paper} read the three-mark centre off its law, and Remark~\ref{a:r:leaf} found the same exchange rate on four rays. This entry states that reading for every ray through the origin of every family, as derived statements, and runs it where \S\S15--17 had not looked: off the centre, at a lean of three marks, and on the one connected list on four marks that had not been read. The receipts are in the script ray\_reading\_s18.py (python-flint, exact arithmetic and ball arithmetic, no quadrature in the procedure).

Throughout, a \emph{ray through the origin} is a ray of Def.~35.1, p.~307, with origin $0$ and an integral direction $(\mvec_1;\mathbf s_1,\mathbf t_1)$ on the balance wall, $s_{1j},t_{1j}\in\mathbb N$, all of whose proper walls are positive, $\kappa_{1,S}>0$ for every nonempty proper $S$. Its values and its stride are
\[
P(k)=P_F(k\mvec_1;k\mathbf s_1,k\mathbf t_1)=\int_{y_{e_0}=0}f^k\,dy\quad(k\ge1),\qquad f=\alpha^{\mvec_1}\prod_jK_j^{-s_{1j}}\Psi_j^{-t_{1j}},
\]
with $\alpha^\star$ the native gauge of the direction on the orthant, where $f$ takes its largest value $f(\alpha^\star)$ there (Thm.~35.10(i), p.~309). The centre ray of Def.~35.9, p.~309, is the case $\mvec_1=\sigma\one$, $\mathbf s_1=\mathbf t_1=\boldsymbol\sigma$.

\subsection*{18.1\quad One pencil, two chains}

\begin{proposition}[The closings of a direction]\label{a:p:dirclose}\derived\
Put $x_j=[\alpha^0]\,(1/f)^j$ for $j\in\mathbb N$, the constant terms of the powers of the Laurent polynomial $1/f=\alpha^{-\mvec_1}\prod_jK_j^{s_{1j}}\Psi_j^{t_{1j}}$, and
\[
\mathrm{Cor}=\sum_\pi\prod_{i=2}^n\kappa_{1,U_i}^{-1},
\]
the sum over the orders $\pi$ of $\cE$, with tails $U_i=\{\pi(i),\dots,\pi(n)\}$. Then $P$ continues meromorphically in $k$, with its poles at negative rational $k$, and
\[
\lim_{k\to-j}(k+j)^{n-1}P(k)=\mathrm{Cor}\cdot x_j\qquad(j\in\mathbb N).
\]
At the centre this is Proposition~\ref{a:p:negcorner}: there $1/f=\prod_j(K_j\Psi_j/\alpha^{\cE})^{\sigma_j}$ and $\kappa_{1,S}=\sum_j\sigma_jc_j(S)$.
\end{proposition}

\begin{proof}
In the sector coordinates of Thm.~8.7(ii), p.~55, $\alpha_{\pi(1)}=\rho$, $\alpha_{\pi(i)}=\rho\,q_2\cdots q_i$, one has $K_j=\rho^{r_j}\prod_iq_i^{d_j(U_i)}k_{j,\pi}$, $\Psi_j=\rho^{n-r_j}\prod_iq_i^{d^\ast_j(U_i)}\psi_{j,\pi}$ and $\alpha^{\mvec_1}=\rho^{m_{1,\cE}}\prod_iq_i^{m_1(U_i)}$, so by the balance of the direction
\[
f=\prod_{i=2}^nq_i^{\kappa_{1,U_i}}\,\varphi_\pi(q),\qquad \varphi_\pi=\prod_jk_{j,\pi}^{-s_{1j}}\psi_{j,\pi}^{-t_{1j}},
\]
and by Thm.~8.10, p.~56, $P(k)=\sum_\pi\int_{[0,1]^{n-1}}\prod_iq_i^{k\kappa_{1,U_i}-1}\varphi_\pi(q)^k\,dq$. The sector letters have nonnegative coefficients and positive constant terms (Lemma~8.5, p.~55), so $(1/\varphi_\pi)^j=\sum_vc_{\pi,v}q^v$ is a polynomial. At $k=-j+\varepsilon$ write $\varphi_\pi^{k}=(1/\varphi_\pi)^j\varphi_\pi^{\varepsilon}$. By Lemma~8.6, p.~55, the term $q^v$ has a pole in $\varepsilon$ of order at most the number of $i$ with $v_i=j\kappa_{1,U_i}$, so the pole of order $n-1$ comes from $v=j\kappa_1$ alone, with the leading coefficient $c_{\pi,j\kappa_1}\prod_i\kappa_{1,U_i}^{-1}$. Finally $c_{\pi,j\kappa_1}=[q^0]\big(\prod_iq_i^{-\kappa_{1,U_i}}/\varphi_\pi\big)^j$ is the constant term of $(1/f)^j$ in the sector coordinates, and $\alpha^v\mapsto\prod_iq_i^{v(U_i)}$ is injective on exponents of degree zero, so it is $x_j$ for every $\pi$.
\end{proof}

\begin{corollary}[The law and its closings]\label{a:c:dirlaw}\derived\
If $\sum_{i=0}^hc_i(k)P(k+i)=0$ identically, then $\sum_{i=0}^{\min(h,j)}c_i(-j)\,x_{j-i}=0$ for every $j\in\mathbb N$, and $c_0$ vanishes to order at least $n-1$ at $0$, with
\[
[k^{n-1}]c_0\cdot\mathrm{Cor}+\sum_{i\ge1}c_i(0)\,P(i)=0 .
\]
\end{corollary}

\begin{proof}
As Corollary~\ref{a:c:law}, with Proposition~\ref{a:p:dirclose} in place of Proposition~\ref{a:p:negcorner}; the second statement is the case $j=0$, $x_0=1$.
\end{proof}

So a law of the values is the recurrence of the integers $x_j$ read backwards, and reading it is exact linear algebra on integers. The converse, that the least recurrence of the $x_j$, reflected, is a law of $P$, is what a certificate of Thm.~35.3, p.~308, settles, as in Proposition~\ref{a:p:certs}.

\begin{remark}[The torus and the chart]\label{a:r:twochains}\derived\
The two sequences are one integrand on two chains. For $|q|<f(\alpha^\star)^{-1}$, $\sum_{k\ge1}P(k)q^{k-1}=\int_{y_{e_0}=0}f/(1-qf)\,dy$ (Thm.~35.10(ii) with $k_0=1$). On the torus $T=\{|\alpha_e|=1\}$, where $|1/f|\le(1/f)(\one)$ because $1/f$ has nonnegative coefficients, expanding in $1/(qf)$ gives for $|q|>f(\one)^{-1}$
\[
\sum_{j\ge0}x_j\,q^{-j-1}=-\frac1{(2\pi i)^{n-1}}\oint_{T}\frac{f}{1-qf}\prod_{e\ne e_0}\frac{d\alpha_e}{\alpha_e}.
\]
The pencil $\prod_jK_j^{s_{1j}}\Psi_j^{t_{1j}}-q\,\alpha^{\mvec_1}$ of Thm.~35.10 carries both: the positive chart gives its solution at $q=0$, whose coefficients are the periods, and the torus its solution at $q=\infty$, whose coefficients are integers for a list. Corollary~\ref{a:c:dirlaw} is the one equation, inhomogeneous by its boundary terms, read at its two ends, and $\mathrm{Cor}$ is the ratio of the two normalizations.
\end{remark}

\subsection*{18.2\quad The slit}

\begin{proposition}[The ray lives on the slit plane]\label{a:p:slit}\derived\
$G(q)=\int_{y_{e_0}=0}f/(1-qf)\,dy$ is holomorphic on $\mathbb C\setminus[f(\alpha^\star)^{-1},\infty)$, and it continues $\sum_{k\ge1}P(k)q^{k-1}$.
\end{proposition}

\begin{proof}
On the chart $0<f\le f(\alpha^\star)$ (Thm.~35.10(i)). For $q$ off the slit the segment $q\cdot(0,f(\alpha^\star)]$ misses $1$, so $|1-qf|$ is bounded below by the distance from $1$ to that segment, locally uniformly in $q$, and $\int f\,dy=P(1)<\infty$; dominated convergence gives holomorphy. On $|q|<f(\alpha^\star)^{-1}$ the integral is the series.
\end{proof}

\begin{theorem}[Positive selection]\label{a:t:select}(i)--(iii) \derived; the count in (iii) \computed\ on every ray below.\
Let $\sum_{i=0}^hc_i(k)P(k+i)=0$ be a law of a ray through the origin, $c_0c_h\ne0$, let $D$ be the largest degree of the $c_i$, and $\chi(\lambda)=\sum_i[k^D]c_i\,\lambda^i$ its symbol, with roots $\lambda_c$; it vanishes at the native value of every nondegenerate native gauge (Cor.~35.6, p.~308), and a root need not be a native value.
\begin{enumerate}[label=(\roman*)]
\item For every solution $(x_k)_{k\ge1}$ of the law, $X(q)=\sum_kx_kq^{k-1}$ satisfies a linear differential equation in $q$ with polynomial coefficients, inhomogeneous by a polynomial that depends linearly on $x_1,\dots,x_h$; in $\mathbb C^\times$ its singular points are the $q_c=1/\lambda_c$. So $X$ continues along every path in $\mathbb C^\times$ that avoids them.
\item The values are the solution whose $X$ is holomorphic on $\mathbb C\setminus[f(\alpha^\star)^{-1},\infty)$ (Proposition~\ref{a:p:slit}): holomorphic at every $q_c$ off the slit, that is at every root $\lambda_c\notin(0,f(\alpha^\star)]$.
\item At a singular point $q_c$, holomorphy of $X$ is $\nu_c$ linear conditions on $(x_1,\dots,x_h)$, $\nu_c$ the number of independent local solutions of the equation that are not holomorphic there, at most the multiplicity of $\lambda_c$ as a root of $\chi$. At a simple root that is the native value of a nondegenerate native gauge it is one: the branch of Thm.~35.11, p.~310, read at that gauge in place of $\alpha^\star$, of exponent $\frac{n-3}2$, with a logarithm for odd $n$. With the corner of Corollary~\ref{a:c:dirlaw}, if $1+\sum_{q_c\text{ off the slit}}\nu_c=h$ and the conditions are independent, the law, the corner and positivity determine the values.
\item Suppose further that every root $\lambda_c\neq f(\alpha^\star)$ that carries a local solution not holomorphic at $q_c$ has $|\lambda_c|>f(\alpha^\star)$; a root in $(-f(\alpha^\star),0)$, or a complex one in the disc $|\lambda|<f(\alpha^\star)$, carries a solution smaller than $P$, and the limits below diverge. Then, for a basis $x^{(0)},\dots,x^{(h-1)}$ of rational solutions fixed by the corner, $P=x^{(0)}+\sum_{i\ge1}M_ix^{(i)}$, and the $M_i$ are the limits as $K\to\infty$ of the solutions of $x^{(0)}_{K+l}+\sum_iM_ix^{(i)}_{K+l}=0$, $0\le l<h-1$, with error $O\big((f(\alpha^\star)/\min_{c\ne\star}|\lambda_c|)^K\big)$ up to powers of $K$.
\end{enumerate}
\end{theorem}

\begin{proof}
(i) Multiply the law by $q^{k-1}$ and sum over $k\ge1$: $k$ acts as $q\partial_q+1$, a shift by $i$ as $q^{-i}$ after removing the first $i$ terms, which leaves a polynomial. The coefficient of the top power of $q\partial_q$ is $\sum_i[k^D]c_i\,q^{-i}$, which vanishes in $\mathbb C^\times$ exactly where $\chi(1/q)=0$. (ii) Proposition~\ref{a:p:slit}. (iii) The inhomogeneity is holomorphic at $q_c$, so the singular part of $X$ there lies in the span of the local homogeneous solutions that are not holomorphic, with coefficients linear in $(x_1,\dots,x_h)$. At a nondegenerate native gauge the value over a steepest-descent disc (Thm.~35.5, p.~308) gives one branch of exponent $\frac{n-1}2-1$; Thm.~35.11 is this computation at $\alpha^\star$, and it is local. (iv) The coefficients of $X$ grow at the rate $|\lambda_c|$ of the singular point nearest $0$ at which $X$ is singular. The rational solutions are singular at the $q_c$ off the slit and $P$ is not, so the combination annihilating their $h-1$ largest rates is $P$, and its own size $f(\alpha^\star)^K$ against the smallest of the others is the error.
\end{proof}

\begin{procedure}[The reading of a ray]\label{a:p:reading}
(1) From the family and the direction write $1/f$ and the walls $\kappa_{1,S}$, and $\mathrm{Cor}$ (Proposition~\ref{a:p:dirclose}). (2) The closings $x_j=[\alpha^0](1/f)^j$, exactly. (3) Their least recurrence by exact linear algebra, reflected (Corollary~\ref{a:c:dirlaw}): the law; its certificate by Thm.~35.3 as in Proposition~\ref{a:p:certs}. (4) The native values, the roots of the symbol, and the native gauges, the critical points of $\log f$: the signs and what a lean moves them to, and families. (5) The seeds: the corner, and a second corner wherever $c_0$ and $c_1$ share a zero (\S15.4). (6) Positivity, Theorem~\ref{a:t:select}(iii)--(iv): the masters as limits of rational solutions, at the rate $f(\alpha^\star)/\min_{c\ne\star}|\lambda_c|$. (7) Integer relations among the masters and against the members already in use; what is left over $\mathbb Q$ are the masters of the ray.
\end{procedure}

\begin{remark}[What sits on the slit]\label{a:r:onslit}
Positivity says nothing at a native value in $(0,f(\alpha^\star))$. For a ray strictly inside its walls $f$ has one critical point on the orthant (Thm.~8.18(ii)--(iii), p.~58), so such a value would sit at a native gauge off the orthant; on every ray strictly inside its walls read in \S\S15--18 there is none. Example~\ref{b:e:Hslit} has one, a conjugate of an irrational $f(\alpha^\star)$; the gauge end reads that ray all the same (\S21). On a ray along walls the slit carries the faces: on the line $1$ of $\binom{[5]}2$ (\S17.1) the native values $\frac13$ and $\frac14$ of Proposition~\ref{a:p:u25roots} lie below $\frac49$, they come from strata at the boundary of the chart, and the constants there are the faces' own, the three-mark centre among them (Theorem~\ref{a:t:u25masters}). That is Thm.~35.15, p.~312, read on the slit.
\end{remark}

\subsection*{18.3\quad Off the centre: a lean of three marks}

The centre is the fixed point of Cremona (Thm.~8.13, p.~57), where both letters take one exponent; a lean gives them two. For $\binom{\cE}1$, $K=\sum_e\alpha_e$ and $\Psi=\alpha^{\cE}\sum_e\alpha_e^{-1}$, and at uniform masses $m\one$ the balance is $nm=s+(n-1)t$.

\begin{proposition}[The lean is the index of the Gramma]\label{a:p:leanindex}\derived; receipt \computed\
For masses $\mvec$ on the balance wall and $\re s,\re t>0$,
\[
\Gamma(s)\Gamma(t)\,P_{\binom{\cE}1}(\mvec;s,t)=\int_0^\infty y^{t-1}\prod_{e\in\cE}\sG_{m_e-t;(1,-1)}(1,y)\,dy ,\qquad \sum_e(m_e-t)=s-t .
\]
At the centre this is Proposition~\ref{a:p:bananapin}; the lean and the offsets move the index slots of the marks' Grammas. In the direction $(2\one;4,1)$ of three marks, $P_{\binom{[3]}1}(2\one;4,1)=\frac16\int_0^\infty\sG_{1;(1,-1)}(1,y)^3\,dy=0.030080945282262304399884274898349847\ldots$ in ball arithmetic.
\end{proposition}

\begin{proof}
On the share side (Proposition~\ref{a:p:share}) the mark $e$ carries the share $a_e$ of the basis $\{e\}$ in $K$ and the share $b_e$ of the cobasis $\cE\setminus e$ in $\Psi$, and the pin at $e$ reads $a_e+\sum_{e'\ne e}b_{e'}=m_e$, that is $a_e=b_e+m_e-t$. So each mark carries $\Gamma(b_e+m_e-t)\Gamma(b_e)$, the share kernel of $\sG_{m_e-t;(1,-1)}(1,y)$ (Prop.~4.2, p.~25), and the Mellin read of the product pins $\sum_eb_e=t$, as in Proposition~\ref{a:p:bananapin}.
\end{proof}

\begin{proposition}[The native gauges of a lean of $\binom{\cE}1$]\label{a:p:leangauges}\derived\
In the direction $(m\one;s,t)$, $s\ne t$, put $u_e=\alpha_e/K(\alpha)$, $v_e=\alpha_e^{-1}/K(\alpha^{-1})$ and $\lambda=\big(K(\alpha)K(\alpha^{-1})\big)^{-1}$. At a native gauge $su_e-tv_e=\frac{s-t}n$ and $u_ev_e=\lambda$ for every $e$, so every $u_e$ is a root of $su^2-\frac{s-t}nu-t\lambda=0$; with $j$ marks at one root $u_-$ and $n-j$ at the other $u_+$,
\[
(n-2j)\,u_+=1-j\,\frac{s-t}{ns},\qquad f=\Big(\prod_eu_e\Big)^{\frac{s-t}n}\lambda^t .
\]
At $s=t$ the two roots are $\pm u_+$ and this is Proposition~\ref{a:p:onlysigns}. For three marks in the direction $(2\one;4,1)$: $j=0$ is the table's gauge, $f=\frac1{243}$; $j=1$ is $\alpha\propto(3,3,-2)$ and its images, $f=-\frac{27}{64}$, the sign $(1,1,-1)$ moved along the real line.
\end{proposition}

\begin{proof}
Nativity is $m=sR_e(\alpha)+tR^\ast_e(\alpha)$ with $R_e=u_e$ and $R^\ast_e=1-v_e$ (Thm.~7.14, p.~50), and $u_ev_e=1/\big(K(\alpha)K(\alpha^{-1})\big)$ for every $e$. The $u_e$ sum to $1$, which fixes $u_+$. For the value, $\alpha^{\one}=K(\alpha)^n\prod_eu_e$ and $K(\alpha^{-1})=1/(K(\alpha)\lambda)$, and the powers of $K(\alpha)$ cancel by the balance.
\end{proof}

\noindent Receipts \computed. The walls of $(2\one;4,1)$ are $2$ on marks and $3$ on pairs, so $\mathrm{Cor}=6/(2\cdot3)=1$. The closings are $x_j=(4j)!\,j!\sum_{b_1+b_2+b_3=j}\prod_f\big(b_f!\,(b_f+j)!\big)^{-1}=1,36,4620,762300,140216076,\dots$; sixty of them have a least recurrence of order $2$ and degree $6$, and reflected it is the law
\begin{gather*}
-3k^2(3k+1)(3k+2)(265k^2+636k+383)\,P(k)\\
+(1721705k^6+7575502k^5+13047276k^4+11068046k^3+4766659k^2+965100k+75600)\,P(k+1)\\
+972(k+1)^2(4k+5)(4k+7)(265k^2+106k+12)\,P(k+2)=0,
\end{gather*}
$P(k)=P_{\binom{[3]}1}(2k\one;4k,k)$, with the symbol $265(243\lambda-1)(64\lambda+27)$: the two native values of Proposition~\ref{a:p:leangauges}, which were read off the list before the closings were computed. The corner gives $12600\,P(1)+68040\,P(2)=383$; positivity gives
\[
P_{\binom{[3]}1}(2\one;4,1)=0.0300809452822623043998842748983498467336738511\ldots,
\]
about two digits per step (the rate $\frac{64}{6561}$), against $0.95$ for the centre, and Proposition~\ref{a:p:leanindex} confirms it. Over $\mathbb Q$,
\[
54\,P_{\binom{[3]}1}(2\one;4,1)=11-4\,\Pth\qquad(380\text{ digits}).
\]
The Cremona partner $(3\one;1,4)$ has the same closings, by $\alpha\mapsto\alpha^{-1}$, and the same law. So the lean is a new coordinate system on the content of the centre: a different law, different native gauges, a different corner, the same master.

The rational coordinates of this ray do not give a gap (Thm.~34.3, p.~300): on $(1,\Pth)$ their denominators beyond $d_k^2$ grow like $e^{8.7k}$, with primes of the apparent factor $265k^2+\cdots$ among them, against the decay $243^{-k}=e^{-5.49k}$ \computed; at the centre the excess is $3^{1.88k}$, against $9^{-k}$.

\subsection*{18.4\quad The fourth list on four marks}

On four marks the connected lists are $\binom{[4]}1$, $\binom{[4]}2$, $\binom{[4]}3$ and the triangle $H$ with one mark doubled; the first two were read in \S15.1, and the third is the dual of the first (Prop.~8.14, p.~57). $H$ is the list of \S17.4, where it was a step toward the triangle with two marks doubled; its own centre was not read there. With the doubled class $\{a,b\}$ and the marks $2,3$ its bases are the pairs other than $\{a,b\}$, and with $A=\alpha_a+\alpha_b$, $B=\alpha_2+\alpha_3$
\[
K=AB+\alpha_2\alpha_3,\qquad \Psi=AB+\alpha_a\alpha_b :
\]
the letters exchange under $\{a,b\}\leftrightarrow\{2,3\}$, so $H$ is self-dual up to that relabeling. The widths are $1$ on marks, on $\{a,b\}$, on $\{2,3\}$ and on triples, and $2$ on the four other pairs (Def.~7.10, p.~49), so at the centre $\mathrm{Cor}=8\cdot1+16\cdot\frac12=16$.

\begin{proposition}[The native gauges of the centre of $H$]\label{a:p:Hgauges}\derived; that there are no others \computed\
The table's gauge, of native value $\frac1{25}$; the sign $(-1,-1,1,1)$, of native value $\frac19$; and, of native value $1$, the three curves
\[
\{A=0,\ \alpha_a^2=\alpha_2\alpha_3\},\qquad \{B=0,\ \alpha_2^2=\alpha_a\alpha_b\},\qquad \{A=B=0\},
\]
the second the image of the first under the self-duality; on them the class letters vanish and $K$, $\Psi$ collapse to monomials, the case of Proposition~\ref{a:p:families}(ii) with one class. Every sign of value $1$ lies on them.
\end{proposition}

\begin{proof}
At the centre nativity is $R_e(\alpha)=R_e(\alpha^{-1})$ (\S15.2). On $A=0$ also $A(\alpha^{-1})=A/(\alpha_a\alpha_b)=0$, so $K(\alpha)=\alpha_2\alpha_3$, $K(\alpha^{-1})=1/(\alpha_2\alpha_3)$ and $R_2=R_3=1$ at $\alpha$ and at $\alpha^{-1}$, while $R_a(\alpha)=\alpha_aB/(\alpha_2\alpha_3)$ and $R_a(\alpha^{-1})=B/\alpha_a$ agree exactly when $B=0$ or $\alpha_a^2=\alpha_2\alpha_3$, and likewise at $b$; there $K(\alpha)K(\alpha^{-1})=1$. The second curve follows by the self-duality. At a sign $K(\eta)=(\eta_a+\eta_b)(\eta_2+\eta_3)+\eta_2\eta_3\in\{\pm5,\pm3,\pm1\}$ (Proposition~\ref{a:p:signsparallel}), and $K(\eta)=\pm1$ forces $\eta_a=-\eta_b$ or $\eta_2=-\eta_3$, which puts the sign on one of the curves. Newton's method from several thousand starts in $(\mathbb C^\times)^3$ finds no critical point of $\log f$ off these.
\end{proof}

\noindent Receipts \computed. By Proposition~\ref{a:p:parallel} the closings are
\[
x_k=\sum_{i+j+l=k}\Big(\frac{k!}{i!\,j!\,l!}\Big)^2\binom{2i+2j}{i+j}=1,5,53,761,12661,229705,4410665,\dots,
\]
and the least recurrence of $160$ of them has order $4$ and degree $7$ (156 equations, 40 unknowns, one solution). Reflected, the law of the centre of $H$, $P(k)=P_H(k\one;k,k)$, is
\begin{gather*}
k^3(k+1)(144k^3+864k^2+1693k+1075)\,P(k)\\
-2(k+1)(2592k^6+19440k^5+56322k^4+80296k^3+60004k^2+23017k+3580)\,P(k+1)\\
+2(21168k^7+211680k^6+884919k^5+2003456k^4\\
\qquad+2646984k^3+2031207k^2+830703k+137260)\,P(k+2)\\
-2(k+2)^2(34848k^5+261360k^4+741842k^3+984642k^2+598948k+127215)\,P(k+3)\\
+225(k+2)^2(k+3)^2(144k^3+432k^2+397k+102)\,P(k+4)=0,
\end{gather*}
with the symbol $144(25\lambda-1)(9\lambda-1)(\lambda-1)^2$: the curves double the root $1$, an integral direction in the exceptional set of Thm.~35.19(i), p.~313. (A search by total size first returns a recurrence of order $5$ and degree $5$: this law composed with a factor of order one, whose symbol has the extra root $\frac{131003}{4771109}$.) The corner relation is $17200-7160P(1)+274520P(2)-1017720P(3)+826200P(4)=0$, and $c_0$, $c_1$ share the zero $k=-1$; the second corner (\S15.4) is rational,
\[
57\,P(1)-807\,P(2)+1050\,P(3)=196\qquad(530\text{ digits}),
\]
the constant term at $k=-1$ being $1176$. Positivity (Theorem~\ref{a:t:select}: three conditions, one at $\frac19$ and two at the double root) gives
\begin{gather*}
P_H(\one;1,1)=4.096236998337882850055553228351268637634670896\ldots,\\
P_H(2\one;2,2)=0.04771649869151311029741208089066366661776262936\ldots .
\end{gather*}

\begin{proposition}[The centre of $H$ on its share side]\label{a:p:Hshare}\derived; receipt \computed\
$H$ is a series root over the parallel cell of $\{a,b\}$ and the two marks left in place (Def.~22.17(5), p.~191), so, as in Proposition~\ref{a:p:triangles},
\[
\Gamma(\sigma)^2P_H(\sigma\one;\sigma,\sigma)=\frac1{2\pi i}\int_{(c)}\Gamma(d)^2\,\Gamma(\sigma-d)^2\,P_{\binom{[2]}1}(d\one;d,d)^2\,dd,\qquad P_{\binom{[2]}1}(d\one;d,d)=\Gst21{d,\ d}{2d},
\]
$0<c<\sigma$. At $\sigma=1$, on $c=\frac12$, the integral is $4.0962369983378828500555532283512686376346709\ldots$ in ball arithmetic: the value of the positive selection, to $44$ digits.
\end{proposition}

\begin{proof}
A series node multiplies the transforms of its cells (Thm.~22.13, p.~188): the two marks give $\sG_{0;(1,-1)}(1,y)^2$, whose Mellin read is $\Gamma(z)^2P_{\binom{[2]}1}(z\one;z,z)$ (Proposition~\ref{a:p:bananapin}), and the parallel cell gives the transform whose Mellin read is $\Gamma(w)^2P_{\binom{[2]}1}((\sigma-w)\one;\sigma-w,\sigma-w)$ (Proposition~\ref{a:p:triangles}); the Mellin read of the product pins the sum of the shares, and $d=\sigma-w$.
\end{proof}

The integrand is $\Gamma(d)^6\Gamma(\sigma-d)^2/\Gamma(2d)^2$. Its poles at $-\mathbb N$ have order four and residues that grow like $16^m$, and to the right the double poles of $\Gamma(\sigma-d)^2$ meet a factor that grows faster than every power: neither closing of Thm.~4.5, p.~26, converges. The centre of $H$ is a fibre integral of stacked Gammas that does not land on a lattice, as the five-mark centre does not (Remark~\ref{a:r:live}).

\noindent The lean. The direction $(2\one;3,1)$ has the walls $2$ on marks, $3$ on $\{a,b\}$, $1$ on $\{2,3\}$, $4$ on the other pairs and $2$ on triples, so $\mathrm{Cor}=\frac13+1+1=\frac73$. On the slice $\alpha_a=\alpha_b=x$, $\alpha_2=\alpha_3=1$ the stride is $x^3/\big((4x+1)^3(x+4)\big)$, whose critical points are the roots of $2x^2-x-6=(x-2)(2x+3)$ \derived: the native gauge of the direction on the orthant is $(2,2,1,1)$, where $K=9$, $\Psi=12$ and $f=\frac4{2187}$, and the sign has moved to $(-\frac32,-\frac32,1,1)$, $f=\frac{27}{2500}$; the curves move to curves of value $-\frac{16}{27}$ \computed. The closings $x_j=[\alpha^{2j\one}]K^{3j}\Psi^j=1,48,10240,3131184,\dots$ have a least recurrence of order $4$ and degree $25$ ($196$ equations, $104$ unknowns, one solution), with the symbol a multiple of $(2187\lambda-4)(2500\lambda-27)(27\lambda+16)^2$; positivity converges at the rate $\frac{10000}{59049}$, about $0.77$ digits per step, and the lean's own second corner is rational: $413667900\,P'(1)+6131684550\,P'(2)+12031031250\,P'(3)=24789199$, $P'(k)=P_H(2k\one;3k,k)$.

\begin{theorem}[One content, two rays]\label{a:t:Hcontent}\computed\
\begin{gather*}
28\,P_H(2\one;3,1)=14-3\,P_H(\one;1,1)-P_H(2\one;2,2),\\
5670000\,P_H(4\one;6,2)=126002\,P_H(\one;1,1)+143398\,P_H(2\one;2,2)-522781,\\
95681250000\,P_H(6\one;9,3)=3048368669-731189098\,P_H(\one;1,1)-1115791502\,P_H(2\one;2,2),
\end{gather*}
to $530$ digits. The centre ray and the lean ray of $H$, with two laws, two corners and two positive selections, land on the one space $\mathbb Q+\mathbb Q\,P_H(\one;1,1)+\mathbb Q\,P_H(2\one;2,2)$.
\end{theorem}

\begin{remark}[What the masters of $H$ are not]\label{a:r:Hnot}\computed\
There is no rational relation among $1$, $P_H(\one;1,1)$, $P_H(2\one;2,2)$ and the members in use, $\Pth$, $\Gamma(\frac12)^4$, $\Tx{3}{\IZ}{0}$, $\Gamma(\frac12)^6/\sqrt3$, $\sqrt3\,\Gamma(\frac12)^2$, $\sqrt3\,\Gamma(\frac12)^2\log3$, $\Gamma(\frac12)^4\log2$, $\Gamma(\frac12)^4\log3$, $\log2$, $\log3$, $\Pth\log2$, $\Pth\log3$, $\Gamma(\frac13)^6/\Gamma(\frac12)^4$ and $\Gamma(\frac12)^8/\Gamma(\frac13)^6$, with coefficients below $10^{24}$ ($420$ digits); nor with $\Pth$, the two masters of $\binom{[5]}1$, the stacked Gamma $X$ at fifteenths of Proposition~\ref{a:p:stackedhalf}, $\Gamma(\frac12)^8/X$, $\Gamma(\frac12)^4$, $\sqrt3\,\Gamma(\frac12)^2$, $\sqrt{15}$, and $\sqrt{15}$ times $X$, $\Gamma(\frac12)^8/X$ and $\Gamma(\frac12)^4$, below $10^{28}$ ($380$ digits). In that framework, over $\mathbb Q$, the two masters of $H$ are new numbers. $H$ and $\binom{[5]}1$ have the same native values $\frac1{25},\frac19,1$ at their centres and different masters: the native values do not decide the numbers.
\end{remark}

\noindent What the tree expects \notyet. Opened at the mark $2$ (Prop.~19.19, p.~144), $H$ has the children $\binom{[3]}1$ and the pair $\{a,b\}$ with the coloop $3$, the strides $X_{H/1}$ and $X_{H\setminus1}$ of \S17.4, and its input is $\Psi_\sigma$ of \S17.4 without the factor $P_{\binom{[2]}1}$. Its lattice values are three-mark leans, in $\mathbb Q+\mathbb Q\,\Pth$, and its slopes should have one free direction (Proposition~\ref{a:p:r4zero}). As for $\binom{[5]}2$ (Theorem~\ref{a:t:u25masters}), the masters over the field of the weights of that closing should be $\Pth$ and one free slope, the weights carrying what Remark~\ref{a:r:Hnot} sees as new.

\subsection*{18.5\quad Two observations}

(i) \emph{The exact pairing belongs to the banana.} Propositions~\ref{a:p:pairing} and~\ref{a:p:edge} give
\[
\mathrm{Cor}\cdot\det H=-\frac{\lambda^\star\chi'(\lambda^\star)}{[k^{n-1}]c_0}=\prod_{c\ne\star}\Big(1-\frac{\lambda^\star}{\lambda_c}\Big)
\]
whenever the closings solve the adjoint of the law and $c_0$ has degree $n-1$; here $H$ in $\det H$ is the Hessian of Proposition~\ref{a:p:edge}. The laws of $\binom{[n]}1$ for $n=2,3,4,5$ (Theorem~\ref{a:t:cert}, \S15.7, Proposition~\ref{a:p:pairing}) are self-adjoint up to sign, $c_i(-k-i)=\pm c_i(k)$ \computed, and the identity is
\[
n!\cdot\frac{2^{n-1}}{n^n}=\prod_{j=1}^{\lceil n/2\rceil-1}\Big(1-\frac{(n-2j)^2}{n^2}\Big)=\prod_{j=1}^{\lceil n/2\rceil-1}\frac{4j(n-j)}{n^2}\quad\derived,
\]
which holds for every $n$: the low end of strength, the corner (Thm.~38.5(i), p.~330), and the high end, the Gaussian at the native gauge (Thm.~38.5(ii)), tied through the native values. On $H$, whose law is not self-adjoint, and on $\binom{[5]}2$ (\S15.5) the closings solve no adjoint and the identity fails.

(ii) \emph{The content belongs to the list.} On $\binom{[3]}1$ and on $H$, two rays each, the masters of the ray are those of the list. The letter and mark relations (Thm.~8.3, p.~54) keep integer data on the balance wall with rational coefficients, so every integer period of a list lies in the span they leave, and each ray samples that span. A ray with the symmetries of the list may sample only their invariant part, which offsets would break. Hypothesis: every ray through the origin strictly inside its walls, without the list's symmetries, spans the whole content. \notyet

\renewcommand{\thesection}{19}
\setcounter{proposition}{0}
\renewcommand{\theproposition}{19.\arabic{proposition}}

\section{The three-mark centre, set as an examination}

Remark~\ref{a:r:paper} read the three-mark centre off its law as a list of ratios, and Procedure~\ref{a:p:reading} states the reading of every ray through the origin. This entry runs the reading once, on that one ray, written as an examination question with its answer key and its marking scheme. A key has to say what a complete solution shows at every step and what each omission costs. So every hypothesis that \S\S15--18 checked once in passing is checked here, at the step that uses it, and the places where a solution can look right and be wrong are named, with the credit they earn. The ratios of Remark~\ref{a:r:paper} become a certified bracket. The receipts are in the script exam\_s19.py (python-flint: exact arithmetic for the solution, ball arithmetic only for the independent check of \S19.5).

\subsection*{19.1\quad The question}

\begin{center}
\fbox{\parbox{0.94\linewidth}{\smallskip
\textbf{Question} (40 points). Let $F=\binom{[3]}1$ be the list on the marks $\{1,2,3\}$ whose bases are the singletons, with unit weights, so that its letters are
\[
K=\alpha_1+\alpha_2+\alpha_3,\qquad \Psi=\alpha_1\alpha_2+\alpha_1\alpha_3+\alpha_2\alpha_3,
\]
and let $\ell=P_F(\one;1,1)$ be the period of its centre. By the reading of its centre ray (Procedure~\ref{a:p:reading}), that is, the law of the ray certified by the relations along a ray (Thm.~35.3, p.~308), its corner, and positive selection (Theorem~\ref{a:t:select}), find a rational number $q$ with
\[
0<\ell-q<10^{-6},
\]
and prove both inequalities.

\smallskip
You may cite the handbook and \S\S15--18 of this notebook. When you cite a statement, check its hypotheses in the case at hand. Exact computation (integers, rationals, polynomials, linear algebra over $\mathbb Q$) may be done by machine; report it so that it can be repeated. Floating-point arithmetic, numerical integration, integer-relation searches and comparison with a known decimal value of $\ell$ may be used as checks. They earn no credit.
\smallskip}}
\end{center}

\subsection*{19.2\quad The answer and the scheme}

\emph{Answer.} $q=a_9=\dfrac{2447622221}{1044248800}=2.3439071426\ldots$, with
\[
\frac{2447622221}{1044248800}<\ell<\frac{6414407843}{2736630400},\qquad b_9-a_9=\frac{6561}{59511099776}<1.11\cdot10^{-7}.
\]
Any rational $q$ with a complete proof of $0<\ell-q<10^{-6}$ earns full credit, whatever bounds the proof uses (other consequences of positivity, such as the determinants of Thm.~35.10(iii), are admissible), and the points of the scheme then go to the corresponding steps of that proof; \S19.4, items 1, 5 and 6, lists the alternatives a grader will meet most often. What the key proves is this.

\begin{theorem}[The three-mark centre, bracketed]\label{a:t:bracket}\derived; the table \computed\ exactly.\
Let $c$ and $r$ be the solutions of
\[
k^2x(k)-(10k^2+10k+3)\,x(k+1)+9(k+1)^2\,x(k+2)=0\qquad(k\ge1)
\]
with $c(1)=1$, $c(2)=\frac13$ and $r(1)=0$, $r(2)=-\frac23$. For every integer $K\ge2$, $c(K)>0$ and $9c(K)>c(K-1)$, and
\begin{gather*}
a_K=-\frac{r(K)}{c(K)}\ <\ \Pth\ <\ b_K=\frac{r(K-1)-9\,r(K)}{9\,c(K)-c(K-1)},\\
b_K-a_K=\frac{2\cdot9^{2-K}}{3(K-1)^2\,c(K)\big(9c(K)-c(K-1)\big)} .
\end{gather*}
\end{theorem}

\begin{center}
\small
\begin{tabular}{@{}r p{0.72\linewidth} r@{}}
\toprule
Step & What is shown & Points\\ \midrule
1 & The ray: the chart and its measure, the stride, the values, the walls; $\ell=P(1)$ & 3\\
2 & The stride: $f\le\frac19$; $\frac19e^{-\varphi}\le f\le e^{-\varphi}$; $\int e^{-k\varphi}dy=6/k^2$; what follows at every real $k>0$ & 5\\
3 & The closings, their recurrence, the reflected candidate law & 3\\
4 & A certificate for Thm.~35.3 and its exact verification & 5\\
5 & The law at every real $k>0$: the integration by parts, with the bounds of Step 2 & 5\\
6 & The native values; what positivity can select & 2\\
7 & The corner: $k^2P(k)\to6$, so $9P(2)=3\ell-6$ & 5\\
8 & The rational solutions, their signs, the exact table & 4\\
9 & The bracket $a_K<\ell<b_K$, the choice of $K$, the answer & 8\\ \midrule
& & 40\\
\bottomrule
\end{tabular}
\end{center}

\subsection*{19.3\quad The key}

Throughout, $e_1,e_2,e_3$ are the elementary symmetric polynomials of $\alpha_1,\alpha_2,\alpha_3$, so that $K=e_1$ and $\Psi=e_2$, and $D_e=\alpha_e\,\partial/\partial\alpha_e$.

\paragraph{Step 1 (3 points). The ray.}
Take the chart $\alpha_3=1$ of the mark $e_0=3$, with $\alpha_1=e^{y_1}$ and $\alpha_2=e^{y_2}$, so that
\[
dy=dy_1\,dy_2=\frac{d\alpha_1\,d\alpha_2}{\alpha_1\alpha_2}.
\]
The centre ray at the strength $\sigma=1$ (Def.~35.9, p.~309) has origin $0$ and direction $(\one;1,1)$. Its stride and its values are
\[
f=\frac{\alpha^{\one}}{K\Psi}=\frac{e_3}{e_1e_2},\qquad P(k)=P_F(k\one;k,k)=\int_{\mathbb R^2}f^k\,dy\qquad(k=1,2,\dots),
\]
so that
\[
\ell=P(1)=\int_0^\infty\!\!\int_0^\infty\frac{d\alpha_1\,d\alpha_2}{(1+\alpha_1+\alpha_2)(\alpha_1+\alpha_2+\alpha_1\alpha_2)} .
\]
For real $k>0$ the same integral defines $P(k)$; Step 2 shows that it converges. The walls of the direction are $\kappa_{1,S}=|S|-d(S)-d^\ast(S)$ (Def.~8.1, p.~53), the width $c(S)$ of Def.~7.10, p.~49. Every nonempty $S$ has rank $r(S)=1$, so for $S$ nonempty and proper $d(S)=|S|-r^\ast(S)=r-r(\cE\setminus S)=0$ and $d^\ast(S)=|S|-r(S)=|S|-1$, and
\[
\kappa_{1,S}=1\qquad(\varnothing\neq S\subsetneq\cE).
\]
All proper walls are positive, the hypothesis of Thm.~35.3 and of Thm.~35.10, p.~309.

\smallskip\noindent\emph{Marking.} 1 point for the stride with the chart and its measure; 1 for $\ell=P(1)$ with the definition of the values cited; 1 for the walls. The measure $d\alpha_1d\alpha_2$ in place of $dy$ makes the integral diverge (the integrand is $\sim1/(\alpha_1+\alpha_2)$ at infinity): $-1$.

\paragraph{Step 2 (5 points). The stride.}
(a) \emph{The gauge.} The identity
\[
e_1e_2-9e_3=\alpha_1(\alpha_2-\alpha_3)^2+\alpha_2(\alpha_3-\alpha_1)^2+\alpha_3(\alpha_1-\alpha_2)^2
\]
holds because both sides equal $\sum_{i\neq j}\alpha_i^2\alpha_j-6e_3$, the left since $e_1e_2=\sum_{i\neq j}\alpha_i^2\alpha_j+3e_3$. So $0<f\le\frac19$ on the orthant, with equality exactly on the diagonal $\alpha_1=\alpha_2=\alpha_3$, which on the chart is the single point $y=0$. This is Thm.~35.10(i) here, with $\alpha^\star=\one$ and $f(\alpha^\star)=\frac19$.

(b) \emph{The decay.} Order the coordinates, $\alpha_{(1)}\ge\alpha_{(2)}\ge\alpha_{(3)}$. Then $\alpha_{(1)}\le e_1\le3\alpha_{(1)}$, $\alpha_{(1)}\alpha_{(2)}\le e_2\le3\alpha_{(1)}\alpha_{(2)}$ and $e_3=\alpha_{(1)}\alpha_{(2)}\alpha_{(3)}$, so
\[
\frac19\,\frac{\alpha_{(3)}}{\alpha_{(1)}}\le f\le\frac{\alpha_{(3)}}{\alpha_{(1)}} .
\]
On the chart $\alpha_{(3)}/\alpha_{(1)}=e^{-\varphi(y)}$ with
\[
\varphi(y)=\max(y_1,y_2,0)-\min(y_1,y_2,0),
\]
and since $0$ and $y_i$ are both among the three entries, $\varphi(y)\ge|y_i|$: $\varphi(y)\ge\|y\|_\infty$.

(c) \emph{The volume.} For real $k>0$,
\[
\int_{\mathbb R^2}e^{-k\varphi(y)}\,dy=\frac6{k^2}.
\]
The lines $y_1=y_2$, $y_1=0$, $y_2=0$ cut the plane into six cones, one for each order of the entries $(y_1,y_2,0)$, and on each cone $\varphi$ is the largest entry less the smallest:
\begin{gather*}
\int_{y_1\ge y_2\ge0}e^{-ky_1}dy=\int_0^\infty y_1e^{-ky_1}dy_1=\frac1{k^2},\qquad
\int_{y_1\ge0\ge y_2}e^{-k(y_1-y_2)}dy=\frac1k\cdot\frac1k,\\
\int_{0\ge y_1\ge y_2}e^{ky_2}dy=\int_0^\infty ue^{-ku}du=\frac1{k^2},
\end{gather*}
and the three cones with $y_1$ and $y_2$ exchanged give the same.

(d) \emph{Consequences}, for real $k>0$:
\[
\frac{6\cdot9^{-k}}{k^2}\le P(k)\le\frac6{k^2},\qquad P(k+1)<\frac19\,P(k),
\]
the second from (a), strict because $f<\frac19$ off the single point $y=0$. And $P$ is continuous on $(0,\infty)$: on $[k_1,\infty)$ the integrands $f^k$ are dominated by $f^{k_1}$, because $0<f<1$, and $f^{k_1}$ is integrable.

\smallskip\noindent\emph{Marking.} (a) 2 points; the identity must be shown or verified by expansion. The two inequalities of the arithmetic and geometric means, $e_1\ge3e_3^{1/3}$ and $e_2\ge3e_3^{2/3}$, are an accepted proof of $f\le\frac19$. The strict inequalities of Step 9 need only $f<\frac19$ on a set of positive measure; without any strictness the bracket is closed at the top, which suffices for the answer. (b) 1 point for the upper side, the decay; citing Thm.~35.10(i) for $f\le C'e^{-c'\|y\|}$ on the chart earns it, and gives convergence at every real $k>0$. 1 point for the lower side of (b) and the volume (c) together; they serve only the corner, and a solution that obtains the corner by a correct citation of Proposition~\ref{a:p:corner} or~\ref{a:p:dirclose} in Step 7 earns this point there. (d) 1 point: convergence, positivity, continuity and $P(k+1)<\frac19P(k)$. Step 2(d) credits the proof of the last inequality and Step 9(b) its use.

\paragraph{Step 3 (3 points). The candidate law.}
Here $1/f=e_1e_2/e_3=(\alpha_1+\alpha_2+\alpha_3)(\alpha_1^{-1}+\alpha_2^{-1}+\alpha_3^{-1})$, and the closings of the direction (Proposition~\ref{a:p:dirclose}) are
\[
x_j=[\alpha^0]\Big(\frac1f\Big)^j=\sum_{a_1+a_2+a_3=j}\Big(\frac{j!}{a_1!\,a_2!\,a_3!}\Big)^2=1,\ 3,\ 15,\ 93,\ 639,\ 4653,\ 35169,\ 272835,\ 2157759,\ \dots,
\]
since the constant term pairs the monomial $\alpha^a$ of $(\sum_e\alpha_e)^j$ with the monomial $\alpha^{-a}$ of $(\sum_e\alpha_e^{-1})^j$, and both have the coefficient $j!/(a_1!a_2!a_3!)$. Exact linear algebra with the ansatz $\sum_{i=0}^2p_i(j)\,x_{j+i}=0$, $\deg p_i\le2$, on $x_0,\dots,x_{15}$ (14 equations, 9 unknowns) has a one-dimensional space of solutions,
\[
(j+2)^2x_{j+2}-(10j^2+30j+23)\,x_{j+1}+9(j+1)^2x_j=0 .
\]
Reflected as in Corollary~\ref{a:c:dirlaw}, $c_i(k)=p_{2-i}(-k-2)$:
\begin{equation}\tag{L}
k^2P(k)-(10k^2+10k+3)\,P(k+1)+9(k+1)^2\,P(k+2)=0 .
\end{equation}
At this point (L) is a candidate. Corollary~\ref{a:c:dirlaw} runs the other way, from a law of $P$ to the recurrence of the closings, and sixteen closings do not prove their own recurrence.

\smallskip\noindent\emph{Marking.} 1 point each for the closings, the recurrence (with the size of the system, or any other exact derivation) and the reflection. Any route to (L) earns the 3 points: the closings are how the law is found, not part of its proof. A solution that takes (L) as proved here keeps these 3 points and earns nothing in Steps 4 and 5.

\paragraph{Step 4 (5 points). The certificate.}
The origin is $0$, so the defects of the origin vanish, and the defect of the direction is (Lemma~35.2(ii), p.~308)
\[
c_e=D_e\log f=1-R_e-R^\ast_e,\qquad R_e=\frac{\alpha_e}{e_1},\qquad R^\ast_e=\frac{\alpha_e(e_1-\alpha_e)}{e_2},
\]
the occupation of $e$ under the basis law and under the cobasis law. Since $\sum_eR_e=1$ and $\sum_eR^\ast_e=2$, $\sum_ec_e=0$. A \emph{certificate} of (L) is a choice of insertions $g_1,g_2,g_3$ (Def.~35.1: regular functions on the letter complement; here rational functions of degree $0$ whose denominators are products of $e_1$ and $e_2$, with coefficients polynomial in $k$) such that
\begin{equation}\tag{C}
k^2-(10k^2+10k+3)\,f+9(k+1)^2f^2=\sum_{e=1}^3\big(D_eg_e+k\,c_e\,g_e\big).
\end{equation}
\emph{Ansatz:} $g_e=N_e/(e_1^2e_2)$, with $N_e$ homogeneous of degree $4$ with coefficients in $\mathbb Q+\mathbb Qk$, $g_e$ the image of $g_1$ under the exchange of the marks $1$ and $e$, and $N_1$ symmetric in $\alpha_2,\alpha_3$. Multiplied by $e_1^3e_2^2$, (C) becomes a polynomial identity, of degree $7$ in $\alpha$ and $2$ in $k$, through
\[
e_1^3e_2^2\big(D_eg_e+k\,c_eg_e\big)=e_1e_2\,D_eN_e-N_e\big(2\alpha_ee_2+\alpha_e(e_1-\alpha_e)e_1\big)+k\,N_e\big(e_1e_2-\alpha_ee_2-\alpha_e(e_1-\alpha_e)e_1\big),
\]
and it is linear in the $18$ coefficients of $N_1$: $99$ equations of rank $10$. One solution is
\begin{equation}\tag{G}
g_1=-\frac{\alpha_1^2\big(k\,\alpha_1(\alpha_2+\alpha_3)+(2+3k)(\alpha_2^2+\alpha_3^2)-(1-k)\,\alpha_2\alpha_3\big)}{e_1^2e_2},
\end{equation}
and within the ansatz the others differ from it by one insertion $h$ added at every mark, $h$ with a symmetric numerator (an $8$-dimensional space). Such an $h$ changes nothing, because $\sum_eD_eh=0$ (Euler's identity in degree $0$) and $\sum_ec_eh=0$.

\emph{Verification.} Write $g_e=g^{(0)}_e+k\,g^{(1)}_e$. The three powers of $k$ in (C) read
\begin{align*}
k^2:&\quad \textstyle\sum_ec_eg^{(1)}_e=1-10f+9f^2=(1-f)(1-9f),\\
k^1:&\quad \textstyle\sum_e\big(D_eg^{(1)}_e+c_eg^{(0)}_e\big)=-10f+18f^2,\\
k^0:&\quad \textstyle\sum_eD_eg^{(0)}_e=-3f+9f^2 .
\end{align*}
Times $e_1^3e_2^2$ each sum over the marks is a symmetric polynomial of degree $7$, and in the basis of products of $e_1,e_2,e_3$ exact expansion gives
\begin{gather*}
\textstyle\sum_ee_1^3e_2^2\,c_eg^{(1)}_e=e_1^3e_2^2-10\,e_1^2e_2e_3+9\,e_1e_3^2,\\
\textstyle\sum_ee_1^3e_2^2\,D_eg^{(1)}_e=-2\,e_1e_2^3-3\,e_1^2e_2e_3+9\,e_1e_3^2,\qquad
\sum_ee_1^3e_2^2\,c_eg^{(0)}_e=2\,e_1e_2^3-7\,e_1^2e_2e_3+9\,e_1e_3^2,\\
\textstyle\sum_ee_1^3e_2^2\,D_eg^{(0)}_e=-3\,e_1^2e_2e_3+9\,e_1e_3^2,
\end{gather*}
which, with $f=e_3/(e_1e_2)$, are the three right sides times $e_1^3e_2^2$. So (C) holds identically in $\alpha$ and $k$. Proposition~\ref{a:p:certs} reported a certificate of this law with the same denominator.

\smallskip\noindent\emph{Marking.} 1 point for the defects $c_e$, in any equivalent form. 2 points for a certificate. Any certificate is accepted; it need not be symmetric, and any denominator is accepted for which Step 5 can be carried out. 2 points for the verification: by exact expansion, by the symmetric-function form above, or by exact evaluation of the cleared identity on a grid large enough to decide it. After clearing, both sides are polynomials of degree at most $7$ in each $\alpha_i$ and $2$ in $k$, homogeneous in $\alpha$, so agreement at $\alpha_3=1$ on $8$ values of $\alpha_1$, $8$ of $\alpha_2$ and $3$ of $k$ (all $192$ points) proves it. Agreement at a few random points is a test, not a proof. To test a certificate other than (G), a grader evaluates both sides of (C) at $\alpha=(1,2,4)$, $k=\frac13$, where the left side is $-\frac{2809}{7203}$ (there $f=\frac4{49}$, and (G) gives $g_1,g_2,g_3=-\frac{85}{1029},-\frac{310}{1029},-\frac{424}{1029}$); a certificate written on the chart, with $g_3=0$ and $D_1,D_2$ read as $\partial_{y_1},\partial_{y_2}$, is tested at the same point $(\alpha_1,\alpha_2)=(\frac14,\frac12)$. A certificate with floating-point coefficients, or verified only in floating point, earns 0 of the 4: the identity has to hold exactly. How the certificate was found does not matter: rational coefficients found by a numerical search and then verified exactly earn full credit.

\paragraph{Step 5 (5 points). The law at every real $k>0$.}
Thm.~35.3 states the relations for all \emph{large} $k$. The law is needed at $k=1,\dots,K-2$ (Step 8) and as $k\to0^+$ (Step 7), so the integration by parts is carried out here at every real $k>0$, with the bounds of Step 2. Fix a real $k>0$.

(a) $D_e(f^k)=k\,c_ef^k$, so the right side of (C), times $f^k$, is $\sum_eD_e(g_ef^k)$.

(b) On the chart $D_1=\partial_{y_1}$ and $D_2=\partial_{y_2}$, and for every function $h$ of degree $0$, $D_3h=-D_1h-D_2h$ (Euler's identity). So
\[
\sum_{e=1}^3D_e(g_ef^k)=\partial_{y_1}H_1+\partial_{y_2}H_2,\qquad H_i=(g_i-g_3)f^k .
\]

(c) \emph{Bounds.} Every monomial of every $N_e$ occurs in the expansion of $e_1^2e_2$: for $N_1$, $\alpha_1^3\alpha_2=\alpha_1^2\cdot\alpha_1\alpha_2$, $\alpha_1^2\alpha_2^2=\alpha_1\alpha_2\cdot\alpha_1\alpha_2$ and $\alpha_1^2\alpha_2\alpha_3=\alpha_1^2\cdot\alpha_2\alpha_3$, with the images under $2\leftrightarrow3$, and the other $N_e$ are images of $N_1$. The coefficients of $e_1^2e_2$ are positive, so $\alpha^v\le e_1^2e_2$ on the orthant for each such monomial, and $|g_e|\le\|N_e\|_1$, the sum of the absolute values of the coefficients. For every mark $i$,
\[
D_ig_e=\frac{D_iN_e}{e_1^2e_2}-g_e\,\big(2R_i+R^\ast_i\big),
\]
where $D_iN_e$ has the monomials of $N_e$ or fewer and $0\le R_i,R^\ast_i\le1$; and $|c_i|\le1$. With Step 2(b) there is, for fixed $k$, a constant $C$ with
\[
|H_i|\le C\,e^{-k\varphi},\qquad |\partial_{y_j}H_i|\le C\,e^{-k\varphi} .
\]

(d) If $H$ is $C^1$ on $\mathbb R^2$, tends to $0$ at infinity and $\partial_{y_1}H$ is integrable, then $\int\partial_{y_1}H\,dy=0$: for every $y_2$, $\int_{-R}^R\partial_{y_1}H\,dy_1=H(R,y_2)-H(-R,y_2)\to0$, and Fubini applies. By (c) and Step 2, $H_1$ and $H_2$ qualify ($e^{-k\varphi}\le e^{-k\|y\|_\infty}\to0$, and $e^{-k\varphi}$ is integrable), so the right side of (C), times $f^k$, integrates to $0$.

(e) The left side of (C), times $f^k$, integrates to the left side of (L), all three integrals convergent by Step 2. Hence (L) holds at every real $k>0$.

\smallskip\noindent\emph{Marking.} 1 point for (a)--(b), including the mark $3$ on its own chart. 2 points for the bounds (c). 1 point for (d). 1 point for the conclusion at every real $k>0$, with the reason it is needed. Citing Thm.~35.3 at $k=1,2,\dots$ and as $k\to0^+$ without (c): 2 of 5. Any complete extension earns full credit; for instance (C) integrated at large integers by Thm.~35.3 as stated (its hypothesis checked in Step 1), then carried to every $k$ with $\re k>0$ by Carlson's theorem: the left side of (L) is holomorphic there, it is bounded by a polynomial in $|k|$ on $\re k\ge1$ because $|P(k)|\le P(\re k)\le6/(\re k)^2$ (Step 2), and Carlson's theorem applied to $z\mapsto F(N+z)$, with $F$ the left side of (L) and $N\ge1$ an integer from which on (L) holds at every integer, makes $F$ vanish identically. Integrating $D_3(g_3f^k)$ by parts ``in $y_3$'' on the chart $\alpha_3=1$, where $y_3$ is not a coordinate: $-1$, unless repaired by Euler's identity or by a certificate with $g_3=0$.

\paragraph{Step 6 (2 points). The native values; what positivity can select.}
The symbol of (L) (Cor.~35.6, p.~308) is
\[
\chi(\lambda)=1-10\lambda+9\lambda^2=(1-\lambda)(1-9\lambda).
\]
The native gauges of the centre ray are the signs $\eta$ with $\eta_3=1$ (Proposition~\ref{a:p:onlysigns}, Thm.~35.12, p.~310), of native value $K(\eta)^{-2}$: the table's gauge $\eta=(1,1,1)=\alpha^\star$, with $K(\eta)=3$ and the value $\frac19=f(\alpha^\star)$; and $\eta=(-1,1,1),(1,-1,1),(-1,-1,1)$, with $K(\eta)=\pm1$ and the value $1$. They are nondegenerate. By Thm.~35.12(ii) the Hessian of $-\log f$ at $\eta$ is $2\big(\kappa_\eta(X_e,X_{e'})\big)_{e,e'\neq3}$, the covariances of the indicators under the signed measure $\mathbb P_\eta(\{e\})=\eta_e/K(\eta)$. Its determinant is $\frac4{27}$ at $(1,1,1)$ and $-4$ at each of the other three signs. So Cor.~35.6, with $g=1$, makes both native values roots of $\chi$, and since $\chi$ has degree $2$ they are all its roots. (Cor.~35.6 runs from native values to roots; a root need not be a native value, as the extra root $\frac{131003}{4771109}$ in \S18.4 shows.)

In Theorem~\ref{a:t:select} the singular points are $q=9$, where the slit $[9,\infty)$ begins, and $q=1$, off the slit. At $q=1$, a simple root and the value of nondegenerate gauges, positivity imposes $\nu=1$ condition, and with the corner $1+\nu=2=h$. The condition is independent, and here that is seen directly: by Step 8, $c(K)\ge3^{1-K}$ and $9c(K)-c(K-1)\ge2c(K-1)\ge2\cdot3^{2-K}$, so the width of Step 9(d) is at most $1/(K-1)^2$ and the bracket closes on one number. So the law, the corner and positivity determine $P$ (Theorem~\ref{a:t:select}(iii)). The rational solutions of (L) have the exponential rate $1$ of the signs and $P$ the rate $\frac19$ of the gauge, so the bracket shrinks by about $\frac19$ per step, $\log_{10}9=0.954$ digits: about seven steps for $10^{-6}$, before constants. The constants decide in Step 9.

\smallskip\noindent\emph{Marking.} 1 point for the native values with their nondegeneracy, matched to the roots of the symbol in the direction Cor.~35.6 gives. ``The roots are native values by Cor.~35.6'' reverses the corollary and earns 0 for this point unless the roots are matched to gauges directly, as above. 1 point for the count $1+\nu=h$ with the independence, or any equivalent statement that the corner supplies the condition positivity does not. The step is not used in the proofs of the bounds; a complete solution without it loses these 2 points only.

\paragraph{Step 7 (5 points). The corner.}
(a) By Step 2(d), $6\cdot9^{-k}\le k^2P(k)\le6$ for real $k>0$, so
\[
\lim_{k\to0^+}k^2P(k)=6 .
\]
(b) Take (L) at real $k\in(0,1)$ (Step 5) and let $k\to0^+$. By the continuity of Step 2(d), $P(k+1)\to P(1)=\ell$ and $P(k+2)\to P(2)$, while $10k^2+10k+3\to3$ and $9(k+1)^2\to9$. So $6-3\ell+9P(2)=0$:
\begin{equation}\tag{S}
P(2)=\frac{\ell-2}3 .
\end{equation}
(c) This is Corollary~\ref{a:c:dirlaw} at $j=0$,
\[
[k^2]c_0\cdot\mathrm{Cor}+c_1(0)P(1)+c_2(0)P(2)=0,\qquad \mathrm{Cor}=\sum_\pi\prod_{i=2}^3\kappa_{1,U_i}^{-1}=3!\cdot1=6
\]
(Proposition~\ref{a:p:dirclose}; for $\binom{\cE}1$ the corner is $n!$, Proposition~\ref{a:p:corner}). The six cones of Step 2(c) are the six orders $\pi$, each contributing $\prod_i\kappa_{1,U_i}^{-1}=1$.

\smallskip\noindent\emph{Marking.} 2 points for (a), 2 for (b), 1 for (c). A correct citation of Proposition~\ref{a:p:corner} or~\ref{a:p:dirclose} in place of (a), with the walls of Step 1 and the sum over the $3!$ orders, earns the 2 points of (a). Errors:
\begin{itemize}
\item Reading (L) at $k=0$ as $0\cdot P(0)-3P(1)+9P(2)=0$. Here $P(0)=\int dy=\infty$, and the product $0\cdot\infty$ is the corner, $6$, not $0$. The reading gives $P(2)=\ell/3$, which is false ($P(2)=0.1146\ldots$, $\ell/3=0.78\ldots$) and contradicts Step 2(d) at once: $P(2)<\frac19P(1)$ would say $\ell/3<\ell/9$. After it every $P(K)$ is $\ell$ times a positive rational, positivity fixes only ratios, and the bracket is empty. 0 for Step 7, and in Steps 8--9 only the points that do not depend on (S): the signs and their use, 2 points.
\item Passing the limit under the integral, $k^2f^k\to0$ pointwise, so ``$k^2P(k)\to0$'': 0 for (a). There is no domination, since $\int dy=\infty$.
\item The wrong power ($kP(k)$, the corner of a ray on two marks) or $\mathrm{Cor}=1$ (one order): 0 for (a); (b) is marked as a consequence.
\item Citing Corollary~\ref{a:c:dirlaw} for (S). Its hypothesis is the law \emph{identically}, as meromorphic functions of $k$; Step 5 gives it at real $k>0$, and the continuation of Proposition~\ref{a:p:dirclose} extends it. Citing the corollary with the law proved only at the integers leaves a gap: $-1$.
\end{itemize}

\paragraph{Step 8 (4 points). The rational solutions and their signs.}
For integers $k\ge1$, $9(k+1)^2\neq0$, and (L) gives
\[
P(k+2)=\frac{(10k^2+10k+3)\,P(k+1)-k^2P(k)}{9(k+1)^2}.
\]
Let $c$ and $r$ be the solutions of (L) with $c(1)=1$, $c(2)=\frac13$ and $r(1)=0$, $r(2)=-\frac23$. From $P(1)=\ell$, (S) and induction,
\[
P(K)=c(K)\,\ell+r(K)\qquad(K\ge1).
\]
\emph{Signs, at every $K$.} Put $x'_j=9^jc(j+1)$. Then (L) at $k=j+1$ reads $(j+2)^2x'_{j+2}=(10j^2+30j+23)\,x'_{j+1}-9(j+1)^2x'_j$, with $x'_0=1$, $x'_1=3$. If $x'_{j+1}\ge3x'_j>0$, then
\[
(j+2)^2x'_{j+2}\ge\big(10j^2+30j+23-3(j+1)^2\big)x'_{j+1}=(7j^2+24j+20)\,x'_{j+1}\ge3(j+2)^2x'_{j+1},
\]
since $7j^2+24j+20-3(j+2)^2=4(j+1)(j+2)\ge0$. So $x'_{j+1}\ge3x'_j>0$ for every $j$, hence $x'_j\ge3^j$, that is,
\[
c(K)\ge3^{1-K}>0,\qquad 9c(K)\ge3c(K-1)>c(K-1)\qquad(K\ge2).
\]
The $x'_j$ obey the recurrence of the closings with the same first two terms, so $x'_j=x_j$ and $c(K)=9^{1-K}x_{K-1}$ wherever that recurrence holds: in the table below by inspection, and for every $j$ by Corollary~\ref{a:c:dirlaw} once (L) holds identically. The proof of the signs does not use this.

\begin{center}
\small
\renewcommand{\arraystretch}{1.12}
\begin{tabular}{@{}r r r r l@{}}
\toprule
$K$ & $c(K)$ & $r(K)$ & $a_K=-r(K)/c(K)$ & \\ \midrule
1 & $1$ & $0$ & $0$ & \\
2 & $1/3$ & $-2/3$ & $2$ & \\
3 & $5/27$ & $-23/54$ & $23/10$ & $2.3$\\
4 & $31/243$ & $-145/486$ & $145/62$ & $2.33870967\ldots$\\
5 & $71/729$ & $-1331/5832$ & $1331/568$ & $2.34330985\ldots$\\
6 & $517/6561$ & $-242353/1312200$ & $242353/103400$ & $2.34383945\ldots$\\
7 & $11723/177147$ & $-5495507/35429400$ & $5495507/2344600$ & $2.34389959\ldots$\\
8 & $10105/177147$ & $-46422941/347208120$ & $46422941/19805800$ & $2.34390638\ldots$\\
9 & $26639/531441$ & $-2447622221/20832487200$ & $2447622221/1044248800$ & $2.34390714\ldots$\\
10 & $5773279/129140163$ & $-530455590781/5062294389600$ & $530455590781/226312536800$ & $2.34390722\ldots$\\
\bottomrule
\end{tabular}
\end{center}
Two checks a grader can run on a student's table: $9^{K-1}c(K)$ is the closing $x_{K-1}$ ($1,3,15,93,\dots$), and the frame of the two solutions (Prop.~35.8, p.~309: $C_{k+1}=\frac{c_0(k)}{c_2(k)}C_k=\frac{k^2}{9(k+1)^2}C_k$, $C_1=-\frac23$) is
\[
c(k)\,r(k+1)-c(k+1)\,r(k)=-\frac23\cdot\frac{9^{1-k}}{k^2} .
\]

\smallskip\noindent\emph{Marking.} 2 points for $P(K)=c(K)\ell+r(K)$ with the initial values from (S). 1 point for the signs $c(K)>0$ and $9c(K)-c(K-1)>0$, at every $K$ as above or at the $K$ used from the exact table. 1 point for the table, exact, through the $K$ used. A decimal table: 0 for the last point.

\paragraph{Step 9 (8 points). The bracket and the answer.}
(a) \emph{Below, by positivity.} $P(K)=c(K)\ell+r(K)>0$, divided by $c(K)$, gives
\[
\ell>a_K=-\frac{r(K)}{c(K)} .
\]
(b) \emph{Above, by the gauge.} By Step 2(d), $P(K)<\frac19P(K-1)$, that is, $c(K)\ell+r(K)<\frac19\big(c(K-1)\ell+r(K-1)\big)$, or $\big(9c(K)-c(K-1)\big)\ell<r(K-1)-9r(K)$, and divided by $9c(K)-c(K-1)$,
\[
\ell<b_K=\frac{r(K-1)-9\,r(K)}{9\,c(K)-c(K-1)} .
\]
(c) \emph{The signs.} Both divisions keep the direction of the inequality: $c(K)>0$ and $9c(K)-c(K-1)>0$ by Step 8.

At $K=2$ this is the first bracket: $P(2)=\frac{\ell-2}3>0$ gives $\ell>2$, and $P(2)<\frac19P(1)$ gives $\ell<3$.

(d) \emph{The width and the answer.} With the frame of Step 8,
\[
b_K-a_K=\frac{r(K-1)\,c(K)-r(K)\,c(K-1)}{c(K)\big(9c(K)-c(K-1)\big)}=\frac{2\cdot9^{2-K}}{3(K-1)^2\,c(K)\big(9c(K)-c(K-1)\big)},
\]
and in exact arithmetic:
\begin{center}
\small
\renewcommand{\arraystretch}{1.12}
\begin{tabular}{@{}r r r r@{}}
\toprule
$K$ & $b_K$ & $b_K-a_K$ & width with $P(K)<9^{1-K}\ell$ instead\\ \midrule
2 & $3$ & $1$ & $1$\\
3 & $19/8$ & $7.50\cdot10^{-2}$ & $1.64\cdot10^{-1}$\\
4 & $61/26$ & $7.44\cdot10^{-3}$ & $2.54\cdot10^{-2}$\\
5 & $3413/1456$ & $7.84\cdot10^{-4}$ & $3.67\cdot10^{-3}$\\
6 & $104539/44600$ & $8.43\cdot10^{-5}$ & $5.04\cdot10^{-4}$\\
7 & $149014/63575$ & $9.17\cdot10^{-6}$ & $6.66\cdot10^{-5}$\\
8 & $909876251/388187800$ & $1.0036\cdot10^{-6}$ & $8.59\cdot10^{-6}$\\
9 & $6414407843/2736630400$ & $1.1025\cdot10^{-7}$ & $1.0863\cdot10^{-6}$\\
10 & $232184895407/99058909600$ & $1.21\cdot10^{-8}$ & $1.35\cdot10^{-7}$\\
\bottomrule
\end{tabular}
\end{center}
At $K=8$ the certified width is $1.0036\cdot10^{-6}>10^{-6}$; at $K=9$ it is $\frac{6561}{59511099776}<1.11\cdot10^{-7}$. So
\[
\frac{2447622221}{1044248800}=a_9<\ell<b_9=\frac{6414407843}{2736630400},
\]
and $q=a_9$ satisfies $0<\ell-q<b_9-a_9<1.11\cdot10^{-7}<10^{-6}$.

\smallskip\noindent\emph{Marking.} (a) 2 points. (b) 2 points: 1 for using $P(K)<\frac19P(K-1)$ (its proof is credited in Step 2(d)), 1 for $b_K$. (c) 1 point for invoking the two signs at the divisions (their proof is credited in Step 8; a sign proved there and used here silently loses nothing). (d) 3 points: 1 for a $K$ at which the width certified by (a)--(c) is below $10^{-6}$, 1 for the exact comparison, 1 for a stated $q$ whose inequalities $0<\ell-q<10^{-6}$ follow from the bracket.

\subsection*{19.4\quad Where a solution that looks right is wrong}

\begin{enumerate}[label=\arabic*.]
\item \emph{The answer one step early.} $a_8=\frac{46422941}{19805800}$ is a correct answer: $\ell-a_8=8.58\cdot10^{-7}$. But the bracket at $K=8$ has width $1.0036\cdot10^{-6}$, so a solution that certifies $a_8$ from $K=8$ alone has an error in it. The upper bound at $K=9$ certifies it: $\ell-a_8<b_9-a_8=8.72\cdot10^{-7}$. Full credit with that; with the bound claimed at $K=8$, 5 of the 8 points of Step 9. The true error falls below the tolerance one step before the certified one does.
\item \emph{The observed error.} Comparing $a_K$ with a decimal value of $\ell$ (from the share side, from a quadrature, or from a closed form recalled elsewhere) proves nothing about $\ell-a_K$: 0 for Step 9(d).
\item \emph{The rate as a bound.} Theorem~\ref{a:t:select}(iv) gives the error $O(9^{-K})$ up to powers of $K$, with no constant: 0 for Step 9(d). Nor is the last step a bound: $a_9-a_8=7.62\cdot10^{-7}$ is smaller than $\ell-a_8=8.58\cdot10^{-7}$, because the $a_K$ increase to $\ell$ with ratio about $\frac19$, so the remaining error is about $\frac98$ of the next step.
\item \emph{An unproved upper bound.} ``$P(K)\le9^{-K}\ell$, since $f\le\frac19$''. The gauge gives $P(K)<\frac19P(K-1)$, and iterated $P(K)<9^{1-K}P(1)$; the factor $9^{-K}$ does not follow from $f\le\frac19$, it would need $\int dy<\infty$. (The inequality is in fact true for $K\ge8$ and false for $K\le7$: $9^KP(K)$ decreases by Step 2(d), and $9^7P(7)/\ell=1.03$, $9^8P(8)/\ell=0.90$. Nothing in such a solution proves it, and a grader who checks it numerically should not be misled.) With it the width at $K=9$ is $1.21\cdot10^{-7}$, while the simple bound actually proved gives $1.09\cdot10^{-6}$ there, so the solution certifies nothing below $10^{-6}$. The stated $q=a_9$ is right and the proof is not: 1 of 2 in Step 9(b), 0 in 9(d), 4 of the 8 points of Step 9.
\item \emph{The simple bound, done right.} $P(K)<9^{1-K}P(1)=9^{1-K}\ell$ gives $\big(c(K)-9^{1-K}\big)\ell<-r(K)$, and $c(K)-9^{1-K}=9^{1-K}(x'_{K-1}-1)>0$ for $K\ge2$ by Step 8, so $\ell<-r(K)/\big(c(K)-9^{1-K}\big)$: valid and cruder (the last column of the table). It needs $K=10$, with $q=a_{10}=\frac{530455590781}{226312536800}$. Full credit, with the sign stated.
\item \emph{Decimals.} $q=2.343907$ is accepted, with the exact comparisons $\frac{2343907}{10^6}<a_9$ and $b_9-\frac{2343907}{10^6}=2.53\cdot10^{-7}$. A decimal above $a_9$, such as $2.3439072$, is a correct answer ($\ell-2.3439072=3.9\cdot10^{-8}$) but needs a new lower bound: $a_{10}=2.34390722\ldots$ supplies it. A decimal rounded above the lower bound in hand: $-1$ in Step 9(d), its last point.
\item \emph{The signs.} Signs of $c(K)$ and of $9c(K)-c(K-1)$ neither proved nor checked: $-2$ in all, the sign point of Step 8 and Step 9(c). They hold here at every $K$ (Step 8), but where a rational solution changes sign the inequality reverses; on the line $1$ of $\binom{[5]}2$, Remark~\ref{a:r:paper} checked the sign only for $m\le80$.
\item \emph{The corner.} Reading the law at $k=0$ with $0\cdot P(0)=0$, or passing the limit $k\to0^+$ under the integral: the approximation is lost (Step 7). The first gives $P(2)=\ell/3$, every $P(K)$ is then $\ell$ times a positive rational, positivity gives only $\ell>0$, and the gauge bound at $K=2$, $\ell/3<\ell/9$, gives $\ell<0$: the student's own bracket is empty, which the grader can point to.
\item \emph{The theorem at small $k$.} Thm.~35.3 cited at $k=1,2,\dots$ or as $k\to0^+$ without the decay of Step 2: Step 5.
\item \emph{A certificate checked at sample points only}: Step 4.
\item \emph{The measure} $d\alpha_1d\alpha_2$ on the chart: Step 1.
\item \emph{A quadrature for $\ell$ or for $P(2)$ used in the solution}: no credit at the step where it is used. It is a check (\S19.5).
\end{enumerate}

\subsection*{19.5\quad Receipts and the independent check}

The script exam\_s19.py prints every number of this entry. Exact: the sum of squares; the closings $x_0,\dots,x_{15}$ and their recurrence (14 equations, 9 unknowns, nullity 1); the reflection; the certificate (G) and the identity (C) after clearing $e_1^3e_2^2$; its three powers of $k$ in the basis of products of $e_1,e_2,e_3$; the rank $10$ of the symmetric ansatz in $18$ unknowns; the monomial check of Step 5(c); the spot value $-\frac{2809}{7203}$; the symbol and the signs; the walls and $\mathrm{Cor}=6$; the table and the brackets to $K=10$; and to $K=16$ the signs with $c(K)\ge3^{1-K}$, the identity $9^{K-1}c(K)=x_{K-1}$, the frame, and the two forms of the width with its bound $1/(K-1)^2$.

\noindent Check \computed, not part of the solution, in ball arithmetic, with the integrals cut where their tails are below $10^{-70}$. On the chart, the integral over $\alpha_2$ is elementary, $\int_0^\infty\frac{d\alpha_2}{(\alpha_2+A)(A\alpha_2+B)}=\frac{\log(A^2/B)}{A^2-B}$ with $A=1+\alpha_1$, $B=\alpha_1$, so
\[
\ell=\int_0^\infty\frac{\log\big((1+x)^2/x\big)}{1+x+x^2}\,dx=2.343907238689458890601562288872277069095\ldots .
\]
On the share side (Proposition~\ref{a:p:bananapin}), $\Gamma(k)^2P(k)=\int_0^\infty y^{k-1}\sG_{0;(1,-1)}(1,y)^3\,dy$ gives the same $\ell$ to $40$ digits, $P(2)-\frac{\ell-2}3=0$ to $47$ digits, and $P(9)=4.8152090557463543143\cdot10^{-9}=c(9)\ell+r(9)$. With this value,
\[
\ell-a_8=8.5773\cdot10^{-7},\qquad \ell-a_9=9.6062\cdot10^{-8},\qquad b_9-\ell=1.4186\cdot10^{-8}.
\]
The upper side is the sharper. Its slack is $\frac19P(K-1)-P(K)$, which is about $P(K)/(K-1)$ because $P(k)$ decays like $9^{-k}/k$ (for three marks the branch of Thm.~35.11, p.~310, is logarithmic), so $b_K-\ell$ is about $\frac9{8(K-1)}(\ell-a_K)$.

\subsection*{19.6\quad What the example shows}

The reading needed three inputs, and each was elementary once it was named. The analytic input is one inequality, $e_1e_2\ge9e_3$, sharpened to $f\asymp e^{-\varphi}$ by the tropical stride $\varphi$: it gives the convergence at every real $k>0$, the corner, and the upper side of the bracket. The algebraic input is one certificate: it gives the law at every real $k>0$, not only at large $k$. The arithmetic is exact. Positivity gives the lower side and the gauge the upper side. On every ray through the origin the first two inputs exist \derived: $f\le f(\alpha^\star)$ (Thm.~35.10(i)), and the tropical stride, whose volume is $\mathrm{Cor}$ (the six cones of Step 2(c) are the six orders of Proposition~\ref{a:p:dirclose}). A certificate exists for every law that follows from the relations (Thm.~35.19, p.~313), but Step 5 needs more than existence: bounds on the $g_e$ and their derivatives against the decay of $f^k$ down to $k\to0^+$. Here they came from the numerators lying among the monomials of $e_1^2e_2$, as the exponent condition did in Proposition~\ref{a:p:certs}; that a certificate with such bounds exists on every ray is \notyet. With $h-1$ masters, positivity of $P(K),\dots,P(K+h-2)$ and the gauge bounds $P(K+i)<f(\alpha^\star)P(K+i-1)$ cut out a polytope in the coordinates of the masters, and Theorem~\ref{a:t:select}(iv) is the rate at which it shrinks \notyet; with all the conditions of the stride's law the rate is Theorem~\ref{b:t:rate}. The width of the bracket is a ratio of integers from the other end of the pencil (Remark~\ref{a:r:twochains}):
\[
b_K-a_K=\frac{2\cdot9^{K-1}}{3(K-1)^2\,x_{K-1}\,(x_{K-1}-x_{K-2})}\qquad\derived,
\]
from $c(K)=9^{1-K}x_{K-1}$, which holds at every $K$ by Corollary~\ref{a:c:dirlaw}: by Step 5 and the continuation of Proposition~\ref{a:p:dirclose}, (L) holds identically.

\renewcommand{\thesection}{20}
\setcounter{proposition}{0}
\renewcommand{\theproposition}{20.\arabic{proposition}}

\section{The law of the stride: eight questions after the examination}

The examination of \S19 read one ray, the three-mark centre, with three inputs: a law certified by the relations along a ray, the corner, and positivity. It closed on the period with a certified bracket. The eight questions below ask what that reading is when it is pointed elsewhere: at every ray, at the corner alone, at the slopes, at the walls, at the other solutions, at the stride itself, at lists that no graph or field realizes, and at the arithmetic. They turn out to be questions about one measure, the law of the stride (\S20.0): positivity is that it is a measure, the corner is its mass near $0$, the slopes are its logarithmic moments, a wall is a singular point of its density inside its support, the other solutions are the other branches of its density, and the rates are the geometry of its support against the native values. \S20.0 sets it up, with three objects that every ray carries: the tropical stride, the two ends of the stride's law, and the two ends of a certificate. \S\S20.1--20.8 state each question as it was posed and answer it.

Throughout, a ray is a ray through the origin in the sense of \S18: origin $0$, an integral direction $(\mvec_1;\mathbf s_1,\mathbf t_1)$ on the balance wall, every proper wall positive, values $P(k)=\int_{y_{e_0}=0}f^k\,dy$, stride $f$ (Def.~35.1, p.~307), native gauge $\alpha^\star$ of the direction on the orthant and $f^\star=f(\alpha^\star)$, the largest value of $f$ there (Thm.~35.10(i), p.~309). The \emph{slit} is $(0,f^\star)$, read in the plane of native values. The laws below that are not those of \S\S15--19 are least recurrences of closings, reflected (Corollary~\ref{a:c:dirlaw}); the converse, that such a recurrence is a law of the values, is what a certificate of Thm.~35.3, p.~308, settles, and none of them is certified yet. Every reading made with one of them is checked against an independent value.

\subsection*{20.0\quad The stride's law}

\begin{definition}[The stride's law]\label{b:d:law}
The \emph{stride's law} of a ray is the image of the Lebesgue input of the chart under the stride, $\mu=f_*\Lch$ (the law of a measure under a map, Def.~1.1(vii), p.~10; the Lebesgue input, Def.~5.4, p.~29), a measure on $(0,\infty)$; $\varrho$ is its density, $\mu(dt)=\varrho(t)\,dt$. With the reciprocal stride $X=1/f$, the image $X_*\Lch$ is the same measure read at $x=1/t$.
\end{definition}

\begin{proposition}[A ray is one Trailmoment on a line, and one Gramma]\label{b:p:law}\derived\
\begin{enumerate}[label=(\roman*)]
\item $\mu$ is carried by $(0,f^\star]$, it has a density, and for every real $k>0$
\[
P(k)=\int_0^{f^\star}t^k\,\mu(dt)=\big(\sT_{k;(1)}[x]\,X_*\Lch\big)\big((0,\infty)\big):
\]
the values of the ray are the Trailmoment on a line of index $k$, one unit exponent and coefficient function $x$, of the transported Lebesgue input (Thm.~5.9, p.~31).
\item Write $X=\sum_iw_i\,\alpha^{t_i}$ on the chart, a Laurent polynomial with positive coefficients $w_i$ and exponents $t_i$ of degree $0$. The Laplace transform of $X_*\Lch$ beyond the cut $0$ (Def.~5.4) is the Gramma of the reciprocal stride at index $0$,
\[
\int e^{-ux}\,X_*\Lch(dx)=\int_{y_{e_0}=0}e^{-uX(e^y)}\,dy=\sG_{\mathbf 0;t}(u\,\mathbf w),
\]
and the Laplace side (Thm.~5.11, p.~31) reads the whole ray off it:
\[
\Gamma(k)P(k)=\int_0^\infty u^{k-1}\sG_{\mathbf 0;t}(u\,\mathbf w)\,du .
\]
\item For $U\sim\mathrm{Breta}^\circ(F;k\mvec_1;k\mathbf s_1,k\mathbf t_1)$ (Thm.~17.1(ii), p.~107), the law of $f(U)$ is $t^k\mu(dt)/P(k)$. So the ray is the exponential family of Thm.~17.4(ii), p.~108, restricted to the direction: sufficient statistic $\log f$, base measure $\mu$, log-partition $\log P$, and
\[
\frac{d}{dk}\log P(k)=\oE\big[\log f(U)\big],\qquad \frac{d^2}{dk^2}\log P(k)=\operatorname{Var}\big[\log f(U)\big]>0 .
\]
\end{enumerate}
\end{proposition}

\begin{proof}
(i) $0<f\le f^\star$ on the orthant with equality only on the ray of $\alpha^\star$ (Thm.~35.10(i)), and $f$ has no critical point on the chart other than $\log\alpha^\star$ there, so the coarea formula gives a density. The integral of $f^k$ is the integral of $t^k$ against the image measure, which is Thm.~5.9 with the coefficient function $x\mapsto x$ and the kernel $\sZ_{k;(1)}(x)=x^{-k}$ (Def.~5.2). (ii) $\int_0^\infty u^{k-1}e^{-uX}\,du=\Gamma(k)X^{-k}$ and Fubini, all terms positive; on the chart $x^{\nu-\one}dx=dy$ at $\nu=\mathbf 0$, so the integral of $e^{-uX}$ is the Gramma of Def.~2.1, p.~13, with the exponents $t_i$ and coefficients $uw_i$. It converges because the direction's walls are positive, which puts $\mathbf 0$ in the interior of the cone of the $t_i$ (Thm.~2.2). (iii) On the balance wall the density of $\mathrm{Breta}^\circ$ in $y$ is proportional to the integrand of the period (Thm.~17.1(i)--(ii)), here $f^k$ on the chart. The two derivatives are those of a log-partition (Thm.~17.4(ii)).
\end{proof}

At the centre of a list, $X=K(\alpha)K(\alpha^{-1})$, and for the banana Prop.~\ref{a:p:bananapin} adds a third reading: since $\int_0^\infty y^{\sigma-1}\sG_{0;(1,-1)}(1,y)\,dy=\Gamma(\sigma)^2$,
\[
\sG_{0;(1,-1)}(1,y)^n=\int_0^{f^\star}\sG_{0;(1,-1)}(1,y/t)\,\mu(dt)\qquad\derived:
\]
the $n$-th power of the Gramma of one mark is a $\mu$-mixture of its own dilations. Every question below is a question about $\mu$.

\begin{proposition}[The corner is the volume of the tropical stride]\label{a:p:tropcorner}\derived; sampled \computed\
Let a ray through the origin of a family with positive weights have every proper wall positive. On the chart put
\[
\varphi(y)=\sum_j\Big(s_{1j}\max_{T\in\cB_j}\langle\one_T,y\rangle+t_{1j}\max_{T\in\cB_j}\langle\one_{\cE\setminus T},y\rangle\Big)-\langle\mvec_1,y\rangle,
\]
the \emph{tropical stride}: $-\log f$ with every letter replaced by its largest monomial and the weights dropped.
\begin{enumerate}[label=(\roman*)]
\item On the cone of an order $\pi$ of $\cE$, $y_{\pi(1)}\ge\dots\ge y_{\pi(n)}$, $\varphi$ is linear, and on the generator $-\one_{U_i}$ of the cone, $U_i$ the tails of $\pi$, it takes the value $\kappa_{1,U_i}$.
\item On the chart
\[
C_-e^{-\varphi}\le f\le C_+e^{-\varphi},\qquad C_-=\prod_j\Big(\sum_T\theta_{j,T}\Big)^{-s_{1j}-t_{1j}},\qquad C_+=\prod_j\big(\min_T\theta_{j,T}\big)^{-s_{1j}-t_{1j}}.
\]
\item For real $k>0$,
\[
\int e^{-k\varphi}dy=\frac{\mathrm{Cor}}{k^{n-1}},\qquad \mathrm{Cor}=\sum_\pi\prod_{i=2}^n\kappa_{1,U_i}^{-1}=(n-1)!\,\mathrm{vol}\{\varphi\le1\},
\]
so $C_-^k\,\mathrm{Cor}\le k^{n-1}P(k)\le C_+^k\,\mathrm{Cor}$ and $k^{n-1}P(k)\to\mathrm{Cor}$.
\item At the centre of a list, $\varphi$ is the support function of $\mathcal P-\mathcal P$, $\mathcal P$ the convex hull of the $\one_T$, $T\in\cB$; so $\{\varphi\le1\}$ is the polar body of $\mathcal P-\mathcal P$ and
\[
\mathrm{Cor}=(n-1)!\;\mathrm{vol}\big((\mathcal P-\mathcal P)^\circ\big).
\]
\end{enumerate}
\end{proposition}

\begin{proof}
(i) Along $\pi$ the largest of the $\langle\one_T,y\rangle$ is taken at the basis built greedily in the order $\pi$ (the exchange property), the same basis on the whole cone, so each maximum is linear there. On $-\one_U$, $\max_T(-|T\cap U|)=-d_j(U)$ and $\max_T(-|U\setminus T|)=-d^\ast_j(U)$ (Def.~7.10, p.~49), so $\varphi(-\one_U)=m_1(U)-\sum_j\big(s_{1j}d_j(U)+t_{1j}d^\ast_j(U)\big)=\kappa_{1,U}$ (Def.~8.1, p.~53). (ii) Each letter lies between its largest monomial times the least weight and the same monomial times the sum of the weights. (iii) Modulo $\one$ the cone of $\pi$ is generated by $-\one_{U_2},\dots,-\one_{U_n}$, a basis of the lattice of the chart, so $\int_{\text{cone}}e^{-k\varphi}dy=\prod_i(k\kappa_{1,U_i})^{-1}$, and homogeneity gives the volume; the squeeze is that of \S19, Steps 2 and 7. (iv) At the centre $1/f=K(\alpha)K(\alpha^{-1})$ (Def.~35.9, p.~309), whose monomials are the $\alpha^{\one_T-\one_{T'}}$.
\end{proof}

So the corner of Proposition~\ref{a:p:dirclose} at $j=0$ and of Proposition~\ref{a:p:corner} is the volume of one polytope, read cone by cone, and the squeeze of \S19 is general. For $\binom{[3]}1$, $\mathcal P-\mathcal P$ is a hexagon whose polar has area $3$, and $\mathrm{Cor}=2!\cdot3=6$. Sampling $\{\varphi\le1\}$ gives $12$, $16$ and $24$ for $\binom{[4]}2$, $H$ and $\binom{[4]}1$, the corners of \S15 and \S18.4.

\begin{proposition}[The two ends of the stride's law]\label{a:p:stridelaw}\derived\
Let $\mu$ be the stride's law (Definition~\ref{b:d:law}). It lives on $(0,f(\alpha^\star)]$, and for real $k>0$ and $q\notin[f(\alpha^\star)^{-1},\infty)$
\begin{gather*}
P(k)=\int t^k\,\mu(dt),\qquad G(q)=\int\frac{t}{1-qt}\,\mu(dt),\\
\mu\big([t,f(\alpha^\star)]\big)\sim\frac{\mathrm{Cor}}{(n-1)!}\Big(\log\frac1t\Big)^{n-1}\quad(t\to0^+).
\end{gather*}
At the centre of a list the level sets $\{f=t\}$ are the hypersurfaces
\[
K(\alpha)\,\Psi(\alpha)=t^{-1}\alpha^{\one}
\]
of degree $n$ in the projective chart, of dimension $n-2$, the same at every strength.
\end{proposition}

\begin{proof}
The first two are the image of $dy$ under $f$ and Thm.~35.10(ii), p.~309. For the third, $\{f\ge t\}$ lies between $\{\varphi\le\log(C_-/t)\}$ and $\{\varphi\le\log(C_+/t)\}$ by Proposition~\ref{a:p:tropcorner}, and $\mathrm{vol}\{\varphi\le s\}=s^{n-1}\mathrm{vol}\{\varphi\le1\}$. At the centre $\deg K+\deg\Psi=r+(n-r)=n=\deg\alpha^{\one}$: the balance of the centre puts the two sides in one degree. At strength $\sigma$ the stride is $f^\sigma$, with the same level sets.
\end{proof}

So the corner is the mass of $\mu$ near $0$, the native value of the table's gauge is the end of its support, and the branch of Thm.~35.11, p.~310, is its shape at that end. On two marks $\mu=2\,dt/\big(t\sqrt{1-4t}\big)$ on $(0,\frac14]$, whose moments are $\Gst21{k,\ k}{2k}$.

\begin{proposition}[The two ends of a certificate]\label{a:p:certends}\derived\
Let $\sum_{i=0}^hc_i(k)f^i=\sum_e\big(D_eg_e+k\,c_eg_e\big)$ certify a law of a ray through the origin, $c_e$ the defects of the direction, $D$ the largest degree of the $c_i$ and $g_e=\sum_{d<D}k^dg^{[d]}_e$. Then
\[
\chi(f)=\sum_ec_e\,g^{[D-1]}_e,\qquad \sum_{i\ge1}c_i(0)\,f^i=\sum_eD_e\,g^{[0]}_e .
\]
The first puts the symbol, read at the stride, in the ideal of the defects: at a native gauge every $c_e$ vanishes, which is Cor.~35.6, p.~308, read off a certificate. The second is a divergence. When the $g^{[0]}_e$ are bounded, integrating it over $\{\varphi\le R\}$ on the chart (with the mark $e_0$ handled as in \S19, Step 5(b)) and letting $R\to\infty$ gives
\[
\sum_{i\ge1}c_i(0)\,P(i)=\lim_{R\to\infty}\big(\text{flux of the field of the }g^{[0]}_e\text{ through }\partial\{\varphi\le R\}\big)=-[k^{n-1}]c_0\cdot\mathrm{Cor},
\]
the last by Corollary~\ref{a:c:dirlaw}: the corner is the flux at infinity of the bottom piece of a certificate.
\end{proposition}

\begin{proof}
Compare the coefficients of $k^D$ and of $k^0$; $c_0(0)=0$ by Corollary~\ref{a:c:dirlaw}. The left side of the second identity is integrable, so its integral over the chart is the limit of its integrals over $\{\varphi\le R\}$, and on these the divergence theorem applies.
\end{proof}

For the three-mark law of \S19 the top piece is $(1-f)(1-9f)$, the bottom piece $-3f+9f^2$, and the flux $-3\ell+9P(2)=-6$.

\begin{proposition}[The corner of two marks is the law of a logarithm]\label{a:p:gumbel}\derived; checked \computed\
For $\binom{[2]}1$, $P(k)=\Gst21{k,\ k}{2k}$, the corner is $2$, and with $X$ of the Gamma law of index $1$,
\[
\log\Big(\frac k2\,P(k)\Big)=2\log\oE[X^k]-\log\oE[X^{2k}]=\sum_{m\ge2}(-1)^m\,(2-2^m)\,\frac{\kappa_m}{m!}\,k^m,
\]
where $\kappa_m=(m-1)!\,\Tx{m}{\IZ}{0}$, $m\ge2$, are the cumulants of $-\log X$, the law that opens Chapter~G.
\end{proposition}

\begin{proof}
$\frac k2P(k)=\Gamma(1+k)^2/\Gamma(1+2k)$ and $\oE[X^k]=\Gamma(1+k)$; the cumulant generating function of $-\log X$ is $\log\Gamma(1-k)$, whose coefficients of order $m\ge2$ are $\Tx{m}{\IZ}{0}/m$. The terms of order one cancel, as Proposition~\ref{a:p:corner} says they do for every list.
\end{proof}

\subsection*{20.1\quad Is every ray examinable?}

\begin{question}[Is every ray examinable?]\label{a:q:exam}
On a ray through the origin with every proper wall positive, do the corners (the Laurent data of Corollary~\ref{a:c:dirlaw} at $k=0,-1,\dots$) and positivity always determine the values, with the masters held in brackets (polytopes, for several masters) whose widths shrink at the rate of Theorem~\ref{a:t:select}(iv)?
\end{question}

\emph{The only thing that can fail is a native value on the slit, where positivity says nothing. Positivity is a moment problem, and the exam used only its first two conditions.}

Hypothesis, as posed: positivity, all of it, examines a ray unless a native value sits on the slit; and a native value sits there only for a reason that can be named.

\begin{proposition}[Positivity is the moment problem of the stride's law]\label{b:p:moment}\derived\
For every polynomial $p\ge0$ on $[0,f^\star]$ and every $K\ge1$,
\[
\int_0^{f^\star}t^Kp(t)\,\mu(dt)>0 ,
\]
and in the coordinates of the masters of Theorem~\ref{a:t:select}(iv) each such condition is one linear inequality. The conditions with $p=q^2$, $p=t\,q^2$, $p=(f^\star-t)q^2$ and $p=t(f^\star-t)q^2$ generate all of them (alias, last). The examination used two: $p=1$, giving $P(K)>0$, and $p=f^\star-t$, giving $P(K)<f^\star P(K-1)$.
\end{proposition}

\begin{proof}
$\mu\ge0$ is carried by $[0,f^\star]$ (Proposition~\ref{b:p:law}(i)); a nonnegative polynomial times $t^K$ integrates to a positive number unless it vanishes. The values are affine in the masters.
\end{proof}

\begin{theorem}[The rate of positivity]\label{b:t:rate}Upper bound \derived\ from the expansion of Thm.~35.5; attained \computed.\
Let every root $\lambda_c\neq f^\star$ of the symbol that carries a local solution not holomorphic at $q_c=1/\lambda_c$ lie off the segment $[0,f^\star]$. Put $w=2\lambda/f^\star-1$ and $e^{g(\lambda)}=|w+\sqrt{w^2-1}|$, the branch with $e^{g}\ge1$. Then the brackets cut out by the conditions of Proposition~\ref{b:p:moment} of degree at most $N$ shrink like $e^{-N\min_cg(\lambda_c)}$ up to powers of $N$. For a real $\lambda_c>f^\star$, with $\rho=f^\star/\lambda_c$,
\[
e^{-g(\lambda_c)}=\frac{1-\sqrt{1-\rho}}{1+\sqrt{1-\rho}}=\frac\rho4+O(\rho^2),
\]
against $\rho$ for the two conditions of \S19.
\end{theorem}

\begin{proof}[Proof of the upper bound]
Take one master, $P=x^{(0)}+M\,x^{(1)}$, and $p=q(t)^2$ with $q(t)=T_m(2t/f^\star-1)$, $T_m$ the polynomial of degree $m$ with $T_m(\cos\vartheta)=\cos m\vartheta$, so that $0\le p\le1$ on $[0,f^\star]$. The condition $\int t^Kp\,\mu>0$ reads $A+BM>0$ with $B=\sum_jp_j\,x^{(1)}(K+j)$, and its error in $M$ is $\int t^Kp\,\mu/|B|\le P(K)/|B|$. By Thm.~35.5(i), p.~308, $x^{(1)}(k)=C\lambda_c^kk^{-(n-1)/2}(1+O(1/k))$ plus terms of smaller rate, $n=|\cE|$, so $B=C\lambda_c^K\int p\,d\nu+\dots$ for a positive measure $\nu$ on $[0,\lambda_c]$ with mass near $\lambda_c$, and $|B|\gtrsim|T_m(w(\lambda_c))|^2\lambda_c^K/\mathrm{poly}(m)$. As $|T_m(w)|\sim\frac12e^{mg(\lambda)}$ and $2m=N$, the error is $e^{-Ng}$ times powers of $N$. The side $p=(f^\star-t)q^2$ bounds $M$ from the other side in the same way. With several masters the same polynomials are used one direction at a time.
\end{proof}

\noindent Receipts \computed, on the ray of \S19: the full conditions of degree at most $N$, read on $P(1),\dots,P(N+1)$, against the two conditions of \S19 at $K=N+1$ (the same values). The interval is cut out exactly, its ends being the roots of Hankel determinants affine in $\ell$:
\begin{center}
\small
\begin{tabular}{@{}r l l l@{}}
\toprule
$N$ & two conditions, $b_K-a_K$ & all conditions & ratio of consecutive widths\\ \midrule
2 & $7.50\cdot10^{-2}$ & $4.73\cdot10^{-2}$ & \\
4 & $7.84\cdot10^{-4}$ & $6.90\cdot10^{-5}$ & $0.0382$\\
8 & $1.10\cdot10^{-7}$ & $9.52\cdot10^{-11}$ & $0.0335$\\
12 & $1.64\cdot10^{-11}$ & $1.04\cdot10^{-16}$ & $0.0320$\\
16 & $2.46\cdot10^{-15}$ & $1.03\cdot10^{-22}$ & $0.0314$\\
22 & $4.58\cdot10^{-21}$ & $9.10\cdot10^{-32}$ & $0.0308$\\
\bottomrule
\end{tabular}
\end{center}
The ratio tends to $e^{-g(1)}=17-12\sqrt2=0.0294\ldots$ ($\rho=\frac19$), slowly, as the powers of $N$ predict: $1.53$ digits of the three-mark centre per value against $0.954$. With the nine values of \S19 the whole moment problem gives $\ell$ to $10^{-10}$ instead of $10^{-7}$.

\begin{proposition}[What positivity cannot see]\label{b:p:blind}\derived\
If two solutions of the law with the same corner differ by $u(k)=\int_0^{f^\star}t^k\nu(dt)$ for a signed measure $\nu$ with $|\nu|\le C\mu$, then every condition of Proposition~\ref{b:p:moment} holds along $P+\varepsilon u$ for $|\varepsilon|\le1/C$: positivity cannot tell them apart. A root $\lambda_c$ on the slit whose local solution is not holomorphic at $q_c$ gives, by Theorem~\ref{a:t:select}(iii), one more solution holomorphic off $[1/f^\star,\infty)$ than the count needs; its jump across the slit is its measure $\nu$, singular at $\lambda_c$ with the exponent $\frac{n-3}2$ of a nondegenerate gauge (Thm.~35.11, p.~310, read at that gauge; \S20.6), so bounded by a multiple of $\varrho$ near $\lambda_c$ when $n\ge4$.
\end{proposition}

\begin{proof}
$\mu\pm\nu/C\ge0$. For the second statement: Theorem~\ref{a:t:select}(ii)--(iii) with a root on the slit leaves the space of solutions holomorphic off the slit of dimension one more than the count; a solution there is a Stieltjes transform of its jump.
\end{proof}

So positivity examines exactly the rays with no seen root in $(0,f^\star)$, and off the segment it does so at the rate of Theorem~\ref{b:t:rate}, complex roots inside the disc $|\lambda|<f^\star$ included (the two conditions of \S19 cannot see those). What remains is whether roots land on the slit.

\begin{proposition}[Where the slit gets its native values]\label{b:p:slitorigin}\derived\
\begin{enumerate}[label=(\roman*)]
\item At the centre of a list no native gauge on the torus $|\alpha_e|=1$, the signs among them, has its native value on the slit: there $K(\alpha^{-1})=\overline{K(\alpha)}$, so $f(\alpha)=|K(\alpha)|^{-2}\ge K(\one)^{-2}=f^\star$, with equality only at $\alpha\propto\one$.
\item If $f^\star$ is irrational, each of its conjugates is a native value: the equations of nativity have rational coefficients, $\alpha^\star$ is algebraic, its conjugates are native gauges, and their values are the conjugates of $f^\star$ (so they are roots of every symbol, Prop.~35.7, p.~309). A real conjugate in $(0,f^\star)$ sits on the slit.
\end{enumerate}
At the centre of a list $f^\star=|\cB|^{-2}$ is rational; at a lean it need not be.
\end{proposition}

\noindent Survey \computed: all native gauges (Newton's method in ball arithmetic from many complex starts) of
\begin{itemize}
\item the three-mark list $\binom{[3]}1$ at the $55$ integral directions with entries at most $4$ and positive walls: no native value in $|\lambda|\le f^\star$ other than $f^\star$;
\item $\binom{[4]}1$ ($144$ directions) and $\binom{[4]}2$ ($243$), with $s_1,t_1\le3$: none;
\item $H$, the triangle with one mark doubled (\S18.4), at the $74$ directions with $s_1,t_1\le3$ up to its symmetries: $10$ have a native value on the slit, every one of them at a gauge where the doubled class is flipped, $\alpha_a,\alpha_b<0$, and every one the conjugate of an irrational $f^\star$ (Proposition~\ref{b:p:slitorigin}(ii)).
\end{itemize}

\begin{example}[A ray strictly inside its walls that positivity cannot examine]\label{b:e:Hslit}\derived\ where marked; the rest \computed.\
Let $H$ carry the class $\{a,b\}$ and the marks $2,3$, with $A=\alpha_a+\alpha_b$, $B=\alpha_2+\alpha_3$, $K=AB+\alpha_2\alpha_3$, $\Psi=AB+\alpha_a\alpha_b$, and take the direction $\mvec_1=(1,1,3,3)$ on $(a,b,2,3)$, $s_1=3$, $t_1=1$. Its walls are $1$ on the marks and on $\{a,b\}$, $3$ or $4$ on the other pairs, $3$ on the triples; $\mathrm{Cor}=\frac{160}{27}$.
\begin{enumerate}[label=(\roman*)]
\item \derived\ On the slice $\alpha=(x,x,1,1)$, fixed by the symmetries $a\leftrightarrow b$, $2\leftrightarrow3$, the stride is $f=x/\big((4x+1)^3(x+4)\big)$, critical at $3x^2+8x-1=0$, so
\begin{gather*}
f^\star=f\Big(\tfrac{\sqrt{19}-4}3\Big)=\frac{7643+1064\sqrt{19}}{1366875}=0.0089846\ldots,\\
f\Big(\tfrac{-\sqrt{19}-4}3\Big)=\frac{7643-1064\sqrt{19}}{1366875}=0.0021985\ldots,
\end{gather*}
the conjugate, on the slit, at the gauge $(-2.786,-2.786,1,1)$, the class flipped.
\item The closings, $x_j=[\alpha^{j(1,1,3,3)}]K^{3j}\Psi^j=\sum_{b=0}^j\binom{3j}{2j+b}\binom jb\binom{2j-2b}{j-b}^2$ (expand both letters in $AB$ and read the two pairs), are $1,13,589,37597,2791949,\dots$; their least recurrence has order $4$ and degree $25$ ($261$ closings), and the law has the symbol $(1366875\lambda^2-15286\lambda+27)(\lambda-1)^2$: $f^\star$, its conjugate, and the double root $1$ of the curves of \S18.4. The count of Theorem~\ref{a:t:select}(iii) is $1$ (corner) $+\,2$ (at most, at $\lambda=1$) $<4$: law, corner and positivity leave one direction, the conjugate's. The procedure of Theorem~\ref{a:t:select}(iv) diverges: at $K=40,80,120$ it returns $P(1)\approx-1.6\cdot10^{22},-4.5\cdot10^{46},-1.3\cdot10^{71}$.
\item \derived\ On its share side, splitting both letters (Cor.~4.4, p.~26) and integrating the two pairs (two Beta corners, the scale between them a pin),
\[
\Gamma(3k)\Gamma(k)\,P(k)=\frac1{2\pi i}\int_{(c)}\frac{\Gamma(v)^6\,\Gamma(3k-v)\,\Gamma(k-v)}{\Gamma(2v)^2}\,dv,\qquad0<c<k,
\]
the integrand of Prop.~\ref{a:p:Hshare} for the centre of $H$ with the strengths of the two letters pulled apart. In ball arithmetic \computed:
\begin{gather*}
P(1)=0.48830319489526888388541464396010900874354031922042897524562\ldots,\\
P(2)=0.00132119179498773222717641714721510065075969610885135498475\ldots,
\end{gather*}
with $P(3)$, $P(4)$, $P(5)$ to the same precision; the law holds at $k=1$ to $60$ digits.
\item The second corner. The coefficients $c_0$ and $c_1$ share the factor $k+1$, and the constant term at $k=-1$ is rational:
\[
\CT_{k=-1}\big(c_0(k)P(k)+c_1(k)P(k+1)\big)=2077616120857361856,
\]
that is $12610378582200\,P(1)-133936231480200\,P(2)+160181043750000\,P(3)=5981712159277$ \computed\ ($30$ digits). With it the count is $1+1+2=4$, and law, corner, second corner and positivity read the values: from $K=20$ on they agree with (iii) to all $60$ digits \computed.
\end{enumerate}
\end{example}

\noindent So for the inputs of \S19, law, corner and positivity, the answer is no, and the reason has a name: positivity loses one direction exactly when a seen root sits on the slit; the roots that can sit there are not at the torus, and the first ones met are the conjugates of an irrational $f^\star$. In the example that direction is fixed by a second corner (\S20.4). The ray itself is examinable. From its gauge end, with the law and the Hessian at $\alpha^\star$ and neither corner nor positivity, it is read to $68$ digits, both corners coming out as outputs (Theorem~\ref{c:t:gauge}, Example~\ref{c:e:slitgauge}). The ``no'' belongs to the inputs, not to the ray.

\subsection*{20.2\quad How much of a period does its corner know?}

\begin{question}[How much of a period does its corner know?]\label{a:q:corner}
Bound $P(1)$ on every ray above and below by multiples of $\mathrm{Cor}$ read off the law at $k=0$ and $f(\alpha^\star)$, and decide how far the corner determines the first value.
\end{question}

\emph{The corner plus positivity already gives a two-sided bound: $\mathrm{Cor}/3<\ell<\mathrm{Cor}/2$ for three marks. At $(\one;0,2)$ on a graph, $P(1)$ is the period of the graph (the note, Corollary~5.5, p.~35), and the corner is the volume already observed to track that period (alias, last). The law explains why: it ties $k\to0$ to $k=1$ with rational coefficients.}

\begin{proposition}[A one-master ray is its corner times a continued fraction of its law]\label{b:p:cf}\derived\
Let the law have order $2$, $c_0(k)=a\,k^{n-1}+O(k^n)$ (Corollary~\ref{a:c:dirlaw}). Then
\[
P(1)=\mathrm{Cor}\cdot\frac{-a}{c_1(0)+c_2(0)\rho_1},\qquad \rho_k=\frac{P(k+1)}{P(k)}=\frac{-c_0(k)}{c_1(k)+c_2(k)\,\rho_{k+1}},
\]
the second the ratio of the minimal solution (Prop.~35.7), a continued fraction in the coefficients of the law. Since $0<\rho_1<f^\star$ (Proposition~\ref{b:p:moment} with $p=1$ and $p=f^\star-t$), $P(1)/\mathrm{Cor}$ lies between $-a/c_1(0)$ and $-a/\big(c_1(0)+f^\star c_2(0)\big)$ whenever both are positive. For three marks these are $\frac13$ and $\frac12$.
\end{proposition}

\begin{proof}
The corner relation (Corollary~\ref{a:c:dirlaw}) is $a\,\mathrm{Cor}+c_1(0)P(1)+c_2(0)P(2)=0$; divide by $P(1)$. The ratio $\rho_k$ of the solution selected by positivity satisfies the law divided by $P(k+1)$.
\end{proof}

With more masters the corner fixes the scale and the law fixes the ratio: the selection of Theorem~\ref{a:t:select}(iv) is homogeneous, so $P(1)/\mathrm{Cor}$ is a function of the law alone \derived. How much of the period the corner knows is the first bracket; what the law adds is that ratio.

\begin{theorem}[The fixture at $(\one;0,2)$, read off its ray]\label{b:t:fixture}The law \computed; the reading \computed; the corner exact.\
Let the fixture be the six-mark list $\mathrm M(K_4)$ (Ledger, \ref{o:MK4}), at the direction $(\one;0,2)$: $f=\alpha^{\one}/\Psi^2$, $P(1)$ its full cobasis lean at unit masses.
\begin{enumerate}[label=(\roman*)]
\item Its walls are $\kappa_{1,S}=2r(S)-|S|>0$ on every nonempty proper $S$, and $\mathrm{Cor}=84$.
\item The closings $x_j=[\alpha^{j\one}]\Psi^{2j}$ are $1,12,756,78960,10451700,1587790512,263964176784,\dots$, and on $31$ of them ($29$ equations, $18$ unknowns, one solution)
\begin{multline*}
(j+2)^5x_{j+2}-2(2j+3)(65j^4+390j^3+885j^2+900j+346)\,x_{j+1}\\
+16(j+1)(2j+1)(2j+3)(4j+3)(4j+5)\,x_j=0 .
\end{multline*}
Reflected (Corollary~\ref{a:c:dirlaw}), the law of the ray is
\begin{multline*}
-k^5P(k)+2(2k+1)(65k^4+130k^3+105k^2+40k+6)\,P(k+1)\\
-16(k+1)(2k+1)(2k+3)(4k+3)(4k+5)\,P(k+2)=0,
\end{multline*}
of symbol $-(1-4\lambda)(1-256\lambda)$: $f^\star=\frac1{256}$ and one other native value, $\frac14$.
\item The corner relation is $-84+12P(1)-720P(2)=0$, so $P(2)=\frac{P(1)-7}{60}$ and every value lies in $\mathbb Q+\mathbb Q P(1)$. The first bracket of Proposition~\ref{b:p:cf} is
\[
\frac{\mathrm{Cor}}{12}=7<P(1)<\frac{64}7=\frac{16\,\mathrm{Cor}}{147},
\]
and the brackets of \S19 shrink by $\frac1{64}$ per step (by Theorem~\ref{b:t:rate}, by $0.0039$ with all conditions). At $K=60$ the reading gives
\[
P(1)=7.21234141895756571239842896906869994458991775404299329075363\ldots,
\]
which is $6\,\Tx3{\IZ}0$ to all $60$ digits \recalled.
\end{enumerate}
\end{theorem}

\noindent So at $(\one;0,2)$ on the fixture the corner, the volume of the tropical stride (Prop.~\ref{a:p:dirclose}), knows the period of the graph within the factor $\frac{64}{49}$ before positivity is used past its first step, and the law makes it exact. The ratios $P(1)/\mathrm{Cor}$ met in this entry:
\begin{center}
\small
\begin{tabular}{@{}l l l l l@{}}
\toprule
ray & $\mathrm{Cor}$ & $P(1)$ & $P(1)/\mathrm{Cor}$ & law\\ \midrule
three-mark centre & $6$ & $2.343907238689\ldots$ & $0.3907$ & order $2$\\
fixture, $(\one;0,2)$ & $84$ & $7.212341418958\ldots$ & $0.0859=\frac1{14}\Tx3{\IZ}0$ & order $2$\\
$N$, $(\one;0,2)$ (\S20.7) & $66$ & $4.756802563758\ldots$ & $0.0721$ & order $6$\\
$N$, centre (\S20.7) & $63$ & $4.578912478882\ldots$ & $0.0727$ & order $11$\\
$V$, $(\one;0,2)$ (\S20.7) & $300$ & $7.12\pm0.01$ (check) & $0.0237$ & \notyet\\
\bottomrule
\end{tabular}
\end{center}

\subsection*{20.3\quad What do the slopes carry?}

\begin{question}[What do the slopes carry?]\label{a:q:slopes}
What numbers are the derivatives of a ray in its strength, the $P^{(j)}(i)$ at the integers and the coefficients of $k^{n-1}P(k)$ at the corner, and how do they enter the content of larger lists?
\end{question}

\emph{The law and the corner tie the derivatives together, $3P'(1)-9P'(2)=-4\ell-12$, with $P'(1)=-8.0751522635316\ldots$, not in $\mathbb Q+\mathbb Q\ell$ below $10^{18}$. In \S17 slopes of a smaller list became masters of $\binom{[5]}2$.}

Throughout this subsection $P(k)=\Pth$ at $k\one$, $\ell=P(1)$, and $\ell^{(j)}=P^{(j)}(1)$.

\begin{proposition}[The slopes are examinable by positivity, one corner per order]\label{b:p:slopes}\derived\
Write $k^2P(k)=\sum_{j\ge0}a_jk^j$ near $k=0$, the corners of every order ($a_0=\mathrm{Cor}=6$, $a_1=0$ by Prop.~\ref{a:p:corner}).
\begin{enumerate}[label=(\roman*)]
\item Differentiating the law (L) of \S19 $j$ times and letting $k\to0$ gives one relation among $P^{(i)}(1)$, $P^{(i)}(2)$, $i\le j$, and $a_j$; at $j=1$ it is the relation of the question, at $j=2$
\[
\tfrac32P''(1)-\tfrac92P''(2)=a_2-10P(1)-10P'(1)+9P(2)+18P'(2).
\]
\item $(-1)^jP^{(j)}(K)=\int t^K\log^j(1/t)\,\mu(dt)$: positive, and at most $f^\star$ times its value at $K-1$. The $j$-th slopes obey the law with an inhomogeneity built from lower slopes, so $P^{(j)}(K)$ is a rational combination of $\ell^{(j)}$ and of lower data, with the coefficient $c(K)$ of $\ell^{(j)}$ the solution of \S19 with $c(1)=1$, $c(2)=\frac13$.
\item So with $a_j$ known, positivity examines $\ell^{(j)}$ exactly as it examined $\ell$, at the same rate. Each order of slope costs one corner, and the first slope costs none. From the gauge end no order costs a corner (Proposition~\ref{c:p:gaugeslopes}).
\end{enumerate}
\end{proposition}

\begin{proof}
(i) Taylor-expand (L) at $k=0$ to order $j$. (ii) Differentiate $\int f^k\,dy$ under the integral, $0<f\le f^\star$. The homogeneous part of the differentiated law is (L), and $P^{(j)}(1)$ is the free initial value, the second being fixed by (i). (iii) As Step 9 of \S19.
\end{proof}

\begin{proposition}[The corners of the three-mark centre]\label{b:p:corners3}Order $2$ \derived; orders $3$ and $4$ \computed.\
\[
k^2P(k)=6-\Gamma(\tfrac12)^4\,k^2+4\,\Tx3{\IZ}0\,k^3+\Big(\tfrac7{60}\Gamma(\tfrac12)^8+\tfrac34\,\ell^2\Big)k^4+a_5k^5+O(k^6),
\]
with $a_5=-40.6602066583709935349082641262349054086879571\ldots$. The term of order two is the sum over the three pairs of the term of order $d$ of the two-mark centre, $\Gst21{d,\ d}{2d}=\frac2d-\frac13\Gamma(\frac12)^4d+O(d^3)$ (Prop.~\ref{a:p:r4zero}).
\end{proposition}

\begin{proof}[Proof of order two]
In the sector coordinates of Prop.~\ref{a:p:dirclose} every order $\pi$ has $\kappa_{1,U_2}=\kappa_{1,U_3}=1$ and the same sector letter, $f=q_2q_3\varphi$ with $\varphi=\big((1+q_2+q_2q_3)(1+q_3+q_2q_3)\big)^{-1}$, $\varphi(0)=1$. Expand $\varphi^k=\sum_mk^m\log^m\varphi/m!$ and split each $\log^m\varphi$ into its value at $0$, its restrictions to the two faces $q_3=0$ and $q_2=0$, and a remainder vanishing on both. In $\int_{[0,1]^2}q_2^{k-1}q_3^{k-1}(\cdot)\,dq$ the value at $0$ gives $k^{-2}$, a face gives $k^{-1}$ times $\int_0^1\big(g(q)-g(0)\big)\,dq/q$, the remainder $O(1)$. Only $m=1$ reaches $k^0$, through the faces, where $\varphi=(1+q)^{-1}$: each contributes $-\int_0^1\log(1+q)\,dq/q$, which is $-\frac1{12}\Gamma(\frac12)^4$ (the finite part at $d=0$ of the face's own two-mark letter, Def.~1.4, p.~11). Two faces, six orders: $a_2=-\Gamma(\frac12)^4$.
\end{proof}

\noindent Receipts \computed. The coefficients $a_j$ by Cauchy's integral on $|k|=0.42$ of the share side $k^2P(k)=k^2\,4^{2-k}\Gamma(k)^{-2}\int_0^\infty x^{2k-1}K_0(x)^3dx$ (alias), the part on $(0,1)$ summed exactly from the expansion of the integrand at $0$: $a_0=6$ and $a_1=0$ to $10^{-45}$, $a_2,a_3,a_4$ as displayed to $44$ digits, the relation for $a_4$ having coefficients $60,7,45$. With $a_2$ the second slope is examinable, and positivity gives
\begin{gather*}
\ell'=-8.07515226353164401026965124146737650134735782164232219222\ldots,\\
\ell''=31.2486570062485056456256488196147038526756\ldots,
\end{gather*}
$\ell'$ to $380$ digits ($K=400$) and $\ell''$ to $39$; a Cauchy integral of the share side around $k=1$ (a check) agrees to $14$ digits for both. There is no rational relation among $\ell'$ and $1$, $\ell$, $\ell\log2$, $\ell\log3$, $\Tx3{\IZ}0$, $\Gamma(\frac12)^4\log2$, $\Gamma(\frac12)^4\log3$, the cubes and mixed cubes of $\log2,\log3$, $\Gamma(\frac12)^6/\sqrt3$, $\frac{\Gamma(1/2)^2}{\sqrt3}$ times $\log^22,\log^23,\log2\log3,\log2,\log3,1$, $\Gamma(\frac12)^4$, $\log2$, $\log3$, $\log^22$, $\log2\log3$, $\log^23$ ($24$ numbers) with coefficients below $10^{13}$ ($350$ digits). With the arc integrals $S_3$ and $S_4$ admitted, $\ell'$ and $\ell''$ have relations of small height (Proposition~\ref{c:p:arcslopes}).

\begin{remark}[What the slopes carry]\label{b:r:slopes}
(i) New numbers, one per order: $\ell'$ is free of any corner datum, and for positivity $\ell''$ needs $a_2$, which the two-mark faces supply; from the gauge end neither needs one, and the new numbers are the arc integrals $S_3$, $S_4$ (Proposition~\ref{c:p:arcslopes}). (ii) The next mark: by Prop.~\ref{a:p:addmark} the four-mark centre is the closing of the three-mark slope input, $\Pfour=-\sum_{m\ge1}P'(m)$, and the five-mark centre that of the four-mark one (Remark~\ref{a:r:live}); slopes are how a list's numbers feed the list with one more mark. (iii) Along a hinge: the free slopes of \S17, $\Phi_1'(0)=2\ell-\ell'-\int\log X_{/}\,dy/X_{\setminus}$ and $f_1(1)=2\ell-\ell'-\int\log X_{H\setminus1}\,dy/X_{H/1}$, both carry $\ell'$, with a logarithmic integral of the children beside it; neither integral is a small rational combination of $1,\ell,\ell',\Tx3{\IZ}0,\Gamma(\frac12)^4,\Gamma(\frac12)^4\log2,\Gamma(\frac12)^4\log3,\ell\log2,\ell\log3,\log2,\log3,\Gamma(\frac12)^2/\sqrt3$, four at a time, coefficients below $2000$ (at the $27$ to $37$ digits known) \computed. (iv) The corners carry the faces and, from order four, the masters: $a_4$ holds $\ell^2$. A product of values written linearly in the slopes of order up to four is a relation of degree two, of the kind asked for in the Ledger, \ref{o:degtwo}.
\end{remark}

\subsection*{20.4\quad Are corners the channels of the tree?}

\begin{question}[Are the corners the channels of the tree?]\label{a:q:channels}
Along a wall, is the inhomogeneity that replaces the corner a value or a slope of the faces, so that every number a smaller list contributes to a larger one enters through a corner of a law; and are the roots of the symbol on the slit exactly the native values of the faces?
\end{question}

\emph{Along a wall the squeeze fails exactly in the face's directions, and the ray carries the faces (Thm.~35.15, p.~312). The question is whether every number a smaller list contributes enters as a corner. This turns ``the tree in one equation'' into a computation.}

The corners are read here at every order: $a_j=[k^j]\,k^{n-1}P(k)$ at $k=0$, and the same at the other poles $k\in-\mathbb N$ and, along a hinge, at the integer points of the split. The leading one is $\mathrm{Cor}$ (Prop.~\ref{a:p:dirclose}); a \emph{second corner} is a constant term at a pole where $c_0$ and $c_1$ share a zero (\S15.4).

\begin{proposition}[The corners of a ray are its faces]\label{b:p:facecorners}\derived\
In the sector coordinates of the proof of Prop.~\ref{a:p:dirclose}, $f=\prod_{i=2}^nq_i^{\kappa_i}\varphi_\pi(q)$ with $\kappa_i=\kappa_{1,U_i}$ and $\varphi_\pi(0)=1$. For $A\subseteq\{2,\dots,n\}$ let $\Delta_A$ restrict a function of $q$ to the face $\{q_i=0,\ i\notin A\}$ and subtract its restrictions to the faces of that face, so that $\Delta_Ah$ vanishes wherever some $q_i=0$, $i\in A$. Then
\[
a_j=\sum_\pi\ \sum_{A}\ \prod_{i\notin A}\kappa_i^{-1}\sum_{m+|\mathbf b|+|A|=j}\frac1{m!}\int_{[0,1]^A}\prod_{i\in A}\frac{(\kappa_i\log q_i)^{b_i}}{b_i!\,q_i}\ \Delta_A\big(\log^m\varphi_\pi\big)\,dq_A .
\]
So the corner of order $j$ is a sum, over the orders and over the faces of the sector cube of dimension at most $j$, of integrals of powers of the logarithm of the stride restricted to those faces. Up to order $n-1$ only proper faces enter, since $\Delta_A(\log^0\varphi_\pi)=0$ for $A\neq\varnothing$; the faces of the cube are the strata where a tail $U_i$ is scaled to zero, on which $\varphi_\pi$ is the product of the bottom letters of the bunch and its room (Lemma~35.14, p.~311). From order $n$ on the whole cube enters as well.
\end{proposition}

\begin{proof}
Thm.~8.10, p.~56, as in the proof of Prop.~\ref{a:p:dirclose}: $P(k)=\sum_\pi\int_{[0,1]^{n-1}}\prod_iq_i^{k\kappa_i-1}\varphi_\pi^k\,dq$. Write $\varphi_\pi^k=\sum_mk^m\log^m\varphi_\pi/m!$, decompose every function of $q$ as $\sum_A\Delta_A$ (inclusion--exclusion over the faces), and integrate: the coordinates outside $A$ give $\prod_{i\notin A}(k\kappa_i)^{-1}$, those in $A$ an integral analytic at $k=0$ whose Taylor coefficients are the logarithmic moments displayed. $\Delta_A(\log^0\varphi)=0$ for $A\neq\varnothing$ and $\log\varphi_\pi(0)=0$, which gives $a_1=0$ for unit weights (Prop.~\ref{a:p:corner}).
\end{proof}

At three marks this is the computation of Proposition~\ref{b:p:corners3}: $a_2$ is the two-mark faces, $-\Gamma(\frac12)^4=3\cdot[d]\Gst21{d,\ d}{2d}$, and $a_3=4\Tx3{\IZ}0$ and $a_4$ involve the whole cube, the second through $\ell^2$. So the corners carry the faces, and from order $n$ on the list itself.

\noindent\emph{The evidence, three kinds of ray.}
\begin{enumerate}[label=(\alph*)]
\item \emph{Strictly inside the walls, no root on the slit} (every centre ray of \S\S15--18). Law, corner and positivity determine the values (Theorem~\ref{a:t:select}(iii)); the higher corners are then consequences, through the slope relations of Proposition~\ref{b:p:slopes}. The faces' numbers reach the values as well: $\Pth$ and $\Gamma(\frac12)^4$ are among the masters of the triangle with two marks doubled (Theorem~\ref{a:t:r4masters}), selected by positivity, while its two corners in $\sigma$ are rational (\S15.4). On such a ray a face number is not confined to the corners.
\item \emph{Along walls} (the inputs of \S17, rays in the split $d$ between two children). The faces' native values sit on the slit (Prop.~\ref{a:p:u25roots}: $\frac13$ and $\frac14$ below $\frac49$), and the faces' numbers are exactly the boundary data of the input's law at the integer points of the split: $\Phi_1(0)=\Pth$, the deletion's centre; $\Phi_1(1)=\Pfour$, the contraction's; at the wall $d=2$ the residue $-6=-6\,P_{\binom{[2]}1}(\one;1,1)^2$ and the finite part $\frac{21}2+\frac34\Pfour$ (Prop.~\ref{a:p:u25wall}); at zero shares the class's $\Gamma(\frac12)^4$ (Prop.~\ref{a:p:r4zero}). Each number a child contributes enters there as a corner of the input.
\item \emph{Strictly inside the walls, a root on the slit} (Example~\ref{b:e:Hslit}). The direction that positivity cannot squeeze is fixed by the second corner, and that corner is rational. From the gauge end no condition is missing, and the second corner is an output (Example~\ref{c:e:slitgauge}).
\end{enumerate}

\begin{hypothesis}[Corners are the channels]\label{b:h:channels}
For every ray of a family, the numbers its faces contribute enter its law only through its corners, at the poles in the strength and, along a hinge, at the integer points of the split; positivity carries them from there into the masters.
\end{hypothesis}

\noindent What it turns into a computation, node by node of the Foreman tree, leaves first: (1) the law of the node's input in the split (or of its ray); (2) its corners, read from the children by Proposition~\ref{b:p:facecorners}; (3) positivity off the slit, or the gauge end (\S21), which needs neither (3) nor (4); (4) the corners on the slit; (5) the closing (summation by parts, Prop.~\ref{a:p:sbp}) up to the node's masters. Steps (1), (3), (5) are those of \S\S15.7 and 17; step (2) is derived above at the strength $0$; at the other poles and at the walls of a split it is the same expansion around $q^{\kappa\cdot j}$ (Prop.~\ref{a:p:dirclose}) and around the wall's face (Thm.~35.15), carried out at three marks only. That the second corner of Example~\ref{b:e:Hslit} is a sum of face numbers, and which: \notyet.

\subsection*{20.5\quad Is a ray's full set of solutions the coaction made visible?}

\begin{question}[Is the period matrix of a ray the face of the coaction?]\label{a:q:matrix}
Are the solutions of a law selected at its other native values and at the torus, the rest of the period matrix of the ray, built from the open periods of the small blocks that the note's Theorem~A attaches to the tree?
\end{question}

\emph{The law's other solutions, picked out at the other native values and at the torus, should be built from the small blocks of the note's Theorem~A. The exam already found one: the dominant rational solution is the closings.}

On the banana the other solutions have integral representations: the lowered settings of Remark~\ref{a:r:lowered}, $x^{(j)}(s)=\int_0^\infty y^{s-1}g^{n-j}h^j\,dy$ with $g=\sG_{0;(1,-1)}(1,\cdot)$ and $h(y)=\sum_my^m/\Gamma(1+m)^2$, the sign classes with $j$ minus signs, of native value $(n-2j)^{-2}$; divided by $\Gamma(s)^2$ they obey the ray law (Theorem~\ref{a:t:cert}), and their corners vanish (Remark~\ref{a:r:lowered}, computed there at five marks; at three and four marks in the proof below).

\begin{proposition}[The other solutions of the bananas]\label{b:p:otherbanana}At $k=1,2$ \computed; at every $k$ \derived\ from the law and the zero corner.\
Let $u_n$ be the rational solution of the ray law of $\binom{[n]}1$ with zero corner and $u_n(1)=1$.
\begin{enumerate}[label=(\roman*)]
\item Three marks: $u_3(k)=9^{1-k}[\alpha^0](K\Psi/\alpha^{\one})^{k-1}$, the closings (\S19.3, Step 8), and
\[
\frac{x^{(1)}(k)}{\Gamma(k)^2}=\frac13\Gst20{\frac13,\ \frac23}{}\;u_3(k).
\]
\item Four marks: $u_4(2)=\frac1{16}$, and $x^{(1)}(k)/\Gamma(k)^2=\frac14\,\Gamma(\frac12)^4\,u_4(k)$.
\item Five marks: the two sign classes span the zero-corner solutions with the coefficients of Prop.~\ref{a:p:stackedhalf}, the stacked Gamma
\[
X=\frac1{120}\Gst84{\frac1{15},\ \frac2{15},\ \frac4{15},\ \frac8{15},\ \frac15,\ \frac25,\ \frac35,\ \frac45}{\frac12,\ \frac12,\ \frac12,\ \frac12},
\]
$\Gamma(\frac12)^8/X$ and $\sqrt{15}/\Gamma(\frac12)^2$.
\end{enumerate}
\end{proposition}

\begin{proof}
At three and four marks the space of solutions with zero corner has dimension one (order $2$, one corner relation), and $x^{(1)}/\Gamma^2$ lies in it; so it is $x^{(1)}(1)\,u_n$. The values at $k=1,2$, by quadrature as a check ($18$ and $30$ digits): $x^{(1)}(1)=\frac{2\pi}{3\sqrt3}$ and $x^{(1)}(2)=\frac13x^{(1)}(1)$ at three marks; $x^{(1)}(1)=\frac{\pi^2}4$ and $x^{(1)}(2)=\frac{\pi^2}{64}$ at four (aliases, last).
\end{proof}

\noindent So the frame of the chart and the sign class (Prop.~35.8, p.~309) is a stacked Gamma of a two-mark block times the frame of \S19: at three marks
\[
P(k)\,\frac{x^{(1)}(k+1)}{\Gamma(k+1)^2}-P(k+1)\,\frac{x^{(1)}(k)}{\Gamma(k)^2}=\frac13\Gst20{\frac13,\ \frac23}{}\cdot\frac23\cdot\frac{9^{1-k}}{k^2}\qquad\derived .
\]

\begin{remark}[The period matrix of a banana ray]\label{b:r:periodmatrix}
The solutions split into the one that carries the corner, the chart, and those with zero corner. The latter are rational solutions times one number each, so their rows of the period matrix are a Gamma product times a rational row; the chart's row is the only one that mixes. That is the shape a coaction has: the pure part is read off the frames, and the mixed entries sit in one row.

Against Theorem~A of the note: Theorem~A is about left factors, periods of the same chain against other forms, and the other solutions are other chains. What can be compared is what the frames are made of. At three and four marks they are Beta corners $\Gst21{a,\ b}{a+b}$ of a two-mark block at masses in $\frac13\mathbb Z$ and $\frac12\mathbb Z$; at five marks the frame is the stacked Gamma $X$ at fifteenths, which is a product of Beta corners of two-mark blocks at fifteenths and halves (Proposition~\ref{c:p:frame}). So the hypothesis holds through five marks in the form ``the other solutions are rational solutions times algebraic multiples of products of Beta corners of two-mark blocks and their inverses''. At the fixture's full lean (Theorem~\ref{b:t:fixture}) the zero-corner solution has $u(2)=\frac1{60}$ and the other native value is $\frac14$; the number that multiplies $u$ there, by analogy a stacked Gamma of weight three: \notyet.
\end{remark}

\begin{proposition}[The five-mark frame is a product of Beta corners]\label{c:p:frame}\derived\
\begin{align*}
X&=\frac1{120}\Gst84{\frac1{15},\ \frac2{15},\ \frac4{15},\ \frac8{15},\ \frac15,\ \frac25,\ \frac35,\ \frac45}{\frac12,\ \frac12,\ \frac12,\ \frac12}\\
&=\frac1{15\sqrt{15}}\,\Gst21{\frac12,\ \frac12}{1}\Gst21{\frac1{15},\ \frac4{15}}{\frac13}\Gst21{\frac2{15},\ \frac8{15}}{\frac23}=8.5702804433744612681\ldots
\end{align*}
\end{proposition}

\begin{proof}
The reflection formula gives $\Gamma(\frac15)\Gamma(\frac45)\Gamma(\frac25)\Gamma(\frac35)=\Gamma(\frac12)^4/(\sin\frac\pi5\sin\frac{2\pi}5)=4\Gamma(\frac12)^4/\sqrt5$, so $X=\Gamma(\frac1{15})\Gamma(\frac2{15})\Gamma(\frac4{15})\Gamma(\frac8{15})/(30\sqrt5)$. Also $\Gamma(\frac13)\Gamma(\frac23)=2\Gamma(\frac12)^2/\sqrt3$. So the right side is $\Gamma(\frac1{15})\Gamma(\frac4{15})\Gamma(\frac2{15})\Gamma(\frac8{15})\cdot\sqrt3/(2\cdot15\sqrt{15})$, the same number.
\end{proof}

\begin{remark}[Balance is not an invariant]\label{c:r:balance}
``Its top entries sum to $3$, its bottom ones to $2$'' is a property of the written monomial. The reflection formula writes the unbalanced $\Gamma(a)\Gamma(1-a)$ (top $1$, bottom $0$) as an algebraic multiple of the balanced $\Gamma(\frac12)^2/\Gamma(1)$, and the proof above is three such moves. A Gamma monomial whose function $F(t)=\sum_an_a\langle ta\rangle$ on $(\mathbb Z/30)^\times$ is constant is algebraic up to a power of $\Gamma(\frac12)^2$ \recalled; the receipt finds $F_X=F$ of the product, identically, the same statement as the proposition. So the hypothesis above holds through five marks in its strong form: the coefficients of the zero-corner solutions are algebraic multiples of products of Beta corners of two-mark blocks and their inverses, for instance $\Gamma(\frac12)^8/X=15\sqrt{15}\,\Gamma(\frac12)^6\big/\Gst21{\frac1{15},\ \frac4{15}}{\frac13}\Gst21{\frac2{15},\ \frac8{15}}{\frac23}$, at thirds (three marks), halves (four), and fifteenths with halves (five). The fifteenths are the orbit $\{1,2,4,8\}$ of doubling in $(\mathbb Z/15)^\times$, as in the top row of $X$. Which block of the note's Theorem~A gives these corners: \notyet.
\end{remark}

\subsection*{20.6\quad What is the law of the stride?}

\begin{question}[What is the law of the stride?]\label{a:q:stride}
Is $\mu$ a member of the family, with its density on $(0,f(\alpha^\star))$ a value of the serif family; and what decides its moments when the native values do not?
\end{question}

\emph{$H$ and $\binom{[5]}1$ share native values but not numbers, and their level sets differ in dimension.}

The law is $\mu$ of Definition~\ref{b:d:law}; its density satisfies the law of the ray read on the other side of the Trailmoment.

\begin{proposition}[The density's law]\label{b:p:density}\derived\
Let $\sum_{i=0}^hc_i(k)P(k+i)=0$ hold identically, and $\theta=t\,d/dt$. Then on $(0,f^\star)$
\[
\Big(\sum_{i=0}^ht^i\,c_i(-1-i-\theta)\Big)\varrho=0,
\]
an equation of order $D=\max_i\deg c_i$ whose leading coefficient is $(-1)^D\chi(t)$, $\chi$ the symbol: its singular points in $\mathbb C^\times$ are the roots of the symbol, the native values. At $t=0$ the exponents are those of $c_0(-1-\theta)$, among them $-1$ with multiplicity $n-1$, which is the corner, $\varrho(t)\sim\mathrm{Cor}\,(-\log t)^{n-2}/\big((n-2)!\,t\big)$. At a simple root carried by a nondegenerate native gauge the exponents are $0,1,\dots,D-2$ and $\frac{n-3}2$ (with a logarithm when $n$ is odd); a family of native gauges of dimension $d$ (Remark~\ref{a:auto:15.16}) moves the last one to $\frac{n-3-d}2$.
\end{proposition}

\begin{proof}
$k\,t^k=\theta\,t^k$, and integrating $\theta$ by parts against $dt$ gives $-1-\theta$; $c_i(k)t^{k+i}=t^{k+i}c_i(\ldots)$ moves $t^i$ through, $\theta t^i=t^i(\theta+i)$. The law holds if and only if $\varrho\one_{(0,f^\star)}$ solves the equation in the sense of distributions (Theorem~\ref{c:t:gauge}(ii)); in particular $\varrho$ solves it on the open interval. The exponents at the native values are those of the branch of Thm.~35.11, read at each gauge; the families are the case of Remark~\ref{a:auto:15.16}.
\end{proof}

\noindent Receipts \computed: the exponents, from the laws of \S\S15--19 and Theorem~\ref{b:t:fixture}.
\begin{center}
\small
\begin{tabular}{@{}l c l l@{}}
\toprule
ray & $n$ & native values & exponents of $\varrho$ beyond $0,\dots,D-2$\\ \midrule
$\binom{[3]}1$, centre & $3$ & $\frac19,\ 1$ & $0$ with logarithm at both\\
$\binom{[4]}1$, centre & $4$ & $\frac1{16},\ \frac14$ & $\frac12$ at both\\
$\binom{[5]}1$, centre & $5$ & $\frac1{25},\ \frac19,\ 1$ & $1$ with logarithm at all three\\
$H$, centre & $4$ & $\frac1{25},\ \frac19,\ 1\ (\text{double})$ & $\frac12$ at $\frac1{25}$ and $\frac19$; $0$ thrice at $1$\\
fixture, $(\one;0,2)$ & $6$ & $\frac1{256},\ \frac14$ & $\frac32$ at both\\
\bottomrule
\end{tabular}
\end{center}
At the native value $1$ of $H$ the indicial polynomial is $s^3(s-1)(s-2)(s-3)(s-4)$: besides the holomorphic exponents, $0$ twice more, the logarithms of the curves of Prop.~\ref{a:p:Hgauges}.

\begin{remark}[$H$ against $\binom{[5]}1$]\label{b:r:Hfive}
The two laws of the stride have the same singular points, $\frac1{25},\frac19,1$, and different exponents at every one of them: $\frac12$ against $1$ with a logarithm at the gauge and at the sign, and at $1$ a triple exponent $0$ from the curves of native gauges against a nondegenerate sign class. ``The level sets differ in dimension'' is the exponent $\frac{n-3}2$, the level sets $\{f=t\}$ having dimension $n-2$, two for $H$ and three for the banana; and the curves at $1$ are the critical locus of dimension one that the banana does not have. The numbers of a ray are the connection data of this one equation between its ends, the corner at $t=0$ and the branch at $f^\star$, so the native values alone do not decide them (Remark~\ref{a:r:Hnot}): the exponents do their part. Whether one law is a transform of the other is taken up in Remark~\ref{c:r:transform} below (Ledger, \ref{o:transform}); such a transform would have to turn the exponent $\frac12$ at two singular points into $1$ with a logarithm while turning a triple $0$ at the third into a single one.
\end{remark}

\begin{remark}[What adding a mark accounts for, and what it does not]\label{c:r:transform}\derived\ from Prop.~\ref{b:p:density} and the display after Prop.~\ref{b:p:law}.\
At a nondegenerate native gauge the exponent beyond $0,\dots,D-2$ is $\frac{n-3}2$, with a logarithm when $n$ is odd; at a family of native gauges of dimension $d$ it is $\frac{n-3-d}2$. A transform from the law of $H$ ($n=4$) to that of $\binom{[5]}1$ ($n=5$) adds a mark. At the gauges $\frac1{25}$ and $\frac19$ it must move $\frac12$ to $1$ with a logarithm, which is exactly what adding a mark does there. These two mismatches of Remark~\ref{b:r:Hfive} are the signature of such a transform; they rule out only the transforms that keep $n$, the rational changes of the unknown, which keep exponents modulo $\mathbb Z$ and logarithms. At $t=1$ adding a mark moves the exponent of the curves of native gauges ($d=1$) from $0$ to $\frac12$, while $\binom{[5]}1$ has $1$ with a logarithm, the exponent of an isolated sign class. The curves would have to become a point as well. Adding a mark does not account for that, and this mismatch stands as a constraint on any transform. On bananas the mark-adding transform is at hand. The display after Prop.~\ref{b:p:law} makes $\sG_{0;(1,-1)}(1,\cdot)^{n}$ a $\mu_n$-mixture of dilations. Applying it once more, $\mu_{n+1}=\int\nu_t\,\mu_n(dt)$, with $\nu_t$ the stride's law of two marks at coefficients $(1,t)$: splitting the two letters of $\sG_{0;(1,-1)}(1,y)\,\sG_{0;(1,-1)}(1,y/t)$ at their sum gives $\int_0^1\sG_{0;(1,-1)}\big(1,y(\frac1w+\frac1{t(1-w)})\big)\frac{dw}{w(1-w)}$. Its density has the exponent $-\frac12$ at its gauge end (Prop.~\ref{b:p:density} at $n=2$): a convolution with a square root, which at a nondegenerate gauge turns a sign into a logarithm and a logarithm into a sign. Whether $H$ and $\binom{[5]}1$ are related by such a transform together with one that shrinks the curves at $t=1$: \notyet.
\end{remark}

So the answer to the sixth question: the law of the stride is the base measure of the ray as an exponential family (Proposition~\ref{b:p:law}(iii)), the input of its Trailmoment on a line (Proposition~\ref{b:p:law}(i)), whose Laplace transform is one Gramma (Proposition~\ref{b:p:law}(ii)) and whose density obeys the law of the ray read backwards (Proposition~\ref{b:p:density}); its singular points are the native values and its exponents the dimensions of the level sets and of the critical loci.

\subsection*{20.7\quad What does a list without a realization make?}

\begin{question}[What does a list without a realization make?]\label{a:q:unreal}
What are the contents, at the centre and at the full lean $(\one;0,2)$, of prime lists that are not lists of spanning trees: the note's list of rank three on six marks whose bases are the triples other than $\{1,2,3\}$, and a list on eight marks that is realized over no field?
\end{question}

\emph{Two candidates: the note's prime rank-3 list on six marks, and an eight-mark list realizable over no field. The second is balanced, with every wall positive, at $(\one;0,2)$.}

Nothing in \S\S15--19 or above uses a realization: the stride, the walls, the corner, the closings, the law and the selection are read off the bases. So both lists make rays, and the question is what the rays hold.

\begin{definition}[The two lists]\label{b:d:NV}
$N$ is the list of rank three on $[6]$ whose bases are all triples but $a=\{1,2,3\}$ (the note, Remark~5.6, p.~37); put $b=\{4,5,6\}$. It is realized by six points of a plane, three of them on a line, and by no graph: contracting the mark $4$ and deleting the mark $6$ leaves $\binom{\{1,2,3,5\}}2$, a minor that the list of a graph never has \recalled. $V$ is the list of rank four on $[8]$ with the pairs $P_1=\{1,2\}$, $P_2=\{3,4\}$, $P_3=\{5,6\}$, $P_4=\{7,8\}$, whose bases are the $65$ four-sets other than the five unions $P_i\cup P_j$ with $\{i,j\}\neq\{3,4\}$; it is realized over no field \recalled.
\end{definition}

\begin{proposition}[The letters and the closings of $N$]\label{b:p:Nclos}\derived\
Let $K_a=\alpha_1+\alpha_2+\alpha_3$ and $\Psi_a=\alpha_1\alpha_2+\alpha_1\alpha_3+\alpha_2\alpha_3$ be the letters of $\binom a1$, $\alpha^a=\alpha_1\alpha_2\alpha_3$, and the same on $b$.
\begin{enumerate}[label=(\roman*)]
\item With $Q=\Psi_aK_b+K_a\Psi_b$,
\[
K_N=Q+\alpha^b,\qquad \Psi_N=Q+\alpha^a .
\]
\item For $i\equiv l\pmod 3$ let $N(i,l)$ be the coefficient of $\alpha^{D\one}$ in $\Psi^iK^l$ on $\binom{[3]}1$, $D=\frac{2i+l}3$: the first closing of the three-mark direction $(D\one;l,i)$. Then
\[
N(i,l)=\sum_w\binom iw\binom l{w+\frac{l-i}3}\binom{2D-2w}{D-w},
\]
and the closings of $N$ are
\begin{align*}
\text{centre:}\quad x_j&=[\alpha^{j\one}](K_N\Psi_N)^j=\sum_{p,q=0}^j\binom jp\binom jq\binom{p+q}{2q-p}\,N(2q-p,2p-q)\,N(2p-q,2q-p),\\
(\one;0,2):\quad x_j&=[\alpha^{j\one}]\Psi_N^{2j}=\sum_{p}\binom{2j}p\binom p{2p-3j}\,N(2p-3j,3j-p)\,N(3j-p,2p-3j),
\end{align*}
$1,19,1567,224821,41485807,\dots$ and $1,18,1494,211860,38486070,\dots$.
\end{enumerate}
\end{proposition}

\begin{proof}
(i) A triple of $[6]$ other than $a$ and $b$ meets $a$ in one or two marks, so its monomial is a mark of $a$ times a pair of $b$ or the reverse; $b$ is a basis and $a$ is not. The cobases are the complements of the bases, all triples but $b$. (ii) In $(K_N\Psi_N)^j$ take $u_1,u_2,u_3$ factors $\Psi_aK_b,K_a\Psi_b,\alpha^b$ from $K_N^j$ and $v_1,v_2,v_3$ factors $K_a\Psi_b,\Psi_aK_b,\alpha^a$ from $\Psi_N^j$. The part on $a$ is $\Psi_a^{\,u_1+v_2}K_a^{\,u_2+v_1}(\alpha^a)^{v_3}$ and the part on $b$ is $\Psi_b^{\,u_2+v_1}K_b^{\,u_1+v_2}(\alpha^b)^{u_3}$, so the two triples separate and each gives an $N(i,l)$ with $i=u_1+v_2$, $l=u_2+v_1$; with $p=j-u_3$, $q=j-v_3$ the degrees force $i=2q-p$, $l=2p-q$, and the sum over $u_1$ is Vandermonde's. The lean is the same with one letter. For $N(i,l)$ write $\Psi=\alpha_1\alpha_2+\alpha_3(\alpha_1+\alpha_2)$, $K=\alpha_3+(\alpha_1+\alpha_2)$ and read the power of $\alpha_3$, then the central binomial coefficient of $\alpha_1+\alpha_2$.
\end{proof}

\noindent So the closings of $N$ are bilinear in first closings of the three-mark list along all its directions, one factor per triple: the classes of the lean, joined at the centre by $i\mapsto i+3$ (the note, Remark~5.6). From them the closings cost $O(j^3)$ operations, against the expansion of a power of a letter in five coordinates.

\begin{theorem}[$N$ at $(\one;0,2)$]\label{b:t:Nlean}The law and the reading \computed; the native gauge \derived.\
\begin{enumerate}[label=(\roman*)]
\item Every wall is positive (the least is $1$) and $\mathrm{Cor}=66$.
\item The closings obey a recurrence of order $6$ and degree $24$ (found on $220$ closings, verified on all $221$). The law of the ray has order $6$, $c_0(k)=k^5(k+1)^3(k+2)(k+3)\,r(k)$ with $r$ irreducible of degree $14$, and symbol
\[
\chi(\lambda)\propto(24\lambda-1)(10800\lambda^2+324\lambda-1)(1728\lambda^3-320\lambda^2+28\lambda+1).
\]
\item \derived\ On the slice $\alpha=(x,x,x,1,1,1)$, the fixed locus of the permutations of $a$ and of $b$, $f=x/(x^2+9x+9)^2$, critical at $x^2+3x-3=0$, so $\alpha^\star=\big(\tfrac{\sqrt{21}-3}2\one_a,\one_b\big)$ and
\[
f^\star=\frac{7\sqrt{21}-27}{1800}=0.0028211277\ldots,
\]
irrational; its conjugate $-\frac{7\sqrt{21}+27}{1800}$ is negative. The other native values are $\frac1{24}$ and the roots of the cubic, $-0.0265231$ and $0.105854\pm0.103023\,i$ (modulus $0.147712$): none lies on $[0,f^\star]$, so positivity examines the ray (\S20.1), and the selection converges at the rate $\big(f^\star/0.0265231\big)^K$.
\item From $K=60$ on the reading is stable; at $K=200$ it is good to $189$ digits, the rate being $10^{-0.973K}$ (the gauge end gives $500$, \S21.4):
\[
P(1)=4.75680256375822944340803728605708480533517181028275669568732\ldots,
\]
and $P(1)/\mathrm{Cor}=0.07207$. Check: tropical sampling gives $4.751\pm0.010$.
\item The relations among $1,P(1),\dots,P(6)$ form a lattice of rank $3$: the values lie in $\mathbb Q+\mathbb QP(1)+\mathbb QP(2)+\mathbb QP(3)$, three masters against the fixture's one (Theorem~\ref{b:t:fixture}), and the content holds the three-mark centre (Proposition~\ref{d:p:Ncontent}(i)--(ii)). The constants searched against first, the \emph{numbers of \S20.7}, were $1$, $\Tx3{\IZ}0$, $\Gamma(\frac12)^4$, $\ell$, $\Gamma(\frac12)^4\log2$, $\Gamma(\frac12)^4\log3$, $\ell\log2$, $\ell\log3$, $\Lambda_{-7}$, $\Lambda_{-8}$, $\Lambda_{-24}$, $\Lambda_{-84}$, $\Gamma(\frac12)^2/\sqrt3$, $\log2$, $\log3$, where $\ell=\Pth$ and $\Lambda_D=\sum_{a=1}^{|D|}\big(\frac Da\big)\Tx2{\Ilat{a+|D|\mathbb Z}}0$ with $\big(\frac Da\big)$ the Kronecker sign, and the seven numbers $1$, $\Tx3{\IZ}0$, $\Gamma(\frac12)^4$, $\varepsilon$, $\Gamma(\frac12)^4\varepsilon$, $\varepsilon^2$, $\varepsilon^3$ with their products with $\sqrt{21}$, and $\log2,\log3,\log7$, where $\varepsilon=\log\frac{5+\sqrt{21}}2$, the logarithm of the unit of the field of $f^\star$.
\end{enumerate}
\end{theorem}

\begin{proof}[Proof of (iii)]
On the slice $\Psi_N=x^3+9x^2+9x$ and $\alpha^\one=x^3$. A critical point of $f$ on the fixed locus of a group of symmetries of $f$ is a critical point of $f$; on the orthant there is one (Thm.~35.10(i), p.~309). With $x^2=3-3x$, $x^2+9x+9=6(x+2)$ and $f^\star=x/\big(36(x+2)^2\big)$.
\end{proof}

\begin{proposition}[$N$ at the centre]\label{b:p:Ncentre}The law of order $11$ \computed\ modulo two primes, that of order $12$ exactly; the native values \derived\ where marked.\
\begin{enumerate}[label=(\roman*)]
\item Every wall is positive and $\mathrm{Cor}=63$; $f^\star=19^{-2}$ at $\alpha^\star=\one$ \derived\ (Proposition~\ref{b:p:slitorigin}(i)).
\item The closings (to $j=800$ modulo $2^{61}-1$, by Proposition~\ref{b:p:Nclos}) obey no recurrence of order at most $10$ and degree at most $70$, one of order $11$ and degree $52$ (one solution on $790$ equations), and one of order $12$ and degree $31$. The symbol of the law of order $11$ is, modulo two primes of $61$ and $62$ bits with rational reconstruction,
\[
\chi(\lambda)\propto(1-361\lambda)(1-25\lambda)(1-9\lambda)(1-\lambda)^5(1+80\lambda)(361\lambda^2-34\lambda-3).
\]
\item \derived\ The native values $\frac1{25},\frac19,1$ are the signs: at the centre every sign vector $\sigma$ is a native gauge with $f(\sigma)=K_N(\sigma)^{-2}$, and $|K_N(\sigma)|$ takes the values $19,5,3,1$ on $2,12,18,32$ sign vectors. The remaining native values are $-\frac1{80}$ and $\frac{17\pm14\sqrt7}{361}$ ($0.14970$ and $-0.05550$), off the torus.
\item No native value lies on $[0,f^\star]$, and $f^\star$ is simple: positivity examines the ray, with the rate $(80/361)^K$.
\item \computed\ The law of order $12$ and degree $31$ exactly: integer coefficients of $1305$ to $1343$ digits, by Chinese remaindering over $160$ primes of $62$ bits (stable from $150$), verified exactly on the $221$ closings computed exactly and modulo a third prime on $801$. Its symbol is that of (ii) times $\lambda-\lambda_s$, with $\lambda_s=0.44992\ldots$ of height $2581$ digits, the root that the extra order adds, off the slit; $c_0$ vanishes to order $5$ at $k=0$, so the corner relation holds.
\item \computed\ Law, corner and positivity read the ray. The readings at $K=150,250,350$ differ by $4.1\cdot10^{-91}$ and $1.5\cdot10^{-156}$, the rate $(80/361)^K$, and to $220$ digits
\[
P(1)=4.57891247888240480638020373022855744921869580872760320929687575\ldots,
\]
$P(1)/\mathrm{Cor}=0.07268$. Check: tropical sampling gives $4.574\pm0.008$.
\item The relations among $1,P(1),\dots,P(12)$ form a lattice of rank $4$: eight masters, and the content holds $\Pth$ and $\Tx3{\IZ}0$ (Proposition~\ref{d:p:Ncontent}(iii)--(iv)); among $1,P(1),\dots,P(10)$ there are two relations, at height $10^{33}$. There is none of height below $10^{12}$ ($210$ digits) among $P(1)$ alone and $1$, $\Tx3{\IZ}0$, $\Gamma(\frac12)^4$, $\ell$, $\Gamma(\frac12)^4\log2$, $\Gamma(\frac12)^4\log3$, $\Gamma(\frac12)^4\log7$, $\ell\log2$, $\ell\log3$, $\Lambda_{-4}$, $\Lambda_{-7}$, $\Lambda_{-8}$, $\sqrt7\,\Gamma(\frac12)^2$, $\log2$, $\log3$, $\log7$.
\end{enumerate}
\end{proposition}

\begin{proposition}[The eight-mark list $V$]\label{b:p:V}\derived\ where marked; the rest \computed.\
\begin{enumerate}[label=(\roman*)]
\item The walls at $(\one;0,2)$, $(\one;1,1)$ and $(\one;2,0)$ are all positive, with values $1,2,3,4$, and $\mathrm{Cor}=300$, $\frac{896}3$, $300$.
\item \derived\ The cobases are the four-sets other than $P_k\cup P_l$ with $\{k,l\}\neq\{1,2\}$. On the slice $\alpha=(x\one_{P_1\cup P_2},\one_{P_3\cup P_4})$, the fixed locus of the swaps inside the pairs and of $P_1\leftrightarrow P_2$, $P_3\leftrightarrow P_4$, $\Psi_V=x(x^3+16x^2+32x+16)$ and $f=x^2/(x^3+16x^2+32x+16)^2$, critical at
\[
x^3+8x^2-8=0,\qquad f=\frac{x^2}{64\,(x+1)^2(x+3)^2}\ \text{there}.
\]
The root $x=0.945667817991$ gives $f^\star=0.000237093045489\ldots$, of degree three; the conjugates, at $x=-1.07480315405$ and $x=-7.87086466394$, are $0.870337682208$ and $0.000864231320476$, both above $f^\star$: off the slit.
\item A census of native gauges (Newton's method from random complex starts: $31$ gauges, $10$ values, and the slice) finds the native values $f^\star$, $0.000818456$, $0.000864231$, $0.004487298$, $-0.013533898$, $0.007758455\pm0.0136\,i$, $0.021107646$, $-0.026238829$, $0.06696851$, $0.870337682$, none in $(0,f^\star)$. It is a lower bound: only a symbol completes it.
\item The closings $x_j=[\alpha^{j\one}]\Psi_V^{2j}$ are $1,64,34080,39638080,65433978400,131819906744064,\dots$ ($j\le9$, by the dense expansion modulo primes, $5$ seconds); $x_1=64$ counts the cobases whose complement is a cobasis, all but $P_1\cup P_2$. They grow like $(f^\star)^{-j}=4217.7^j$.
\item Check: tropical sampling gives $P(1)=7.124\pm0.009$, so $P(1)/\mathrm{Cor}\approx0.0237$.
\item The law: \notyet. The dense expansion of $\Psi_V^{2j}$ in seven coordinates grows like $j^7$ per closing, and a law of the size of $N$'s needs some hundreds of closings; what $V$ lacks is the structure of Proposition~\ref{b:p:Nclos}(ii), a splitting of the letters along classes that would make the closings a short sum.
\end{enumerate}
\end{proposition}

\begin{remark}[What the two lists make]\label{b:r:NV}
(i) Rays like any other. The machinery has no step at which a realization could enter, and none does: $N$ is examinable at both directions, by the symbols, and read at $(\one;0,2)$; $V$ has its walls, its corners, its gauge and its first closings, and fails only for want of a law. (ii) What the realization changed was arithmetic, not existence. A list of a graph whose marks are alike has $\alpha^\star=\one$ and a rational $f^\star$ at every balanced direction; $N$ at its lean and $V$ have marks of two kinds, $\alpha^\star$ is irrational, and by Proposition~\ref{b:p:slitorigin}(ii) the conjugates of $f^\star$ are native values. For both lists they fall off the slit, below $0$ for $N$, above $f^\star$ for $V$; so the mechanism of Example~\ref{b:e:Hslit} is present but does not bite. (iii) Numbers of its own and of its minors: $N$'s lean has three masters, its content holds $\Pth$, and two masters are new (Proposition~\ref{d:p:Ncontent}, Remark~\ref{d:r:minors}). The fixture at the same direction has one master and the number $6\Tx3{\IZ}0$. So the list without a graph realization carries, at the direction where graphs carry their periods, a ray with a six-dimensional law, a minor's number, and numbers not among the ones the graphs of \S\S15--19 produce. (iv) At the centre, where $N$ is its own and only factor and an atom of its Foreman tree (the note, Remark~5.6), the law has order $11$, eight masters, and native values in $\mathbb Q$ and $\mathbb Q(\sqrt7)$; the period of the atom is read off its ray, $4.5789124788\ldots$, and its content holds the three-mark and the four-mark centres, with six masters new. The price is arithmetic: the law that could be made exact has coefficients of $1300$ digits, against $33$ at the lean.
\end{remark}

The values of $N$ read from its gauge end (\S21.4: $500$ digits at the lean, $520$ at the centre) settle its masters and its content.

\begin{proposition}[The content of $N$]\label{d:p:Ncontent}\computed: lattice reduction on the gauge-end values, $480$ digits at the lean, $520$ at the centre. Every relation below has residual under $10^{-570}$; the generic vectors of the same lattices have residuals above $10^{-480}$.
\begin{enumerate}[label=(\roman*)]
\item At $(\one;0,2)$, the rational relations among $1,P(1),\dots,P(6)$ form a lattice of rank $3$; the next vectors have height $10^{113}$, the generic size. The lattice contains the corner relation and the law's constant term at $k=-1$, where $c_0$ and $c_1$ share $(k+1)^3$: $\CT_{k=-1}\big(c_0(k)P(k)+c_1(k)P(k+1)\big)$ is an integer (it is $3314356313112576$ for the law normalized as the reflection of the primitive recurrence of the closings). The constant terms at $k=-2$ and $k=-3$, where $c_0$ and $c_1$ share $k+2$ and $k+3$, are not rational: their relations are not in the lattice. The origin of the third relation: \notyet. So the values lie in $\mathbb Q+\mathbb QP(1)+\mathbb QP(2)+\mathbb QP(3)$: three masters, not five.
\item The content of the lean holds the three-mark centre:
\[
3906\,P(1)-506772\,P(2)+10837800\,P(3)=15263+1047\,\Pth ,
\]
residual $10^{-662}$; with the positivity values of \S20, good to $189$ digits, the residual is $10^{-185}$. With all the numbers of \S20.7 added, no further relation below heights $10^{26}$ and $10^{23}$.
\item At the centre, the relations among $1,P(1),\dots,P(12)$ form a lattice of rank $4$: the law at $k=1$, the corner, and two more ($c_0$ and $c_1$ share $(k+1)^2$, $k+2$, $k+3$). Among $1,P(1),\dots,P(10)$ there are two, among $1,P(1),\dots,P(9)$ none below height $10^{52}$. So eight masters, not ten.
\item The content of the centre holds the three-mark and the four-mark centre: among $1$, $P(1),\dots,P(8)$, $\Pth$ and $\Zt$ the relations form a lattice of rank $2$, heights $1.4\cdot10^{25}$ and $1.6\cdot10^{25}$. Two further searches, one adding $\Gamma(\frac12)^4$, $\Gamma(\frac12)^8$, $(\Pth)^2$ and $\Gamma(\frac12)^4\Pth$, the other $\Lambda_{-7}$, $\sqrt7\,\Gamma(\frac12)^2$, $\log2$, $\log3$ and $\log7$, find no third relation below height $10^{32}$.
\end{enumerate}
\end{proposition}

\begin{remark}[Whose numbers]\label{d:r:minors}
$N$ has the minors $\binom{\{1,2,3,6\}}1=N/\{4,5\}$ and $\binom{\{1,2,3,5\}}2=(N/4)\setminus6$. The four-mark rank-one list is rational at both full leans and $7\Zt$ at its centre (the note, \S1 and Example~6.3). The four-mark rank-two list is $\Pth$ at its centre (the note, Example~6.2), and also at its full lean \computed: at $(\one;0,2)$ it has a law of order $2$, $f^\star=\frac1{36}$, $\mathrm{Cor}=12$, and its value from the gauge end is $\Pth$ to $3\cdot10^{-308}$. So on $N$, at both directions, the content holds the same-direction periods of these two minors, and nothing of the four-mark rank-one list where that list is rational; besides them, $2$ masters at the lean and $6$ at the centre are new. That minors' numbers enter is not a rule: the centre of $H$ has the content $\mathbb Q+\mathbb QP_H(\one;1,1)+\mathbb QP_H(2\one;2,2)$ (\S18.4), and there is no relation among $1$, those two values and $\Pth$ below height $10^{94}$ ($380$ digits, the values from the gauge end), though $H$ has the three-mark minors $\binom{[3]}1$ and $\binom{[3]}2$, whose centres are $\Pth$; with $\Zt$ and $\Gamma(\frac12)^4$ added, none below $10^{63}$. $N$ is prime (the note, Remark~5.6); $H$ is not, its parallel pair having width one. Which minors' numbers enter a ray's content: \notyet\ beyond these two lists.
\end{remark}

\noindent So the answer to the seventh question: a list without a realization makes rays, laws and corners exactly as a graph does, with irrational native gauges at a lean as the one structural difference met. On $N$ it makes the numbers of two of its minors, in the pattern their own full leans and centres dictate, and new ones besides: two at the lean and six at the centre, unrecognized in the bases tried. On $V$ it makes everything but the law.

\subsection*{20.8\quad Which ray of a content is best?}

\begin{question}[Which ray of a content is best?]\label{a:q:arith}
Over all rays that carry one content (directions, leans, offsets, other lists with the same content), what is the frontier of the rate of positive selection against the growth of the denominators of the rational solutions, and does any ray reach a gap (Thm.~34.3, p.~300)?
\end{question}

\emph{This is the long shot: rate against denominator growth, over every ray that carries $\mathbb Q+\mathbb Q\ell$.}

The \emph{content} of a ray is the $\mathbb Q$-span of its values (\S18.5(ii)). A ray of order two with content $\mathbb Q+\mathbb Q\theta$ is a sequence of rational approximations to $\theta$, and rays of one content compete.

\begin{proposition}[The merit of a one-master ray]\label{b:p:merit}\derived\
Let the law have order $2$, let $c$ and $r$ be its solutions with $c(1)=1$ and zero corner, and with $r(1)=0$ and corner $\mathrm{Cor}$, so that $P=r+P(1)\,c$; let $D_K$ be the least common denominator of $c(K)$ and $r(K)$. Then $D_KP(K)=D_Kr(K)+P(1)\,D_Kc(K)$ is a linear form in $1$ and $P(1)$ with integer coefficients, positive, of size $D_K(f^\star)^K$ up to powers of $K$. So if the \emph{merit}
\[
-\log f^\star-\limsup_K\tfrac1K\log D_K
\]
is positive, $P(1)$ is irrational.
\end{proposition}

\begin{proof}
$P(K)>0$ (Proposition~\ref{b:p:moment}, $p=1$) and $P(K)=(f^\star)^K$ times a power of $K$ up to a constant (Laplace's method at $\alpha^\star$ on the chart, Thm.~35.5(i), p.~308). If $P(1)=a/b$, then $b\,D_KP(K)$ is a positive integer, which is impossible once it is below $1$ \recalled.
\end{proof}

\begin{proposition}[Families of rays on a banana]\label{b:p:families}\derived\ from Prop.~\ref{a:p:leanindex}.\
For $\binom\cE1$ and a balanced direction $(\mvec;s,t)$ put $c_e=m_e-t$, so $\sum_ec_e=s-t$. Then
\[
\Gamma(sk)\Gamma(tk)\,P(k)=\int_0^\infty y^{\frac{(s+t)k}2-1}\prod_{e\in\cE}\sG_{c_ek;(1,-1)}(\sqrt y,\sqrt y)\,dy ,
\]
which depends on the direction only through $S=s+t$ and the multiset $\{|c_e|\}$. Two directions with the same $S$ and the same $\{|c_e|\}$ are a \emph{family}: their values are in the ratio
\[
\frac{P_{(\mvec';s',t')}(k)}{P_{(\mvec;s,t)}(k)}=\Gst21{s'k,\ t'k}{Sk}^{-1}\Gst21{sk,\ tk}{Sk},
\]
two Beta corners at the same total, rational at integer $k$; so a family has one content. The duality $(\mvec;s,t)\leftrightarrow((s+t)\one-\mvec;t,s)$ is $c\mapsto-c$, and the multiples $(c'\one;c',c')$ of the centre are the centre ray read at $c'k$.
\end{proposition}

\begin{proof}
Prop.~\ref{a:p:leanindex} at $(k\mvec;ks,kt)$ gives $\Gamma(sk)\Gamma(tk)P(k)=\int_0^\infty y^{tk-1}\prod_e\sG_{c_ek;(1,-1)}(1,y)\,dy$, and shifting the line of Def.~2.1 by $\frac12\log y$ gives $\sG_{\nu;(1,-1)}(1,y)=y^{\nu/2}\sG_{\nu;(1,-1)}(\sqrt y,\sqrt y)$; the powers of $y$ add up to $y^{tk-1+(s-t)k/2}$. At equal coefficients the Gramma of one mark is even in its index (reflect the line).
\end{proof}

\noindent Receipts \computed: all $123$ balanced directions of $\binom{[3]}1$ with $1\le s,t\le5$, positive walls and entries in decreasing order (the full leans have rational values), each with $70$ closings and a search to order $4$ and degree $24$. A law of order two was found on $17$ directions and of order three on $7$; on every one of the $17$, $|c_e|$ is constant, on every one of the $7$ it takes two values, and every direction with $|c_e|$ constant has a law of order two. The $17$ fall into nine families, five of them the centre read at $c'k$, all of content $\mathbb Q+\mathbb Q\ell$, and the merit is constant on each family (growth of $D_K$ between $K=200$ and $K=240$):
\begin{center}
\footnotesize
\setlength{\tabcolsep}{4pt}
\begin{tabular}{@{}l l l l r r r l@{}}
\toprule
$\{|c_e|\}$, $S$ & directions $(\mvec;s,t)$ & $f^\star$ & other & $-\log f^\star$ & growth & merit & $P(1)$\\ \midrule
$0$, $2c'$ & $(c'\one;c',c')$, $c'=1$ & $1/9$ & $1$ & $2.197$ & $3.742$ & $-1.545$ & $\ell$\\
 & \qquad $c'=2$ & $1/81$ & $1$ & $4.394$ & $8.605$ & $-4.211$ & $(\ell-2)/3$\\
 & \qquad $c'=3$ & $1/729$ & $1$ & $6.592$ & $12.781$ & $-6.189$ & $(10\ell-23)/54$\\
 & \qquad $c'=4$ & $1/6561$ & $1$ & $8.789$ & $16.906$ & $-8.117$ & $(62\ell-145)/486$\\
 & \qquad $c'=5$ & $1/59049$ & $1$ & $10.986$ & $20.599$ & $-9.613$ & $(568\ell-1331)/5832$\\[2pt]
$1$, $5$ & $(3,3,3;1,4)$, $(2,2,2;4,1)$ & $1/243$ & $-27/64$ & $5.493$ & $10.059$ & $-4.566$ & $(11-4\ell)/54$\\
 & $(4,2,2;2,3)$, $(3,3,1;3,2)$ & $64/6561$ & $-1$ & $4.630$ & $9.196$ & $-4.566$ & $(11-4\ell)/18$\\[2pt]
$1$, $7$ & $(4,4,4;2,5)$, $(3,3,3;5,2)$ & $1/2187$ & $-1728/3125$ & $7.690$ & $13.384$ & $-5.694$ & $(19-8\ell)/108$\\
 & $(5,3,3;3,4)$, $(4,4,2;4,3)$ & $3125/3779136$ & $-1$ & $7.098$ & $12.792$ & $-5.694$ & $(19-8\ell)/54$\\[2pt]
$1$, $9$ & $(6,4,4;4,5)$, $(5,5,3;5,4)$ & $1/12500$ & $-1$ & $9.433$ & $16.532$ & $-7.098$ & $(751-320\ell)/2916$\\[2pt]
$2$, $8$ & $(7,3,3;3,5)$, $(5,5,1;5,3)$ & $823543/553584375$ & $1$ & $6.511$ & $15.251$ & $-8.741$ & $(80\ell-181)/648$\\
\bottomrule
\end{tabular}
\end{center}
The growth column is the quotient $(\log D_{240}-\log D_{200})/40$, not the limit superior of $\frac1K\log D_K$ that Proposition~\ref{b:p:merit} uses; at the centre that quotient moves between $3.74$ and $4.34$ from window to window, and $\frac1K\log D_K$ gives the merit $-1.9$ (Proposition~\ref{c:p:merit3}). The other rows have not been recomputed. Within a family the values are in the ratio of Proposition~\ref{b:p:families} at every $k$, exactly, for $k\le30$:
\[
\frac{P_{(4,2,2;2,3)}(k)}{P_{(3,3,3;1,4)}(k)}=\frac{\Gamma(4k)\Gamma(k)}{\Gamma(3k)\Gamma(2k)}=3,7,\tfrac{33}2,39,\dots,\qquad \frac{P_{(5,3,3;3,4)}(k)}{P_{(4,4,4;2,5)}(k)}=\frac{\Gamma(5k)\Gamma(2k)}{\Gamma(4k)\Gamma(3k)}=2,\tfrac{18}5,\tfrac{13}2,\dots;
\] the native values move by the ratio of the rates, $\frac{64}{27}$ and $\frac{3125}{1728}$, and the merit does not move. The seven rays of order three form three families in the same way, with symbols that agree up to that scaling.

\begin{remark}[The fixture's content]\label{b:r:fixmerit}
At $(\one;0,2)$ the fixture has content $\mathbb Q+\mathbb Q\Tx3{\IZ}0$ (Theorem~\ref{b:t:fixture}(iii)), $-\log f^\star=\log256=5.545$, and $D_K$ grows by $7.6$ to $8.4$ per step for $120\le K\le240$: merit about $-2.5$ \computed. Its denominators are $4K$-smooth: at $K=60$ the primes in $(K,2K]$ enter squared and those in $(2K,4K]$ once, and at $K=200,240$ the largest primes are $797$ and $953$. The classical recurrence of order two for the same number has merit $\log(17+12\sqrt2)-3=+0.526$ (alias, last) \recalled. So the ray of the fixture is a different law for that content, and it loses about three per step.
\end{remark}

\begin{proposition}[The centre's rate is spent on the scale of its closings]\label{c:p:merit3}(i), (ii) \derived; (iii) \computed.\
\begin{enumerate}[label=(\roman*)]
\item At the scale of the closings the law of \S19 is the recurrence of the closings:
\[
9k^2Q(k)-(10k^2+10k+3)\,Q(k+1)+(k+1)^2Q(k+2)=0 .
\]
\item The zero-corner solution $c$ of the law ($c(1)=1$, $c(2)=\frac13$) is $c(K)=x_{K-1}/9^{K-1}$, with $x_j=\sum_{a+b+c=j}\big(\frac{j!}{a!\,b!\,c!}\big)^2$ the closings.
\item Write $P(K)=r(K)+\ell\,c(K)$ and let $D_K$ be the least common denominator of $r(K)$ and $c(K)$. For $K\le702$,
\[
v_3(D_K)=2(K-1)-v_3(x_{K-1}),\qquad v_3(x_j)\le13,
\]
the part of $D_K$ prime to $3$ divides $\mathrm{lcm}(1,\dots,K)^2$, and $\frac1K\log D_K$ lies between $4.10$ and $4.16$ at $K=200,300,\dots,700$.
\end{enumerate}
So the rate $\log9=2.197$ is spent exactly on the $3$-part of $D_K$, up to the bounded $v_3(x_{K-1})$, and the merit is minus the growth of the part prime to $3$: about $-1.9$ ($-1.91$ to $-1.96$ at those $K$). The $-1.545$ of the table above is the quotient $(\log D_{240}-\log D_{200})/40$; over windows of $40$ steps that quotient gives values from $-1.55$ to $-2.15$. At the scale where the zero-corner solution is integral, $9^{K-1}P(K)$ decays only like a power of $K$: the rate of this ray is zero.
\end{proposition}

\begin{proof}
(i) Put $P(k)=9^{1-k}Q(k)$ in $k^2P(k)-(10k^2+10k+3)P(k+1)+9(k+1)^2P(k+2)=0$ and multiply by $9^k$. With $Q(k)=x_{k-1}$ it is the closings' recurrence $N^2x_N=(10N^2-10N+3)x_{N-1}-9(N-1)^2x_{N-2}$ at $N=k+1$. (ii) Both sides solve the law and agree at $K=1,2$ ($x_0=1$, $x_1=3$).
\end{proof}

\begin{remark}[The merit, sorted]\label{c:r:merit21}
For this ray the negative merit is absolute, a property of its law. Within the families of Prop.~\ref{b:p:families} it is constant (above), so no search there could change it. For the content $\mathbb Q+\mathbb Q\ell$ the question is whether $\ell$ is irrational, which is open (alias, last). Relative to the law: Rem.~\ref{b:r:fixmerit} has one content on two laws with merits $-2.5$ and $+0.526$. A ray of positive merit for $\ell$ needs a zero-corner solution whose denominators grow more slowly than the rate; at the centre they are the rate.
\end{remark}

\noindent So the answer to the eighth question, on the content $\mathbb Q+\mathbb Q\ell$: the rays of a banana come in families (Proposition~\ref{b:p:families}), the merit is a property of the family, and the best family met is the centre's, read at step one, with merit about $-1.9$ ($-1.545$ in the window $200\le K\le240$): the ray of the examination. No ray met has positive merit, and at the centre the reason is structural: the rate is spent on the scale of the closings (Proposition~\ref{c:p:merit3}), which are the torus value read backwards (Corollary~\ref{d:c:three}(ii)). Whether some ray of $\mathbb Q+\mathbb Q\ell$ has positive merit is whether $\ell$ is irrational by this route: open.

\renewcommand{\thesection}{21}
\setcounter{proposition}{0}
\renewcommand{\theproposition}{21.\arabic{proposition}}

\section{The gauge end}

The examination of \S19 read a ray with three inputs, the law, the corner and positivity, and \S20 found where that reading stops. This section reads a ray from the other end of its stride's law, the native gauge. There the law and the Hessian at $\alpha^\star$ fix the density, and the density, continued along the open support, is the ray: neither the corner nor positivity enters, and nothing that sits at the native values in between matters. \S21.1 states and proves this for every number of marks; \S21.2 shows that the germ at the gauge is local data at $\alpha^\star$; \S21.3 shows that a law must carry the dimension of its chart at every nondegenerate native value; \S21.4 reads every ray met so far with one program; \S21.5 reads the slopes and finds the arc integrals in them.

\subsection*{21.1\quad The gauge end, for every number of marks}

Throughout, a ray is strictly inside its walls, $n=|\cE|$, the chart has dimension $n-1$, $\alpha^\star$ is the native gauge on the orthant and $f^\star=f(\alpha^\star)$, and $H$ is the Hessian of $-\log f$ at $\log\alpha^\star$ in the coordinates of the chart. The density of the stride's law (Definition~\ref{b:d:law}) is $\varrho$, and $\sum_ic_i(k)P(k+i)=0$ is a law of the values with density operator $\mathcal L=\sum_it^i\,c_i(-1-i-\theta)=\sum_{j=0}^Da_j(t)\,\partial_t^j$, $\theta=t\,d/dt$ (Proposition~\ref{b:p:density}), of order $D=\max_i\deg c_i$, whose leading coefficient is $a_D=(-1)^Dt^D\chi(t)$, $\chi$ the symbol. Put $s=f^\star-t$, $b_j(s)=(-1)^ja_j(f^\star-s)$ (the operator in $s$), $b_{j,l}=[s^l]b_j$, and $\one_+=\one_{s>0}$, the chart's side of the gauge.

\begin{theorem}[The gauge end]\label{c:t:gauge}\label{d:t:gauge}\derived\
\begin{enumerate}[label=(\roman*)]
\item $H$ is positive definite, $\varrho$ is real-analytic on $(0,f^\star)$, and near $f^\star$
\[
\varrho(t)=(f^\star-t)^{\frac{n-3}2}\,g(f^\star-t),\qquad g\ \text{holomorphic at }0,\qquad g(0)=\frac{\big(2\Gamma(\frac12)^2/f^\star\big)^{\frac{n-1}2}}{\Gamma\big(\frac{n-1}2\big)\sqrt{\det H}} .
\]
\item The law holds if and only if $\varrho\,\one_{(0,f^\star)}$ solves $\mathcal L u=0$ in the sense of distributions on $(0,\infty)$.
\item Let $f^\star$ be a simple root of the symbol. Then $s=0$ is a regular singular point of $\mathcal L$, with exponents $0,1,\dots,D-2$ and one more, $\mu=D-1-b_{D-1,0}/b_{D,1}$, and the local solutions $u$ that \emph{extend by zero}, meaning that $u\one_+$ solves $\mathcal L$ across $s=0$, form a line. If $\mu\notin\{0,\dots,D-2\}$ it is spanned by the solution $s^\mu(1+O(s))$; if $\mu\in\{0,\dots,D-2\}$, by the holomorphic solution that vanishes to order exactly $\mu$, and never by one with a logarithm. The same holds on the other side, with $\one_{s<0}$ and $|s|^\mu$.
\item $\mu=\frac{n-3}2$.
\item The germ of $\varrho$ spans the line of (iii). So $\varrho$ is the continuation along $(0,f^\star)$ of the solution on that line normalized by (i), and
\[
P^{(j)}(k)=\int_0^{f^\star}t^k\log^jt\;\varrho(t)\,dt\qquad(j\ge0).
\]
\end{enumerate}
Neither a corner nor positivity enters, nor any condition on the native values in $(0,f^\star)$ or on the parity of $n$.
\end{theorem}

\begin{proof}
(i) On the chart, $-\log f$ is a sum of logarithms of Laurent polynomials with positive coefficients in $e^y$, minus a linear form. Its Hessian at any point is a sum of covariances of exponents under positive weights, so it is semidefinite. A null direction $v$ would make $\langle v,\cdot\rangle$ constant on the exponents of each letter, so $f$ would change by an exponential factor along $v$: unbounded if the factor is not $1$, against $f\le f^\star$; constant if it is, against $P(k)<\infty$. So $-\log f$ is strictly convex, $\log\alpha^\star$ is its only critical point (Thm.~35.10(i), p.~309), $-\log f\to\infty$ at infinity of the chart, and every set $\{f\ge t\}$, $t>0$, is compact. For $0<t<f^\star$ the level set $\{f=t\}$ is compact and smooth, and $\varrho(t)=\int_{f=t}d\sigma/|\nabla f|$ is real-analytic there. At $\log\alpha^\star$ the Morse lemma gives analytic coordinates $w$ with $\log(f^\star/f)=\frac12|w|^2$ and $dy=J(w)\,dw$, $J(0)=(\det H)^{-1/2}$. With $\sigma=\log(f^\star/t)$, $\mu([t,f^\star])=\int_{|w|^2\le2\sigma}J\,dw$; the odd spherical moments of $J$ vanish, so this is $\sigma^{\frac{n-1}2}$ times a power series in $\sigma$ with leading coefficient $J(0)\,2^{\frac{n-1}2}\,\Gamma(\frac12)^{n-1}/\Gamma(\frac{n+1}2)$. As $\sigma$ is holomorphic in $f^\star-t$ with $\sigma=(f^\star-t)/f^\star+\dots$, differentiating in $t$ gives the form and the constant.

(ii) For a distribution $T$ on $(0,\infty)$ carried by $(0,f^\star]$ and of moderate growth at $0$, $\int t^k\,\theta T=-(k+1)\int t^kT$ for $\operatorname{Re}k$ large. So $\int t^k\,\mathcal L(\varrho\one_{(0,f^\star)})=\sum_ic_i(k)P(k+i)$ (Prop.~\ref{b:p:law}(i)), and a distribution carried by $(0,f^\star]$ is zero exactly when these moments vanish on a half-plane.

(iii) $a_D$ has a simple zero at $f^\star$, so after dividing by a unit $s^{D-1}\mathcal L=\sum_{m\ge0}s^mR_m(\theta_s)$, $\theta_s=s\,\partial_s$, with $R_m(x)=\sum_jb_{j,m-D+1+j}\,x(x-1)\cdots(x-j+1)$, and
\[
R_0(x)=x(x-1)\cdots(x-D+2)\,\big(b_{D,1}(x-D+1)+b_{D-1,0}\big).
\]
Now let $u=\sum_ru_rs^r$ be holomorphic. The coefficient of $s^N$ in $s^{D-1}\mathcal Lu$ is $\sum_mR_m(N-m)u_{N-m}$, and for $N\le D-2$ every term vanishes: the term of $R_m$ coming from $b_j$ carries $j$ falling factors and needs $N-m\ge j$, while $m\ge D-1-j$. So $u_0,\dots,u_{D-2}$ are free, and from $N=D-1$ on $R_0(N)\neq0$ fixes the rest. The distribution $\mathcal L(u\one_+)$ is carried by $s=0$ and has order at most $D-2$, $\mathcal L(u\one_+)=\sum_{r=0}^{D-2}\gamma_r\,\delta^{(D-2-r)}$. The two leading terms $b_{D,1}s\,\partial^D+b_{D-1,0}\,\partial^{D-1}$ send $s^r\one_+$ to $r!\,b_{D,1}(r-\mu)\,\delta^{(D-2-r)}$, since $\partial^{D-1}(s^r\one_+)=r!\,\delta^{(D-2-r)}$ and $s\,\delta^{(m)}=-m\,\delta^{(m-1)}$; every other term raises the power of $s$ by more than it lowers the order of differentiation, and lands on point masses of lower order. So
\[
\gamma_r=r!\,b_{D,1}\,(r-\mu)\,u_r+\sum_{i<r}\kappa_{r,i}u_i ,
\]
a triangular system with diagonal $r!\,b_{D,1}(r-\mu)$. If $\mu\notin\{0,\dots,D-2\}$ it forces $u=0$. The solution $s^\mu\Phi$ does extend by zero: $\lambda\mapsto s_+^\lambda$ is a holomorphic family of distributions, the Frobenius family with a parameter gives $\mathcal L(s_+^\lambda\Phi_\lambda)=R_0(\lambda)\,v(s)\,s_+^{\lambda-D+1}$ with $v$ a unit, and at $\lambda=\mu$ the factor $R_0$ vanishes while $s_+^{\lambda-D+1}$ is finite. If $\mu\in\{0,\dots,D-2\}$, the system forces $u_0=\dots=u_{\mu-1}=0$, leaves $u_\mu$ free (zero diagonal, vanishing lower terms) and fixes $u_{\mu+1},\dots,u_{D-2}$: a line. The remaining solution has the form $s^\mu\log s\,(1+O(s))$ plus a holomorphic one. For the two leading terms, $L_0=b_{D,1}s\,\partial^D+b_{D-1,0}\,\partial^{D-1}$, one has $L_0(s_+^\lambda)=R_0(\lambda)\,s_+^{\lambda-D+1}$ for every $\lambda$, both sides holomorphic families of distributions. At $\lambda=\mu$, $R_0$ has a double zero and $s_+^{\lambda-D+1}$ a simple pole whose residue is a nonzero multiple of $\delta^{(D-2-\mu)}$; differentiating in $\lambda$, $L_0(s_+^\mu\log s_+)=\frac12R_0''(\mu)\operatorname{Res}\cdot\delta^{(D-2-\mu)}\neq0$. The rest of $\mathcal L$, and the rest of the logarithmic solution, reach only the levels $r>\mu$. At the levels $r<\mu$ the holomorphic part is forced to vanish, and then at the level $r=\mu$ nothing cancels the logarithm's point mass. The other side is the same computation with $s\mapsto-s$.

The count of the holomorphic solutions used $R_0(N)\neq0$ for $N\ge D-1$, that is, $\mu$ not an integer $\ge D-1$. For a law of the values this always holds. The exponent $-1$ has multiplicity $n-1$ at $t=0$ (Prop.~\ref{b:p:density}), so $D\ge n-1$, and the germ of $\varrho$, a local solution by (ii), has the leading exponent $\frac{n-3}2\le D-2$ by (i). For even $n$ that exponent is a half-integer, so it is an indicial root, and $\mu$ is the only candidate. For odd $n$ the germ is holomorphic and extends by zero, and if $\mu$ were not in $\{0,\dots,D-2\}$ the triangular system would force its Taylor coefficients up to order $D-2$ to vanish, among them $g(0)\neq0$.

(iv) By (ii), $\varrho\one_{(0,f^\star)}$ solves $\mathcal L$ across $f^\star$, so its germ lies on the line of (iii); by (i) its leading exponent is $\frac{n-3}2$.

(v) By (iv) the germ of $\varrho$ is on the line, and (i) fixes its scale. By (i) $\varrho$ is analytic on $(0,f^\star)$, so it is the continuation of its germ, and the values and slopes are Prop.~\ref{b:p:law}(i) differentiated under the integral, $0<f\le f^\star$.
\end{proof}


\begin{remark}[The line from the law alone]\label{c:r:lawline}
The line of (iii) can also be reached without looking at exponents, an argument that bounds its dimension when $f^\star$ is a multiple root (Corollary~\ref{c:c:multiple}). Let $u$ be a germ at $f^\star$ that extends by zero, and $\psi$ a cutoff, $1$ on $[f^\star-\delta,f^\star]$ and $0$ below $f^\star-2\delta$. By the identity in the proof of (ii), $M_u(k)=\int t^ku\,\one_{t<f^\star}\psi\,dt$ satisfies $\sum_ic_i(k)M_u(k+i)=\int t^k\,\mathcal L(u\one_{t<f^\star}\psi)\,dt$. The distribution $\mathcal L(u\one_{t<f^\star}\psi)$ is carried by $[f^\star-2\delta,f^\star-\delta]$, where $\psi$ is not constant, so this defect is $O\big((f^\star-\delta)^k\big)$ times a power of $k$. The germ $u$ is a finite sum of $(f^\star-t)^\lambda\log^j(f^\star-t)$ times holomorphic functions. So Laplace's method at the endpoint gives $M_u(k)$ an expansion $(f^\star)^k\sum c\,k^{-\lambda-1-m}\log^{j'}k$, whose leading term is not zero unless $u$ is. That expansion solves the law term by term, because the defect is smaller than every term.

Now count the formal solutions $(f^\star)^kB(k)$. Put $B(k+i)=\sum_r\frac{i^r}{r!}B^{(r)}(k)$ into the law, and write $\chi(\lambda)=\sum_i[k^D]c_i\,\lambda^i$ for the symbol and $\chi_1(\lambda)=\sum_i[k^{D-1}]c_i\,\lambda^i$. The leading order is $\chi(f^\star)=0$. The next is
\[
f^\star\chi'(f^\star)\,k\,B'(k)+\chi_1(f^\star)\,B(k)=0 .
\]
As $\chi'(f^\star)\neq0$, this forces $B=k^\rho(a_0+a_1k^{-1}+\dots)$ with $\rho=-\chi_1(f^\star)/(f^\star\chi'(f^\star))$ and no logarithm. Each $a_j$ is then fixed by the order at which it enters, with the factor $-jf^\star\chi'(f^\star)\neq0$. So the expansions of two germs extending by zero are proportional; their difference has the zero expansion, so it is zero.
\end{remark}


\begin{remark}[Where the parity was]\label{c:r:parity}\label{d:r:parity}
Only in how the line looks. For odd $n$ the density is holomorphic at $f^\star$ from the left, $(f^\star-t)^{\frac{n-3}2}g$, and cut there: the line of (iii) lies inside the $D-1$ holomorphic solutions, where the exponent cannot single it out, and the logarithm of Proposition~\ref{b:p:density} belongs to the other solution, the one the density avoids. For even $n$ the line is the half-integer class, since finite parts of half-integer powers carry no point masses. One mechanism covers both: the zeros of the indicial polynomial against the poles of $s_+^{\lambda-D+1}$. To compute the line, write $\mathcal L(s^\nu)=\sum_dT_d(\nu)s^{\nu+d}$: then $T_{1-D}$ is the indicial polynomial and $T_d(m)=0$ whenever $m\ge0$ and $m+d\le-1$. For a holomorphic $h=\sum h_ms^m$, $\mathcal L(h\,\one_+)$ is a sum of point masses at $s=0$: at $\nu=m$ the zero of $T_d$ meets the simple pole of $s_+^{\nu+d}$ at $\nu+d=-j$, which leaves a point mass of order $j-1$. They vanish exactly when
\[
\sum_{m=0}^{D-1-j}h_m\,T'_{-m-j}(m)=0,\qquad j=1,\dots,D-1,
\]
linear conditions on the first $D-1$ Taylor coefficients, which with the recurrence of the Taylor coefficients give the germ (Example~\ref{c:e:fivegauge}). A first proof of (iii) went through finite parts of half-integer powers alone, and so was stated for even $n$.
\end{remark}


\noindent So the corners are outputs: the expansion of $\varrho$ at $t=0$, with the exponent $-1$ of multiplicity $n-1$ (Prop.~\ref{b:p:density}), is the connection of the line at $f^\star$ with the solutions at $0$, and the corner relations are the law at $k=0,-1,\dots$ evaluated on the values the gauge end produces.

\begin{remark}[What positivity sees and what it does not]\label{c:r:sees}
The extra solution of Prop.~\ref{b:p:blind} is $u(k)=\int t^k\nu(dt)$, with $\nu$ the jump across the slit, carried by $[0,\lambda_c]$ and of the form $(\lambda_c-t)_+^{\frac{n-3}2}$ times a holomorphic function at $\lambda_c$. Three inputs cannot see it. The law cannot: $u$ solves it, so $\nu$ is itself a solution in the sense of distributions (Theorem~\ref{c:t:gauge}(ii) read for $\nu$). The corner cannot: $u$ has none. Positivity cannot: $|\nu|\le C\varrho$ (Prop.~\ref{b:p:blind}). But $\varrho+\varepsilon\nu$ is not analytic at $\lambda_c$. So the inputs of \S19, read as conditions on candidates, cut out a cone with a flat direction. The gauge end does not choose among candidates; it constructs $\varrho$, and analyticity on the open support is as much a property of $\mu=f_*\Lch$ as $\mu\ge0$ is. For every ray strictly inside its walls, the ``no'' of \S20.1 belongs to the inputs, not to the ray (with Corollary~\ref{c:c:multiple} when $f^\star$ is a multiple root).
\end{remark}


\subsection*{21.2\quad The germ at the gauge end is local data at $\alpha^\star$}

Theorem~\ref{c:t:gauge}(i) uses one number of the native gauge, $\det H$. The whole germ is local data there.

\begin{proposition}[The germ at the gauge end is local data at $\alpha^\star$]\label{c:p:jet}\derived; the check \computed.\
Expand $h=\log f$ at $\log\alpha^\star$ in the coordinates $y$ of the chart:
\[
h=\log f^\star-\tfrac12H_{ab}\,y^ay^b+\tfrac16h_{abc}\,y^ay^by^c+\tfrac1{24}h_{abcd}\,y^ay^by^cy^d+\dots
\]
\begin{enumerate}[label=(\roman*)]
\item At $\alpha^\star$ each letter $L=\sum_tw_t\alpha^t$ puts the law $\sum_tw_t(\alpha^\star)^t\delta_t/L(\alpha^\star)$ on its exponents. The jet of $h$ is the strengths times the cumulants of these laws: $H$ from the covariances (as in Example~\ref{c:e:slitgauge}(iii)), $-h_{abc}$ from the third cumulants, $-h_{abcd}$ from the fourth, and so on.
\item As $k\to\infty$,
\[
P(k)=(f^\star)^k\Big(\frac{2\Gamma(\frac12)^2}{k}\Big)^{\frac{n-1}2}\frac{1}{\sqrt{\det H}}\Big(1+\frac{c_1}k+\frac{c_2}{k^2}+\dots\Big),
\]
where, with $G=H^{-1}$,
\[
c_1=\tfrac18h_{abcd}G^{ab}G^{cd}+\tfrac18h_{abc}h_{def}G^{ab}G^{cd}G^{ef}+\tfrac1{12}h_{abc}h_{def}G^{ad}G^{be}G^{cf},
\]
and each $c_j$ is a finite expression in the jet.
\item The Taylor coefficients of $g$ in Theorem~\ref{c:t:gauge}(i) are read off (ii) one at a time. The term $(f^\star-t)^{\frac{n-3}2+m}$ of $\varrho$ contributes $(f^\star)^{k+\frac{n-1}2+m}B(k+1,\frac{n-1}2+m)$, whose expansion starts at $k^{-\frac{n-1}2-m}$ with the coefficient $\Gamma(\frac{n-1}2+m)\neq0$.
\end{enumerate}
\end{proposition}

\begin{proof}
(i) $-h=\sum_L\mathrm{str}_L\log L(e^y)-\langle\mvec,y\rangle$, and $\log L(\alpha^\star e^y)-\log L(\alpha^\star)$ is the cumulant generating function of the law in (i). (ii) Laplace's method: $e^{k(h-\log f^\star)}=e^{-\frac k2H_{ab}y^ay^b}(1+\dots)$; expand the cubic and quartic parts and take Gaussian moments with covariance $G/k$. The three terms of $c_1$ are the quartic term and the two ways of pairing two cubic terms. (iii) $B(k+1,b)=\Gamma(b)\,k^{-b}\big(1+O(k^{-1})\big)$.
\end{proof}

\begin{corollary}[Theorem~\ref{c:t:gauge}(iii) and (v) without a hypothesis on the symbol]\label{c:c:multiple}
The germs at $f^\star$ that extend by zero form a space of dimension at most the multiplicity of $f^\star$ as a root of the symbol. The germ of $\varrho$ is the one in it whose expansion is that of Proposition~\ref{c:p:jet}(ii). When the root is simple the space is a line and $g(0)$ alone fixes the germ; otherwise finitely many $c_j$ do.
\end{corollary}

\begin{proof}
The argument of Remark~\ref{c:r:lawline} maps this space injectively into the formal solutions of the law at the rate $f^\star$, of which there are as many independent ones as the multiplicity of the root \recalled. An injective map from a finite-dimensional space is inverted on its image by finitely many coordinates.
\end{proof}

\noindent Check \computed. Compare $c_1$ from the cumulants with $a_1/a_0$ for the formal solution $(f^\star)^kk^\rho(a_0+a_1k^{-1}+\dots)$ of the law. For the five-mark centre both are $\frac12$, to $30$ digits. For Example~\ref{b:e:Hslit} both are $0.332601611985598167464224028657\ldots$ So the large-$k$ expansion of the law at $f^\star$, the germ of $\varrho$, and the jet of the stride at $\alpha^\star$ are one object. The density equation is needed only to carry the germ along $(0,f^\star)$. No law met has a multiple root at $f^\star$ (the table of \S21.3), so Corollary~\ref{c:c:multiple} has not yet been needed on an example.


\subsection*{21.3\quad The law knows the dimension of its chart}

Theorem~\ref{c:t:gauge}(iv) is a test that a law guessed from closings must pass: at the native value $f^\star$ the indicial root beyond $0,\dots,D-2$ is $\frac{n-3}2$, a number that the closings, integers indexed by one strength, do not contain. \computed:
\begin{center}
\small
\begin{tabular}{@{}l c c l@{}}
\toprule
ray & $n$ & $D$ & simple roots of the symbol, and $\mu$ at each\\ \midrule
three-mark centre & $3$ & $2$ & $\frac19,\ 1$: $\mu=0$\\
the $24$ three-mark rays of \S20.8 & $3$ & $\le18$ & all $55$ simple roots: $\mu=0$\\
four-mark centre & $4$ & $3$ & $\frac1{16},\ \frac14$: $\frac12$\\
$H$, centre & $4$ & $7$ & $\frac1{25},\ \frac19$: $\frac12$; $1$ double\\
$H$ at $(1,1,3,3;3,1)$ (Ex.~\ref{b:e:Hslit}) & $4$ & $25$ & $f^\star$ and its conjugate on the slit: $\frac12$; $1$ double\\
five-mark centre & $5$ & $4$ & $\frac1{25},\ \frac19,\ 1$: $1$\\
fixture, $(\one;0,2)$ & $6$ & $5$ & $\frac1{256},\ \frac14$: $\frac32$\\
$N$, $(\one;0,2)$ & $6$ & $24$ & $f^\star$, its negative conjugate, $\frac1{24}$, the cubic's three roots: $\frac32$\\
$N$, centre, law of order $12$ & $6$ & $31$ & $\frac1{361},\frac1{25},\frac19,-\frac1{80},\frac{17\pm14\sqrt7}{361}$: $\frac32$; $1$ five-fold;\\
 & & & $\lambda_s=0.4499215354534\ldots$: $16.8575639440931\ldots$\\
\bottomrule
\end{tabular}
\end{center}
Every simple root that is a native value carries $\frac{n-3}2$: $f^\star$, the conjugate on the slit, negative and complex native values alike. The multiple roots are the families of native gauges (Remark~\ref{a:auto:15.16}). The one root that is not a native value, $\lambda_s$, which the left multiple of order $12$ adds on $N$'s centre (Prop.~\ref{b:p:Ncentre}(v)), carries an exponent that is not a half-integer, and the law itself flags it. At $f^\star$ the value is \derived\ (Theorem~\ref{c:t:gauge}(iv)); at the other simple roots it is the exponent of the vanishing cycle of Thm.~35.11, p.~310, read at that gauge, and \computed\ on every law here. So a law read off closings carries, at every nondegenerate native value, the dimension $n-1$ of the chart; and at $t=0$ it carries $n$ again, as the multiplicity $n-1$ of the exponent $-1$ (Prop.~\ref{b:p:density}).


\subsection*{21.4\quad One program, every $n$}

Receipts \computed. One program reads every ray: law $\to$ density operator $\to$ the line at $f^\star$ (in exact rationals where $f^\star$ is rational) $\to$ its series $\to$ Beta sums at $K,\dots,K+h-1$ $\to$ the law backwards to $k=1$. Its inputs are the law, $\alpha^\star$ and $\det H$; the corner and positivity are not used.
\begin{center}
\small
\setlength{\tabcolsep}{4pt}
\begin{tabular}{@{}l c c r l l l@{}}
\toprule
ray & $n$ & $\mu$ & $D$ & $P(1)$, gauge end & against & difference\\ \midrule
three-mark centre & $3$ & $0$ & $2$ & $2.34390723868946\ldots$ & \S19 & $6\cdot10^{-495}$\\
four-mark centre & $4$ & $\frac12$ & $3$ & $8.41439832211716\ldots$ & $7\Zt$ & $2\cdot10^{-492}$\\
$H$, centre & $4$ & $\frac12$ & $7$ & $4.09623699833788\ldots$ & \S18.4 & all $43$ digits\\
$H$, $(1,1,3,3;3,1)$ & $4$ & $\frac12$ & $25$ & $0.48830319489527\ldots$ & Ex.~\ref{b:e:Hslit}(iii) & all $59$ digits\\
five-mark centre & $5$ & $1$ & $4$ & $39.9455880119962\ldots$ & \S15 & all $50$ digits\\
fixture, $(\one;0,2)$ & $6$ & $\frac32$ & $5$ & $7.21234141895757\ldots$ & $6\Zt$ & $2\cdot10^{-489}$\\
$N$, $(\one;0,2)$ & $6$ & $\frac32$ & $24$ & $4.75680256375823\ldots$ & Thm.~\ref{b:t:Nlean} & $2\cdot10^{-189}$\\
$N$, centre & $6$ & $\frac32$ & $31$ & $4.57891247888240\ldots$ & Prop.~\ref{b:p:Ncentre} & $5\cdot10^{-222}$\\
\bottomrule
\end{tabular}
\end{center}
The independent values are the closed form of \S19; seven and six times $\Zt$ \recalled; the share-side integrals of \S18.4 and Ex.~\ref{b:e:Hslit}(iii); the value printed in \S15; and the positivity readings of \S20.

Readings of the same ray from two depths agree to $4\cdot10^{-387}$ (five marks, $K=260$ and $360$), $5\cdot10^{-410}$ ($H$), $7\cdot10^{-504}$ and $3\cdot10^{-542}$ ($N$). The $H$ row at $(1,1,3,3;3,1)$ passes the native value on the slit, $0.0021985\ldots$, with no special step: the density is analytic there, and the backward run only has to carry the factor $(f^\star/0.0021985)^K$ in its digits. Example~\ref{c:e:slitgauge} below lowered the order of the density equation with a second recurrence of order $19$; the law's own equation, of order $25$, does the same.

\begin{example}[Example~\ref{b:e:Hslit} read from its gauge end]\label{c:e:slitgauge}\computed; the Hessian \derived.\
Take $H$ at $(1,1,3,3;3,1)$, $n=4$, as in Ex.~\ref{b:e:Hslit}.
\begin{enumerate}[label=(\roman*)]
\item The closings have a recurrence of order $19$ and degree $10$ besides that of order $4$ and degree $25$: both exact, by modular nullspaces over $12$ and $8$ primes of $62$ bits, Chinese remaindering and rational reconstruction (coefficients up to $65$ and $36$ digits), and exact at all $401$ and $416$ shifts that the $420$ computed closings allow. Reflected, the first gives a density equation of order $10$ and degree $19$ in $t$ that annihilates $\sum_Nx_Nt^{N-1}$ exactly. Its symbol is
\[
(1366875\,t^2-15286\,t+27)(t-1)^2A(t),\qquad\deg A=15,
\]
and the symbol of the law of order $4$ is the first two factors alone, so $A$ carries no native value.
\item Exponents. At $t=0$: $-1$, $-\frac23$, $-\frac13$, each three times, and $0$ once. At $f^\star$: $0,1,\dots,8$ and $\frac12$. $f^\star$ is a simple root of the symbol, so Theorem~\ref{c:t:gauge}(iii) applies; for this even $n$ the line is the half-integer class, and $\mu=\frac12$ as (iv) requires.
\item \derived\ The Hessian on the chart $\alpha_3=1$, coordinates $(a,b,2)$, is $3\,\mathrm{Cov}_K+\mathrm{Cov}_\Psi$ at $\alpha^\star=(x_0,x_0,1,1)$, $x_0=\frac{\sqrt{19}-4}3$:
\[
\det H=\frac{376096\sqrt{19}-1614848}{151875}=0.161425231735790355\ldots,\qquad g(0)=\frac{2\Gamma(\frac12)^2\,(2/f^\star)^{3/2}}{\sqrt{\det H}} .
\]
\item The series $g(s)=\sum_me_ms^m$ from the recurrence of $\mathcal L$ at $f^\star$:
\begin{center}
\small
\begin{tabular}{@{}r l l@{}}
\toprule
$m$ & $e_m(f^\star)^m$ & $e_m(f^\star-f^{\star\prime})^m$\\ \midrule
$50$ & $2.4538\cdot10^5$ & $0.19761$\\
$400$ & $3.9255\cdot10^5$ & $6.9444\cdot10^{-44}$\\
$1600$ & $5.0797\cdot10^5$ & $4.9749\cdot10^{-190}$\\
\bottomrule
\end{tabular}
\end{center}
with $f^{\star\prime}=0.0021985\ldots$ the conjugate on the slit. The first column grows like $\log^2m$: the radius of $g$ is $f^\star$, the distance to $t=0$, not $f^\star-f^{\star\prime}$. The density continues analytically through the native value on the slit, as Theorem~\ref{c:t:gauge}(i) says it must.
\item Values. Beta sums $P(K)=\sum_me_m(f^\star)^{K+m+3/2}B(K+1,m+\frac32)$ at $K=80,\dots,83$, then the law of order $4$ backwards:
\begin{gather*}
P(1)=0.48830319489526888388541464396010900874354031922042897524562141926647\ldots,\\
P(2)=0.0013211917949877322271764171472151006507596961088513549847545827159\ldots
\end{gather*}
The readings from $K=80$ and $K=120$ differ by $4\cdot10^{-89}$. The share side of Ex.~\ref{b:e:Hslit}(iii), recomputed here independently (the Mellin--Barnes line at $\operatorname{Re}v=\frac12$), agrees to all $62$ digits it was computed to.
\item Both corners as outputs. The law at $k=0$, not used above, gives
\[
\mathrm{Cor}=-\frac{\sum_{i=1}^4c_i(0)P(i)}{[k^3]\,c_0}=5.925925925\ldots=\frac{160}{27}
\]
to $40$ digits. The second corner of Ex.~\ref{b:e:Hslit}(iv),
\[
12610378582200\,P(1)-133936231480200\,P(2)+160181043750000\,P(3)-5981712159277,
\]
is $-9.2\cdot10^{-101}$ from $K=120$.
\end{enumerate}
What went in: the density equation, $\alpha^\star$, $\det H$. What did not: the corner, the second corner, positivity, the share side. The backward run of the law amplifies the slit direction by $(f^\star/f^{\star\prime})^K=10^{0.611K}$, which is why Theorem~\ref{a:t:select}(iv) diverges on this ray. The gauge end hands the law values with no slit component, so the run only has to carry enough digits.
\end{example}

\noindent The same reading on the three-mark centre \computed: $n=3$, $\varrho$ holomorphic at $\frac19$ with $\varrho(\frac19)=2\Gamma(\frac12)^2/(f^\star\sqrt{\det H})=27\sqrt3\,\Gamma(\frac12)^2$ ($\det H=\frac4{27}$, \S19), Beta sums at $K=300$ with $2200$ terms, the law of \S19 backwards. It gives $\ell$ to $390$ digits, and the corner as an output, $3P(1)-9P(2)=6.000\ldots$ to $40$ digits. Here $D=2$, $\mu=0$, and the line is the holomorphic solution at $\frac19$.

\begin{example}[The five-mark centre read from its gauge end: $n$ odd, $\mu$ an integer]\label{c:e:fivegauge}\label{d:e:five}The line \derived; the reading \computed; the Hessian \derived.\
\begin{enumerate}[label=(\roman*)]
\item The closings $x_j=\sum_{a_1+\dots+a_5=j}\big(\frac{j!}{a_1!\cdots a_5!}\big)^2=1,5,45,545,7885,\dots$ have a least recurrence of order $3$ and degree $4$, exact at all $157$ shifts of $160$ closings:
\begin{multline*}
(j+3)^4x_{j+3}-(35j^4+350j^3+1323j^2+2240j+1433)\,x_{j+2}\\
+(j+2)^2(259j^2+1036j+1062)\,x_{j+1}-225(j+1)^2(j+2)^2x_j=0 .
\end{multline*}
Reflected, the density equation has order $4$, symbol $(t-1)(9t-1)(25t-1)$ up to sign, and the exponent $-1$ four times at $t=0$.
\item At $f^\star=\frac1{25}$ the indicial polynomial is $384\,\nu(\nu-1)^2(\nu-2)$: $D=4$, and $\mu=\frac{n-3}2=1$ falls on $1$. On the holomorphic solutions $u_0,u_1,u_2$ are free, a space of dimension $3$. The line of Theorem~\ref{c:t:gauge}(iii) has $u_0=0$ and, by the one condition at the level $r=2$, $u_2=\frac{175}4\,u_1$, exactly: the germ vanishes to first order, as Theorem~\ref{c:t:gauge}(i) says it must, and its Taylor coefficients are rational, $u_3=\frac{48125}{32}u_1$, $u_4=\frac{35703125}{768}u_1,\dots$
\item The law has the formal solution $25^{-k}k^{\rho}(1+\dots)$ at $\frac1{25}$ with $\rho=-2=-\frac{n-1}2$.
\item \derived\ On the chart $\alpha_5=1$, $H=2(\frac15I-\frac1{25}J)$ with $J$ the matrix of ones, so $\det H=2^4/5^5$ (Proposition~\ref{a:p:pairing}), and $u_1=g(0)=\big(50\,\Gamma(\frac12)^2\big)^2/\sqrt{\det H}=15625\sqrt5\,\Gamma(\frac12)^4$.
\item Beta sums at $K=150,151,152$, then the law backwards:
\[
P(1)=39.945588011996245698944028498423995891849782924769859572366\ldots,
\]
the same from $K=250$ and from $K=360$. The share side,
\[
\Gamma(k)^2P(k)=\int_0^\infty y^{k-1}\sG_{0;(1,-1)}(1,y)^5\,dy
\]
(Prop.~\ref{a:p:leanindex}, as in the proof of Proposition~\ref{b:p:families}), at $k=1$ by quadrature, agrees to all $59$ digits computed, and \S15.4 has the value to all $50$ of its digits. The corner comes out as an output: $\mathrm{Cor}=-\sum_ic_i(0)P(i)/[k^4]c_0=120.000\ldots$
\end{enumerate}
Here the exponent alone could not find the germ; extension by zero did.
\end{example}


\subsection*{21.5\quad The slopes at the arc}

\begin{proposition}[Slopes from the gauge end]\label{c:p:gaugeslopes}\derived\
Under Theorem~\ref{c:t:gauge}, $P^{(j)}(K)=\int_0^{f^\star}t^K\log^jt\,\varrho(t)\,dt$ for every $j\ge0$, with $\varrho$ continued along $(0,f^\star)$. When $\sum_me_ms^m$ converges on $[0,f^\star)$ this is, for $K$ large,
\[
P^{(j)}(K)=\sum_{m\ge0}e_m\,\frac{\partial^j}{\partial a^j}\Big[(f^\star)^{a+m+\frac{n-1}2}\,B\big(a+1,m+\tfrac{n-1}2\big)\Big]_{a=K}.
\]
The law differentiated $j$ times carries $P^{(j)}$ from $K$ back to $1$. No corner datum enters at any order.
\end{proposition}

\begin{proof}
$\int_0^{f^\star}t^a(f^\star-t)^b\,dt=(f^\star)^{a+b+1}B(a+1,b+1)$, and $\partial_a^jt^a=t^a\log^jt$. The differentiated law is the law with an inhomogeneity built from lower slopes (Prop.~\ref{b:p:slopes}(ii)).
\end{proof}

\noindent For the next statement put $P(k)=\Pth$ at $k\one$, $\ell^{(j)}=P^{(j)}(1)$, and read the ray at the scale of its closings, $Q(k)=(f^\star)^{1-k}P(k)=9^{k-1}P(k)$. Write
\[
\Lambda_{-3}(s)=\sum_{a=1}^{2}\Big(\frac{-3}a\Big)\Tx{s}{\Ilat{a+3\mathbb Z}}0,\qquad
S_j=-2^{1-j}\int_0^{3/4}\log^{j-1}(4u)\,\frac{du}{\sqrt{u(1-u)}} .
\]
$\Lambda_{-3}$ is the Zamma sum with the Kronecker sign of $-3$ (as in \S20.7). With $B(p,q;x)=\int_0^xu^{p-1}(1-u)^{q-1}du$, the Beta corner cut at $x$,
\[
S_j=-2^{1-j}\,\frac{\partial^{j-1}}{\partial p^{j-1}}\Big[4^{p-\frac12}B\big(p,\tfrac12;\tfrac34\big)\Big]_{p=\frac12},
\]
a slope of the cut corner. With $u=\sin^2\frac\vartheta2$ it is $-\int_0^{2\pi/3}\log^{j-1}(2\sin\frac\vartheta2)\,d\vartheta$, an integral over the arc from $1$ to $e^{2\pi i/3}$; we call the $S_j$ the \emph{arc integrals}. So $S_1=-\frac23\Gamma(\frac12)^2$, and $S_2=\frac{\sqrt3}2\Lambda_{-3}(2)$ \derived\ (expand $-\log(2\sin\frac\vartheta2)=\sum_{m\ge1}\cos(m\vartheta)/m$).

\begin{proposition}[The first two slopes of the three-mark centre]\label{c:p:arcslopes}\computed\
\begin{align*}
Q(1)&=\ell=3\Lambda_{-3}(2)=2\sqrt3\,S_2,\\
Q'(1)&=\ell'+2\ell\log3=\tfrac{27}2\Lambda_{-3}(3)+4\sqrt3\,S_3,\\
Q''(1)&=\ell''+4\ell'\log3+4\ell\log^23=-\tfrac13\Gamma(\tfrac12)^4\,\ell-27\Lambda_{-3}(4)+\tfrac{32}{\sqrt3}\Big(S_4-\Gamma(\tfrac12)^2\,\Tx3{\IZ}0\Big),
\end{align*}
with $\Lambda_{-3}(3)=\frac{4\,\Gamma(1/2)^6}{81\sqrt3}$ \recalled, and
\[
S_3=-2.14476721256949439148264211222\ldots,\qquad S_4=5.95066824413978493075781012901\ldots
\]
\end{proposition}

\noindent Receipts. $\ell,\ell',\ell''$ to $390$ digits from the gauge end (Proposition~\ref{c:p:gaugeslopes} at $K=300$; the slope relation of Prop.~\ref{b:p:slopes}(i), $3P'(1)-9P'(2)+4\ell+12$, is $1.6\cdot10^{-395}$). They agree with \S20.3's values to every digit printed there ($57$ and $42$ significant digits); there positivity read $\ell'$ with no corner and $\ell''$ with the corner $a_2$, here neither used one. $S_3$ by the trilogarithm formula and by quadrature, $S_4$ by quadrature, both to $400$ digits. Integer relations by reduction of lattices at $370$ digits: $\ell'$ against $29$ numbers (the $24$ of \S20.3 together with $\ell\Gamma(\frac12)^2$, $\ell\Gamma(\frac12)^2/\sqrt3$, $\Gamma(\frac12)^6$, $S_3$, $S_3/\sqrt3$) returns the relation above with coefficients $-3,-6,2,36$ and residual $3\cdot10^{-391}$. $Q''(1)$ against $23$ numbers of weight at most four returns coefficients $-3,-1,-96,-81,96$, residual $2\cdot10^{-390}$; written for $\ell''$ itself, with $\ell'\log3$ and $\ell\log^23$ moved across, the relation has height $144$. Neither $S_3$ nor $S_4$ lies in the span of the other numbers: the smallest relations have heights $10^{12}$ and $10^{14}$, the size expected when there is none.

\begin{remark}[What the slopes carry, sorted]\label{c:r:slopes21}
(i) \S20.3's negative was a property of the basis: none of its numbers was an arc integral of weight three ($\ell=2\sqrt3\,S_2$ and $\Gamma(\frac12)^2/\sqrt3=-\frac{\sqrt3}2S_1$ were there). Rem.~\ref{b:r:slopes}(i)'s ``new numbers, one per order'' is borne out, and the numbers have names: order one brings $S_3$, order two brings $S_4$, each beside a value $\Lambda_{-3}(j+2)$ of the same weight. (ii) ``At the price of one corner'' belongs to the positivity reading (Proposition~\ref{c:p:gaugeslopes}). (iii) The scale $(f^\star)^{1-k}$ removes every $\log3$: $Q$ is the ray read where its zero-corner solution is integral (Proposition~\ref{c:p:merit3}). (iv) \derived\ from Proposition~\ref{c:p:arcslopes}: the arc parts through order two are those of one function,
\begin{multline*}
Q(1+\varepsilon)=\frac{\sqrt3}{2\varepsilon}\Big[B\big(\tfrac12,\tfrac12;\tfrac34\big)-16^\varepsilon B\big(\tfrac12+2\varepsilon,\tfrac12;\tfrac34\big)\Big]+\tfrac{27}2\Lambda_{-3}(3)\,\varepsilon\\
-\tfrac12\Big[\Gamma(\tfrac12)^4\Lambda_{-3}(2)+27\Lambda_{-3}(4)+\tfrac{32}{\sqrt3}\Gamma(\tfrac12)^2\,\Tx3{\IZ}0\Big]\varepsilon^2+O(\varepsilon^3),
\end{multline*}
where $B(\frac12,\frac12;\frac34)=\frac23\Gamma(\frac12)^2$. The bracket is $\sum_{j\ge1}(4\varepsilon)^jS_{j+1}/j!$, because $16^\varepsilon u^{2\varepsilon}=(4u)^{2\varepsilon}$; so the arc part is the cut corner of the definition of $S_j$, read at $p=\frac12+2\varepsilon$. Why the cut sits at $\frac34$ is answered on the torus side of the same law: the arc runs to the zero $(1,e^{2\pi i/3},e^{4\pi i/3})$ of the letter $K=\alpha_1+\alpha_2+\alpha_3$ on the torus, where one mark's torus average changes form (Corollary~\ref{d:c:three}(iv)). Whether $Q$ near $k=1$ is a cut Beta corner at $\frac34$ plus Gamma corners: \notyet.
\end{remark}

\renewcommand{\thesection}{22}
\setcounter{proposition}{0}
\renewcommand{\theproposition}{22.\arabic{proposition}}

\section{The torus: two chains through one gauge}

At the centre of a list the native gauge is $\one$, a point of the orthant and of the torus, and the law of the values is also a law of the torus value. \S22.1 sets the two chains side by side: their densities are the two cut solutions on the line of Theorem~\ref{c:t:gauge}(iii), one on each side of the gauge. \S22.2 ties them at three marks by a connection with trigonometric coefficients, which reads the corner as the torus's residue at the zeros of the letter, the zero-corner solution as the torus read backwards, and the arc integrals of \S21.5 as the torus's log-moments of the letter.

\subsection*{22.1\quad Two chains through one gauge}

At the centre of a list, $\alpha^\star=\one$ lies on the orthant and on the torus $\{|\alpha_e|=1\}$. Since $\Psi(\alpha)=\alpha^{\one}K(\alpha^{-1})$ for every list, the stride on the torus is $f=|K(\alpha)|^{-2}$, real and at least $f^\star$ (Prop.~\ref{b:p:slitorigin}(i)).

\begin{definition}[The torus value]\label{d:d:torus}
At the centre, the \emph{torus value} is $T(k)=\int f^k\prod_{e\neq e_0}\frac{d\vartheta_e}{2\pi}$ over the torus $\alpha_e=e^{i\vartheta_e}$, $\alpha_{e_0}=1$, and the \emph{torus law} $\nu_T$ is the image of the normalized Haar measure under $f$, carried by $[f^\star,\infty)$.
\end{definition}

\begin{proposition}[The torus is the second chain]\label{d:p:torus}\derived\
At the centre of a list with a law:
\begin{enumerate}[label=(\roman*)]
\item $T$ converges on a half-plane $\operatorname{Re}k<k_T$ with $k_T>0$, $T(-j)=x_j$ are the closings, and $T$ solves the law there.
\item $\nu_T$ has a density $\varrho_T$; $\varrho_T\one_{(f^\star,\infty)}$ solves the density equation across $f^\star$; and near $f^\star$
\[
\varrho(t)=g(0)\,(f^\star-t)^{\frac{n-3}2}\,\Phi(f^\star-t),\qquad \varrho_T(t)=(2\pi)^{1-n}\,g(0)\,(t-f^\star)^{\frac{n-3}2}\,\Phi(f^\star-t),
\]
with one real-analytic $\Phi$, $\Phi(0)=1$: the chart's density and the torus's are the two cut solutions on the line of Theorem~\ref{c:t:gauge}(iii), one on each side of the gauge.
\item The corners read the torus: by Prop.~\ref{a:p:dirclose} the pole of $P$ at $k=-j$ has the leading coefficient $\mathrm{Cor}\cdot T(-j)$.
\end{enumerate}
\end{proposition}

\begin{proof}
(i) $f^k$ is bounded on the torus for $\operatorname{Re}k\le0$, and $|K|^{-2k}$ is integrable near the zeros of $K$ for small $\operatorname{Re}k>0$. The constant term of $(1/f)^j$ is its torus average. The relations along a ray (Thm.~35.3, p.~308) are exact forms; on the torus, which has no boundary, they integrate to zero wherever $f^k$ is smooth, which is $\operatorname{Re}k\ll0$, and the law continues to the half-plane. (ii) The argument of Thm.~\ref{c:t:gauge}(ii) on $[f^\star,\infty)$, with Mellin transforms on a left half-plane, and Theorem~\ref{c:t:gauge}(iii) on the side $s<0$. On the torus near $\one$, $\log f(e^{i\vartheta})=\log f^\star+\frac12\vartheta^{\mathsf T}H\vartheta+O(\vartheta^3)$ with the same $H$, so the Morse constant is that of Thm.~\ref{c:t:gauge}(i) with $d\vartheta/(2\pi)^{n-1}$ in place of $dy$. (iii) Prop.~\ref{a:p:dirclose} with (i).
\end{proof}

\noindent Receipts at three marks \computed\ (checks, by quadrature). The torus law of three marks has a closed density (alias, last); at the gauge it gives $\varrho_T(\frac19+)=\frac{27\sqrt3}{4\Gamma(1/2)^2}=(2\pi)^{-2}\cdot27\sqrt3\,\Gamma(\frac12)^2=(2\pi)^{-2}\varrho(\frac19-)$. And the moments over $(0,\frac19)$ of $(2\pi)^2$ times its continuation through $\frac19$ are $P(1)$, $P(2)$, $P(3)$ to $19$ digits: at three marks the chart's density is the torus's density continued across the gauge.

\subsection*{22.2\quad Three marks: the connection, and where the arc comes from}

Put $S(k)=x^{(1)}(k)/\Gamma(k)^2$, the sign-class solution of Prop.~\ref{b:p:otherbanana} (the lowered setting with one minus sign, Remark~\ref{a:r:lowered}).

\begin{theorem}[The connection of chart, torus and sign class]\label{d:t:connection}(i)--(ii) \computed; (iii) \derived.\
At the three-mark centre:
\begin{enumerate}[label=(\roman*)]
\item $S(k)=\frac13\Gamma(\frac13)\Gamma(\frac23)\,9^{1-k}\,T(1-k)$: the sign class is the torus read backwards.
\item The chart is a combination of the torus at $k$ and at $1-k$:
\[
P(k)=2\,\Gamma(\tfrac12)^4\Big(T(k)+\frac{\cot(\pi k)}{\sqrt3}\,9^{1-k}\,T(1-k)\Big),\qquad\text{equivalently}\qquad T=\frac{P}{2\Gamma(\frac12)^4}-\frac{3\cot(\pi k)}{2\Gamma(\frac12)^2}\,S .
\]
\item $T$ has a simple pole at $k=1$ with residue $-1/\big(\sqrt3\,\Gamma(\frac12)^2\big)$. It comes from the two zeros $(1,\omega,\omega^2)$, $(1,\omega^2,\omega)$ of $K$ on the torus, $\omega=e^{2\pi i/3}$, at each of which $|K|^2$ is a quadratic form of determinant $\frac34$ in $\vartheta$.
\end{enumerate}
\end{theorem}

\noindent Receipts. The coefficients of $T=\alpha P+\beta S$, solved from $k$ and $k+1$ at $k=0.1,0.2,\dots,0.9$ ($P$ from the gauge end, $T$ by its one-mark form below, $S$ by its integral; $14$ digits): $\alpha=0.0506605918212=1/(2\Gamma(\frac12)^4)$ at all nine, and $\beta\tan(\pi k)=-0.47746482927=-3/(2\Gamma(\frac12)^2)$; (i) at four points to $14$ digits. (iii): near a zero, $\int|K|^{-2k}$ over a disc is $\frac{2\pi}{\sqrt{3/4}}\int r^{1-2k}dr$, and the two zeros and the factor $(2\pi)^{-2}$ give the residue.

\begin{corollary}[What the connection carries]\label{d:c:three}\derived\ from Theorem~\ref{d:t:connection}; the moments $m_j$ \computed.\
\begin{enumerate}[label=(\roman*)]
\item \emph{The corner is the torus's residue at the zeros of the letter.} Near $k=0$, (ii) and (iii) give $\frac{2\Gamma(1/2)^4}{\sqrt3}\cdot\frac1{\pi k}\cdot9\cdot\frac1{\sqrt3\,\Gamma(1/2)^2k}=\frac6{k^2}$: $\mathrm{Cor}=6$. At $k=-j$ the same reading with Prop.~\ref{a:p:dirclose} gives $\operatorname{Res}_{k=1+j}T=-\frac{3\sqrt3}{\Gamma(1/2)^2}\,9^{-1-j}\,T(-j)$: the torus value's poles carry its own values at the negative integers.
\item \emph{The zero-corner solution is the torus read backwards.} $u_3(k)=9^{1-k}T(1-k)$, which at the integers is $9^{1-k}x_{k-1}$, the solution $c(K)$ of Prop.~\ref{c:p:merit3}. So the rate of the chart is the normalization of a chain that does not decay, and every rational solution of the law is that chain read backwards or carries the corner; this is why the merit of this law is negative.
\item \emph{The slopes at $k=1$ are torus data.} Write $T(-\varepsilon)=\sum_j\frac{(2\varepsilon)^j}{j!}m_j$, $m_j=\int\log^j|K|$ over the torus, and $T(1+\varepsilon)=-\frac{1}{\sqrt3\Gamma(1/2)^2\varepsilon}+T_0+T_1\varepsilon+\dots$ Then (ii) at $k=1+\varepsilon$ gives
\begin{align*}
\Pth&=2\Gamma(\tfrac12)^4\,T_0+\tfrac{4\Gamma(1/2)^2}{\sqrt3}\,(m_1-\log3),\\
P'(1)+2\Pth\log3&=2\Gamma(\tfrac12)^4\,T_1+\tfrac{4\Gamma(1/2)^2}{\sqrt3}\,(m_2+\log^23)-\tfrac{2\Gamma(1/2)^6}{3\sqrt3},
\end{align*}
with $m_1=\frac{\sqrt3}{4\Gamma(1/2)^2}\Pth$ and $m_2=\frac3{\Gamma(1/2)^2}S_3+\frac{\Gamma(1/2)^4}4$ \computed\ ($30$ digits; $S_3$ the arc integral of \S21.5). The first line gives $T_0=\frac{2\log3}{\sqrt3\,\Gamma(1/2)^2}$, and the second, with Prop.~\ref{c:p:arcslopes}, gives $T_1=\frac{\Gamma(1/2)^2}{6\sqrt3}-\frac{2\log^23}{\sqrt3\,\Gamma(1/2)^2}$; computed directly from the one-mark form of $T$ near $k=1$ they agree to $9$ and $6$ digits (a check). So the term $4\sqrt3\,S_3$ of Prop.~\ref{c:p:arcslopes} is $\frac{4\Gamma(1/2)^2}{\sqrt3}m_2-\frac{\Gamma(1/2)^6}{\sqrt3}$, the torus's second log-moment of the letter, and the rest of that relation is the torus's Laurent data at its pole, made of $\Gamma(\frac12)^2$, $\sqrt3$ and $\log3$.
\item \emph{Why the cut sits at $\frac34$.} Integrate one mark of the letter over its circle: with $A=\alpha_1+\alpha_2$,
\[
\int_0^{2\pi}|A+e^{i\vartheta_3}|^{-2k}\frac{d\vartheta_3}{2\pi}=\begin{cases}\sum_{m\ge0}\big(\frac{\Gamma(k+m)}{\Gamma(k)\,m!}\big)^2|A|^{2m},&|A|<1,\\[3pt] |A|^{-2k}\sum_{m\ge0}\big(\frac{\Gamma(k+m)}{\Gamma(k)\,m!}\big)^2|A|^{-2m},&|A|>1,\end{cases}
\]
two Beta-type series of one mark. On $\alpha_1=1$, $\alpha_2=e^{i\vartheta}$, $|A|^2=4\cos^2\frac\vartheta2$, and the two regimes meet at $|A|=1$: $\vartheta=\frac{2\pi}3$, $\sin^2\frac\vartheta2=\frac34$, the angle of the zero $(1,\omega,\omega^2)$ of $K$. So $T(k)$ is an integral over $\vartheta\in(0,\pi)$ split at the arc; its log-moments $m_j$ are arc integrals ($m_1=\frac1\pi\int_0^{2\pi/3}\log(2\cos\frac\vartheta2)\,d\vartheta$); and its pole at $k=1$ sits at the arc's endpoint. The cut Beta corner of Remark~\ref{c:r:slopes21}(iv) is this arc.
\end{enumerate}
\end{corollary}

\noindent So the arc integrals that \S21.5 found in the slopes of the chart are the torus side of the same law, and the connection carries them across.

\clearpage
\renewcommand{\thesection}{L}

\section{Ledger}

The running state of \S\S15--22 in one place. L.1 sorts every negative and every hypothesis met; L.2 is what is open, merged from the eight lists of the separate entries; L.3 logs every correction and where it now stands; L.4 gives the old numbers against the new; L.5 says where the receipts are. A later entry updates this ledger rather than adding a list of its own.

\subsection*{L.1\quad What is absolute and what is relative}

One test throughout: change the arena and see whether the negative survives. A negative is \emph{absolute} (a hypothesis \emph{genuine}) when it is a property of the ray or of the number; it is \emph{relative} when it is a property of the inputs read, the basis searched, the route of a proof, or the digits. A hypothesis that only reflects the route of a proof should go when the route changes; one that stays should have a reason that can be said in a sentence.

{\small
\renewcommand{\arraystretch}{1.15}
\begin{longtable}{@{}>{\raggedright}p{4.5cm} >{\raggedright}p{2.3cm} >{\raggedright\arraybackslash}p{8.2cm}@{}}
\toprule
statement, and where & sorted & what decides it\\ \midrule
\endfirsthead
\toprule
statement, and where & sorted & what decides it\\ \midrule
\endhead
\multicolumn{3}{@{}l}{\emph{The negatives of \S20}}\\[2pt]
``The answer to the first question is no'' (\S20.1, Ex.~\ref{b:e:Hslit}) & relative to the inputs & Theorem~\ref{c:t:gauge} reads every ray strictly inside its walls, for every $n$ (a multiple root at $f^\star$: Corollary~\ref{c:c:multiple}). Example~\ref{c:e:slitgauge} reads Ex.~\ref{b:e:Hslit} to $68$ digits with both corners as outputs, and \S21.4 reads every ray met.\\
Positivity cannot see a root on the slit (Prop.~\ref{b:p:blind}) & true; relative & Positivity is $\mu\ge0$. The slit direction is not analytic at $\lambda_c$; the density of $\mu$ is analytic on $(0,f^\star)$ (Remark~\ref{c:r:sees}).\\
Each order of slope costs a corner (Prop.~\ref{b:p:slopes}(iii), Rem.~\ref{b:r:slopes}(i)) & relative to the inputs & From the gauge end no order costs one (Proposition~\ref{c:p:gaugeslopes}).\\
No relation for $\ell'$ among the numbers of \S20.3 & relative to the basis & With the arc integrals $S_3$, $S_4$ the slopes $\ell'$, $\ell''$ have relations of heights $36$ and $144$ (Proposition~\ref{c:p:arcslopes}), and the arc integrals are the torus's log-moments of the letter (Corollary~\ref{d:c:three}(iii)).\\
The five-mark frame is no product of Beta corners (Rem.~\ref{b:r:periodmatrix}, first form) & false & Proposition~\ref{c:p:frame}. Balance is a property of a written monomial, not of a number (Remark~\ref{c:r:balance}).\\
A corner must supply the condition that positivity misses on the slit (Hyp.~\ref{b:h:channels}, last clause, first form) & not needed & The second corner of Ex.~\ref{b:e:Hslit} is an output of the gauge end (Example~\ref{c:e:slitgauge}(vi)).\\
A transform from $H$ to $\binom{[5]}1$ must turn $\frac12$ into $1$ with a logarithm, and a triple $0$ into a single one (Rem.~\ref{b:r:Hfive}) & partly relative & At $\frac1{25}$ and $\frac19$ that is what adding a mark does. At $t=1$ the curve of native gauges of $H$ must also become a point, and that part stands (Remark~\ref{c:r:transform}).\\
$N$: five masters at the lean, ten at the centre, numbers recognized by nothing (Thm.~\ref{b:t:Nlean}, Prop.~\ref{b:p:Ncentre}, first form) & wrong: the digits & The searches read digits the values did not have. From the gauge end: three masters and eight; the contents hold $\Pth$, the centre's also $\Zt$ (Proposition~\ref{d:p:Ncontent}, Remark~\ref{d:r:minors}). Two masters at the lean and six at the centre are new (L.2, \ref{o:NV}).\\
No law of $V$ (Prop.~\ref{b:p:V}) & computational & A law from a certificate or from a splitting of $\Psi_V$; then Theorem~\ref{c:t:gauge} reads $V$.\\
No ray of positive merit (\S20.8) & structural for the ray; open for the content & At the centre the rate $\log9$ is spent on the $3$-part of the denominators, and the merit is about $-1.9$ (Proposition~\ref{c:p:merit3}); the zero-corner solution is the torus read backwards (Corollary~\ref{d:c:three}(ii)). For $\mathbb Q+\mathbb Q\ell$ the question is whether $\ell$ is irrational. A merit belongs to a law, not to a content (Rem.~\ref{b:r:fixmerit}).\\
The procedure of Thm.~\ref{a:t:select}(iv) needs $|\lambda_c|>f^\star$ & true; absolute for that procedure & Theorem~\ref{c:t:gauge} does not use the procedure.\\[4pt]
\multicolumn{3}{@{}l}{\emph{The hypotheses of the readings}}\\[2pt]
$n$ even (the first proof of the gauge end) & route & Extension by zero separates the germ for every $n$ (Theorem~\ref{c:t:gauge}(iii), Remark~\ref{c:r:parity}); the five-mark centre is read (Example~\ref{c:e:fivegauge}).\\
The exponent at $f^\star$ is a check to be run (Ex.~\ref{c:e:slitgauge}(ii), first form) & route & It is forced, $\mu=\frac{n-3}2$ (Theorem~\ref{c:t:gauge}(iv)), and so a free test of every guessed law at every nondegenerate native value (\S21.3).\\
$f^\star$ a simple root of the symbol & route & The whole germ is local data at $\alpha^\star$ (Proposition~\ref{c:p:jet}, Corollary~\ref{c:c:multiple}). No law met has a multiple root at $f^\star$.\\
The series at $f^\star$ converges on $[0,f^\star)$ & route & $\varrho$ is analytic on $(0,f^\star)$; a complex singularity near $f^\star$ limits one series, not the continuation.\\
No root on the slit (\S20.1); $|\lambda_c|>f^\star$ & route & Those readings see only $\mu\ge0$, or the large-$k$ expansion truncated, and a solution carried below $f^\star$ is smaller than every term of that expansion. The germ is that expansion summed, and its continuation sees everything.\\
Rays through the origin & route & The gauge end sees only $\alpha^\star$: for $\int f^k\varphi\,dy$ with $\varphi>0$ analytic, $g(0)$ gains the factor $\varphi(\alpha^\star)$ and the jet gains $\log\varphi$.\\
Strictly inside the walls & genuine & On a wall the maximum of $f$ runs off to a face and $P(k)$ diverges; the reading moves to the face. That is the tree, not a defect.\\
Balance; integral directions & genuine & The chart integral needs degree $0$, and the law needs integer shifts of the exponents. With irrational exponents the level sets are not algebraic and there is no law.\\
A law of the values, certified & genuine & Closings are periods over the torus, values over the orthant; a closing recurrence is a law of the values when the orthant period solves it, which is what a certificate checks (L.2, \ref{o:certs}). At a centre the torus value solves the law (Proposition~\ref{d:p:torus}(i)), which is why closings find it.\\
Native values on the slit exist & genuine; costless & They are conjugates of an irrational $f^\star$ (Prop.~\ref{b:p:slitorigin}(ii)).\\
The merit of the centre is negative & genuine for that ray & The values decay exactly as fast as the closings grow (Proposition~\ref{c:p:merit3}), and the closings are a chain that does not decay (Corollary~\ref{d:c:three}(ii)).\\
The mismatch at $t=1$ between $H$ and $\binom{[5]}1$ & genuine for mark-adding & The doubled mark gives $H$ a curve of native gauges; the banana has an isolated sign class (Remark~\ref{c:r:transform}).\\
The cut at $\frac34$ in the slopes (Rem.~\ref{c:r:slopes21}(iv), first form: unexplained) & explained & One mark's torus average changes form at $|A|=1$, the end of the arc at the zero of the letter (Corollary~\ref{d:c:three}(iv)).\\
$N$'s new masters; Hyp.~\ref{b:h:channels} beyond three marks; which minors' numbers enter a content & not yet explained & Questions of basis or structure. The slopes and the torus point to integrals over the zero loci of the letters before more constants of the kinds already tried (L.2, \ref{o:corners}, \ref{o:NV}).\\
\bottomrule
\end{longtable}}

\noindent What stands as absolute: native values on the slit exist, and cost nothing; a law of the values is needed and must be certified, though every guessed law is now tested for free, at every native value, against the dimension of its chart; at $t=1$, the curve of native gauges of $H$, which adding a mark does not remove; the walls, where the reading moves to a face; and the arithmetic of a content, which no ray of the families met settles (whether $\ell$ is irrational is open). Everything else posed as a negative or a hypothesis in \S\S18--21 was a property of the inputs, the basis, the route of a proof, or the digits.

\subsection*{L.2\quad Open}

Each item ends with the old list items it replaces: $15.8$, $16.4$, $17.7$, $18.6$ and $20.9$ were the lists ``Open'' of \S\S15--18 and of the answers of \S20, $19$ the list at the end of \S19, $21.7$ and $22.8$ those of \S\S21--22 as first written.

\begin{enumerate}[label=\textbf{O\arabic*},ref=O\arabic*,leftmargin=*,itemsep=4pt]
\item\label{o:certs} \emph{Certificates} (Thm.~35.3, p.~308) for the laws read from closings: the banana for every $n$ by an operator proof, which would derive every banana ray law from the stacked Gamma, the law of each list produced by the one before (at each $n\le6$ the same $A_s$ appeared; from five to six marks the symbol of the five-mark law vanishes at its native value $1$, $L_5^\ast$ raises degrees by less than four, and the block of six marks is built differently), and one certificate with $s$ a variable, as Theorem~\ref{a:t:cert} has for the banana; a series root, where the input is a product and the laws have orders $6$ and $5$; the laws of \S\S18.3--18.4, among them the lean $(2\one;4,1)$ and $H$'s law of order $4$; the fixture at $(\one;0,2)$, the ray of Example~\ref{b:e:Hslit}, $N$ at both directions, and $V$; the converse of Corollary~\ref{a:c:dirlaw} in general. Each with bounds on the pieces down to $k\to0^+$, as \S19, Step 5, has for three marks (the four-mark banana ray needs the same bounds, with its own tropical stride), and with the corner as the boundary term of the certificate at $k=0$, where $g^{(0)}_e$ does not vanish at infinity, instead of a limit. Certified so far: the lines of Proposition~\ref{a:p:certs} and the three-mark law of \S19. Theorem~\ref{c:t:gauge} needs a law of the values; every law guessed from closings passes the test of \S21.3. \notyet\ {\small[15.8(i), 16.4(iii), 17.7(i), 18.6(i), 18.6(vii), 19(i), 19(iii), 20.9(i), 21.7(i), 22.8(i)]}
\item\label{o:brackets} \emph{Certified brackets.} Positivity polytopes for the two masters of $H$ (\S18.4) and for the lean $(2\one;4,1)$, whose rational solutions have unproved signs; from the gauge end, tail bounds for the series at $f^\star$ (majorant series), which would make every reading of \S21.4 a bracket. \notyet\ {\small[19(i), 19(ii), 21.7(ii)]}
\item\label{o:slit} \emph{The slit.} Which algebraic native values sit on $(0,f^\star)$ for a ray strictly inside its walls, beyond the conjugates of an irrational $f^\star$ (Proposition~\ref{b:p:slitorigin}); along walls, whether they are exactly the faces' native values, carrying the faces' masters. That there are none strictly inside walls, the first half of the old 18.6(ii), is false (Example~\ref{b:e:Hslit}). For reading a ray the question is moot (Theorem~\ref{c:t:gauge}); it stays a question about native values. \notyet\ {\small[18.6(ii), 20.9(ii)]}
\item\label{o:rates} \emph{Rates.} The lower bound of Theorem~\ref{b:t:rate} on every ray; it is attained on the three-mark centre. \notyet\ {\small[20.9(iii)]}
\item\label{o:slopes} \emph{Slopes.} The three-mark centre near $k=1$: a closed form of $Q(1+\varepsilon)$ (Remark~\ref{c:r:slopes21}(iv)), whether it is a cut Beta corner at $\frac34$ plus Gamma corners, and the slopes of every order against $S_{j+2}$ and $\Lambda_{-3}(j+2)$; the corner $a_5$ in the family's terms, and the relation of degree two carried by $a_4$ (Remark~\ref{b:r:slopes}(iv)). Through order two the slopes are identified (Proposition~\ref{c:p:arcslopes}) and placed on the torus (Corollary~\ref{d:c:three}(iii)). \notyet\ {\small[20.9(iv), 21.7(iv)]}
\item\label{o:corners} \emph{Second corners and channels.} The second corners from the list: the boundary term at $\sigma=-1$ as a constant term of the continued period (Thm.~1.6), assembled from its Laurent coefficients of orders $-(n-1)$, $-(n-2)$, $-(n-3)$ at $\sigma=0$ and $\sigma=-1$: the leading ones are the corner and the corner times $|\cB|$ (Proposition~\ref{a:p:negcorner}), the next vanishes at $0$ (Proposition~\ref{a:p:corner}) and the ray law fixes it at $-1$, and on the share side the third ones are Laurent data of the Gammas at the lattice points of the closing, heads of the counting input and Trailmoments of the counting input cut at the data (Prop.~5.13, p.~32; Prop.~5.14, p.~33), so that the rationality of the two boundary terms says that the Trailmoments cancel and the heads remain; which lists put $(\sigma+1)^2$ into both $c_0$ and $c_1$ (in \S15.4 the two triangles $\binom{[3]}2$ with parallel classes, not $\binom{[4]}2$ with a mark doubled); the second corner of Example~\ref{b:e:Hslit} as a sum of face numbers; Proposition~\ref{b:p:facecorners} at the other poles and at the walls of a split; Hypothesis~\ref{b:h:channels} beyond three marks. Every second corner met is an output of the gauge end (Example~\ref{c:e:slitgauge}(vi)); what is open is its structure. \notyet\ {\small[15.8(v), 20.9(v)]}
\item\label{o:coaction} \emph{The coaction.} The block of the note's Theorem~A that gives the Beta corners of the five-mark frame (Proposition~\ref{c:p:frame}, Remark~\ref{c:r:balance}); the number that multiplies the zero-corner solution at the fixture's full lean (Remark~\ref{b:r:periodmatrix}); the period matrix of the centre of $H$, that is, the solutions selected at its other native gauges, and whether they are stacked Gammas, as the lowered settings of $\binom{[5]}1$ are (Proposition~\ref{a:p:stackedhalf}). \notyet\ {\small[18.6(vi), 20.9(vi)]}
\item\label{o:transform} \emph{A transform between $H$ and $\binom{[5]}1$.} Whether the two laws are related by adding a mark together with a transform that shrinks the curves of native gauges of $H$ at $t=1$ to a point; their frames (Prop.~35.8, p.~309). The local exponents are known (Proposition~\ref{b:p:density}, Remark~\ref{b:r:Hfive}), and adding a mark accounts for $\frac1{25}$ and $\frac19$ (Remark~\ref{c:r:transform}). \notyet\ {\small[18.6(v), 20.9(vii), 21.7(vi)]}
\item\label{o:connection} \emph{Connections between the chains.} The connection data of the density between $t=0$ and $t=f^\star$, which carry the numbers (Remark~\ref{b:r:Hfive}); off the bananas, what replaces the identity $\mathrm{Cor}\cdot\det H=\prod_{c\neq\star}(1-\lambda^\star/\lambda_c)$ of \S18.5(i), which ties the two ends of strength. At three marks Theorem~\ref{d:t:connection} gives them, computed; open are a derivation inside the family (the chart, the torus and the sign class are three real forms of one pencil, and the route is the rotation $y\mapsto-y$ of the one-mark Gramma $\sG_{0;(1,-1)}(1,y)$ on the share side, Proposition~\ref{a:p:bananapin}); the connection at four and five marks, with periodic coefficients, and what replaces the arc at four marks, where $K$ vanishes on the torus along a curve; the second chain at a lean, where the stride on the torus through $\alpha^\star$ is not real and the second real form through the gauge is not a torus. \notyet\ {\small[20.9(vii), 21.7(iii), 22.8(iii), 22.8(iv), 22.8(vii)]}
\item\label{o:NV} \emph{$N$ and $V$.} The law of $V$, by a splitting of $\Psi_V$ along classes or by a certificate, after which Theorem~\ref{c:t:gauge} reads it; the law of order $11$ of $N$'s centre, exactly; the two new masters at $N$'s lean and the six at its centre, against $N$'s torus (the log-moments of $K_N$ and the Laurent data of its torus value) and the arc integrals of its letters' zeros; the third relation of Proposition~\ref{d:p:Ncontent}(i); which minors' numbers enter a ray's content, beyond the two lists of Remark~\ref{d:r:minors}. \notyet\ {\small[20.9(viii), 21.7(v), 22.8(v), 22.8(vi)]}
\item\label{o:content} \emph{The content of a list.} Whether every ray strictly inside its walls without the list's symmetries spans the whole content (the hypothesis of \S18.5(ii)): an offset ray, a ray breaking the symmetries $a\leftrightarrow b$, $2\leftrightarrow3$ of $H$, and an offset ray of $N$. \notyet\ {\small[18.6(iii), 20.9(viii)]}
\item\label{o:Hhinge} \emph{$H$ at its hinge.} The masters of $H$ through its hinge at a mark, as in \S\S17.1--17.3, with the weights; what the tree expects for $H$ (the end of \S18.4). \notyet\ {\small[18.6(iv)]}
\item\label{o:merit} \emph{Arithmetic.} A ray of positive merit; a ray of content $\mathbb Q+\mathbb Q\ell$ whose zero-corner solution is not the closings at the scale $(f^\star)^{k-1}$ (Remark~\ref{c:r:merit21}, Corollary~\ref{d:c:three}(ii)); families on lists other than bananas, from a share side as in Proposition~\ref{b:p:families}; the denominators of the brackets of the whole moment problem, which are rational at every degree (\S20.1); the other rows of the table of \S20.8 recomputed with $\frac1K\log D_K$. \notyet\ {\small[20.9(ix), 21.7(vii)]}
\item\label{o:ell} \emph{The three-mark centre inside the family.} Reduce the live Trailmoment of the receipts of Proposition~\ref{a:p:addmark} at $n=2$ to the Trailmoment of the odd periodic density modulo $3$ at index $2$, inside the family; Thm.~5.17 does not reach it, the parity being mismatched. The identity $\ell=3\Lambda_{-3}(2)$ is computed to $390$ digits (Proposition~\ref{c:p:arcslopes}), and on the torus $m_1=\frac{\sqrt3}{4\Gamma(1/2)^2}\ell$ (Corollary~\ref{d:c:three}(iii)). \notyet\ {\small[15.8(iii)]}
\item\label{o:block} \emph{The direction $x-1$ of the block.} $\Tx{0}{(x-1)\,p\,\IZ}{0}=(q\partial_q)G_0(q)|_{q=1}=\int dy/(X-1)^2=0.16106970227667485885\ldots$, which the ray reaches only through the value term of its second master; expanding $h$ as in Proposition~\ref{a:p:cosetmark}, it is the lowered setting with $h$ replaced by $y\,\partial_yh$, $\int_0^\infty g^4\,y\,h'(y)\,dy$ \derived. Which member of the family it is, and whether a lean of the five-mark list reads it as a period. \notyet\ {\small[15.8(ii)]}
\item\label{o:degtwo} \emph{Relations of degree two; pairings.} The last minor of \S14.17, $M_5(1)x^{(1)}(2)-M_5(2)x^{(1)}(1)=4.34000293153566\ldots$, is the minimal solution of the exterior square of the ray law, of native value $\frac1{225}$ \derived, and it is the masters over the stacked Gamma (Remark~\ref{a:r:found}(iii)), so a relation of degree two fixing it would tie the masters to $x^{(1)}(1)$. Whether pairings of the sign classes with closings of the share side twisted by the signs fix such relations, beside the exact pairing of Proposition~\ref{a:p:pairing}, a relation of degree one between solutions fixed by the residue at zero shares; the products of values that the corner $a_4$ writes linearly in the slopes (Remark~\ref{b:r:slopes}(iv)) are relations of this kind. The pairing for the other lists: the partner of the period in the adjoint law, when the closings are not. \notyet\ {\small[15.8(vi), 15.8(vii)]}
\item\label{o:leans} \emph{Lean rays.} At a lean with uniform masses only $\eta=\one$ stays native for $\binom{\cE}1$; at a sign the native equation reads $(m-t)K(\eta)=(s-t)\eta_e$ for every $e$ \derived. The gauges that move off the signs, and their values. \notyet\ {\small[15.8(viii)]}
\item\label{o:families} \emph{Families of native gauges.} In both cases of Proposition~\ref{a:p:families} the letter collapses at the gauge to the letter of a minor that leaves directions free: a class whose contraction letter vanishes is free, and a class whose own letter vanishes leaves a deletion with coloops, whose relative scales are free. A rule for every list, and the exponents $-3$ and $-\frac52$ of Remark~\ref{a:auto:15.16}. The multiple roots of the symbols met are the families (\S21.3), and at a family of dimension $d$ the density has the exponent $\frac{n-3-d}2$ (Remark~\ref{c:r:transform}). \notyet\ {\small[15.8(ix)]}
\item\label{o:MK4} \emph{$\mathrm M(K_4)$.} The six-mark atom is self-dual only up to a relabelling of its marks, so Cor.~22.8(iii) does not split it (as the handbook notes after Cor.~22.8). Its centre ray, read opened at a mark into $K_4\setminus e$ and $K_4/e$, the latter the triangle with two marks doubled of \S17.4. Its full lean $(\one;0,2)$ is the fixture of Theorem~\ref{b:t:fixture}. \notyet\ {\small[15.8(x), 16.4(iv), 17.7(vi)]}
\item\label{o:bounds17} \emph{The boundaries of \S17.} The finite part $\frac{21}2+\frac34\Pfour$ at the wall $d=2$ from the collapse of Cor.~8.8 as a regularized lean of the four-mark banana, and the third relation of Proposition~\ref{a:p:u25free} from the next order; the relation of Proposition~\ref{a:p:r4zero} that carries $\Gamma(\frac12)^4$, from the next Laurent order at zero shares; Theorem~\ref{a:t:r4masters}(i) and (ii) as boundaries, (i) tying the lines $1$ and $2$ with the coefficient $8=|\cB|$, and which insertion reads it (a mean share as in Proposition~\ref{a:p:block}?); the other weights $A_s$, $B_s$, $R_s$, and $C_s$ for $s\ge2$, of $\binom{[5]}2$, as periods of the surface $e_2^2=e_1e_3$ at its $\mathbb Q(\sqrt{-5})$ point. \notyet\ {\small[17.7(ii)--(v)]}
\end{enumerate}

\noindent\emph{Closed since they were posed.}
\begin{itemize}[leftmargin=*,itemsep=1pt]
\item The three masters of $\binom{[5]}2$, and the law in $m$ of the ray of $\binom{[4]}2$ with its closing: \S\S17.1--17.3 (Propositions~\ref{a:p:u25law} and~\ref{a:p:u25clos}, Theorem~\ref{a:t:u25masters}). {\small[15.8(iv), 16.4(i)]}
\item Two marks doubled, closed, with its second corner as a boundary: \S17.4 (Theorem~\ref{a:t:r4masters}, Remark~\ref{a:r:closingfound}(ii)). {\small[16.4(ii)]}
\item Certificates for the lines used in \S17 (Proposition~\ref{a:p:certs}). {\small[17.7(i), first half]}
\item The local exponents of $H$ and $\binom{[5]}1$ (Proposition~\ref{b:p:density}, Remark~\ref{b:r:Hfive}). {\small[18.6(v), first half]}
\item A list on six marks that is not a graph, with native gauges off the signs and off the torus: $N$ (\S20.7). Its certificate is in \ref{o:certs}. {\small[18.6(vii)]}
\item The eight questions of \S20: each is answered at the head of \S\S20.1--20.8; what stays open of them is above.
\item The gauge end for odd $n\ge5$: Theorem~\ref{c:t:gauge}, Example~\ref{c:e:fivegauge}. {\small[21.6(iii) of the version of \S21 that \S22 cites]}
\item Why the arc ends where $K$ vanishes on the torus: Corollary~\ref{d:c:three}(iv). {\small[21.7(iii), first half]}
\item A multiple root at $f^\star$: Corollary~\ref{c:c:multiple}; no example met. {\small[22.8(ii)]}
\item The masters of $N$: three at the lean and eight at the centre, not five and ten (Proposition~\ref{d:p:Ncontent}). {\small[20.9(viii), in part]}
\end{itemize}

\subsection*{L.3\quad Corrections}

Every correction made by a later entry, with the entry that made it. All are applied in place in this file; the old wording survives only here.

\begin{enumerate}[label=\textbf{C\arabic*},ref=C\arabic*,leftmargin=*,itemsep=3pt]
\item {\small[17.6(i)]} \S15.4, the masters column, counts over $\mathbb Q$. Over $\mathbb Q+\mathbb Q\Pth$ the list $\binom{[5]}2$ has two masters, and over $\mathbb Q+\mathbb Q\Pth+\mathbb Q\Gamma(\frac12)^4$ the triangle with two marks doubled has one (Theorems~\ref{a:t:u25masters} and~\ref{a:t:r4masters}).
\item {\small[17.6(ii)]} Remark~\ref{a:auto:16.9}: the three-mark centre enters every row of \S15.4 but the banana, row four included, through the contraction of $H$ at its single mark; not only the rows whose trees carry a three-mark piece.
\item {\small[17.6(iii)]} Remark~\ref{a:auto:16.8}: two of the three masters of $\binom{[5]}2$ are slope Trailmoments of the deletion's ray, and the third is $\Pth$; not all three.
\item {\small[19.7(i)]} The receipts of Proposition~\ref{a:p:certs}: Thm.~35.3 as stated gives the relations for large $k$; down to $k\to0^+$ it needs bounds on the pieces, which \S19, Step 5, supplies for three marks. Not ``applies as stated''.
\item {\small[19.7(ii)]} Remark~\ref{a:r:paper}: the ratios $-r/c$ are lower bounds at every $q$, since $c(q)>0$ at every $q$ (\S19, Step 8), not only at the $q$ computed.
\item {\small[19.7(iii)]} Theorem~\ref{a:t:select}: a root of the symbol need not be a native value (the extra root $\frac{131003}{4771109}$ of \S18.4). Its (ii) speaks of roots, and the count of one condition in (iii) is for a simple root that is the value of a nondegenerate native gauge.
\item {\small[20.10(i)]} The open item 18.6(ii), first half (``no native value in $(0,f^\star)$ strictly inside the walls''): false; Example~\ref{b:e:Hslit} has one, the conjugate of an irrational $f^\star$. Now \ref{o:slit}.
\item {\small[20.10(ii)]} Remark~\ref{a:r:onslit} stands as written; a sentence pointing to Example~\ref{b:e:Hslit} and to \S21 is added.
\item {\small[20.10(iii)]} Theorem~\ref{a:t:select}(iv) needs $|\lambda_c|>f^\star$ for every root $\lambda_c\neq f^\star$ that carries a local solution not holomorphic at $q_c$, not only $\lambda_c$ off the slit; added to its hypotheses. The whole moment problem needs only $\lambda_c\notin[0,f^\star]$ (Theorem~\ref{b:t:rate}).
\item {\small[20.10(iv)]} The open items 18.6(vii) and 18.6(v) are answered in part: by $N$ (\S20.7) and by the local exponents (Proposition~\ref{b:p:density}). The rest is in \ref{o:certs} and \ref{o:transform}.
\item {\small[21.8(i)]} Remark~\ref{b:r:periodmatrix}: ``which is not balanced and is no product of Beta corners \derived'' is false (Proposition~\ref{c:p:frame}, Remark~\ref{c:r:balance}, which now follow it in \S20.5).
\item {\small[21.8(ii)]} \S20.1, last paragraph: ``So the answer to the first question is no'' holds for the inputs of \S19 only (Theorem~\ref{c:t:gauge}).
\item {\small[21.8(iii)]} Hypothesis~\ref{b:h:channels}: its last clause (``wherever a seen root sits on the slit a corner of that ray supplies the missing condition'') is not needed, and is removed (Example~\ref{c:e:slitgauge}).
\item {\small[21.8(iv)]} Proposition~\ref{b:p:slopes}(iii) and Remark~\ref{b:r:slopes}(i): ``each order of slope costs one corner'' holds for positivity only (Proposition~\ref{c:p:gaugeslopes}).
\item {\small[21.8(v)]} Proposition~\ref{b:p:density}, proof: ``the boundary terms at $0$ and $f^\star$ cancel by themselves'' is replaced by Theorem~\ref{c:t:gauge}(ii).
\item {\small[21.8(vi)]} \S20.8, the table: the growth column is the quotient $(\log D_{240}-\log D_{200})/40$, not $\limsup\frac1K\log D_K$; at the centre the merit is about $-1.9$, not $-1.545$ (Proposition~\ref{c:p:merit3}). The other rows are in \ref{o:merit}.
\item {\small[21.8(vii)]} \S20.3: the list announced as $25$ numbers has $24$.
\item {\small[22.9(i)]} Theorem~\ref{b:t:Nlean}(iv): the positivity reading at $K=200$ is good to $189$ digits, not $400$; its rate is $10^{-0.973K}$.
\item {\small[22.9(ii)]} Theorem~\ref{b:t:Nlean}(v) and Remark~\ref{b:r:NV}(iii): ``five masters, independent over $\mathbb Q$ with $1$ as far as height $10^{64}$'' is false; there are three, and the content holds $\Pth$ (Proposition~\ref{d:p:Ncontent}(i)--(ii)).
\item {\small[22.9(iii)]} Proposition~\ref{b:p:Ncentre}(vii) and Remark~\ref{b:r:NV}(iv): ``ten masters'' is false; there are eight, and the content holds $\Pth$ and $\Zt$ (Proposition~\ref{d:p:Ncontent}(iii)--(iv)). ``No relation among $1,P(1),\dots,P(10)$ below $10^{18}$'' was true and too weak: two exist at height $10^{33}$.
\item {\small[22.9(iv)]} \S20.7, the answer to the seventh question: revised.
\item {\small[22.9(v)]} The gauge end ``for even $n$'' (its first form): the parity belongs to that proof; Theorem~\ref{c:t:gauge} covers every $n$ (Remark~\ref{c:r:parity}).
\item {\small[22.9(vi)]} Example~\ref{c:e:slitgauge}(ii) and the first sorting of \S21: the exponent at $f^\star$ is forced (Theorem~\ref{c:t:gauge}(iv)); it is not a check.
\item {\small[22.9(vii)]} Remark~\ref{c:r:slopes21}(iv): ``why the cut sits at $\frac34$: \notyet'' is answered by Corollary~\ref{d:c:three}(iv).
\item {\small[this file]} Theorem~\ref{c:t:gauge}(iii), proof: the count of holomorphic solutions assumed that $\mu$ is not an integer $\ge D-1$. That this holds for every law of the values is now proved, in the last paragraph of that proof.
\item {\small[this file]} Remarks~\ref{a:auto:16.8} and~\ref{a:auto:16.9}: three tags \notyet\ that \S17 answered are replaced by the answers (Propositions~\ref{a:p:u25law} and~\ref{a:p:u25clos}; Theorem~\ref{a:t:r4masters}; Remark~\ref{a:r:closingfound}(ii)). \S19.6: the rate of the polytope now points to Theorem~\ref{b:t:rate}.
\item {\small[this file]} Remarks~\ref{b:r:periodmatrix} and~\ref{b:r:Hfive}: a tag \notyet\ each, restated by the remarks that now follow them (Remarks~\ref{c:r:balance} and~\ref{c:r:transform}), is folded into those.
\end{enumerate}

\subsection*{L.4\quad Old numbers and new}

Numbers $15.x$--$19.x$ are unchanged. The entries that carried $20.x$--$22.x$ are merged and renumbered; their numbered statements are listed here by their printed numbers in the entry where they first appeared. \emph{Questions} is the entry of 23 October ending the file of \S\S15--20; \emph{Answers} is the entry of 23 October answering it; \S21 is the entry of 24 October in its last form, with the numbers of the earlier form that \S22 cites in brackets where they differ; \S22 is the entry of 24 October. A statement merged into another points to the one that stands.

{\small
\renewcommand{\arraystretch}{1.08}
\setlength{\tabcolsep}{3pt}
\begin{longtable}{@{}l@{\,$\to$\,}l@{\qquad\qquad}l@{\,$\to$\,}l@{\qquad\qquad}l@{\,$\to$\,}l@{}}
\toprule
\multicolumn{6}{@{}l}{old number $\to$ number in this file}\\ \midrule
\endfirsthead
\toprule
\multicolumn{6}{@{}l}{old number $\to$ number in this file}\\ \midrule
\endhead
\multicolumn{6}{@{}l}{\emph{Questions (\S20 of 23 October)}}\\[1pt]
Prop.~20.1 & 20.3 & Prop.~20.2 & 20.4 & Prop.~20.3 & 20.5\\
Prop.~20.4 & 20.6 & Q.~20.5 & 20.7 & Q.~20.6 & 20.13\\
Q.~20.7 & 20.16 & Q.~20.8 & 20.20 & Q.~20.9 & 20.23\\
Q.~20.10 & 20.28 & Q.~20.11 & 20.32 & Q.~20.12 & 20.41\\[3pt]
\multicolumn{6}{@{}l}{\emph{Answers (\S20 of 23 October)}}\\[1pt]
Def.~20.1 & 20.1 & Prop.~20.2 & 20.2 & Prop.~20.3 & 20.8\\
Thm.~20.4 & 20.9 & Prop.~20.5 & 20.10 & Prop.~20.6 & 20.11\\
Ex.~20.7 & 20.12 & Prop.~20.8 & 20.14 & Thm.~20.9 & 20.15\\
Prop.~20.10 & 20.17 & Prop.~20.11 & 20.18 & Rem.~20.12 & 20.19\\
Prop.~20.13 & 20.21 & Hyp.~20.14 & 20.22 & Prop.~20.15 & 20.24\\
Rem.~20.16 & 20.25 & Prop.~20.17 & 20.29 & Rem.~20.18 & 20.30\\
Def.~20.19 & 20.33 & Prop.~20.20 & 20.34 & Thm.~20.21 & 20.35\\
Prop.~20.22 & 20.36 & Prop.~20.23 & 20.37 & Rem.~20.24 & 20.38\\
Prop.~20.25 & 20.42 & Prop.~20.26 & 20.43 & Rem.~20.27 & 20.44\\[3pt]
\multicolumn{6}{@{}l}{\emph{\S21 (24 October; in brackets the numbers of the earlier form, which \S22 cites)}}\\[1pt]
Thm.~21.1 & 21.1 & Rem.~21.2 & 21.3 & Rem.~21.3 [21.2] & 21.4\\
Ex.~21.4 [21.3] & 21.7 & Ex.~21.5 & 21.8 & Prop.~21.6 [21.4] & 21.9\\
Prop.~21.7 [21.5] & 21.10 & Rem.~21.8 [21.6] & 21.11 & Prop.~21.9 [21.7] & 20.26\\
Rem.~21.10 [21.8] & 20.27 & Prop.~21.11 [21.9] & 20.45 & Rem.~21.12 [21.10] & 20.46\\
Rem.~21.13 [21.11] & 20.31 & Prop.~21.14 & 21.5 & Cor.~21.15 & 21.6\\
Rem.~21.16 & L.1 & \multicolumn{2}{l}{} & \multicolumn{2}{l}{}\\[3pt]
\multicolumn{6}{@{}l}{\emph{\S22 (24 October)}}\\[1pt]
Thm.~22.1 & 21.1 & Rem.~22.2 & 21.3 & Ex.~22.3 & 21.8\\
Prop.~22.4 & 20.39 & Rem.~22.5 & 20.40 & Def.~22.6 & 22.1\\
Prop.~22.7 & 22.2 & Thm.~22.8 & 22.3 & Cor.~22.9 & 22.4\\[3pt]
\bottomrule
\end{longtable}}

\noindent Merged statements: Theorem~22.1 of 24 October is Theorem~\ref{c:t:gauge}(iii)--(iv), Remark~22.2 is Remark~\ref{c:r:parity}, and Example~22.3 is Example~\ref{c:e:fivegauge}; the proof by moments of Theorem~21.1(iii) of 24 October is Remark~\ref{c:r:lawline}; the sorting tables (\S21.0, Remark~21.16, \S22.0 of 24 October), the lists ``Open'' and ``Corrections'' of every entry, and the sections ``Sorted again'' and ``Where \S21 stands'' are L.1--L.3. Proposition~20.2 of the questions entry, ``The values are the moments of the law of the stride'', is Proposition~\ref{a:p:stridelaw}, now titled ``The two ends of the stride's law''.


\subsection*{L.5\quad Receipts}

One archive, \texttt{receipts\_s15\_s22.zip}, with a folder per section and a \texttt{MANIFEST.txt} that maps each numbered statement to the script that prints it. \S\S15--17 were printed by \texttt{centre\_ray.py}, which was delivered with those entries and is not in the archive; \S18 by \texttt{ray\_reading\_s18.py}; \S19 by \texttt{exam\_s19.py}; the questions of \S20 by \texttt{questions\_s20.py}; the answers of \S20 by the receipts of that entry, whose toolkit and laws read from closings are carried in \texttt{s22/s20/}; \S21.1--21.2, Example~\ref{c:e:slitgauge}, the three-mark reading, \S21.5 and the propositions of \S20 that came from the entry of 24 October (Propositions~\ref{c:p:frame} and~\ref{c:p:merit3}) by \texttt{receipts\_s21.py}; \S21.3, \S21.4, Example~\ref{c:e:fivegauge}(ii), Proposition~\ref{d:p:Ncontent}, Remark~\ref{d:r:minors} and \S22 by the scripts of \texttt{s22/}. Every script uses python-flint and the standard library only, and in every entry the checks (tropical sampling, quadrature) are never part of a reading.

\section*{Aliases, last}
\addcontentsline{toc}{section}{Aliases, last}

\noindent\emph{Constants and functions.} $\Gamma(\frac12)^2=\pi$; $\Tx m{\IZ}0=\zeta(m)$, so $\Zt=\zeta(3)$; $\ell=\Pth=3L(\chi_{-3},2)$ \recalled, which agrees with the check of \S19 to $40$ digits and with the gauge end to $390$; $\Pfour=7\zeta(3)$; $\Lambda_D=L(\chi_D,2)$, the Kronecker sign being the quadratic character of discriminant $D$, and $\Lambda_{-3}(s)=L(\chi_{-3},s)$, so $\Lambda_{-3}(3)=\frac{4\pi^3}{81\sqrt3}$, which $\Gamma(\frac12)^6/\sqrt3$ stood for in \S18; $S_j=\mathrm{Ls}_j(\frac{2\pi}3)$, the log-sine integrals, and $S_2=\mathrm{Cl}_2(\frac{2\pi}3)$, the Clausen value; $B(p,q;x)$ is the incomplete Beta function; the $\varepsilon$ of \S20.7 is the regulator of $\mathbb Q(\sqrt{21})$. $\sG_{\nu;(1,-1)}(1,y)=2y^{\nu/2}K_\nu(2\sqrt y)$, so $\sG_{\nu;(1,-1)}(a,a)=2K_\nu(2a)$, the leans are moments of products of Bessel functions $K_\nu$ of integer orders, and Proposition~\ref{a:p:leanindex} reads the lean as the order of the Bessel functions. In Bessel form the three-mark centre is $P(k)=4^{2-k}\Gamma(k)^{-2}\int_0^\infty x^{2k-1}K_0(x)^3dx$, the equal-mass two-loop vacuum (sunrise) integral with propagator powers $k$ in dimension $2k$; log-sine integrals at $\frac{2\pi}3$ are the constants of the $\varepsilon$-expansions of that integral \recalled.

\noindent\emph{Lists.} $\binom{[n]}1$ is the banana; $H$ is the cycle matroid of a triangle with one edge doubled; $\binom{[4]}1$ and $\binom{[4]}2$ are $U_{1,4}$ and $U_{2,4}$; $\mathrm M(K_4)$ is the cycle matroid of $K_4$, the wheel with three spokes. $N$ is the matroid of six points of a plane with exactly three of them collinear, $U_{3,6}$ with one basis removed; it has a $U_{2,4}$ minor, so it is neither binary nor graphic (Tutte). $V$, the eight-mark list, is the V\'amos matroid, realizable over no field. At the full lean $(\one;0,2)$ of a graph the corner is the Hepp bound and $P(1)$ the Feynman period, $6\zeta(3)$ for the fixture; the observed correlation between the two, with the conjecture that equal Hepp bounds mean equal periods, is the classical side of Question~\ref{a:q:corner} (Panzer); tropical sampling is tropical Monte Carlo over Hepp sectors.

\noindent\emph{\S15.} $x^{(1)}(1)=8\int_0^\infty tI_0(t)K_0(t)^4\,dt$, whose evaluation by the Gammas at fifteenths is the classical one through the singular modulus of $\mathbb Q(\sqrt{-15})$ \recalled; Theorem~\ref{a:t:cert} gives the recurrences of the moments $\int_0^\infty t^{2s-1}K_0(t)^n\,dt$, $n\le7$, by adding marks. The direction $x-1$ of the block (\ref{o:block}) is $4\int_0^\infty t^2I_1(t)K_0(t)^4\,dt$, to $25$ digits \computed.

\noindent\emph{\S17.} $C_1$ is the Chowla--Selberg period of $\mathbb Q(\sqrt{-5})$; $X_{\setminus}-X_{/}=(e_2^2-e_1e_3)/e_4$, so $x_1=\int dy/(X_{\setminus}-X_{/})$ is a period of the quartic surface $e_2^2=e_1e_3$.

\noindent\emph{\S18.} The $x_j$ are the classical periods of the Laurent polynomial $1/f$, and the law is the Picard--Fuchs operator of its pencil read on coefficients; Theorem~\ref{a:t:select}(iv) is the mechanism of Ap\'ery limits, the selection read off the slit (Pincherle's minimal solution, Perron's theorem).

\noindent\emph{\S19.} The closings are the sums of squares of trinomial coefficients, the even moments of the planar walk of three unit steps; (L) is an Ap\'ery-like recurrence, and Step 9 is Ap\'ery's limit made two-sided by positivity (Pincherle's minimal solution, read with signs); the certificate is a creative-telescoping certificate for the integral; $\varphi$ is the tropicalization of $-\log f$, and $\int e^{-k\varphi}dy=6/k^2$ its volume.

\noindent\emph{\S20.} For $\binom{\cE}1$ the body $\mathcal P-\mathcal P$ of Proposition~\ref{a:p:tropcorner} is the root polytope of type $A_{n-1}$. The stride's law $\mu$ is the pushforward of Lebesgue measure under $f$, the density of states of the stride, and its density a period of the fibres; the values are its moments, and $\log P$ is the log-partition function of an exponential family. At the centre the level sets are anticanonical hypersurfaces of $\mathbb P^{n-1}$ (elliptic curves, K3 surfaces and Calabi--Yau threefolds for $n=3,4,5$), and the equation in $t$ is their Picard--Fuchs equation; for three marks the pencil $(x+y+z)(xy+yz+zx)=\lambda xyz$ is modular of level $6$. Proposition~\ref{b:p:density} is the Mellin transform of the law, and its exponents are the Frobenius exponents. The top piece of a certificate is a Nullstellensatz certificate for the ideal of critical points and the whole a creative-telescoping certificate. Proposition~\ref{b:p:moment} is the Hausdorff moment problem on $[0,f^\star]$, generated by the four families of the Markov--Luk\'acs theorem; $e^{-g}$ in Theorem~\ref{b:t:rate} is $e^{-G}$ for the Green's function $G$ of $\mathbb C\setminus[0,f^\star]$ with pole at infinity (the growth of Chebyshev polynomials); Proposition~\ref{b:p:blind} is the indeterminacy of a moment problem along a Stieltjes transform. Proposition~\ref{a:p:gumbel} is the Gumbel law of Chapter~G. The coaction of Theorem~A is Brown's, and the group acting on rays is the mechanism of Rhin and Viola. Proposition~\ref{b:p:families} writes $\Gamma(sk)\Gamma(tk)P(k)$ as a Mellin transform of a product of Bessel functions of orders $k|c_e|$. In \S20.5, the frame $X$: $\Gamma(\frac1{15})\Gamma(\frac2{15})\Gamma(\frac4{15})\Gamma(\frac8{15})$ is a Chowla--Selberg period of $\mathbb Q(\sqrt{-15})$; the $F$ of Remark~\ref{c:r:balance} is the Koblitz--Ogus function, that a constant $F$ makes a Gamma monomial algebraic up to a power of $\pi$ is Deligne's theorem, and the converse is the Rohrlich--Lang conjecture. In \S20.6, Remark~\ref{c:r:transform}: at a nondegenerate gauge Katz's middle convolution with the quadratic character exchanges the local monodromy $-1$ with a unipotent block of size two, as adding a mark does; on a unipotent block of size $m\ge2$ it gives the sign with a block of size $m-1$, so it does not produce the exponent $1$ at $t=1$. In \S20.8: Proposition~\ref{b:p:merit} is Ap\'ery's criterion, and Remark~\ref{b:r:fixmerit} compares with Ap\'ery's recurrence for $\zeta(3)$, of characteristic roots $17\pm12\sqrt2$; the recurrence of Proposition~\ref{c:p:merit3}(i) is Zagier's sporadic case C, $x_j=\sum_k\binom jk^2\binom{2k}k$, and $\ell$ its Ap\'ery limit up to a rational factor; the irrationality of $L(\chi_{-3},2)$ is open \recalled.

\noindent\emph{\S21.} $\varrho$ is the period of the vanishing cycle of the Morse critical point $\alpha^\star$, and Theorem~\ref{c:t:gauge}(i) is Laplace's method with the Morse lemma. Theorem~\ref{c:t:gauge}(iii) is the Frobenius method at a regular singular point together with the extension of a solution by zero; the line is the image of the variation $T-1$ of the local monodromy at $f^\star$ (Picard--Lefschetz), and the forced exponent of (iv) is the local monodromy of a Morse point in dimension $n-1$: $T$ is a reflection when $n$ is even and a transvection when $n$ is odd, which is all Remark~\ref{c:r:parity} sees of the parity. The formal solution at a simple root is Birkhoff's, Laplace's method at an endpoint is Watson's lemma, and the reading is the numerical solution of a connection problem by Frobenius series. Proposition~\ref{c:p:jet}: (i) says that the Taylor series of $\log f$ at $\alpha^\star$ is the cumulant series of the letters' exponent laws; (ii) is Laplace's method with Wick's theorem; the germ at $f^\star$ is the Borel sum of the large-$k$ expansion; the count of formal solutions at a multiple root in Corollary~\ref{c:c:multiple} is the Birkhoff--Adams theory of linear recurrences. The share side of Ex.~\ref{b:e:Hslit} is a Mellin--Barnes integral, and that of Example~\ref{c:e:fivegauge} is $16\int_0^\infty xK_0(x)^5\,dx$, a Bessel moment.

\noindent\emph{\S22.} The torus value is the zeta Mahler measure of the letter (Akatsuka). At three marks $T(k)=W_3(-2k)$, the moments of the three-step uniform random walk; its density is that of the walk's distance in closed form, $\frac{2\sqrt3}\pi\frac x{3+x^2}\,{}_2F_1\big(\frac13,\frac23;1;\frac{x^2(9-x^2)^2}{(3+x^2)^3}\big)$ with $x=t^{-1/2}$ (Borwein, Straub, Wan, Zudilin); the closings are Zagier's sporadic sequence C. The $m_j$ are higher Mahler measures of $1+x+y$: $m_1$ is Smyth's $\frac{3\sqrt3}{4\pi}L(\chi_{-3},2)$, and $m_2=\frac3\pi\mathrm{Ls}_3(\frac{2\pi}3)+\frac{\pi^2}4$ (Borwein, Straub). The one-mark average in Corollary~\ref{d:c:three}(iv) is Jensen's formula with the hypergeometric function ${}_2F_1(k,k;1;\cdot)$. Theorem~\ref{d:t:connection} is, in Bessel form, an identity between the moments of $K_0(x)^3$, of $J_0(x)^3$ and of $K_0(x)^2I_0(x)$ (a Wick rotation).

\end{document}
